\documentclass[14pt,a4paper,oneside]{extbook}
\usepackage[utf8]{vietnam}
\usepackage{amsmath,amssymb,mathrsfs,fancyhdr,enumerate,multirow,makecell,venndiagram}
\usepackage{tkz-euclide,tikz-3dplot,tikz,tkz-tab}
\usetikzlibrary{arrows,calc,intersections,angles,quotes,backgrounds,shapes.geometric,patterns,through}
\usepackage{yhmath}
\usepackage{pgfplots}
\pgfplotsset{compat=1.9}
\usepgfplotslibrary{fillbetween}
\usepackage[tikz]{bclogo}
\usepackage[top=1.8cm, bottom=1.8cm, left=1.8cm, right=1.8cm, headheight=16pt]{geometry}
\usepackage[hidelinks,unicode]{hyperref}
\usepackage{qrcode}
\usepackage[many]{tcolorbox}

\usepackage[dethi]{ex_test}
\providecommand{\shortans}[1]{\par\vspace{1.5mm}\noindent\textbf{Đáp số: } \framebox[3.2cm]{\phantom{\rule{0pt}{1.8ex}}}\par\vspace{2mm}}

\def\colorEX{\color{blue}}
\everymath{\color{blue}}

\usepackage{fontawesome}
\usepackage{framed}
\usepackage{polynom}

\pagestyle{fancy}
\fancyhf{}
\fancyhead[L]{\small\itshape\color{gray} Tuyển tập Bài tập Toán 2 theo từng bài}
\fancyhead[R]{\small\bfseries\color{blue!80!black} Lớp Toán Thầy Lĩnh -- loptoanthaylinh.vn}
\fancyfoot[C]{\bfseries\thepage}
\renewcommand{\headrulewidth}{0.4pt}

\setcounter{secnumdepth}{4}
\setcounter{tocdepth}{2}
\usepackage{titlesec}
\usepackage{titletoc}

\makeatletter
\renewcommand\thesection{\@arabic\c@section}
\makeatother

\titleformat{\section}{\centering\Large\bfseries\color{blue!80!black}}{Chủ đề~\thesection.}{1ex}{
    \setcounter{ex}{0}
    \MakeUppercase
}
\titleformat{\subsection}{\large\bfseries\color{blue!70!black}}{}{0ex}{}

% KHUNG TIÊU ĐỀ PHIẾU BÀI TẬP KÈM QR CODE
\newcommand{\TaoPhieuBaiTap}[4]{%
  \refstepcounter{subsection}%
  \addcontentsline{toc}{subsection}{Bài #2. #3}%
  \setcounter{ex}{0}%
  \noindent
  \begin{tcolorbox}[
    colback=white,
    colframe=blue!85!black,
    arc=2.5mm,
    boxrule=1pt,
    left=2.5mm, right=2.5mm, top=1.8mm, bottom=1.8mm
  ]
    \begin{minipage}[c]{0.27\textwidth}
      \centering
      \begin{tcolorbox}[colback=blue!8, colframe=blue!80!black, boxrule=0.6pt, arc=1.5mm, halign=center, top=1.2mm, bottom=1.2mm, left=1mm, right=1mm]
        {\bfseries\color{blue!95!black}\small LỚP TOÁN THẦY LĨNH}\\[1.5pt]
        {\scriptsize\color{blue!70!black}\faGlobe~loptoanthaylinh.vn}\\[1.5pt]
        {\scriptsize\bfseries\color{blue!85!black}\faPhone~Hotline: 0866976878}
      \end{tcolorbox}
    \end{minipage}%
    \hfill
    \begin{minipage}[c]{0.56\textwidth}
      \centering
      {\footnotesize\bfseries\color{blue!80!black}\MakeUppercase{#1}}\\[2pt]
      {\large\bfseries\color{red!85!black}\MakeUppercase{BÀI #2: #3}}\\[2pt]
      {\scriptsize\itshape Phiếu Bài Tập Luyện Tập -- Hệ Thống Bài Tập Theo Buổi}
    \end{minipage}%
    \hfill
    \begin{minipage}[c]{0.14\textwidth}
      \centering
      \qrcode[height=1.35cm]{#4}\\[2pt]
      {\tiny\bfseries\color{blue!90!black}\faQrcode~Quét QR học}
    \end{minipage}
  \end{tcolorbox}
  \vspace{2mm}
}

\begin{document}

\tableofcontents
\newpage


% =========================================================================
% Ôn tập và bổ sung
% =========================================================================
\section{Ôn tập và bổ sung}

\TaoPhieuBaiTap{Chủ đề 1 -- Ôn tập và bổ sung}{01}{Ôn tập các số đến 100}{https://loptoanthaylinh.vn/bai-tap/toan-2/bai-01}
\begin{ex}
Số liền sau của số 99 là số nào?
\shortans{100}
\loigiai{
Số liền sau của 99 là 100.
}
\end{ex}

\begin{ex}
Số liền trước của số 100 là số nào?
\shortans{99}
\loigiai{
Số liền trước của 100 là 99.
}
\end{ex}

\begin{ex}
Số gồm 7 chục và 4 đơn vị là số nào?
\shortans{74}
\loigiai{
Số gồm 7 chục và 4 đơn vị là 74.
}
\end{ex}

\begin{ex}
Số 85 có chữ số hàng chục là bao nhiêu?
\shortans{8}
\loigiai{
Chữ số hàng chục của 85 là 8.
}
\end{ex}

\begin{ex}
Số 62 có chữ số hàng đơn vị là bao nhiêu?
\shortans{2}
\loigiai{
Chữ số hàng đơn vị của 62 là 2.
}
\end{ex}

\begin{ex}
Số lớn nhất có hai chữ số là số nào?
\shortans{99}
\loigiai{
Số lớn nhất có hai chữ số là 99.
}
\end{ex}

\begin{ex}
Số bé nhất có hai chữ số là số nào?
\shortans{10}
\loigiai{
Số bé nhất có hai chữ số là 10.
}
\end{ex}

\begin{ex}
Trong các số $45, 83, 38, 79$, số nào lớn nhất?
\shortans{83}
\loigiai{
Số lớn nhất là 83.
}
\end{ex}

\begin{ex}
Trong các số $62, 59, 70, 54$, số nào bé nhất?
\shortans{54}
\loigiai{
Số bé nhất là 54.
}
\end{ex}

\begin{ex}
Điền số thích hợp: $59 < \dots < 61$.
\shortans{60}
\loigiai{
Số tự nhiên ở giữa 59 và 61 là 60.
}
\end{ex}

\begin{ex}
Có bao nhiêu số tròn chục có hai chữ số?
\shortans{9}
\loigiai{
Các số tròn chục từ 10 đến 90 gồm 9 số.
}
\end{ex}

\begin{ex}
Có bao nhiêu số có hai chữ số giống nhau?
\shortans{9}
\loigiai{
Gồm 9 số: 11, 22, 33, 44, 55, 66, 77, 88, 99.
}
\end{ex}

\begin{ex}
Sắp xếp các số $35, 82, 19, 64$ từ bé đến lớn, số đứng ở vị trí thứ hai là số nào?
\shortans{35}
\loigiai{
Thứ tự: 19, 35, 64, 82. Số thứ hai là 35.
}
\end{ex}

\begin{ex}
Sắp xếp các số $71, 48, 93, 65$ từ lớn đến bé, số đứng ở vị trí thứ ba là số nào?
\shortans{65}
\loigiai{
Thứ tự: 93, 71, 65, 48. Số thứ ba là 65.
}
\end{ex}

\begin{ex}
Từ 20 đến 30 có tất cả bao nhiêu số tự nhiên?
\shortans{11}
\loigiai{
Có $(30 - 20) + 1 = 11$ số.
}
\end{ex}

\begin{ex}
Tìm số lớn nhất có hai chữ số mà tổng hai chữ số bằng 9.
\shortans{90}
\loigiai{
Số lớn nhất là 90 vì $9 + 0 = 9$ và có hàng chục lớn nhất.
}
\end{ex}

\begin{ex}
Tìm số bé nhất có hai chữ số mà tổng hai chữ số bằng 9.
\shortans{18}
\loigiai{
Số bé nhất có hàng chục là 1, hàng đơn vị là 8 ($1 + 8 = 9$), vậy là 18.
}
\end{ex}

\begin{ex}
Tính tổng chữ số hàng chục và hàng đơn vị của số 47.
\shortans{11}
\loigiai{
Tổng là: $4 + 7 = 11$.
}
\end{ex}

\begin{ex}
Hiệu giữa chữ số hàng chục và hàng đơn vị của số 82 là bao nhiêu?
\shortans{6}
\loigiai{
Hiệu là: $8 - 2 = 6$.
}
\end{ex}

\begin{ex}
Lớp 2A có 35 học sinh, lớp 2B có 38 học sinh. Lớp 2B nhiều hơn lớp 2A bao nhiêu học sinh?
\shortans{3}
\loigiai{
Nhiều hơn: $38 - 35 = 3$ học sinh.
}
\end{ex}

\begin{ex}
Có bao nhiêu số có một chữ số?
\shortans{10}
\loigiai{
Từ 0 đến 9 gồm 10 số.
}
\end{ex}

\begin{ex}
Có bao nhiêu số có hai chữ số?
\shortans{90}
\loigiai{
Từ 10 đến 99 gồm 90 số.
}
\end{ex}

\begin{ex}
Số tròn chục liền sau số 50 là số nào?
\shortans{60}
\loigiai{
Số tròn chục liền sau 50 là 60.
}
\end{ex}

\begin{ex}
Số tròn chục liền trước số 80 là số nào?
\shortans{70}
\loigiai{
Số tròn chục liền trước 80 là 70.
}
\end{ex}

\begin{ex}
Số nào trừ đi 1 thì được số lớn nhất có một chữ số?
\shortans{10}
\loigiai{
Số lớn nhất có 1 chữ số là 9. Số cần tìm: $9 + 1 = 10$.
}
\end{ex}

\newpage

\TaoPhieuBaiTap{Chủ đề 1 -- Ôn tập và bổ sung}{02}{Tia số. Số liền trước, số liền sau}{https://loptoanthaylinh.vn/bai-tap/toan-2/bai-02}
\begin{ex}
Gốc của tia số biểu diễn số tự nhiên nào?
\shortans{0}
\loigiai{
Gốc tia số bắt đầu từ số 0.
}
\end{ex}

\begin{ex}
Trên tia số, điểm đứng liền sau số 14 là số nào?
\shortans{15}
\loigiai{
Số liền sau số 14 là 15.
}
\end{ex}

\begin{ex}
Trên tia số, điểm đứng liền trước số 28 là số nào?
\shortans{27}
\loigiai{
Số liền trước số 28 là 27.
}
\end{ex}

\begin{ex}
Tìm số liền sau của số 69.
\shortans{70}
\loigiai{
Số liền sau số 69 là 70.
}
\end{ex}

\begin{ex}
Tìm số liền trước của số 90.
\shortans{89}
\loigiai{
Số liền trước số 90 là 89.
}
\end{ex}

\begin{ex}
Tìm số liền trước của số liền trước số 35.
\shortans{33}
\loigiai{
Số liền trước 35 là 34, liền trước 34 là 33.
}
\end{ex}

\begin{ex}
Tìm số liền sau của số liền sau số 48.
\shortans{50}
\loigiai{
Số liền sau 48 là 49, liền sau 49 là 50.
}
\end{ex}

\begin{ex}
Trên tia số, số nào nằm chính giữa số 30 và số 32?
\shortans{31}
\loigiai{
Số nằm chính giữa là 31.
}
\end{ex}

\begin{ex}
Trên tia số, khoảng cách từ vạch số 5 đến vạch số 12 là bao nhiêu đơn vị?
\shortans{7}
\loigiai{
Khoảng cách: $12 - 5 = 7$ đơn vị.
}
\end{ex}

\begin{ex}
Trên tia số, khoảng cách từ vạch số 20 đến vạch số 35 là bao nhiêu đơn vị?
\shortans{15}
\loigiai{
Khoảng cách: $35 - 20 = 15$ đơn vị.
}
\end{ex}

\begin{ex}
Điểm A trên tia số ở vạch 8, điểm B ở vạch 14. Hỏi từ điểm A phải bước thêm bao nhiêu bước đơn vị để đến điểm B?
\shortans{6}
\loigiai{
Số bước: $14 - 8 = 6$ bước.
}
\end{ex}

\begin{ex}
Một chú thỏ đang ở vạch số 15 trên tia số, thỏ nhảy tiến về phía trước 4 bước. Chú thỏ dừng lại ở vạch số mấy?
\shortans{19}
\loigiai{
Dừng lại ở: $15 + 4 = 19$.
}
\end{ex}

\begin{ex}
Một chú ếch đang ở vạch số 24 trên tia số, ếch nhảy lùi lại 6 bước. Chú ếch dừng lại ở vạch số mấy?
\shortans{18}
\loigiai{
Dừng lại ở: $24 - 6 = 18$.
}
\end{ex}

\begin{ex}
Tìm tổng của số liền trước số 20 và số liền sau số 20.
\shortans{40}
\loigiai{
Số liền trước là 19, liền sau là 21. Tổng: $19 + 21 = 40$.
}
\end{ex}

\begin{ex}
Tìm hiệu giữa số liền sau số 50 và số liền trước số 50.
\shortans{2}
\loigiai{
Số liền sau là 51, liền trước là 49. Hiệu: $51 - 49 = 2$.
}
\end{ex}

\begin{ex}
Số liền sau của số lớn nhất có hai chữ số là số nào?
\shortans{100}
\loigiai{
Số lớn nhất có hai chữ số là 99, số liền sau là 100.
}
\end{ex}

\begin{ex}
Số liền trước của số bé nhất có hai chữ số là số nào?
\shortans{9}
\loigiai{
Số bé nhất có hai chữ số là 10, số liền trước là 9.
}
\end{ex}

\begin{ex}
Trên tia số, giữa số 42 và số 47 có bao nhiêu số tự nhiên?
\shortans{4}
\loigiai{
Các số ở giữa: 43, 44, 45, 46. Có 4 số.
}
\end{ex}

\begin{ex}
Trên tia số từ vạch 0 đến vạch 10 có bao nhiêu đoạn đơn vị?
\shortans{10}
\loigiai{
Có đúng 10 đoạn đơn vị.
}
\end{ex}

\begin{ex}
Điểm M biểu diễn số 50. Điểm N nằm bên phải điểm M và cách M 8 đơn vị. Điểm N biểu diễn số nào?
\shortans{58}
\loigiai{
Điểm N biểu diễn: $50 + 8 = 58$.
}
\end{ex}

\begin{ex}
Điểm P biểu diễn số 75. Điểm Q nằm bên trái điểm P và cách P 15 đơn vị. Điểm Q biểu diễn số nào?
\shortans{60}
\loigiai{
Điểm Q biểu diễn: $75 - 15 = 60$.
}
\end{ex}

\begin{ex}
Số nào có số liền sau là 80?
\shortans{79}
\loigiai{
Số đó là: $80 - 1 = 79$.
}
\end{ex}

\begin{ex}
Số nào có số liền trước là 63?
\shortans{64}
\loigiai{
Số đó là: $63 + 1 = 64$.
}
\end{ex}

\begin{ex}
Có bao nhiêu số tự nhiên lớn hơn 15 và không vượt quá 20?
\shortans{5}
\loigiai{
Các số: 16, 17, 18, 19, 20. Có 5 số.
}
\end{ex}

\begin{ex}
Tổng của số liền trước số 10 và số liền sau số 10 là bao nhiêu?
\shortans{20}
\loigiai{
Tổng: $9 + 11 = 20$.
}
\end{ex}

\newpage

\TaoPhieuBaiTap{Chủ đề 1 -- Ôn tập và bổ sung}{03}{Các thành phần của phép cộng, phép trừ}{https://loptoanthaylinh.vn/bai-tap/toan-2/bai-03}
\begin{ex}
Trong phép tính $32 + 15 = 47$, số 47 được gọi là gì? Tính tổng của 32 và 15.
\shortans{47}
\loigiai{
Tổng của hai số là 47.
}
\end{ex}

\begin{ex}
Trong phép tính $25 + 40 = 65$, số 25 được gọi là số hạng thứ nhất. Số hạng thứ hai bằng bao nhiêu?
\shortans{40}
\loigiai{
Số hạng thứ hai là 40.
}
\end{ex}

\begin{ex}
Trong phép tính $58 - 23 = 35$, số bị trừ bằng bao nhiêu?
\shortans{58}
\loigiai{
Số bị trừ là 58.
}
\end{ex}

\begin{ex}
Trong phép tính $76 - 34 = 42$, số trừ bằng bao nhiêu?
\shortans{34}
\loigiai{
Số trừ là 34.
}
\end{ex}

\begin{ex}
Trong phép tính $89 - 50 = 39$, hiệu bằng bao nhiêu?
\shortans{39}
\loigiai{
Hiệu là 39.
}
\end{ex}

\begin{ex}
Biết hai số hạng là 24 và 35. Tính tổng của hai số đó.
\shortans{59}
\loigiai{
Tổng: $24 + 35 = 59$.
}
\end{ex}

\begin{ex}
Biết số bị trừ là 67, số trừ là 25. Tính hiệu của phép trừ.
\shortans{42}
\loigiai{
Hiệu: $67 - 25 = 42$.
}
\end{ex}

\begin{ex}
Biết tổng bằng 58, một số hạng bằng 20. Tìm số hạng còn lại.
\shortans{38}
\loigiai{
Số hạng còn lại: $58 - 20 = 38$.
}
\end{ex}

\begin{ex}
Biết tổng bằng 76, một số hạng bằng 32. Tìm số hạng còn lại.
\shortans{44}
\loigiai{
Số hạng còn lại: $76 - 32 = 44$.
}
\end{ex}

\begin{ex}
Biết số bị trừ là 84, hiệu là 31. Tìm số trừ.
\shortans{53}
\loigiai{
Số trừ: $84 - 31 = 53$.
}
\end{ex}

\begin{ex}
Biết số trừ là 26, hiệu là 43. Tìm số bị trừ.
\shortans{69}
\loigiai{
Số bị trừ: $43 + 26 = 69$.
}
\end{ex}

\begin{ex}
Nếu số hạng thứ nhất tăng thêm 5 đơn vị và giữ nguyên số hạng thứ hai thì tổng tăng thêm bao nhiêu đơn vị?
\shortans{5}
\loigiai{
Tổng tăng thêm đúng 5 đơn vị.
}
\end{ex}

\begin{ex}
Nếu số bị trừ tăng thêm 8 đơn vị và giữ nguyên số trừ thì hiệu tăng thêm bao nhiêu đơn vị?
\shortans{8}
\loigiai{
Hiệu tăng thêm 8 đơn vị.
}
\end{ex}

\begin{ex}
Nếu giữ nguyên số bị trừ và số trừ tăng thêm 4 đơn vị thì hiệu giảm đi bao nhiêu đơn vị?
\shortans{4}
\loigiai{
Hiệu giảm đi 4 đơn vị.
}
\end{ex}

\begin{ex}
Tổng của hai số bằng 60, nếu một số hạng bằng 60 thì số hạng kia bằng bao nhiêu?
\shortans{0}
\loigiai{
Số hạng kia là: $60 - 60 = 0$.
}
\end{ex}

\begin{ex}
Hiệu của hai số bằng 0, biết số bị trừ là 45. Số trừ bằng bao nhiêu?
\shortans{45}
\loigiai{
Hai số bằng nhau thì hiệu bằng 0, nên số trừ là 45.
}
\end{ex}

\begin{ex}
Hiệu của hai số bằng 38, biết số trừ là 0. Số bị trừ bằng bao nhiêu?
\shortans{38}
\loigiai{
Số bị trừ: $38 + 0 = 38$.
}
\end{ex}

\begin{ex}
Tìm số bị trừ biết số trừ là số bé nhất có hai chữ số và hiệu là 52.
\shortans{62}
\loigiai{
Số trừ là 10. Số bị trừ: $52 + 10 = 62$.
}
\end{ex}

\begin{ex}
Tìm số hạng thứ hai biết tổng là số lớn nhất có hai chữ số và số hạng thứ nhất là 70.
\shortans{29}
\loigiai{
Tổng là 99. Số hạng thứ hai: $99 - 70 = 29$.
}
\end{ex}

\begin{ex}
Tìm hiệu biết số bị trừ là số tròn chục lớn nhất có hai chữ số và số trừ là 30.
\shortans{60}
\loigiai{
Số bị trừ là 90. Hiệu: $90 - 30 = 60$.
}
\end{ex}

\begin{ex}
Hai số có tổng bằng 88. Nếu số hạng thứ nhất là 55 thì số hạng thứ hai là bao nhiêu?
\shortans{33}
\loigiai{
Số hạng thứ hai: $88 - 55 = 33$.
}
\end{ex}

\begin{ex}
Tìm số trừ biết số bị trừ là 79 và hiệu là 46.
\shortans{33}
\loigiai{
Số trừ: $79 - 46 = 33$.
}
\end{ex}

\begin{ex}
Trong một phép cộng có hai số hạng bằng nhau và tổng bằng 18. Số hạng đó bằng bao nhiêu?
\shortans{9}
\loigiai{
Vì $9 + 9 = 18$ nên số hạng là 9.
}
\end{ex}

\begin{ex}
Trong một phép cộng có hai số hạng bằng nhau và tổng bằng 16. Số hạng đó bằng bao nhiêu?
\shortans{8}
\loigiai{
Vì $8 + 8 = 16$ nên số hạng là 8.
}
\end{ex}

\begin{ex}
Một phép trừ có số bị trừ bằng 50, số trừ bằng 50. Hiệu bằng bao nhiêu?
\shortans{0}
\loigiai{
Hiệu là: $50 - 50 = 0$.
}
\end{ex}

\newpage

\TaoPhieuBaiTap{Chủ đề 1 -- Ôn tập và bổ sung}{04}{Hơn, kém nhau bao nhiêu đơn vị}{https://loptoanthaylinh.vn/bai-tap/toan-2/bai-04}
\begin{ex}
Bạn Lan có 15 bông hoa, bạn Mai có 10 bông hoa. Bạn Lan nhiều hơn bạn Mai bao nhiêu bông hoa?
\shortans{5}
\loigiai{
Nhiều hơn: $15 - 10 = 5$ bông hoa.
}
\end{ex}

\begin{ex}
Bác nông dân có 38 con gà và 24 con vịt. Số gà nhiều hơn số vịt bao nhiêu con?
\shortans{14}
\loigiai{
Nhiều hơn: $38 - 24 = 14$ con.
}
\end{ex}

\begin{ex}
Lớp 2A có 18 bạn nữ và 15 bạn nam. Số bạn nữ nhiều hơn số bạn nam bao nhiêu bạn?
\shortans{3}
\loigiai{
Nhiều hơn: $18 - 15 = 3$ bạn.
}
\end{ex}

\begin{ex}
Đoạn thẳng AB dài 28 cm, đoạn thẳng CD dài 20 cm. Đoạn thẳng CD ngắn hơn đoạn thẳng AB bao nhiêu xăng-ti-mét?
\shortans{8}
\loigiai{
Ngắn hơn: $28 - 20 = 8$ cm.
}
\end{ex}

\begin{ex}
Cây cau cao 12 mét, cây ổi cao 4 mét. Cây cau cao hơn cây ổi bao nhiêu mét?
\shortans{8}
\loigiai{
Cao hơn: $12 - 4 = 8$ mét.
}
\end{ex}

\begin{ex}
Anh 16 tuổi, em 9 tuổi. Anh nhiều hơn em bao nhiêu tuổi?
\shortans{7}
\loigiai{
Nhiều hơn: $16 - 9 = 7$ tuổi.
}
\end{ex}

\begin{ex}
Sau 5 năm nữa, anh vẫn nhiều hơn em bao nhiêu tuổi?
\shortans{7}
\loigiai{
Hiệu số tuổi không đổi theo thời gian, anh vẫn hơn em 7 tuổi.
}
\end{ex}

\begin{ex}
Mẹ 36 tuổi, con 8 tuổi. Mẹ nhiều hơn con bao nhiêu tuổi?
\shortans{28}
\loigiai{
Nhiều hơn: $36 - 8 = 28$ tuổi.
}
\end{ex}

\begin{ex}
Bình có 45 viên bi, Nam có 32 viên bi. Nam ít hơn Bình bao nhiêu viên bi?
\shortans{13}
\loigiai{
Ít hơn: $45 - 32 = 13$ viên bi.
}
\end{ex}

\begin{ex}
Quyển sách Toán dày 84 trang, quyển vở bài tập dày 50 trang. Quyển sách dày hơn quyển vở bao nhiêu trang?
\shortans{34}
\loigiai{
Dày hơn: $84 - 50 = 34$ trang.
}
\end{ex}

\begin{ex}
Một cửa hàng buổi sáng bán được 65 chiếc xe đạp, buổi chiều bán được 42 chiếc xe đạp. Buổi sáng bán nhiều hơn buổi chiều bao nhiêu chiếc xe đạp?
\shortans{23}
\loigiai{
Nhiều hơn: $65 - 42 = 23$ chiếc.
}
\end{ex}

\begin{ex}
Thùng thứ nhất đựng 48 lít dầu, thùng thứ hai đựng 35 lít dầu. Thùng thứ hai ít hơn thùng thứ nhất bao nhiêu lít dầu?
\shortans{13}
\loigiai{
Ít hơn: $48 - 35 = 13$ lít dầu.
}
\end{ex}

\begin{ex}
Bạn An cao 125 cm, bạn Bình cao 120 cm. Bạn An cao hơn bạn Bình bao nhiêu xăng-ti-mét?
\shortans{5}
\loigiai{
Cao hơn: $125 - 120 = 5$ cm.
}
\end{ex}

\begin{ex}
Bao gạo to nặng 50 kg, bao gạo nhỏ nặng 20 kg. Bao to nặng hơn bao nhỏ bao nhiêu ki-lô-gam?
\shortans{30}
\loigiai{
Nặng hơn: $50 - 20 = 30$ kg.
}
\end{ex}

\begin{ex}
Tổ Một gấp được 46 ngôi sao, tổ Hai gấp được 33 ngôi sao. Tổ Hai gấp kém tổ Một bao nhiêu ngôi sao?
\shortans{13}
\loigiai{
Kém: $46 - 33 = 13$ ngôi sao.
}
\end{ex}

\begin{ex}
Tìm số lớn biết số bé là 25 và số lớn hơn số bé 14 đơn vị.
\shortans{39}
\loigiai{
Số lớn là: $25 + 14 = 39$.
}
\end{ex}

\begin{ex}
Tìm số bé biết số lớn là 68 và số lớn hơn số bé 26 đơn vị.
\shortans{42}
\loigiai{
Số bé là: $68 - 26 = 42$.
}
\end{ex}

\begin{ex}
Hai số hơn kém nhau 15 đơn vị, biết số lớn là 70. Số bé là bao nhiêu?
\shortans{55}
\loigiai{
Số bé là: $70 - 15 = 55$.
}
\end{ex}

\begin{ex}
Hai số hơn kém nhau 18 đơn vị, biết số bé là 31. Số lớn là bao nhiêu?
\shortans{49}
\loigiai{
Số lớn là: $31 + 18 = 49$.
}
\end{ex}

\begin{ex}
Trong bãi xe có 54 ô tô và 30 xe máy. Số ô tô nhiều hơn số xe máy bao nhiêu chiếc?
\shortans{24}
\loigiai{
Nhiều hơn: $54 - 30 = 24$ chiếc.
}
\end{ex}

\begin{ex}
Đàn gà có 27 con gà mái và 5 con gà trống. Số gà mái nhiều hơn số gà trống bao nhiêu con?
\shortans{22}
\loigiai{
Nhiều hơn: $27 - 5 = 22$ con.
}
\end{ex}

\begin{ex}
Ngăn trên có 45 cuốn sách, ngăn dưới có 32 cuốn sách. Ngăn trên nhiều hơn ngăn dưới bao nhiêu cuốn sách?
\shortans{13}
\loigiai{
Nhiều hơn: $45 - 32 = 13$ cuốn sách.
}
\end{ex}

\begin{ex}
Bác thợ mộc cưa thanh gỗ dài 95 cm thành hai đoạn: đoạn thứ nhất dài 50 cm. Đoạn thứ nhất dài hơn đoạn thứ hai bao nhiêu xăng-ti-mét?
\shortans{5}
\loigiai{
Đoạn 2 dài: $95 - 50 = 45$ cm. Chênh lệch: $50 - 45 = 5$ cm.
}
\end{ex}

\begin{ex}
Số 78 hơn số 43 bao nhiêu đơn vị?
\shortans{35}
\loigiai{
Hơn: $78 - 43 = 35$ đơn vị.
}
\end{ex}

\begin{ex}
Số 24 kém số 69 bao nhiêu đơn vị?
\shortans{45}
\loigiai{
Kém: $69 - 24 = 45$ đơn vị.
}
\end{ex}

\newpage

\TaoPhieuBaiTap{Chủ đề 1 -- Ôn tập và bổ sung}{05}{Bảng cộng và bảng trừ (có nhớ)}{https://loptoanthaylinh.vn/bai-tap/toan-2/bai-05}
\begin{ex}
Tính: $35 + 24$.
\shortans{59}
\loigiai{
Ta có $35 + 24 = 59$.
}
\end{ex}

\begin{ex}
Tính: $68 - 45$.
\shortans{23}
\loigiai{
Ta có $68 - 45 = 23$.
}
\end{ex}

\begin{ex}
Tính: $42 + 6$.
\shortans{48}
\loigiai{
Ta có $42 + 6 = 48$.
}
\end{ex}

\begin{ex}
Tính: $79 - 7$.
\shortans{72}
\loigiai{
Ta có $79 - 7 = 72$.
}
\end{ex}

\begin{ex}
Tính: $50 + 36$.
\shortans{86}
\loigiai{
Ta có $50 + 36 = 86$.
}
\end{ex}

\begin{ex}
Tính: $94 - 40$.
\shortans{54}
\loigiai{
Ta có $94 - 40 = 54$.
}
\end{ex}

\begin{ex}
Tính giá trị biểu thức: $32 + 15 - 24$.
\shortans{23}
\loigiai{
Ta có $32 + 15 = 47$, rồi $47 - 24 = 23$.
}
\end{ex}

\begin{ex}
Tính giá trị biểu thức: $87 - 35 + 21$.
\shortans{73}
\loigiai{
Ta có $87 - 35 = 52$, rồi $52 + 21 = 73$.
}
\end{ex}

\begin{ex}
Tính giá trị biểu thức: $20 + 30 + 40$.
\shortans{90}
\loigiai{
Ta có $20 + 30 = 50$, rồi $50 + 40 = 90$.
}
\end{ex}

\begin{ex}
Tính giá trị biểu thức: $90 - 40 - 20$.
\shortans{30}
\loigiai{
Ta có $90 - 40 = 50$, rồi $50 - 20 = 30$.
}
\end{ex}

\begin{ex}
Điền số thích hợp: $24 + \dots = 67$.
\shortans{43}
\loigiai{
Số cần điền: $67 - 24 = 43$.
}
\end{ex}

\begin{ex}
Điền số thích hợp: $89 - \dots = 34$.
\shortans{55}
\loigiai{
Số cần điền: $89 - 34 = 55$.
}
\end{ex}

\begin{ex}
Điền số thích hợp: $\dots + 31 = 75$.
\shortans{44}
\loigiai{
Số cần điền: $75 - 31 = 44$.
}
\end{ex}

\begin{ex}
Điền số thích hợp: $\dots - 23 = 52$.
\shortans{75}
\loigiai{
Số cần điền: $52 + 23 = 75$.
}
\end{ex}

\begin{ex}
Lớp 2A có 18 học sinh nam và 21 học sinh nữ. Hỏi lớp 2A có tất cả bao nhiêu học sinh?
\shortans{39}
\loigiai{
Cả lớp có: $18 + 21 = 39$ học sinh.
}
\end{ex}

\begin{ex}
Một cửa hàng nhập về 78 gói bánh, đã bán được 46 gói. Cửa hàng còn lại bao nhiêu gói bánh?
\shortans{32}
\loigiai{
Còn lại: $78 - 46 = 32$ gói bánh.
}
\end{ex}

\begin{ex}
Bác Ba thu hoạch được 54 bắp ngô nếp và 45 bắp ngô ngọt. Bác thu hoạch được tất cả bao nhiêu bắp ngô?
\shortans{99}
\loigiai{
Tổng số bắp: $54 + 45 = 99$ bắp ngô.
}
\end{ex}

\begin{ex}
Đàn gà có 68 con, bác nông dân bán đi 25 con. Đàn gà còn lại bao nhiêu con?
\shortans{43}
\loigiai{
Còn lại: $68 - 25 = 43$ con gà.
}
\end{ex}

\begin{ex}
Đoạn dây thứ nhất dài 45 cm, đoạn thứ hai dài 34 cm. Cả hai đoạn dây dài bao nhiêu xăng-ti-mét?
\shortans{79}
\loigiai{
Tổng chiều dài: $45 + 34 = 79$ cm.
}
\end{ex}

\begin{ex}
Sợi dây đồng dài 88 cm, cắt bớt đi 35 cm. Sợi dây đồng còn lại dài bao nhiêu xăng-ti-mét?
\shortans{53}
\loigiai{
Còn lại: $88 - 35 = 53$ cm.
}
\end{ex}

\begin{ex}
Số nào cộng với 32 thì được 76?
\shortans{44}
\loigiai{
Số đó là: $76 - 32 = 44$.
}
\end{ex}

\begin{ex}
Số nào trừ đi 25 thì được 41?
\shortans{66}
\loigiai{
Số đó là: $41 + 25 = 66$.
}
\end{ex}

\begin{ex}
Lấy 98 trừ đi số tròn chục lớn nhất có hai chữ số thì được kết quả là bao nhiêu?
\shortans{8}
\loigiai{
Số tròn chục lớn nhất có hai chữ số là 90. Hiệu: $98 - 90 = 8$.
}
\end{ex}

\begin{ex}
Tính tổng của số bé nhất có hai chữ số và số 75.
\shortans{85}
\loigiai{
Số bé nhất có hai chữ số là 10. Tổng: $10 + 75 = 85$.
}
\end{ex}

\begin{ex}
Điền số thích hợp vào chỗ chấm: $45 + 13 - \dots = 20$.
\shortans{38}
\loigiai{
Ta có $45 + 13 = 58$, vậy số cần điền là $58 - 20 = 38$.
}
\end{ex}

\newpage

\TaoPhieuBaiTap{Chủ đề 1 -- Ôn tập và bổ sung}{06}{Luyện tập chung}{https://loptoanthaylinh.vn/bai-tap/toan-2/bai-06}
\begin{ex}
Số gồm 9 chục và 9 đơn vị là số nào?
\shortans{99}
\loigiai{
Số đó là 99.
}
\end{ex}

\begin{ex}
Số liền sau của số 79 là số nào?
\shortans{80}
\loigiai{
Số liền sau của 79 là 80.
}
\end{ex}

\begin{ex}
Số liền trước của số 60 là số nào?
\shortans{59}
\loigiai{
Số liền trước của 60 là 59.
}
\end{ex}

\begin{ex}
Tính: $43 + 35$.
\shortans{78}
\loigiai{
Ta có $43 + 35 = 78$.
}
\end{ex}

\begin{ex}
Tính: $86 - 52$.
\shortans{34}
\loigiai{
Ta có $86 - 52 = 34$.
}
\end{ex}

\begin{ex}
Trong phép trừ $75 - 23 = 52$, số bị trừ là bao nhiêu?
\shortans{75}
\loigiai{
Số bị trừ là 75.
}
\end{ex}

\begin{ex}
Trong phép cộng $34 + 45 = 79$, tổng là bao nhiêu?
\shortans{79}
\loigiai{
Tổng là 79.
}
\end{ex}

\begin{ex}
Bạn Hà có 24 cái kẹo, bạn Nga có 11 cái kẹo. Bạn Hà nhiều hơn bạn Nga bao nhiêu cái kẹo?
\shortans{13}
\loigiai{
Nhiều hơn: $24 - 11 = 13$ cái kẹo.
}
\end{ex}

\begin{ex}
Đoạn thẳng MN dài 48 cm, đoạn thẳng PQ dài 35 cm. Đoạn MN dài hơn đoạn PQ bao nhiêu xăng-ti-mét?
\shortans{13}
\loigiai{
Dài hơn: $48 - 35 = 13$ cm.
}
\end{ex}

\begin{ex}
Điền số thích hợp: $41 < \dots < 43$.
\shortans{42}
\loigiai{
Số ở giữa là 42.
}
\end{ex}

\begin{ex}
Tính giá trị biểu thức: $54 + 23 - 31$.
\shortans{46}
\loigiai{
Ta có $54 + 23 = 77$, rồi $77 - 31 = 46$.
}
\end{ex}

\begin{ex}
Tính giá trị biểu thức: $89 - 46 + 15$.
\shortans{58}
\loigiai{
Ta có $89 - 46 = 43$, rồi $43 + 15 = 58$.
}
\end{ex}

\begin{ex}
Tìm số tròn chục lớn nhất có hai chữ số.
\shortans{90}
\loigiai{
Số đó là 90.
}
\end{ex}

\begin{ex}
Tìm số tròn chục bé nhất có hai chữ số.
\shortans{10}
\loigiai{
Số đó là 10.
}
\end{ex}

\begin{ex}
Hiệu của số tròn chục lớn nhất và số tròn chục bé nhất có hai chữ số là bao nhiêu?
\shortans{80}
\loigiai{
Hiệu: $90 - 10 = 80$.
}
\end{ex}

\begin{ex}
Một cửa hàng có 65 bao xi măng, đã bán được 32 bao. Cửa hàng còn lại bao nhiêu bao xi măng?
\shortans{33}
\loigiai{
Còn lại: $65 - 32 = 33$ bao.
}
\end{ex}

\begin{ex}
Thùng cam có 42 quả, thùng quýt có 36 quả. Cả hai thùng có tất cả bao nhiêu quả?
\shortans{78}
\loigiai{
Tổng số quả: $42 + 36 = 78$ quả.
}
\end{ex}

\begin{ex}
Tìm một số biết rằng số đó cộng với 25 thì được 68.
\shortans{43}
\loigiai{
Số đó là: $68 - 25 = 43$.
}
\end{ex}

\begin{ex}
Tìm một số biết rằng số đó trừ đi 34 thì được 45.
\shortans{79}
\loigiai{
Số đó là: $45 + 34 = 79$.
}
\end{ex}

\begin{ex}
Trên tia số, điểm biểu diễn số 45 cách điểm biểu diễn số 60 bao nhiêu đơn vị?
\shortans{15}
\loigiai{
Khoảng cách: $60 - 45 = 15$ đơn vị.
}
\end{ex}

\begin{ex}
Tìm tổng của số lớn nhất có một chữ số và số bé nhất có hai chữ số.
\shortans{19}
\loigiai{
Số lớn nhất có 1 chữ số là 9, số bé nhất có 2 chữ số là 10. Tổng: $9 + 10 = 19$.
}
\end{ex}

\begin{ex}
Hai số có hiệu bằng 24, nếu số trừ là 35 thì số bị trừ là bao nhiêu?
\shortans{59}
\loigiai{
Số bị trừ: $24 + 35 = 59$.
}
\end{ex}

\begin{ex}
Hai số có tổng bằng 97, nếu một số hạng là 54 thì số hạng kia là bao nhiêu?
\shortans{43}
\loigiai{
Số hạng kia: $97 - 54 = 43$.
}
\end{ex}

\begin{ex}
Bác An năm nay 58 tuổi, chú Bình 36 tuổi. Bác An nhiều hơn chú Bình bao nhiêu tuổi?
\shortans{22}
\loigiai{
Nhiều hơn: $58 - 36 = 22$ tuổi.
}
\end{ex}

\begin{ex}
Điền số thích hợp: $30 + 40 + \dots = 95$.
\shortans{25}
\loigiai{
Ta có $30 + 40 = 70$, vậy số cần điền là $95 - 70 = 25$.
}
\end{ex}

\newpage


% =========================================================================
% Phép cộng, phép trừ trong phạm vi 20
% =========================================================================
\section{Phép cộng, phép trừ trong phạm vi 20}

\TaoPhieuBaiTap{Chủ đề 2 -- Phép cộng, phép trừ trong phạm vi 20}{07}{Phép cộng có tổng là số tròn chục}{https://loptoanthaylinh.vn/bai-tap/toan-2/bai-07}
\begin{ex}
Tính kết quả của phép tính: $9 + 2$.
\shortans{11}
\loigiai{
Tách 2 thành 1 và 1: $9 + 1 = 10$, rồi $10 + 1 = 11$.
}
\end{ex}

\begin{ex}
Tính kết quả của phép tính: $9 + 3$.
\shortans{12}
\loigiai{
Ta có $9 + 3 = 12$.
}
\end{ex}

\begin{ex}
Tính kết quả của phép tính: $9 + 4$.
\shortans{13}
\loigiai{
Ta có $9 + 4 = 13$.
}
\end{ex}

\begin{ex}
Tính kết quả của phép tính: $9 + 5$.
\shortans{14}
\loigiai{
Ta có $9 + 5 = 14$.
}
\end{ex}

\begin{ex}
Tính kết quả của phép tính: $9 + 6$.
\shortans{15}
\loigiai{
Ta có $9 + 6 = 15$.
}
\end{ex}

\begin{ex}
Tính kết quả của phép tính: $9 + 7$.
\shortans{16}
\loigiai{
Ta có $9 + 7 = 16$.
}
\end{ex}

\begin{ex}
Tính kết quả của phép tính: $9 + 8$.
\shortans{17}
\loigiai{
Ta có $9 + 8 = 17$.
}
\end{ex}

\begin{ex}
Tính kết quả của phép tính: $9 + 9$.
\shortans{18}
\loigiai{
Ta có $9 + 9 = 18$.
}
\end{ex}

\begin{ex}
Tính kết quả của phép tính: $8 + 3$.
\shortans{11}
\loigiai{
Tách 3 thành 2 và 1: $8 + 2 = 10$, rồi $10 + 1 = 11$.
}
\end{ex}

\begin{ex}
Tính kết quả của phép tính: $8 + 4$.
\shortans{12}
\loigiai{
Ta có $8 + 4 = 12$.
}
\end{ex}

\begin{ex}
Tính kết quả của phép tính: $8 + 5$.
\shortans{13}
\loigiai{
Ta có $8 + 5 = 13$.
}
\end{ex}

\begin{ex}
Tính kết quả của phép tính: $8 + 6$.
\shortans{14}
\loigiai{
Ta có $8 + 6 = 14$.
}
\end{ex}

\begin{ex}
Tính kết quả của phép tính: $8 + 7$.
\shortans{15}
\loigiai{
Ta có $8 + 7 = 15$.
}
\end{ex}

\begin{ex}
Tính kết quả của phép tính: $8 + 8$.
\shortans{16}
\loigiai{
Ta có $8 + 8 = 16$.
}
\end{ex}

\begin{ex}
Tính kết quả của phép tính: $7 + 4$.
\shortans{11}
\loigiai{
Tách 4 thành 3 và 1: $7 + 3 = 10$, rồi $10 + 1 = 11$.
}
\end{ex}

\begin{ex}
Tính kết quả của phép tính: $7 + 5$.
\shortans{12}
\loigiai{
Ta có $7 + 5 = 12$.
}
\end{ex}

\begin{ex}
Tính kết quả của phép tính: $7 + 6$.
\shortans{13}
\loigiai{
Ta có $7 + 6 = 13$.
}
\end{ex}

\begin{ex}
Tính kết quả của phép tính: $7 + 7$.
\shortans{14}
\loigiai{
Ta có $7 + 7 = 14$.
}
\end{ex}

\begin{ex}
Tính kết quả của phép tính: $6 + 5$.
\shortans{11}
\loigiai{
Tách 5 thành 4 và 1: $6 + 4 = 10$, rồi $10 + 1 = 11$.
}
\end{ex}

\begin{ex}
Tính kết quả của phép tính: $6 + 6$.
\shortans{12}
\loigiai{
Ta có $6 + 6 = 12$.
}
\end{ex}

\begin{ex}
Tính kết quả của phép tính: $6 + 7$.
\shortans{13}
\loigiai{
Ta có $6 + 7 = 13$.
}
\end{ex}

\begin{ex}
Tính kết quả của phép tính: $6 + 8$.
\shortans{14}
\loigiai{
Ta có $6 + 8 = 14$.
}
\end{ex}

\begin{ex}
Tính kết quả của phép tính: $6 + 9$.
\shortans{15}
\loigiai{
Ta có $6 + 9 = 15$.
}
\end{ex}

\begin{ex}
Mai có 9 bông hoa màu đỏ và 5 bông hoa màu vàng. Mai có tất cả bao nhiêu bông hoa?
\shortans{14}
\loigiai{
Mai có: $9 + 5 = 14$ bông hoa.
}
\end{ex}

\begin{ex}
Trong sân có 8 con gà mái và 7 con gà con. Trong sân có tất cả bao nhiêu con gà?
\shortans{15}
\loigiai{
Có tất cả: $8 + 7 = 15$ con gà.
}
\end{ex}

\newpage

\TaoPhieuBaiTap{Chủ đề 2 -- Phép cộng, phép trừ trong phạm vi 20}{08}{Phép cộng (có nhớ) trong phạm vi 20}{https://loptoanthaylinh.vn/bai-tap/toan-2/bai-08}
\begin{ex}
Tính: $9 + 8$.
\shortans{17}
\loigiai{
Ta có $9 + 8 = 17$.
}
\end{ex}

\begin{ex}
Tính: $8 + 9$.
\shortans{17}
\loigiai{
Phép cộng có tính giao hoán: $8 + 9 = 9 + 8 = 17$.
}
\end{ex}

\begin{ex}
Tính: $7 + 8$.
\shortans{15}
\loigiai{
Ta có $7 + 8 = 15$.
}
\end{ex}

\begin{ex}
Tính: $8 + 7$.
\shortans{15}
\loigiai{
Ta có $8 + 7 = 15$.
}
\end{ex}

\begin{ex}
Tính: $6 + 8$.
\shortans{14}
\loigiai{
Ta có $6 + 8 = 14$.
}
\end{ex}

\begin{ex}
Tính: $8 + 6$.
\shortans{14}
\loigiai{
Ta có $8 + 6 = 14$.
}
\end{ex}

\begin{ex}
Tính: $5 + 6$.
\shortans{11}
\loigiai{
Ta có $5 + 6 = 11$.
}
\end{ex}

\begin{ex}
Tính: $5 + 7$.
\shortans{12}
\loigiai{
Ta có $5 + 7 = 12$.
}
\end{ex}

\begin{ex}
Tính: $5 + 8$.
\shortans{13}
\loigiai{
Ta có $5 + 8 = 13$.
}
\end{ex}

\begin{ex}
Tính: $5 + 9$.
\shortans{14}
\loigiai{
Ta có $5 + 9 = 14$.
}
\end{ex}

\begin{ex}
Tính: $4 + 7$.
\shortans{11}
\loigiai{
Ta có $4 + 7 = 11$.
}
\end{ex}

\begin{ex}
Tính: $4 + 8$.
\shortans{12}
\loigiai{
Ta có $4 + 8 = 12$.
}
\end{ex}

\begin{ex}
Tính: $4 + 9$.
\shortans{13}
\loigiai{
Ta có $4 + 9 = 13$.
}
\end{ex}

\begin{ex}
Tính: $3 + 8$.
\shortans{11}
\loigiai{
Ta có $3 + 8 = 11$.
}
\end{ex}

\begin{ex}
Tính: $3 + 9$.
\shortans{12}
\loigiai{
Ta có $3 + 9 = 12$.
}
\end{ex}

\begin{ex}
Tính: $2 + 9$.
\shortans{11}
\loigiai{
Ta có $2 + 9 = 11$.
}
\end{ex}

\begin{ex}
Tính giá trị biểu thức: $9 + 4 + 3$.
\shortans{16}
\loigiai{
Ta có $9 + 4 = 13$, rồi $13 + 3 = 16$.
}
\end{ex}

\begin{ex}
Tính giá trị biểu thức: $8 + 5 + 2$.
\shortans{15}
\loigiai{
Ta có $8 + 2 = 10$, rồi $10 + 5 = 15$.
}
\end{ex}

\begin{ex}
Tính giá trị biểu thức: $7 + 6 + 3$.
\shortans{16}
\loigiai{
Ta có $7 + 3 = 10$, rồi $10 + 6 = 16$.
}
\end{ex}

\begin{ex}
Tính giá trị biểu thức: $6 + 8 + 4$.
\shortans{18}
\loigiai{
Ta có $6 + 4 = 10$, rồi $10 + 8 = 18$.
}
\end{ex}

\begin{ex}
Điền số thích hợp: $7 + \dots = 16$.
\shortans{9}
\loigiai{
Số cần điền là: $16 - 7 = 9$.
}
\end{ex}

\begin{ex}
Điền số thích hợp: $8 + \dots = 13$.
\shortans{5}
\loigiai{
Số cần điền là: $13 - 8 = 5$.
}
\end{ex}

\begin{ex}
Điền số thích hợp: $\dots + 6 = 15$.
\shortans{9}
\loigiai{
Số cần điền là: $15 - 6 = 9$.
}
\end{ex}

\begin{ex}
Điền số thích hợp: $\dots + 8 = 14$.
\shortans{6}
\loigiai{
Số cần điền là: $14 - 8 = 6$.
}
\end{ex}

\begin{ex}
Tổng của hai số có một chữ số lớn nhất là bao nhiêu?
\shortans{18}
\loigiai{
Số lớn nhất có một chữ số là 9. Tổng: $9 + 9 = 18$.
}
\end{ex}

\newpage

\TaoPhieuBaiTap{Chủ đề 2 -- Phép cộng, phép trừ trong phạm vi 20}{09}{Phép trừ (có nhớ) trong phạm vi 20}{https://loptoanthaylinh.vn/bai-tap/toan-2/bai-09}
\begin{ex}
Tổ Em có 8 bạn, cô giáo chuyển thêm 5 bạn nữa vào tổ. Hỏi tổ Em có tất cả bao nhiêu bạn?
\shortans{13}
\loigiai{
Tổ có tất cả: $8 + 5 = 13$ bạn.
}
\end{ex}

\begin{ex}
Trên cành cây có 9 chú chim đang đậu, có 4 chú chim khác bay đến cùng đậu. Hỏi trên cành cây có tất cả bao nhiêu chú chim?
\shortans{13}
\loigiai{
Có tất cả: $9 + 4 = 13$ chú chim.
}
\end{ex}

\begin{ex}
Bình có 7 viên bi, bố cho thêm Bình 6 viên bi nữa. Hỏi Bình có tất cả bao nhiêu viên bi?
\shortans{13}
\loigiai{
Bình có: $7 + 6 = 13$ viên bi.
}
\end{ex}

\begin{ex}
Trong rổ có 8 quả cam, mẹ bỏ thêm vào rổ 7 quả cam nữa. Hỏi trong rổ có tất cả bao nhiêu quả cam?
\shortans{15}
\loigiai{
Trong rổ có: $8 + 7 = 15$ quả cam.
}
\end{ex}

\begin{ex}
Đàn vịt đang bơi dưới ao có 9 con, sau đó có thêm 8 con vịt nữa nhảy xuống ao. Hỏi dưới ao có bao nhiêu con vịt?
\shortans{17}
\loigiai{
Dưới ao có: $9 + 8 = 17$ con vịt.
}
\end{ex}

\begin{ex}
Bác Ba nuôi 15 con thỏ, bác đã bán bớt đi 6 con. Hỏi bác Ba còn lại bao nhiêu con thỏ?
\shortans{9}
\loigiai{
Bác còn lại: $15 - 6 = 9$ con thỏ.
}
\end{ex}

\begin{ex}
Một cành hoa có 14 bông hoa, gió làm rụng bớt 5 bông. Trên cành còn lại bao nhiêu bông hoa?
\shortans{9}
\loigiai{
Còn lại: $14 - 5 = 9$ bông hoa.
}
\end{ex}

\begin{ex}
Lan có 16 cái kẹo, Lan cho các bạn 7 cái kẹo. Lan còn lại bao nhiêu cái kẹo?
\shortans{9}
\loigiai{
Lan còn lại: $16 - 7 = 9$ cái kẹo.
}
\end{ex}

\begin{ex}
Trong bãi xe có 18 chiếc xe máy, lúc sau có 9 chiếc xe máy rời khỏi bãi. Trong bãi còn lại bao nhiêu chiếc xe máy?
\shortans{9}
\loigiai{
Còn lại: $18 - 9 = 9$ chiếc xe máy.
}
\end{ex}

\begin{ex}
Một đàn gà có 13 con, mẹ đã nhốt bớt 4 con vào chuồng. Ngoài sân còn lại bao nhiêu con gà?
\shortans{9}
\loigiai{
Còn lại ngoài sân: $13 - 4 = 9$ con gà.
}
\end{ex}

\begin{ex}
Lớp 2A có 12 bạn tham gia đội ca múa, sau đó có thêm 6 bạn đăng ký tham gia. Đội ca múa có tất cả bao nhiêu bạn?
\shortans{18}
\loigiai{
Đội có: $12 + 6 = 18$ bạn.
}
\end{ex}

\begin{ex}
Trên xe buýt có 17 hành khách, đến trạm dừng có 8 hành khách xuống xe và không có ai lên. Trên xe buýt còn lại bao nhiêu hành khách?
\shortans{9}
\loigiai{
Còn lại: $17 - 8 = 9$ hành khách.
}
\end{ex}

\begin{ex}
Một giá sách có 15 quyển truyện tranh, bạn Nam mượn đi 8 quyển. Trên giá còn lại bao nhiêu quyển truyện?
\shortans{7}
\loigiai{
Còn lại: $15 - 8 = 7$ quyển.
}
\end{ex}

\begin{ex}
Hộp bánh có 11 chiếc bánh, bé ăn bớt 3 chiếc. Hộp bánh còn lại bao nhiêu chiếc?
\shortans{8}
\loigiai{
Còn lại: $11 - 3 = 8$ chiếc bánh.
}
\end{ex}

\begin{ex}
Hồ nước có 9 bông hoa sen nở, ngày hôm sau nở thêm 7 bông hoa sen nữa. Hỏi hồ nước có tất cả bao nhiêu bông hoa sen nở?
\shortans{16}
\loigiai{
Có tất cả: $9 + 7 = 16$ bông hoa.
}
\end{ex}

\begin{ex}
Bác thợ may có một mảnh vải dài 16 mét, bác đã cắt may áo hết 9 mét vải. Mảnh vải còn lại dài bao nhiêu mét?
\shortans{7}
\loigiai{
Còn lại: $16 - 9 = 7$ mét.
}
\end{ex}

\begin{ex}
Mẹ mua về 14 quả trứng, mẹ dùng 6 quả để làm bánh rán. Mẹ còn lại bao nhiêu quả trứng?
\shortans{8}
\loigiai{
Còn lại: $14 - 6 = 8$ quả trứng.
}
\end{ex}

\begin{ex}
Tổ Một trồng được 8 cây hoa, cô giáo yêu cầu trồng thêm 8 cây hoa nữa. Hỏi sau khi trồng thêm tổ Một có bao nhiêu cây hoa?
\shortans{16}
\loigiai{
Có tất cả: $8 + 8 = 16$ cây hoa.
}
\end{ex}

\begin{ex}
Có 12 bạn nhỏ đang chơi trốn tìm, có 5 bạn chạy về nhà ăn cơm. Còn lại bao nhiêu bạn đang chơi?
\shortans{7}
\loigiai{
Còn lại: $12 - 5 = 7$ bạn.
}
\end{ex}

\begin{ex}
Một rổ táo có 13 quả, bé mang biếu bà 5 quả táo. Trong rổ còn lại bao nhiêu quả táo?
\shortans{8}
\loigiai{
Trong rổ còn lại: $13 - 5 = 8$ quả táo.
}
\end{ex}

\begin{ex}
Hùng có 6 viên bi, Hùng được thưởng thêm 7 viên bi nữa. Hùng có tất cả bao nhiêu viên bi?
\shortans{13}
\loigiai{
Hùng có: $6 + 7 = 13$ viên bi.
}
\end{ex}

\begin{ex}
Vườn có 9 cây chuối, bố trồng thêm 6 cây chuối nữa. Vườn có bao nhiêu cây chuối?
\shortans{15}
\loigiai{
Có tất cả: $9 + 6 = 15$ cây chuối.
}
\end{ex}

\begin{ex}
Đàn cá vàng có 11 con, bác bán bớt 7 con. Đàn cá còn lại bao nhiêu con?
\shortans{4}
\loigiai{
Còn lại: $11 - 7 = 4$ con cá.
}
\end{ex}

\begin{ex}
Chùm bóng bay có 14 quả, bị bay mất 8 quả. Chùm bóng còn lại bao nhiêu quả?
\shortans{6}
\loigiai{
Còn lại: $14 - 8 = 6$ quả.
}
\end{ex}

\begin{ex}
Một hộp bút có 9 chiếc, cô giáo bỏ thêm 9 chiếc bút nữa vào hộp. Hộp có tất cả bao nhiêu chiếc bút?
\shortans{18}
\loigiai{
Có tất cả: $9 + 9 = 18$ chiếc bút.
}
\end{ex}

\newpage

\TaoPhieuBaiTap{Chủ đề 2 -- Phép cộng, phép trừ trong phạm vi 20}{10}{Bảng cộng (có nhớ) trong phạm vi 20}{https://loptoanthaylinh.vn/bai-tap/toan-2/bai-10}
\begin{ex}
Tính: $9 + 6$.
\shortans{15}
\loigiai{
Ta có $9 + 6 = 15$.
}
\end{ex}

\begin{ex}
Tính: $8 + 7$.
\shortans{15}
\loigiai{
Ta có $8 + 7 = 15$.
}
\end{ex}

\begin{ex}
Tính: $7 + 5$.
\shortans{12}
\loigiai{
Ta có $7 + 5 = 12$.
}
\end{ex}

\begin{ex}
Tính: $6 + 6$.
\shortans{12}
\loigiai{
Ta có $6 + 6 = 12$.
}
\end{ex}

\begin{ex}
Tính: $8 + 8$.
\shortans{16}
\loigiai{
Ta có $8 + 8 = 16$.
}
\end{ex}

\begin{ex}
Tính: $9 + 4$.
\shortans{13}
\loigiai{
Ta có $9 + 4 = 13$.
}
\end{ex}

\begin{ex}
Tính giá trị biểu thức: $9 + 3 + 4$.
\shortans{16}
\loigiai{
Ta có $9 + 3 = 12$, rồi $12 + 4 = 16$.
}
\end{ex}

\begin{ex}
Tính giá trị biểu thức: $8 + 4 + 3$.
\shortans{15}
\loigiai{
Ta có $8 + 4 = 12$, rồi $12 + 3 = 15$.
}
\end{ex}

\begin{ex}
Tính giá trị biểu thức: $7 + 7 + 2$.
\shortans{16}
\loigiai{
Ta có $7 + 7 = 14$, rồi $14 + 2 = 16$.
}
\end{ex}

\begin{ex}
Tính giá trị biểu thức: $6 + 5 + 4$.
\shortans{15}
\loigiai{
Ta có $6 + 4 = 10$, rồi $10 + 5 = 15$.
}
\end{ex}

\begin{ex}
Điền số thích hợp: $8 + \dots = 17$.
\shortans{9}
\loigiai{
Số cần điền là: $17 - 8 = 9$.
}
\end{ex}

\begin{ex}
Điền số thích hợp: $7 + \dots = 13$.
\shortans{6}
\loigiai{
Số cần điền là: $13 - 7 = 6$.
}
\end{ex}

\begin{ex}
Điền số thích hợp: $\dots + 9 = 14$.
\shortans{5}
\loigiai{
Số cần điền là: $14 - 9 = 5$.
}
\end{ex}

\begin{ex}
Điền số thích hợp: $\dots + 8 = 11$.
\shortans{3}
\loigiai{
Số cần điền là: $11 - 8 = 3$.
}
\end{ex}

\begin{ex}
So sánh: Kết quả của $9 + 5$ lớn hơn kết quả của $8 + 5$ bao nhiêu đơn vị?
\shortans{1}
\loigiai{
Ta có $9 + 5 = 14$, $8 + 5 = 13$. Chênh lệch: $14 - 13 = 1$.
}
\end{ex}

\begin{ex}
Hai số có một chữ số có tổng bằng 15, biết một số là 9. Số còn lại là bao nhiêu?
\shortans{6}
\loigiai{
Số còn lại là: $15 - 9 = 6$.
}
\end{ex}

\begin{ex}
Hai số có một chữ số có tổng bằng 17. Tích của hai số đó là bao nhiêu nếu hai số đó là 9 và 8? Tính tổng của 9 và 8.
\shortans{17}
\loigiai{
Tổng là: $9 + 8 = 17$.
}
\end{ex}

\begin{ex}
Bác nông dân có 8 con gà và 9 con vịt. Hỏi bác nông dân có tất cả bao nhiêu con gà và vịt?
\shortans{17}
\loigiai{
Bác có: $8 + 9 = 17$ con.
}
\end{ex}

\begin{ex}
Lan có 7 cái nhãn vở hoa và 6 cái nhãn vở hoạt hình. Lan có tất cả bao nhiêu cái nhãn vở?
\shortans{13}
\loigiai{
Lan có: $7 + 6 = 13$ cái nhãn vở.
}
\end{ex}

\begin{ex}
Trong hộp có 6 chiếc bút chì màu xanh và 8 chiếc bút chì màu đỏ. Trong hộp có bao nhiêu chiếc bút chì?
\shortans{14}
\loigiai{
Trong hộp có: $6 + 8 = 14$ chiếc bút chì.
}
\end{ex}

\begin{ex}
Một người nuôi 9 con bồ câu trắng và 7 con bồ câu nâu. Người đó nuôi tất cả bao nhiêu con bồ câu?
\shortans{16}
\loigiai{
Tổng số bồ câu: $9 + 7 = 16$ con.
}
\end{ex}

\begin{ex}
Tìm một số biết lấy số đó trừ đi 8 thì được 4.
\shortans{12}
\loigiai{
Số đó là: $4 + 8 = 12$.
}
\end{ex}

\begin{ex}
Tìm một số biết lấy số đó trừ đi 9 thì được 7.
\shortans{16}
\loigiai{
Số đó là: $7 + 9 = 16$.
}
\end{ex}

\begin{ex}
Tính: $9 + 9 - 10$.
\shortans{8}
\loigiai{
Ta có $9 + 9 = 18$, rồi $18 - 10 = 8$.
}
\end{ex}

\begin{ex}
Tính: $8 + 6 - 5$.
\shortans{9}
\loigiai{
Ta có $8 + 6 = 14$, rồi $14 - 5 = 9$.
}
\end{ex}

\newpage

\TaoPhieuBaiTap{Chủ đề 2 -- Phép cộng, phép trừ trong phạm vi 20}{11}{Bảng trừ (có nhớ) trong phạm vi 20}{https://loptoanthaylinh.vn/bai-tap/toan-2/bai-11}
\begin{ex}
Tính kết quả của phép tính: $11 - 2$.
\shortans{9}
\loigiai{
Ta có $11 - 2 = 9$.
}
\end{ex}

\begin{ex}
Tính kết quả của phép tính: $11 - 3$.
\shortans{8}
\loigiai{
Ta có $11 - 3 = 8$.
}
\end{ex}

\begin{ex}
Tính kết quả của phép tính: $11 - 4$.
\shortans{7}
\loigiai{
Ta có $11 - 4 = 7$.
}
\end{ex}

\begin{ex}
Tính kết quả của phép tính: $11 - 5$.
\shortans{6}
\loigiai{
Ta có $11 - 5 = 6$.
}
\end{ex}

\begin{ex}
Tính kết quả của phép tính: $11 - 6$.
\shortans{5}
\loigiai{
Ta có $11 - 6 = 5$.
}
\end{ex}

\begin{ex}
Tính kết quả của phép tính: $11 - 7$.
\shortans{4}
\loigiai{
Ta có $11 - 7 = 4$.
}
\end{ex}

\begin{ex}
Tính kết quả của phép tính: $11 - 8$.
\shortans{3}
\loigiai{
Ta có $11 - 8 = 3$.
}
\end{ex}

\begin{ex}
Tính kết quả của phép tính: $11 - 9$.
\shortans{2}
\loigiai{
Ta có $11 - 9 = 2$.
}
\end{ex}

\begin{ex}
Tính kết quả của phép tính: $12 - 3$.
\shortans{9}
\loigiai{
Ta có $12 - 3 = 9$.
}
\end{ex}

\begin{ex}
Tính kết quả của phép tính: $12 - 4$.
\shortans{8}
\loigiai{
Ta có $12 - 4 = 8$.
}
\end{ex}

\begin{ex}
Tính kết quả của phép tính: $12 - 5$.
\shortans{7}
\loigiai{
Ta có $12 - 5 = 7$.
}
\end{ex}

\begin{ex}
Tính kết quả của phép tính: $12 - 6$.
\shortans{6}
\loigiai{
Ta có $12 - 6 = 6$.
}
\end{ex}

\begin{ex}
Tính kết quả của phép tính: $12 - 7$.
\shortans{5}
\loigiai{
Ta có $12 - 7 = 5$.
}
\end{ex}

\begin{ex}
Tính kết quả của phép tính: $12 - 8$.
\shortans{4}
\loigiai{
Ta có $12 - 8 = 4$.
}
\end{ex}

\begin{ex}
Tính kết quả của phép tính: $12 - 9$.
\shortans{3}
\loigiai{
Ta có $12 - 9 = 3$.
}
\end{ex}

\begin{ex}
Tính kết quả của phép tính: $13 - 4$.
\shortans{9}
\loigiai{
Ta có $13 - 4 = 9$.
}
\end{ex}

\begin{ex}
Tính kết quả của phép tính: $13 - 5$.
\shortans{8}
\loigiai{
Ta có $13 - 5 = 8$.
}
\end{ex}

\begin{ex}
Tính kết quả của phép tính: $13 - 6$.
\shortans{7}
\loigiai{
Ta có $13 - 6 = 7$.
}
\end{ex}

\begin{ex}
Tính kết quả của phép tính: $13 - 7$.
\shortans{6}
\loigiai{
Ta có $13 - 7 = 6$.
}
\end{ex}

\begin{ex}
Tính kết quả của phép tính: $13 - 8$.
\shortans{5}
\loigiai{
Ta có $13 - 8 = 5$.
}
\end{ex}

\begin{ex}
Tính kết quả của phép tính: $13 - 9$.
\shortans{4}
\loigiai{
Ta có $13 - 9 = 4$.
}
\end{ex}

\begin{ex}
Có 12 quả táo, Mai ăn hết 5 quả. Hỏi còn lại bao nhiêu quả táo?
\shortans{7}
\loigiai{
Còn lại: $12 - 5 = 7$ quả táo.
}
\end{ex}

\begin{ex}
Một đàn gà có 11 con, 4 con chạy vào chuồng. Còn lại bao nhiêu con ngoài sân?
\shortans{7}
\loigiai{
Còn lại: $11 - 4 = 7$ con gà.
}
\end{ex}

\begin{ex}
Bình có 13 viên kẹo, Bình chia cho em 6 viên kẹo. Bình còn lại bao nhiêu viên kẹo?
\shortans{7}
\loigiai{
Bình còn lại: $13 - 6 = 7$ viên kẹo.
}
\end{ex}

\begin{ex}
Tổ Em có 11 bạn học sinh, trong đó có 6 bạn nữ. Hỏi tổ Em có bao nhiêu bạn nam?
\shortans{5}
\loigiai{
Số bạn nam: $11 - 6 = 5$ bạn.
}
\end{ex}

\newpage

\TaoPhieuBaiTap{Chủ đề 2 -- Phép cộng, phép trừ trong phạm vi 20}{12}{Bài toán nhiều hơn, ít hơn một số đơn vị}{https://loptoanthaylinh.vn/bai-tap/toan-2/bai-12}
\begin{ex}
Tính: $14 - 5$.
\shortans{9}
\loigiai{
Ta có $14 - 5 = 9$.
}
\end{ex}

\begin{ex}
Tính: $14 - 6$.
\shortans{8}
\loigiai{
Ta có $14 - 6 = 8$.
}
\end{ex}

\begin{ex}
Tính: $14 - 7$.
\shortans{7}
\loigiai{
Ta có $14 - 7 = 7$.
}
\end{ex}

\begin{ex}
Tính: $14 - 8$.
\shortans{6}
\loigiai{
Ta có $14 - 8 = 6$.
}
\end{ex}

\begin{ex}
Tính: $14 - 9$.
\shortans{5}
\loigiai{
Ta có $14 - 9 = 5$.
}
\end{ex}

\begin{ex}
Tính: $15 - 6$.
\shortans{9}
\loigiai{
Ta có $15 - 6 = 9$.
}
\end{ex}

\begin{ex}
Tính: $15 - 7$.
\shortans{8}
\loigiai{
Ta có $15 - 7 = 8$.
}
\end{ex}

\begin{ex}
Tính: $15 - 8$.
\shortans{7}
\loigiai{
Ta có $15 - 8 = 7$.
}
\end{ex}

\begin{ex}
Tính: $15 - 9$.
\shortans{6}
\loigiai{
Ta có $15 - 9 = 6$.
}
\end{ex}

\begin{ex}
Tính: $16 - 7$.
\shortans{9}
\loigiai{
Ta có $16 - 7 = 9$.
}
\end{ex}

\begin{ex}
Tính: $16 - 8$.
\shortans{8}
\loigiai{
Ta có $16 - 8 = 8$.
}
\end{ex}

\begin{ex}
Tính: $16 - 9$.
\shortans{7}
\loigiai{
Ta có $16 - 9 = 7$.
}
\end{ex}

\begin{ex}
Tính: $17 - 8$.
\shortans{9}
\loigiai{
Ta có $17 - 8 = 9$.
}
\end{ex}

\begin{ex}
Tính: $17 - 9$.
\shortans{8}
\loigiai{
Ta có $17 - 9 = 8$.
}
\end{ex}

\begin{ex}
Tính: $18 - 9$.
\shortans{9}
\loigiai{
Ta có $18 - 9 = 9$.
}
\end{ex}

\begin{ex}
Tính giá trị biểu thức: $15 - 7 + 4$.
\shortans{12}
\loigiai{
Ta có $15 - 7 = 8$, rồi $8 + 4 = 12$.
}
\end{ex}

\begin{ex}
Tính giá trị biểu thức: $16 - 8 + 5$.
\shortans{13}
\loigiai{
Ta có $16 - 8 = 8$, rồi $8 + 5 = 13$.
}
\end{ex}

\begin{ex}
Tính giá trị biểu thức: $17 - 9 + 6$.
\shortans{14}
\loigiai{
Ta có $17 - 9 = 8$, rồi $8 + 6 = 14$.
}
\end{ex}

\begin{ex}
Tính giá trị biểu thức: $14 - 8 + 7$.
\shortans{13}
\loigiai{
Ta có $14 - 8 = 6$, rồi $6 + 7 = 13$.
}
\end{ex}

\begin{ex}
Điền số thích hợp: $15 - \dots = 8$.
\shortans{7}
\loigiai{
Số cần điền là: $15 - 8 = 7$.
}
\end{ex}

\begin{ex}
Điền số thích hợp: $16 - \dots = 9$.
\shortans{7}
\loigiai{
Số cần điền là: $16 - 9 = 7$.
}
\end{ex}

\begin{ex}
Điền số thích hợp: $\dots - 8 = 7$.
\shortans{15}
\loigiai{
Số cần điền là: $7 + 8 = 15$.
}
\end{ex}

\begin{ex}
Điền số thích hợp: $\dots - 9 = 6$.
\shortans{15}
\loigiai{
Số cần điền là: $6 + 9 = 15$.
}
\end{ex}

\begin{ex}
Mẹ mua 16 quả cam, cả nhà ăn hết 9 quả. Mẹ còn lại bao nhiêu quả cam?
\shortans{7}
\loigiai{
Còn lại: $16 - 9 = 7$ quả cam.
}
\end{ex}

\begin{ex}
Bác nông dân có 17 con vịt, bác bán đi 8 con vịt. Bác còn lại bao nhiêu con vịt?
\shortans{9}
\loigiai{
Còn lại: $17 - 8 = 9$ con vịt.
}
\end{ex}

\newpage

\TaoPhieuBaiTap{Chủ đề 2 -- Phép cộng, phép trừ trong phạm vi 20}{13}{Luyện tập chung}{https://loptoanthaylinh.vn/bai-tap/toan-2/bai-13}
\begin{ex}
Lan có 9 bông hoa, Mai có nhiều hơn Lan 4 bông hoa. Hỏi Mai có bao nhiêu bông hoa?
\shortans{13}
\loigiai{
Mai có: $9 + 4 = 13$ bông hoa.
}
\end{ex}

\begin{ex}
Hùng có 8 viên bi, Dũng có nhiều hơn Hùng 7 viên bi. Hỏi Dũng có bao nhiêu viên bi?
\shortans{15}
\loigiai{
Dũng có: $8 + 7 = 15$ viên bi.
}
\end{ex}

\begin{ex}
Lớp 2A có 16 bạn nữ, số bạn nam ít hơn số bạn nữ 3 bạn. Hỏi lớp 2A có bao nhiêu bạn nam?
\shortans{13}
\loigiai{
Số bạn nam: $16 - 3 = 13$ bạn.
}
\end{ex}

\begin{ex}
Vườn nhà bác An có 14 cây bưởi, số cây cam ít hơn số cây bưởi là 6 cây. Hỏi vườn nhà bác An có bao nhiêu cây cam?
\shortans{8}
\loigiai{
Số cây cam: $14 - 6 = 8$ cây.
}
\end{ex}

\begin{ex}
Bình có 15 quyển truyện, Nam có ít hơn Bình 7 quyển truyện. Hỏi Nam có bao nhiêu quyển truyện?
\shortans{8}
\loigiai{
Nam có: $15 - 7 = 8$ quyển truyện.
}
\end{ex}

\begin{ex}
Năm nay anh 13 tuổi, em ít hơn anh 5 tuổi. Hỏi năm nay em bao nhiêu tuổi?
\shortans{8}
\loigiai{
Tuổi em là: $13 - 5 = 8$ tuổi.
}
\end{ex}

\begin{ex}
Tổ Một gấp được 9 chiếc thuyền, tổ Hai gấp được nhiều hơn tổ Một 6 chiếc thuyền. Hỏi tổ Hai gấp được bao nhiêu chiếc thuyền?
\shortans{15}
\loigiai{
Tổ Hai gấp được: $9 + 6 = 15$ chiếc.
}
\end{ex}

\begin{ex}
Một cây cau cao 11 mét, cây ổi thấp hơn cây cau 4 mét. Hỏi cây ổi cao bao nhiêu mét?
\shortans{7}
\loigiai{
Cây ổi cao: $11 - 4 = 7$ mét.
}
\end{ex}

\begin{ex}
Đoạn thẳng AB dài 8 cm, đoạn thẳng CD dài hơn đoạn thẳng AB là 6 cm. Hỏi đoạn thẳng CD dài bao nhiêu xăng-ti-mét?
\shortans{14}
\loigiai{
Đoạn CD dài: $8 + 6 = 14$ cm.
}
\end{ex}

\begin{ex}
Đoạn thẳng MN dài 15 cm, đoạn thẳng PQ ngắn hơn đoạn thẳng MN là 8 cm. Hỏi đoạn thẳng PQ dài bao nhiêu xăng-ti-mét?
\shortans{7}
\loigiai{
Đoạn PQ dài: $15 - 8 = 7$ cm.
}
\end{ex}

\begin{ex}
Thùng to đựng 18 lít nước, thùng nhỏ đựng ít hơn thùng to 9 lít nước. Hỏi thùng nhỏ đựng bao nhiêu lít nước?
\shortans{9}
\loigiai{
Thùng nhỏ đựng: $18 - 9 = 9$ lít nước.
}
\end{ex}

\begin{ex}
Bao gạo nặng 17 kg, bao ngô nhẹ hơn bao gạo 8 kg. Hỏi bao ngô nặng bao nhiêu ki-lô-gam?
\shortans{9}
\loigiai{
Bao ngô nặng: $17 - 8 = 9$ kg.
}
\end{ex}

\begin{ex}
Bác thợ mộc cưa thanh gỗ thứ nhất dài 12 dm, thanh gỗ thứ hai ngắn hơn thanh thứ nhất 5 dm. Hỏi thanh gỗ thứ hai dài bao nhiêu đề-xi-mét?
\shortans{7}
\loigiai{
Thanh thứ hai dài: $12 - 5 = 7$ dm.
}
\end{ex}

\begin{ex}
Mẹ mua 7 quả táo, mẹ mua số quả cam nhiều hơn số quả táo là 8 quả. Hỏi mẹ mua bao nhiêu quả cam?
\shortans{15}
\loigiai{
Mẹ mua số cam là: $7 + 8 = 15$ quả.
}
\end{ex}

\begin{ex}
Buổi sáng cửa hàng bán được 9 chiếc cặp, buổi chiều bán được nhiều hơn buổi sáng 5 chiếc cặp. Buổi chiều bán được bao nhiêu chiếc cặp?
\shortans{14}
\loigiai{
Buổi chiều bán được: $9 + 5 = 14$ chiếc.
}
\end{ex}

\begin{ex}
Một cửa hàng ngày thứ nhất bán được 16 bao gạo, ngày thứ hai bán được ít hơn ngày thứ nhất 7 bao gạo. Hỏi ngày thứ hai bán được bao nhiêu bao gạo?
\shortans{9}
\loigiai{
Ngày thứ hai bán: $16 - 7 = 9$ bao gạo.
}
\end{ex}

\begin{ex}
Hà có 14 cái bút chì màu, Mai có ít hơn Hà 6 cái bút chì màu. Hỏi Mai có bao nhiêu cái bút chì màu?
\shortans{8}
\loigiai{
Mai có: $14 - 6 = 8$ cái bút chì màu.
}
\end{ex}

\begin{ex}
An có 8 chiếc tem thư, Bình có nhiều hơn An 9 chiếc tem thư. Hỏi Bình có bao nhiêu chiếc tem thư?
\shortans{17}
\loigiai{
Bình có: $8 + 9 = 17$ chiếc tem thư.
}
\end{ex}

\begin{ex}
Ngăn sách trên có 13 quyển sách, ngăn dưới có ít hơn ngăn trên 5 quyển sách. Hỏi ngăn dưới có bao nhiêu quyển sách?
\shortans{8}
\loigiai{
Ngăn dưới có: $13 - 5 = 8$ quyển sách.
}
\end{ex}

\begin{ex}
Bạn Việt có 11 viên bi, bạn Nam có nhiều hơn bạn Việt 7 viên bi. Hỏi bạn Nam có bao nhiêu viên bi?
\shortans{18}
\loigiai{
Nam có: $11 + 7 = 18$ viên bi.
}
\end{ex}

\begin{ex}
Chị hái được 15 bông hoa, em hái được ít hơn chị 6 bông hoa. Hỏi em hái được bao nhiêu bông hoa?
\shortans{9}
\loigiai{
Em hái được: $15 - 6 = 9$ bông hoa.
}
\end{ex}

\begin{ex}
Đàn vịt có 16 con, đàn gà ít hơn đàn vịt 9 con. Hỏi đàn gà có bao nhiêu con?
\shortans{7}
\loigiai{
Đàn gà có: $16 - 9 = 7$ con.
}
\end{ex}

\begin{ex}
Bác nông dân nuôi 8 con cừu, số con dê nhiều hơn số con cừu là 8 con. Hỏi bác có bao nhiêu con dê?
\shortans{16}
\loigiai{
Bác có: $8 + 8 = 16$ con dê.
}
\end{ex}

\begin{ex}
Một chiếc thước kẻ dài 17 cm, chiếc bút bi ngắn hơn thước kẻ 5 cm. Chiếc bút bi dài bao nhiêu xăng-ti-mét?
\shortans{12}
\loigiai{
Chiếc bút bi dài: $17 - 5 = 12$ cm.
}
\end{ex}

\begin{ex}
Một quyển truyện có 18 chương, quyển sách thứ hai có ít hơn quyển truyện 9 chương. Quyển sách thứ hai có bao nhiêu chương?
\shortans{9}
\loigiai{
Quyển thứ hai có: $18 - 9 = 9$ chương.
}
\end{ex}

\newpage

\TaoPhieuBaiTap{Chủ đề 2 -- Phép cộng, phép trừ trong phạm vi 20}{14}{Đề kiểm tra chủ đề}{https://loptoanthaylinh.vn/bai-tap/toan-2/bai-14}
\begin{ex}
Tính: $9 + 7$.
\shortans{16}
\loigiai{
Ta có $9 + 7 = 16$.
}
\end{ex}

\begin{ex}
Tính: $16 - 9$.
\shortans{7}
\loigiai{
Ta có $16 - 9 = 7$.
}
\end{ex}

\begin{ex}
Tính: $8 + 6$.
\shortans{14}
\loigiai{
Ta có $8 + 6 = 14$.
}
\end{ex}

\begin{ex}
Tính: $14 - 8$.
\shortans{6}
\loigiai{
Ta có $14 - 8 = 6$.
}
\end{ex}

\begin{ex}
Tính: $7 + 5$.
\shortans{12}
\loigiai{
Ta có $7 + 5 = 12$.
}
\end{ex}

\begin{ex}
Tính: $12 - 7$.
\shortans{5}
\loigiai{
Ta có $12 - 7 = 5$.
}
\end{ex}

\begin{ex}
Tính giá trị biểu thức: $9 + 5 - 6$.
\shortans{8}
\loigiai{
Ta có $9 + 5 = 14$, rồi $14 - 6 = 8$.
}
\end{ex}

\begin{ex}
Tính giá trị biểu thức: $13 - 7 + 8$.
\shortans{14}
\loigiai{
Ta có $13 - 7 = 6$, rồi $6 + 8 = 14$.
}
\end{ex}

\begin{ex}
Tính giá trị biểu thức: $15 - 8 + 9$.
\shortans{16}
\loigiai{
Ta có $15 - 8 = 7$, rồi $7 + 9 = 16$.
}
\end{ex}

\begin{ex}
Tính giá trị biểu thức: $8 + 8 - 9$.
\shortans{7}
\loigiai{
Ta có $8 + 8 = 16$, rồi $16 - 9 = 7$.
}
\end{ex}

\begin{ex}
Điền số thích hợp: $7 + 8 = 9 + \dots$
\shortans{6}
\loigiai{
Vế trái: $7 + 8 = 15$. Số cần điền: $15 - 9 = 6$.
}
\end{ex}

\begin{ex}
Điền số thích hợp: $16 - 7 = 14 - \dots$
\shortans{5}
\loigiai{
Vế trái: $16 - 7 = 9$. Số cần điền: $14 - 9 = 5$.
}
\end{ex}

\begin{ex}
Điền số thích hợp: $13 - 6 = \dots - 8$.
\shortans{15}
\loigiai{
Vế trái: $13 - 6 = 7$. Số cần điền: $7 + 8 = 15$.
}
\end{ex}

\begin{ex}
Điền số thích hợp: $8 + 5 = \dots + 4$.
\shortans{9}
\loigiai{
Vế trái: $8 + 5 = 13$. Số cần điền: $13 - 4 = 9$.
}
\end{ex}

\begin{ex}
Một đàn gà có 15 con, trong đó có 7 con gà trống. Hỏi đàn gà có bao nhiêu con gà mái?
\shortans{8}
\loigiai{
Số gà mái: $15 - 7 = 8$ con.
}
\end{ex}

\begin{ex}
Lan hái được 8 bông hoa hồng, Mai hái được nhiều hơn Lan 5 bông hoa hồng. Hỏi Mai hái được bao nhiêu bông hoa hồng?
\shortans{13}
\loigiai{
Mai hái được: $8 + 5 = 13$ bông hoa.
}
\end{ex}

\begin{ex}
Hùng có 17 viên bi, Hùng cho Dũng 9 viên bi. Hùng còn lại bao nhiêu viên bi?
\shortans{8}
\loigiai{
Hùng còn lại: $17 - 9 = 8$ viên bi.
}
\end{ex}

\begin{ex}
Bình có 14 cái nhãn vở, Minh có ít hơn Bình 6 cái nhãn vở. Hỏi Minh có bao nhiêu cái nhãn vở?
\shortans{8}
\loigiai{
Minh có: $14 - 6 = 8$ cái nhãn vở.
}
\end{ex}

\begin{ex}
Cả hai bạn Lan và Mai có 16 cái kẹo, biết Lan có 9 cái kẹo. Hỏi Mai có bao nhiêu cái kẹo?
\shortans{7}
\loigiai{
Mai có: $16 - 9 = 7$ cái kẹo.
}
\end{ex}

\begin{ex}
Hiệu của 18 và số lớn nhất có một chữ số là bao nhiêu?
\shortans{9}
\loigiai{
Số lớn nhất có 1 chữ số là 9. Hiệu: $18 - 9 = 9$.
}
\end{ex}

\begin{ex}
Tổng của 8 và số liền sau của 6 là bao nhiêu?
\shortans{15}
\loigiai{
Số liền sau của 6 là 7. Tổng: $8 + 7 = 15$.
}
\end{ex}

\begin{ex}
Hiệu của 14 và số liền trước của 8 là bao nhiêu?
\shortans{7}
\loigiai{
Số liền trước của 8 là 7. Hiệu: $14 - 7 = 7$.
}
\end{ex}

\begin{ex}
Bác thợ rèn thanh sắt dài 15 cm, bác cắt bớt một đoạn dài 8 cm. Đoạn sắt còn lại dài bao nhiêu xăng-ti-mét?
\shortans{7}
\loigiai{
Còn lại: $15 - 8 = 7$ cm.
}
\end{ex}

\begin{ex}
Một cuộn len dài 16 dm, dùng đan khăn hết 7 dm. Cuộn len còn lại dài bao nhiêu đề-xi-mét?
\shortans{9}
\loigiai{
Còn lại: $16 - 7 = 9$ dm.
}
\end{ex}

\begin{ex}
Điền số thích hợp vào chỗ chấm: $17 - 9 = 2 + \dots$
\shortans{6}
\loigiai{
Vế trái: $17 - 9 = 8$. Số cần điền: $8 - 2 = 6$.
}
\end{ex}

\newpage


% =========================================================================
% Làm quen với khối lượng, dung tích
% =========================================================================
\section{Làm quen với khối lượng, dung tích}

\TaoPhieuBaiTap{Chủ đề 3 -- Làm quen với khối lượng, dung tích}{15}{Ki-lô-gam}{https://loptoanthaylinh.vn/bai-tap/toan-2/bai-15}
\begin{ex}
Ki-lô-gam được viết tắt bằng mấy chữ cái trong bảng chữ cái tiếng Anh? (viết tắt là kg).
\shortans{2}
\loigiai{
Viết tắt là kg, gồm 2 chữ cái k và g.
}
\end{ex}

\begin{ex}
Tính: $5\text{ kg} + 3\text{ kg}$. Kết quả là bao nhiêu ki-lô-gam?
\shortans{8}
\loigiai{
Ta có: $5 + 3 = 8$ kg.
}
\end{ex}

\begin{ex}
Tính: $14\text{ kg} + 5\text{ kg}$. Kết quả là bao nhiêu ki-lô-gam?
\shortans{19}
\loigiai{
Ta có: $14 + 5 = 19$ kg.
}
\end{ex}

\begin{ex}
Tính: $26\text{ kg} + 12\text{ kg}$. Kết quả là bao nhiêu ki-lô-gam?
\shortans{38}
\loigiai{
Ta có: $26 + 12 = 38$ kg.
}
\end{ex}

\begin{ex}
Tính: $45\text{ kg} - 20\text{ kg}$. Kết quả là bao nhiêu ki-lô-gam?
\shortans{25}
\loigiai{
Ta có: $45 - 20 = 25$ kg.
}
\end{ex}

\begin{ex}
Tính: $18\text{ kg} - 9\text{ kg}$. Kết quả là bao nhiêu ki-lô-gam?
\shortans{9}
\loigiai{
Ta có: $18 - 9 = 9$ kg.
}
\end{ex}

\begin{ex}
Tính: $35\text{ kg} - 15\text{ kg}$. Kết quả là bao nhiêu ki-lô-gam?
\shortans{20}
\loigiai{
Ta có: $35 - 15 = 20$ kg.
}
\end{ex}

\begin{ex}
Tính giá trị biểu thức: $12\text{ kg} + 8\text{ kg} - 5\text{ kg}$.
\shortans{15}
\loigiai{
Ta có $12 + 8 = 20$, rồi $20 - 5 = 15$ kg.
}
\end{ex}

\begin{ex}
Tính giá trị biểu thức: $25\text{ kg} - 10\text{ kg} + 7\text{ kg}$.
\shortans{22}
\loigiai{
Ta có $25 - 10 = 15$, rồi $15 + 7 = 22$ kg.
}
\end{ex}

\begin{ex}
Trên đĩa cân thăng bằng: đĩa bên trái có quả dưa hấu, đĩa bên phải có 2 quả cân loại 2 kg và 1 quả cân loại 1 kg. Quả dưa hấu nặng bao nhiêu ki-lô-gam?
\shortans{5}
\loigiai{
Quả dưa nặng: $2 + 2 + 1 = 5$ kg.
}
\end{ex}

\begin{ex}
Đĩa bên trái có một túi bột mì và quả cân 1 kg, đĩa bên phải có quả cân 5 kg. Hai đĩa cân thăng bằng. Túi bột mì nặng bao nhiêu ki-lô-gam?
\shortans{4}
\loigiai{
Túi bột mì nặng: $5 - 1 = 4$ kg.
}
\end{ex}

\begin{ex}
Bao gạo to nặng 25 kg, bao gạo nhỏ nặng 15 kg. Cả hai bao gạo nặng bao nhiêu ki-lô-gam?
\shortans{40}
\loigiai{
Cả hai bao nặng: $25 + 15 = 40$ kg.
}
\end{ex}

\begin{ex}
Bao gạo to nặng hơn bao gạo nhỏ bao nhiêu ki-lô-gam?
\shortans{10}
\loigiai{
Nặng hơn: $25 - 15 = 10$ kg.
}
\end{ex}

\begin{ex}
Con ngỗng nặng 6 kg, con vịt nặng 2 kg. Con ngỗng nặng hơn con vịt bao nhiêu ki-lô-gam?
\shortans{4}
\loigiai{
Nặng hơn: $6 - 2 = 4$ kg.
}
\end{ex}

\begin{ex}
Con ngỗng và con vịt nặng tất cả bao nhiêu ki-lô-gam?
\shortans{8}
\loigiai{
Cả hai con nặng: $6 + 2 = 8$ kg.
}
\end{ex}

\begin{ex}
Con chó nặng 12 kg, con mèo nặng 4 kg. Con chó nặng hơn con mèo bao nhiêu ki-lô-gam?
\shortans{8}
\loigiai{
Con chó nặng hơn: $12 - 4 = 8$ kg.
}
\end{ex}

\begin{ex}
Mẹ mua về 18 kg gạo, gia đình đã ăn hết 9 kg gạo. Còn lại bao nhiêu ki-lô-gam gạo?
\shortans{9}
\loigiai{
Còn lại: $18 - 9 = 9$ kg.
}
\end{ex}

\begin{ex}
Một thùng cam nặng 32 kg, người ta lấy ra 10 kg cam. Thùng cam còn lại nặng bao nhiêu ki-lô-gam?
\shortans{22}
\loigiai{
Còn lại: $32 - 10 = 22$ kg.
}
\end{ex}

\begin{ex}
Điền số thích hợp: $17\text{ kg} + \dots\text{ kg} = 25\text{ kg}$.
\shortans{8}
\loigiai{
Số cần điền là: $25 - 17 = 8$.
}
\end{ex}

\begin{ex}
Điền số thích hợp: $48\text{ kg} - \dots\text{ kg} = 23\text{ kg}$.
\shortans{25}
\loigiai{
Số cần điền là: $48 - 23 = 25$.
}
\end{ex}

\begin{ex}
Điền số thích hợp: $\dots\text{ kg} + 15\text{ kg} = 39\text{ kg}$.
\shortans{24}
\loigiai{
Số cần điền là: $39 - 15 = 24$.
}
\end{ex}

\begin{ex}
Bé Nam cân nặng 23 kg, bé Mai cân nặng 21 kg. Cả hai bạn cân nặng tất cả bao nhiêu ki-lô-gam?
\shortans{44}
\loigiai{
Cả hai bạn nặng: $23 + 21 = 44$ kg.
}
\end{ex}

\begin{ex}
Một con cừu mẹ nặng 35 kg, cừu con nặng 12 kg. Cừu mẹ nặng hơn cừu con bao nhiêu ki-lô-gam?
\shortans{23}
\loigiai{
Nặng hơn: $35 - 12 = 23$ kg.
}
\end{ex}

\begin{ex}
Bác nông dân thu hoạch được 50 kg bí đỏ, bác đã bán đi 30 kg. Bác còn lại bao nhiêu ki-lô-gam bí đỏ?
\shortans{20}
\loigiai{
Còn lại: $50 - 30 = 20$ kg.
}
\end{ex}

\begin{ex}
Có 3 gói kẹo: gói thứ nhất nặng 2 kg, gói thứ hai nặng 3 kg, gói thứ ba nặng 4 kg. Ba gói nặng tất cả bao nhiêu ki-lô-gam?
\shortans{9}
\loigiai{
Cả ba gói nặng: $2 + 3 + 4 = 9$ kg.
}
\end{ex}

\newpage

\TaoPhieuBaiTap{Chủ đề 3 -- Làm quen với khối lượng, dung tích}{16}{Lít}{https://loptoanthaylinh.vn/bai-tap/toan-2/bai-16}
\begin{ex}
Lít được viết tắt bằng chữ cái nào trong bảng chữ cái? Viết tắt có mấy chữ cái?
\shortans{1}
\loigiai{
Viết tắt là chữ l (chỉ gồm 1 chữ cái).
}
\end{ex}

\begin{ex}
Tính: $4\text{ }l + 5\text{ }l$. Kết quả là bao nhiêu lít?
\shortans{9}
\loigiai{
Ta có: $4 + 5 = 9$ lít.
}
\end{ex}

\begin{ex}
Tính: $12\text{ }l + 7\text{ }l$. Kết quả là bao nhiêu lít?
\shortans{19}
\loigiai{
Ta có: $12 + 7 = 19$ lít.
}
\end{ex}

\begin{ex}
Tính: $35\text{ }l + 20\text{ }l$. Kết quả là bao nhiêu lít?
\shortans{55}
\loigiai{
Ta có: $35 + 20 = 55$ lít.
}
\end{ex}

\begin{ex}
Tính: $18\text{ }l - 8\text{ }l$. Kết quả là bao nhiêu lít?
\shortans{10}
\loigiai{
Ta có: $18 - 8 = 10$ lít.
}
\end{ex}

\begin{ex}
Tính: $47\text{ }l - 23\text{ }l$. Kết quả là bao nhiêu lít?
\shortans{24}
\loigiai{
Ta có: $47 - 23 = 24$ lít.
}
\end{ex}

\begin{ex}
Tính: $50\text{ }l - 30\text{ }l$. Kết quả là bao nhiêu lít?
\shortans{20}
\loigiai{
Ta có: $50 - 30 = 20$ lít.
}
\end{ex}

\begin{ex}
Tính giá trị biểu thức: $15\text{ }l + 5\text{ }l - 8\text{ }l$.
\shortans{12}
\loigiai{
Ta có $15 + 5 = 20$, rồi $20 - 8 = 12$ lít.
}
\end{ex}

\begin{ex}
Tính giá trị biểu thức: $30\text{ }l - 10\text{ }l + 15\text{ }l$.
\shortans{35}
\loigiai{
Ta có $30 - 10 = 20$, rồi $20 + 15 = 35$ lít.
}
\end{ex}

\begin{ex}
Thùng thứ nhất có 16 lít nước ngọt, thùng thứ hai có 12 lít nước ngọt. Cả hai thùng có tất cả bao nhiêu lít nước ngọt?
\shortans{28}
\loigiai{
Cả hai thùng có: $16 + 12 = 28$ lít.
}
\end{ex}

\begin{ex}
Thùng thứ nhất nhiều hơn thùng thứ hai bao nhiêu lít nước ngọt?
\shortans{4}
\loigiai{
Nhiều hơn: $16 - 12 = 4$ lít.
}
\end{ex}

\begin{ex}
Can dầu to có 25 lít dầu, rót sang can nhỏ 10 lít dầu. Can dầu to còn lại bao nhiêu lít dầu?
\shortans{15}
\loigiai{
Còn lại: $25 - 10 = 15$ lít.
}
\end{ex}

\begin{ex}
Một bể nước có 80 lít nước, sau khi dùng để tưới cây thì bể còn 50 lít nước. Đã dùng bao nhiêu lít nước để tưới cây?
\shortans{30}
\loigiai{
Đã dùng: $80 - 50 = 30$ lít nước.
}
\end{ex}

\begin{ex}
Trong bình có 2 lít nước mơ, mẹ pha thêm 3 lít nước lọc vào bình. Trong bình có tất cả bao nhiêu lít nước?
\shortans{5}
\loigiai{
Có tất cả: $2 + 3 = 5$ lít.
}
\end{ex}

\begin{ex}
Bác bảo vệ mua 2 can xăng, mỗi can 5 lít. Hỏi bác mua tất cả bao nhiêu lít xăng?
\shortans{10}
\loigiai{
Bác mua: $5 + 5 = 10$ lít xăng.
}
\end{ex}

\begin{ex}
Một xô nước đựng được 10 lít nước, một chậu nước đựng được 15 lít nước. Cả xô và chậu đựng được bao nhiêu lít nước?
\shortans{25}
\loigiai{
Đựng được: $10 + 15 = 25$ lít.
}
\end{ex}

\begin{ex}
Điền số thích hợp: $14\text{ }l + \dots\text{ }l = 20\text{ }l$.
\shortans{6}
\loigiai{
Số cần điền là: $20 - 14 = 6$.
}
\end{ex}

\begin{ex}
Điền số thích hợp: $38\text{ }l - \dots\text{ }l = 25\text{ }l$.
\shortans{13}
\loigiai{
Số cần điền là: $38 - 25 = 13$.
}
\end{ex}

\begin{ex}
Điền số thích hợp: $\dots\text{ }l + 22\text{ }l = 45\text{ }l$.
\shortans{23}
\loigiai{
Số cần điền là: $45 - 22 = 23$.
}
\end{ex}

\begin{ex}
Một máy giặt lần giặt thứ nhất dùng hết 35 lít nước, lần giặt thứ hai dùng hết 30 lít nước. Cả hai lần giặt dùng hết bao nhiêu lít nước?
\shortans{65}
\loigiai{
Dùng hết: $35 + 30 = 65$ lít nước.
}
\end{ex}

\begin{ex}
Một thùng sữa chua có 24 lít sữa, đã bán được 14 lít. Thùng còn lại bao nhiêu lít sữa chua?
\shortans{10}
\loigiai{
Còn lại: $24 - 14 = 10$ lít.
}
\end{ex}

\begin{ex}
Có 4 chai nước khoáng, mỗi chai chứa 2 lít nước. Cả 4 chai chứa tất cả bao nhiêu lít nước?
\shortans{8}
\loigiai{
Chứa: $2 + 2 + 2 + 2 = 8$ lít.
}
\end{ex}

\begin{ex}
Một téc nước mini có 95 lít nước, đã xả bớt 42 lít để cọ rửa sân. Téc nước còn lại bao nhiêu lít?
\shortans{53}
\loigiai{
Còn lại: $95 - 42 = 53$ lít.
}
\end{ex}

\begin{ex}
Bình nước nhà bạn An có 18 lít, hôm nay dùng hết 6 lít. Bình còn lại bao nhiêu lít nước?
\shortans{12}
\loigiai{
Còn lại: $18 - 6 = 12$ lít.
}
\end{ex}

\begin{ex}
Tính: $10\text{ }l + 20\text{ }l + 30\text{ }l - 15\text{ }l$.
\shortans{45}
\loigiai{
Ta có $10 + 20 + 30 = 60$, rồi $60 - 15 = 45$ lít.
}
\end{ex}

\newpage

\TaoPhieuBaiTap{Chủ đề 3 -- Làm quen với khối lượng, dung tích}{17}{Thực hành cân đo}{https://loptoanthaylinh.vn/bai-tap/toan-2/bai-17}
\begin{ex}
Để cân quả mít, người ta đặt lên đĩa cân bên kia các quả cân: 5 kg, 2 kg và 1 kg thì cân thăng bằng. Quả mít nặng bao nhiêu ki-lô-gam?
\shortans{8}
\loigiai{
Quả mít nặng: $5 + 2 + 1 = 8$ kg.
}
\end{ex}

\begin{ex}
Cân một con gà cùng quả cân 2 kg ở đĩa trái thăng bằng với quả cân 5 kg ở đĩa phải. Con gà nặng bao nhiêu ki-lô-gam?
\shortans{3}
\loigiai{
Con gà nặng: $5 - 2 = 3$ kg.
}
\end{ex}

\begin{ex}
Mẹ đong nước vào 3 chiếc chai loại 1 lít và 2 chiếc can loại 5 lít. Mẹ đã đong được tất cả bao nhiêu lít nước?
\shortans{13}
\loigiai{
Mẹ đong được: $(1 + 1 + 1) + (5 + 5) = 13$ lít.
}
\end{ex}

\begin{ex}
Một chiếc ca đựng được 1 lít nước, một chiếc bình đựng được 5 ca nước đó. Chiếc bình đựng được bao nhiêu lít nước?
\shortans{5}
\loigiai{
Chiếc bình đựng: $1 \times 5 = 5$ lít.
}
\end{ex}

\begin{ex}
Muốn đong được 8 lít nước bằng ca 1 lít thì cần đong bao nhiêu lần ca đầy nước?
\shortans{8}
\loigiai{
Cần đong đúng 8 lần.
}
\end{ex}

\begin{ex}
Một bao xi măng nặng 50 kg, bác thợ xây đã dùng 30 kg để trát tường. Bao xi măng còn lại bao nhiêu ki-lô-gam?
\shortans{20}
\loigiai{
Còn lại: $50 - 30 = 20$ kg.
}
\end{ex}

\begin{ex}
Một chiếc ba lô đựng sách vở nặng 4 kg, khi bỏ bớt 1 kg sách ra thì ba lô nặng bao nhiêu ki-lô-gam?
\shortans{3}
\loigiai{
Nặng: $4 - 1 = 3$ kg.
}
\end{ex}

\begin{ex}
Hai bao ngô cân nặng lần lượt là 23 kg và 25 kg. Cả hai bao ngô nặng tất cả bao nhiêu ki-lô-gam?
\shortans{48}
\loigiai{
Hai bao nặng: $23 + 25 = 48$ kg.
}
\end{ex}

\begin{ex}
Một can 20 lít chứa đầy nước mắm, người ta rót vào 4 chai, mỗi chai 2 lít. Can nước mắm còn lại bao nhiêu lít?
\shortans{12}
\loigiai{
Rót ra: $2 + 2 + 2 + 2 = 8$ lít. Còn lại: $20 - 8 = 12$ lít.
}
\end{ex}

\begin{ex}
Bác nông dân có 60 lít sữa bò, bác đem bán ở chợ 40 lít. Bác còn lại bao nhiêu lít sữa bò?
\shortans{20}
\loigiai{
Bác còn lại: $60 - 40 = 20$ lít.
}
\end{ex}

\begin{ex}
Cân nặng của bạn Bình là 26 kg, bạn Nam nặng hơn bạn Bình 3 kg. Bạn Nam nặng bao nhiêu ki-lô-gam?
\shortans{29}
\loigiai{
Nam nặng: $26 + 3 = 29$ kg.
}
\end{ex}

\begin{ex}
Bạn Lan nặng 24 kg, bạn Mai nhẹ hơn bạn Lan 2 kg. Bạn Mai nặng bao nhiêu ki-lô-gam?
\shortans{22}
\loigiai{
Mai nặng: $24 - 2 = 22$ kg.
}
\end{ex}

\begin{ex}
Một chiếc ấm điện đun được 2 lít nước. Sau khi đun 3 ấm đầy nước như thế, ta có tất cả bao nhiêu lít nước sôi?
\shortans{6}
\loigiai{
Có tất cả: $2 + 2 + 2 = 6$ lít.
}
\end{ex}

\begin{ex}
Một bình lọc nước có dung tích 15 lít, đã có sẵn 9 lít nước. Cần đổ thêm bao nhiêu lít nước nữa để đầy bình?
\shortans{6}
\loigiai{
Cần đổ thêm: $15 - 9 = 6$ lít.
}
\end{ex}

\begin{ex}
Một giỏ hoa quả nặng 7 kg, chiếc giỏ không nặng 1 kg. Hoa quả trong giỏ nặng bao nhiêu ki-lô-gam?
\shortans{6}
\loigiai{
Hoa quả nặng: $7 - 1 = 6$ kg.
}
\end{ex}

\begin{ex}
Mẹ mua 5 kg khoai tây và 4 kg cà rốt. Mẹ mua tất cả bao nhiêu ki-lô-gam củ?
\shortans{9}
\loigiai{
Mẹ mua: $5 + 4 = 9$ kg.
}
\end{ex}

\begin{ex}
Một can xăng có 10 lít, đổ vào bình xe máy 3 lít. Trong can còn lại bao nhiêu lít xăng?
\shortans{7}
\loigiai{
Còn lại: $10 - 3 = 7$ lít.
}
\end{ex}

\begin{ex}
Một con lợn con nặng 15 kg, sau một tháng tăng thêm 8 kg. Con lợn con lúc này nặng bao nhiêu ki-lô-gam?
\shortans{23}
\loigiai{
Nặng: $15 + 8 = 23$ kg.
}
\end{ex}

\begin{ex}
Một bồn hoa cần tưới 25 lít nước, bạn An đã xách được 15 lít nước tưới. An còn phải xách thêm bao nhiêu lít nước nữa?
\shortans{10}
\loigiai{
An cần thêm: $25 - 15 = 10$ lít.
}
\end{ex}

\begin{ex}
Trên cân đĩa: đĩa trái có quả bưởi và 1 quả cân 2 kg, đĩa phải có 2 quả cân loại 2 kg. Quả bưởi nặng bao nhiêu ki-lô-gam?
\shortans{2}
\loigiai{
Đĩa phải nặng 4 kg. Quả bưởi nặng: $4 - 2 = 2$ kg.
}
\end{ex}

\begin{ex}
Một gói bánh nặng 500 gam (tức nửa kg). Hai gói bánh như thế nặng đúng bao nhiêu ki-lô-gam?
\shortans{1}
\loigiai{
Hai gói nặng: $500\text{ g} + 500\text{ g} = 1000\text{ g} = 1$ kg.
}
\end{ex}

\begin{ex}
Một bể cá chứa 45 lít nước, người ta xả bớt 15 lít nước bẩn và thêm vào 15 lít nước sạch. Trong bể cá có bao nhiêu lít nước?
\shortans{45}
\loigiai{
Số lít nước trong bể không đổi: 45 lít.
}
\end{ex}

\begin{ex}
Con bê cân nặng 34 kg, con dê cân nặng 22 kg. Con bê nặng hơn con dê bao nhiêu ki-lô-gam?
\shortans{12}
\loigiai{
Con bê nặng hơn: $34 - 22 = 12$ kg.
}
\end{ex}

\begin{ex}
Có 5 can nước loại 2 lít. Tổng số lít nước là bao nhiêu?
\shortans{10}
\loigiai{
Tổng số lít: $2 \times 5 = 10$ lít.
}
\end{ex}

\begin{ex}
Một xô nước nặng 8 kg tính cả vỏ xô. Biết vỏ xô nặng 1 kg, hỏi nước trong xô nặng bao nhiêu ki-lô-gam?
\shortans{7}
\loigiai{
Nước trong xô nặng: $8 - 1 = 7$ kg.
}
\end{ex}

\newpage

\TaoPhieuBaiTap{Chủ đề 3 -- Làm quen với khối lượng, dung tích}{18}{Luyện tập chung}{https://loptoanthaylinh.vn/bai-tap/toan-2/bai-18}
\begin{ex}
Tính: $25\text{ kg} + 14\text{ kg}$.
\shortans{39}
\loigiai{
Ta có: $25 + 14 = 39$ kg.
}
\end{ex}

\begin{ex}
Tính: $48\text{ kg} - 23\text{ kg}$.
\shortans{25}
\loigiai{
Ta có: $48 - 23 = 25$ kg.
}
\end{ex}

\begin{ex}
Tính: $32\text{ }l + 16\text{ }l$.
\shortans{48}
\loigiai{
Ta có: $32 + 16 = 48$ lít.
}
\end{ex}

\begin{ex}
Tính: $57\text{ }l - 35\text{ }l$.
\shortans{22}
\loigiai{
Ta có: $57 - 35 = 22$ lít.
}
\end{ex}

\begin{ex}
Tính giá trị biểu thức: $18\text{ kg} + 12\text{ kg} - 10\text{ kg}$.
\shortans{20}
\loigiai{
Ta có $18 + 12 = 30$, rồi $30 - 10 = 20$ kg.
}
\end{ex}

\begin{ex}
Tính giá trị biểu thức: $45\text{ }l - 20\text{ }l + 15\text{ }l$.
\shortans{40}
\loigiai{
Ta có $45 - 20 = 25$, rồi $25 + 15 = 40$ lít.
}
\end{ex}

\begin{ex}
Điền số thích hợp: $28\text{ kg} + \dots\text{ kg} = 49\text{ kg}$.
\shortans{21}
\loigiai{
Số cần điền là: $49 - 28 = 21$.
}
\end{ex}

\begin{ex}
Điền số thích hợp: $65\text{ }l - \dots\text{ }l = 30\text{ }l$.
\shortans{35}
\loigiai{
Số cần điền là: $65 - 30 = 35$.
}
\end{ex}

\begin{ex}
Điền số thích hợp: $\dots\text{ kg} - 14\text{ kg} = 35\text{ kg}$.
\shortans{49}
\loigiai{
Số cần điền là: $35 + 14 = 49$.
}
\end{ex}

\begin{ex}
Một con chó nặng 14 kg, con mèo nhẹ hơn con chó 10 kg. Con mèo nặng bao nhiêu ki-lô-gam?
\shortans{4}
\loigiai{
Con mèo nặng: $14 - 10 = 4$ kg.
}
\end{ex}

\begin{ex}
Cả con chó và con mèo nặng tất cả bao nhiêu ki-lô-gam?
\shortans{18}
\loigiai{
Cả hai con nặng: $14 + 4 = 18$ kg.
}
\end{ex}

\begin{ex}
Thùng thứ nhất có 38 lít dầu, thùng thứ hai có ít hơn thùng thứ nhất 15 lít dầu. Thùng thứ hai có bao nhiêu lít dầu?
\shortans{23}
\loigiai{
Thùng thứ hai có: $38 - 15 = 23$ lít.
}
\end{ex}

\begin{ex}
Cả hai thùng có tất cả bao nhiêu lít dầu?
\shortans{61}
\loigiai{
Cả hai thùng có: $38 + 23 = 61$ lít.
}
\end{ex}

\begin{ex}
Bao thóc thứ nhất nặng 45 kg, bao thóc thứ hai nặng nhiều hơn bao thứ nhất 12 kg. Bao thóc thứ hai nặng bao nhiêu ki-lô-gam?
\shortans{57}
\loigiai{
Bao thứ hai nặng: $45 + 12 = 57$ kg.
}
\end{ex}

\begin{ex}
Một can chứa đầy 20 lít nước. Sau khi dùng 8 lít, người ta lại đổ thêm vào can 5 lít nữa. Can nước lúc này có bao nhiêu lít?
\shortans{17}
\loigiai{
Còn lại: $20 - 8 + 5 = 17$ lít.
}
\end{ex}

\begin{ex}
Một quả dưa hấu nặng 4 kg, một nải chuối nặng 3 kg. Cả hai thứ nặng bao nhiêu ki-lô-gam?
\shortans{7}
\loigiai{
Cả hai thứ nặng: $4 + 3 = 7$ kg.
}
\end{ex}

\begin{ex}
Hai thùng nước mắm có tổng cộng 55 lít, thùng thứ nhất có 30 lít. Thùng thứ hai có bao nhiêu lít nước mắm?
\shortans{25}
\loigiai{
Thùng thứ hai có: $55 - 30 = 25$ lít.
}
\end{ex}

\begin{ex}
Hiệu cân nặng giữa bao gạo 50 kg và bao ngô 35 kg là bao nhiêu ki-lô-gam?
\shortans{15}
\loigiai{
Hiệu: $50 - 35 = 15$ kg.
}
\end{ex}

\begin{ex}
Bác nông dân có 70 lít mật ong, bác đã bán đi 40 lít mật ong. Bác còn lại bao nhiêu lít mật ong?
\shortans{30}
\loigiai{
Còn lại: $70 - 40 = 30$ lít.
}
\end{ex}

\begin{ex}
Một cửa hàng ngày thứ nhất bán được 26 kg đường, ngày thứ hai bán được nhiều hơn ngày thứ nhất 13 kg đường. Ngày thứ hai bán được bao nhiêu ki-lô-gam đường?
\shortans{39}
\loigiai{
Ngày thứ hai bán: $26 + 13 = 39$ kg.
}
\end{ex}

\begin{ex}
Một bình nước có 15 lít nước, người ta rót ra 3 ca, mỗi ca 2 lít. Trong bình còn lại bao nhiêu lít nước?
\shortans{9}
\loigiai{
Rót ra: $2 + 2 + 2 = 6$ lít. Còn lại: $15 - 6 = 9$ lít.
}
\end{ex}

\begin{ex}
Bé Nam nặng 25 kg, khi đeo cặp sách 3 kg lên lưng thì cả người và cặp nặng bao nhiêu ki-lô-gam?
\shortans{28}
\loigiai{
Nặng: $25 + 3 = 28$ kg.
}
\end{ex}

\begin{ex}
Trên cân đĩa, bên trái có 1 bao gạo và 1 quả cân 2 kg, bên phải có 2 quả cân loại 5 kg thì cân thăng bằng. Bao gạo nặng bao nhiêu ki-lô-gam?
\shortans{8}
\loigiai{
Đĩa phải: $5 + 5 = 10$ kg. Bao gạo nặng: $10 - 2 = 8$ kg.
}
\end{ex}

\begin{ex}
Một chiếc thùng rỗng nặng 2 kg, khi đổ vào thùng 15 lít nước (mỗi lít nước nặng 1 kg) thì cả thùng và nước nặng bao nhiêu ki-lô-gam?
\shortans{17}
\loigiai{
Cả thùng nặng: $15 + 2 = 17$ kg.
}
\end{ex}

\begin{ex}
Tính: $10\text{ kg} + 20\text{ kg} + 30\text{ kg} - 15\text{ kg}$.
\shortans{45}
\loigiai{
Ta có $10 + 20 + 30 = 60$, rồi $60 - 15 = 45$ kg.
}
\end{ex}

\newpage


% =========================================================================
% Phép cộng, phép trừ (có nhớ) trong phạm vi 100
% =========================================================================
\section{Phép cộng, phép trừ (có nhớ) trong phạm vi 100}

\TaoPhieuBaiTap{Chủ đề 4 -- Phép cộng, phép trừ (có nhớ) trong phạm vi 100}{19}{Phép cộng (có nhớ) số có hai chữ số với số có một chữ số}{https://loptoanthaylinh.vn/bai-tap/toan-2/bai-19}
\begin{ex}
Tính kết quả của phép tính: $27 + 5$.
\shortans{32}
\loigiai{
Ta có $7 + 5 = 12$, viết 2 nhớ 1; $2 + 1 = 3$. Kết quả là 32.
}
\end{ex}

\begin{ex}
Tính kết quả của phép tính: $38 + 4$.
\shortans{42}
\loigiai{
Ta có $38 + 4 = 42$.
}
\end{ex}

\begin{ex}
Tính kết quả của phép tính: $49 + 3$.
\shortans{52}
\loigiai{
Ta có $49 + 3 = 52$.
}
\end{ex}

\begin{ex}
Tính kết quả của phép tính: $56 + 6$.
\shortans{62}
\loigiai{
Ta có $56 + 6 = 62$.
}
\end{ex}

\begin{ex}
Tính kết quả của phép tính: $65 + 8$.
\shortans{73}
\loigiai{
Ta có $65 + 8 = 73$.
}
\end{ex}

\begin{ex}
Tính kết quả của phép tính: $74 + 7$.
\shortans{81}
\loigiai{
Ta có $74 + 7 = 81$.
}
\end{ex}

\begin{ex}
Tính kết quả của phép tính: $83 + 9$.
\shortans{92}
\loigiai{
Ta có $83 + 9 = 92$.
}
\end{ex}

\begin{ex}
Tính kết quả của phép tính: $18 + 7$.
\shortans{25}
\loigiai{
Ta có $18 + 7 = 25$.
}
\end{ex}

\begin{ex}
Tính kết quả của phép tính: $29 + 6$.
\shortans{35}
\loigiai{
Ta có $29 + 6 = 35$.
}
\end{ex}

\begin{ex}
Tính kết quả của phép tính: $36 + 8$.
\shortans{44}
\loigiai{
Ta có $36 + 8 = 44$.
}
\end{ex}

\begin{ex}
Tính kết quả của phép tính: $47 + 9$.
\shortans{56}
\loigiai{
Ta có $47 + 9 = 56$.
}
\end{ex}

\begin{ex}
Tính kết quả của phép tính: $58 + 5$.
\shortans{63}
\loigiai{
Ta có $58 + 5 = 63$.
}
\end{ex}

\begin{ex}
Tính kết quả của phép tính: $69 + 4$.
\shortans{73}
\loigiai{
Ta có $69 + 4 = 73$.
}
\end{ex}

\begin{ex}
Tính kết quả của phép tính: $77 + 6$.
\shortans{83}
\loigiai{
Ta có $77 + 6 = 83$.
}
\end{ex}

\begin{ex}
Tính kết quả của phép tính: $88 + 7$.
\shortans{95}
\loigiai{
Ta có $88 + 7 = 95$.
}
\end{ex}

\begin{ex}
Điền số thích hợp: $28 + \dots = 34$.
\shortans{6}
\loigiai{
Số cần điền là: $34 - 28 = 6$.
}
\end{ex}

\begin{ex}
Điền số thích hợp: $45 + \dots = 52$.
\shortans{7}
\loigiai{
Số cần điền là: $52 - 45 = 7$.
}
\end{ex}

\begin{ex}
Điền số thích hợp: $\dots + 8 = 43$.
\shortans{35}
\loigiai{
Số cần điền là: $43 - 8 = 35$.
}
\end{ex}

\begin{ex}
Điền số thích hợp: $\dots + 9 = 72$.
\shortans{63}
\loigiai{
Số cần điền là: $72 - 9 = 63$.
}
\end{ex}

\begin{ex}
Lan có 28 cái kẹo, mẹ cho thêm Lan 5 cái kẹo nữa. Lan có tất cả bao nhiêu cái kẹo?
\shortans{33}
\loigiai{
Lan có tất cả: $28 + 5 = 33$ cái kẹo.
}
\end{ex}

\begin{ex}
Lớp 2A có 29 học sinh, có thêm 3 bạn mới chuyển vào lớp. Hỏi lớp 2A có tất cả bao nhiêu học sinh?
\shortans{32}
\loigiai{
Lớp 2A có: $29 + 3 = 32$ học sinh.
}
\end{ex}

\begin{ex}
Bác nông dân nuôi 36 con vịt, bác mua thêm 8 con vịt nữa. Bác có tất cả bao nhiêu con vịt?
\shortans{44}
\loigiai{
Bác có: $36 + 8 = 44$ con vịt.
}
\end{ex}

\begin{ex}
Đoạn dây dài 47 cm, nối thêm một đoạn dài 8 cm. Cả đoạn dây dài bao nhiêu xăng-ti-mét?
\shortans{55}
\loigiai{
Đoạn dây dài: $47 + 8 = 55$ cm.
}
\end{ex}

\begin{ex}
Một cây bưởi có 58 quả, hôm nay có thêm 7 quả chín. Hỏi có bao nhiêu quả bưởi chín nếu trước đó có 58 quả chín?
\shortans{65}
\loigiai{
Tổng số quả là: $58 + 7 = 65$.
}
\end{ex}

\begin{ex}
Tính giá trị biểu thức: $37 + 5 + 8$.
\shortans{50}
\loigiai{
Ta có $37 + 5 = 42$, rồi $42 + 8 = 50$.
}
\end{ex}

\newpage

\TaoPhieuBaiTap{Chủ đề 4 -- Phép cộng, phép trừ (có nhớ) trong phạm vi 100}{20}{Phép cộng (có nhớ) số có hai chữ số với số có hai chữ số}{https://loptoanthaylinh.vn/bai-tap/toan-2/bai-20}
\begin{ex}
Tính kết quả của phép tính: $26 + 15$.
\shortans{41}
\loigiai{
Ta có $6 + 5 = 11$, viết 1 nhớ 1; $2 + 1 + 1 = 4$. Kết quả là 41.
}
\end{ex}

\begin{ex}
Tính kết quả của phép tính: $37 + 24$.
\shortans{61}
\loigiai{
Ta có $37 + 24 = 61$.
}
\end{ex}

\begin{ex}
Tính kết quả của phép tính: $48 + 35$.
\shortans{83}
\loigiai{
Ta có $48 + 35 = 83$.
}
\end{ex}

\begin{ex}
Tính kết quả của phép tính: $59 + 23$.
\shortans{82}
\loigiai{
Ta có $59 + 23 = 82$.
}
\end{ex}

\begin{ex}
Tính kết quả của phép tính: $18 + 26$.
\shortans{44}
\loigiai{
Ta có $18 + 26 = 44$.
}
\end{ex}

\begin{ex}
Tính kết quả của phép tính: $29 + 34$.
\shortans{63}
\loigiai{
Ta có $29 + 34 = 63$.
}
\end{ex}

\begin{ex}
Tính kết quả của phép tính: $47 + 45$.
\shortans{92}
\loigiai{
Ta có $47 + 45 = 92$.
}
\end{ex}

\begin{ex}
Tính kết quả của phép tính: $38 + 48$.
\shortans{86}
\loigiai{
Ta có $38 + 48 = 86$.
}
\end{ex}

\begin{ex}
Tính kết quả của phép tính: $56 + 27$.
\shortans{83}
\loigiai{
Ta có $56 + 27 = 83$.
}
\end{ex}

\begin{ex}
Tính kết quả của phép tính: $64 + 19$.
\shortans{83}
\loigiai{
Ta có $64 + 19 = 83$.
}
\end{ex}

\begin{ex}
Tính kết quả của phép tính: $45 + 38$.
\shortans{83}
\loigiai{
Ta có $45 + 38 = 83$.
}
\end{ex}

\begin{ex}
Tính kết quả của phép tính: $28 + 54$.
\shortans{82}
\loigiai{
Ta có $28 + 54 = 82$.
}
\end{ex}

\begin{ex}
Tính kết quả của phép tính: $19 + 68$.
\shortans{87}
\loigiai{
Ta có $19 + 68 = 87$.
}
\end{ex}

\begin{ex}
Tính kết quả của phép tính: $33 + 49$.
\shortans{82}
\loigiai{
Ta có $33 + 49 = 82$.
}
\end{ex}

\begin{ex}
Tính kết quả của phép tính: $55 + 36$.
\shortans{91}
\loigiai{
Ta có $55 + 36 = 91$.
}
\end{ex}

\begin{ex}
Điền số thích hợp: $28 + \dots = 52$.
\shortans{24}
\loigiai{
Số cần điền là: $52 - 28 = 24$.
}
\end{ex}

\begin{ex}
Điền số thích hợp: $37 + \dots = 65$.
\shortans{28}
\loigiai{
Số cần điền là: $65 - 37 = 28$.
}
\end{ex}

\begin{ex}
Điền số thích hợp: $\dots + 19 = 45$.
\shortans{26}
\loigiai{
Số cần điền là: $45 - 19 = 26$.
}
\end{ex}

\begin{ex}
Điền số thích hợp: $\dots + 38 = 71$.
\shortans{33}
\loigiai{
Số cần điền là: $71 - 38 = 33$.
}
\end{ex}

\begin{ex}
Lớp 2A có 28 học sinh, lớp 2B có 27 học sinh. Cả hai lớp có tất cả bao nhiêu học sinh?
\shortans{55}
\loigiai{
Cả hai lớp có: $28 + 27 = 55$ học sinh.
}
\end{ex}

\begin{ex}
Đàn gà thứ nhất có 38 con, đàn gà thứ hai có 25 con. Cả hai đàn gà có bao nhiêu con?
\shortans{63}
\loigiai{
Cả hai đàn gà có: $38 + 25 = 63$ con.
}
\end{ex}

\begin{ex}
Bác nông dân gặt được 46 bao thóc nếp và 37 bao thóc tẻ. Bác gặt được tất cả bao nhiêu bao thóc?
\shortans{83}
\loigiai{
Tất cả có: $46 + 37 = 83$ bao.
}
\end{ex}

\begin{ex}
Cuộn dây thứ nhất dài 39 mét, cuộn thứ hai dài 28 mét. Cả hai cuộn dài bao nhiêu mét?
\shortans{67}
\loigiai{
Hai cuộn dài: $39 + 28 = 67$ mét.
}
\end{ex}

\begin{ex}
Một cửa hàng bán được 48 kg đường vào buổi sáng và 35 kg đường vào buổi chiều. Cả ngày bán được bao nhiêu ki-lô-gam đường?
\shortans{83}
\loigiai{
Cả ngày bán được: $48 + 35 = 83$ kg.
}
\end{ex}

\begin{ex}
Tính giá trị biểu thức: $27 + 18 + 35$.
\shortans{80}
\loigiai{
Ta có $27 + 18 = 45$, rồi $45 + 35 = 80$.
}
\end{ex}

\newpage

\TaoPhieuBaiTap{Chủ đề 4 -- Phép cộng, phép trừ (có nhớ) trong phạm vi 100}{21}{Phép trừ (có nhớ) số có hai chữ số cho số có một chữ số}{https://loptoanthaylinh.vn/bai-tap/toan-2/bai-21}
\begin{ex}
Tính: $36 + 9$.
\shortans{45}
\loigiai{
Ta có $36 + 9 = 45$.
}
\end{ex}

\begin{ex}
Tính: $48 + 7$.
\shortans{55}
\loigiai{
Ta có $48 + 7 = 55$.
}
\end{ex}

\begin{ex}
Tính: $27 + 18$.
\shortans{45}
\loigiai{
Ta có $27 + 18 = 45$.
}
\end{ex}

\begin{ex}
Tính: $35 + 29$.
\shortans{64}
\loigiai{
Ta có $35 + 29 = 64$.
}
\end{ex}

\begin{ex}
Tính: $58 + 36$.
\shortans{94}
\loigiai{
Ta có $58 + 36 = 94$.
}
\end{ex}

\begin{ex}
Tính: $49 + 43$.
\shortans{92}
\loigiai{
Ta có $49 + 43 = 92$.
}
\end{ex}

\begin{ex}
Tính giá trị biểu thức: $28 + 14 + 18$.
\shortans{60}
\loigiai{
Ta có $28 + 14 = 42$, rồi $42 + 18 = 60$.
}
\end{ex}

\begin{ex}
Tính giá trị biểu thức: $37 + 25 + 18$.
\shortans{80}
\loigiai{
Ta có $37 + 25 = 62$, rồi $62 + 18 = 80$.
}
\end{ex}

\begin{ex}
Tính giá trị biểu thức: $46 + 19 + 25$.
\shortans{90}
\loigiai{
Ta có $46 + 19 = 65$, rồi $65 + 25 = 90$.
}
\end{ex}

\begin{ex}
Điền số thích hợp: $38 + 25 = 40 + \dots$
\shortans{23}
\loigiai{
Vế trái: $38 + 25 = 63$. Số cần điền: $63 - 40 = 23$.
}
\end{ex}

\begin{ex}
Điền số thích hợp: $47 + 16 = 50 + \dots$
\shortans{13}
\loigiai{
Vế trái: $47 + 16 = 63$. Số cần điền: $63 - 50 = 13$.
}
\end{ex}

\begin{ex}
Điền số thích hợp: $29 + \dots = 62$.
\shortans{33}
\loigiai{
Số cần điền: $62 - 29 = 33$.
}
\end{ex}

\begin{ex}
Điền số thích hợp: $56 + \dots = 81$.
\shortans{25}
\loigiai{
Số cần điền: $81 - 56 = 25$.
}
\end{ex}

\begin{ex}
Hà có 27 chiếc kẹp tóc, mẹ mua thêm cho Hà 15 chiếc nữa. Hà có tất cả bao nhiêu chiếc kẹp tóc?
\shortans{42}
\loigiai{
Hà có: $27 + 15 = 42$ chiếc kẹp tóc.
}
\end{ex}

\begin{ex}
Trong thư viện trường có 48 cuốn truyện tranh và 36 cuốn sách khoa học. Có tất cả bao nhiêu cuốn sách và truyện?
\shortans{84}
\loigiai{
Có tất cả: $48 + 36 = 84$ cuốn.
}
\end{ex}

\begin{ex}
Một đội trồng rừng đã trồng được 57 cây thông và 28 cây bạch đàn. Đội đã trồng được bao nhiêu cây?
\shortans{85}
\loigiai{
Đội trồng được: $57 + 28 = 85$ cây.
}
\end{ex}

\begin{ex}
Thùng thứ nhất đựng 38 lít dầu, thùng thứ hai đựng nhiều hơn thùng thứ nhất 17 lít dầu. Thùng thứ hai đựng bao nhiêu lít dầu?
\shortans{55}
\loigiai{
Thùng thứ hai đựng: $38 + 17 = 55$ lít.
}
\end{ex}

\begin{ex}
Bao gạo cân nặng 46 kg, bao ngô nặng hơn bao gạo 9 kg. Bao ngô cân nặng bao nhiêu ki-lô-gam?
\shortans{55}
\loigiai{
Bao ngô nặng: $46 + 9 = 55$ kg.
}
\end{ex}

\begin{ex}
Tìm một số biết lấy số đó trừ đi 27 thì được 38.
\shortans{65}
\loigiai{
Số đó là: $38 + 27 = 65$.
}
\end{ex}

\begin{ex}
Tìm một số biết lấy số đó trừ đi 19 thì được 56.
\shortans{75}
\loigiai{
Số đó là: $56 + 19 = 75$.
}
\end{ex}

\begin{ex}
Tổng của số liền trước 48 và số liền sau 25 là bao nhiêu?
\shortans{73}
\loigiai{
Số liền trước 48 là 47, liền sau 25 là 26. Tổng: $47 + 26 = 73$.
}
\end{ex}

\begin{ex}
Một đoàn tàu chở hàng có 29 toa hàng khô và 18 toa hàng lạnh. Đoàn tàu có tất cả bao nhiêu toa hàng?
\shortans{47}
\loigiai{
Đoàn tàu có: $29 + 18 = 47$ toa.
}
\end{ex}

\begin{ex}
Tính: $19 + 29 + 39$.
\shortans{87}
\loigiai{
Ta có $19 + 29 = 48$, rồi $48 + 39 = 87$.
}
\end{ex}

\begin{ex}
Hai số có hiệu bằng 18, số trừ là 35. Số bị trừ là bao nhiêu?
\shortans{53}
\loigiai{
Số bị trừ: $18 + 35 = 53$.
}
\end{ex}

\begin{ex}
Tính tổng của số lớn nhất có hai chữ số khác nhau và số nhỏ nhất có một chữ số khác 0 (là 98 và 1).
\shortans{99}
\loigiai{
Tổng là: $98 + 1 = 99$.
}
\end{ex}

\newpage

\TaoPhieuBaiTap{Chủ đề 4 -- Phép cộng, phép trừ (có nhớ) trong phạm vi 100}{22}{Phép trừ (có nhớ) số có hai chữ số cho số có hai chữ số}{https://loptoanthaylinh.vn/bai-tap/toan-2/bai-22}
\begin{ex}
Tính kết quả của phép tính: $32 - 5$.
\shortans{27}
\loigiai{
Ta có 2 không trừ được 5, lấy $12 - 5 = 7$, nhớ 1; $3 - 1 = 2$. Kết quả là 27.
}
\end{ex}

\begin{ex}
Tính kết quả của phép tính: $41 - 4$.
\shortans{37}
\loigiai{
Ta có $41 - 4 = 37$.
}
\end{ex}

\begin{ex}
Tính kết quả của phép tính: $53 - 6$.
\shortans{47}
\loigiai{
Ta có $53 - 6 = 47$.
}
\end{ex}

\begin{ex}
Tính kết quả của phép tính: $62 - 7$.
\shortans{55}
\loigiai{
Ta có $62 - 7 = 55$.
}
\end{ex}

\begin{ex}
Tính kết quả của phép tính: $74 - 8$.
\shortans{66}
\loigiai{
Ta có $74 - 8 = 66$.
}
\end{ex}

\begin{ex}
Tính kết quả của phép tính: $85 - 9$.
\shortans{76}
\loigiai{
Ta có $85 - 9 = 76$.
}
\end{ex}

\begin{ex}
Tính kết quả của phép tính: $91 - 5$.
\shortans{86}
\loigiai{
Ta có $91 - 5 = 86$.
}
\end{ex}

\begin{ex}
Tính kết quả của phép tính: $24 - 8$.
\shortans{16}
\loigiai{
Ta có $24 - 8 = 16$.
}
\end{ex}

\begin{ex}
Tính kết quả của phép tính: $35 - 7$.
\shortans{28}
\loigiai{
Ta có $35 - 7 = 28$.
}
\end{ex}

\begin{ex}
Tính kết quả của phép tính: $46 - 9$.
\shortans{37}
\loigiai{
Ta có $46 - 9 = 37$.
}
\end{ex}

\begin{ex}
Tính kết quả của phép tính: $50 - 6$.
\shortans{44}
\loigiai{
Ta có $50 - 6 = 44$.
}
\end{ex}

\begin{ex}
Tính kết quả của phép tính: $70 - 8$.
\shortans{62}
\loigiai{
Ta có $70 - 8 = 62$.
}
\end{ex}

\begin{ex}
Tính kết quả của phép tính: $90 - 4$.
\shortans{86}
\loigiai{
Ta có $90 - 4 = 86$.
}
\end{ex}

\begin{ex}
Tính kết quả của phép tính: $82 - 7$.
\shortans{75}
\loigiai{
Ta có $82 - 7 = 75$.
}
\end{ex}

\begin{ex}
Tính kết quả của phép tính: $63 - 8$.
\shortans{55}
\loigiai{
Ta có $63 - 8 = 55$.
}
\end{ex}

\begin{ex}
Điền số thích hợp: $43 - \dots = 36$.
\shortans{7}
\loigiai{
Số cần điền là: $43 - 36 = 7$.
}
\end{ex}

\begin{ex}
Điền số thích hợp: $61 - \dots = 55$.
\shortans{6}
\loigiai{
Số cần điền là: $61 - 55 = 6$.
}
\end{ex}

\begin{ex}
Điền số thích hợp: $\dots - 8 = 45$.
\shortans{53}
\loigiai{
Số cần điền là: $45 + 8 = 53$.
}
\end{ex}

\begin{ex}
Điền số thích hợp: $\dots - 9 = 74$.
\shortans{83}
\loigiai{
Số cần điền là: $74 + 9 = 83$.
}
\end{ex}

\begin{ex}
Mẹ có 31 quả trứng, mẹ dùng 6 quả để nướng bánh. Mẹ còn lại bao nhiêu quả trứng?
\shortans{25}
\loigiai{
Mẹ còn lại: $31 - 6 = 25$ quả trứng.
}
\end{ex}

\begin{ex}
Vườn nhà bà có 42 con gà, bão làm thất lạc mất 5 con gà. Hỏi vườn nhà bà còn lại bao nhiêu con gà?
\shortans{37}
\loigiai{
Còn lại: $42 - 5 = 37$ con gà.
}
\end{ex}

\begin{ex}
Một sợi dây dài 54 cm, cắt bớt 7 cm. Đoạn dây còn lại dài bao nhiêu xăng-ti-mét?
\shortans{47}
\loigiai{
Còn lại: $54 - 7 = 47$ cm.
}
\end{ex}

\begin{ex}
Một thùng nước có 60 lít, người ta múc ra 8 lít nước. Trong thùng còn lại bao nhiêu lít nước?
\shortans{52}
\loigiai{
Trong thùng còn lại: $60 - 8 = 52$ lít nước.
}
\end{ex}

\begin{ex}
Bao thóc nặng 75 kg, người ta lấy ra 9 kg thóc để xay xát gạo. Bao thóc còn lại nặng bao nhiêu ki-lô-gam?
\shortans{66}
\loigiai{
Còn lại: $75 - 9 = 66$ kg.
}
\end{ex}

\begin{ex}
Tính giá trị biểu thức: $52 - 7 - 5$.
\shortans{40}
\loigiai{
Ta có $52 - 7 = 45$, rồi $45 - 5 = 40$.
}
\end{ex}

\newpage

\TaoPhieuBaiTap{Chủ đề 4 -- Phép cộng, phép trừ (có nhớ) trong phạm vi 100}{23}{Luyện tập phép cộng, phép trừ có nhớ}{https://loptoanthaylinh.vn/bai-tap/toan-2/bai-23}
\begin{ex}
Tính kết quả của phép tính: $42 - 17$.
\shortans{25}
\loigiai{
Ta có 2 không trừ được 7, lấy $12 - 7 = 5$, nhớ 1; $1 + 1 = 2$, rồi $4 - 2 = 2$. Kết quả là 25.
}
\end{ex}

\begin{ex}
Tính kết quả của phép tính: $53 - 28$.
\shortans{25}
\loigiai{
Ta có $53 - 28 = 25$.
}
\end{ex}

\begin{ex}
Tính kết quả của phép tính: $61 - 35$.
\shortans{26}
\loigiai{
Ta có $61 - 35 = 26$.
}
\end{ex}

\begin{ex}
Tính kết quả của phép tính: $74 - 46$.
\shortans{28}
\loigiai{
Ta có $74 - 46 = 28$.
}
\end{ex}

\begin{ex}
Tính kết quả của phép tính: $85 - 59$.
\shortans{26}
\loigiai{
Ta có $85 - 59 = 26$.
}
\end{ex}

\begin{ex}
Tính kết quả của phép tính: $92 - 64$.
\shortans{28}
\loigiai{
Ta có $92 - 64 = 28$.
}
\end{ex}

\begin{ex}
Tính kết quả của phép tính: $50 - 24$.
\shortans{26}
\loigiai{
Ta có $50 - 24 = 26$.
}
\end{ex}

\begin{ex}
Tính kết quả của phép tính: $70 - 38$.
\shortans{32}
\loigiai{
Ta có $70 - 38 = 32$.
}
\end{ex}

\begin{ex}
Tính kết quả của phép tính: $80 - 45$.
\shortans{35}
\loigiai{
Ta có $80 - 45 = 35$.
}
\end{ex}

\begin{ex}
Tính kết quả của phép tính: $90 - 57$.
\shortans{33}
\loigiai{
Ta có $90 - 57 = 33$.
}
\end{ex}

\begin{ex}
Tính kết quả của phép tính: $66 - 29$.
\shortans{37}
\loigiai{
Ta có $66 - 29 = 37$.
}
\end{ex}

\begin{ex}
Tính kết quả của phép tính: $73 - 36$.
\shortans{37}
\loigiai{
Ta có $73 - 36 = 37$.
}
\end{ex}

\begin{ex}
Tính kết quả của phép tính: $84 - 48$.
\shortans{36}
\loigiai{
Ta có $84 - 48 = 36$.
}
\end{ex}

\begin{ex}
Tính kết quả của phép tính: $95 - 67$.
\shortans{28}
\loigiai{
Ta có $95 - 67 = 28$.
}
\end{ex}

\begin{ex}
Tính kết quả của phép tính: $63 - 39$.
\shortans{24}
\loigiai{
Ta có $63 - 39 = 24$.
}
\end{ex}

\begin{ex}
Điền số thích hợp: $52 - \dots = 27$.
\shortans{25}
\loigiai{
Số cần điền là: $52 - 27 = 25$.
}
\end{ex}

\begin{ex}
Điền số thích hợp: $71 - \dots = 38$.
\shortans{33}
\loigiai{
Số cần điền là: $71 - 38 = 33$.
}
\end{ex}

\begin{ex}
Điền số thích hợp: $\dots - 26 = 47$.
\shortans{73}
\loigiai{
Số cần điền là: $47 + 26 = 73$.
}
\end{ex}

\begin{ex}
Điền số thích hợp: $\dots - 35 = 28$.
\shortans{63}
\loigiai{
Số cần điền là: $28 + 35 = 63$.
}
\end{ex}

\begin{ex}
Một cửa hàng có 65 mét vải hoa, đã bán được 28 mét vải. Cửa hàng còn lại bao nhiêu mét vải hoa?
\shortans{37}
\loigiai{
Còn lại: $65 - 28 = 37$ mét vải.
}
\end{ex}

\begin{ex}
Đàn vịt có 72 con, mẹ đem bán ở chợ 35 con. Đàn vịt còn lại bao nhiêu con?
\shortans{37}
\loigiai{
Còn lại: $72 - 35 = 37$ con vịt.
}
\end{ex}

\begin{ex}
Trường tiểu học có 80 bạn dự thi học sinh giỏi, trong đó có 43 bạn nữ. Hỏi có bao nhiêu bạn nam dự thi?
\shortans{37}
\loigiai{
Số bạn nam: $80 - 43 = 37$ bạn.
}
\end{ex}

\begin{ex}
Kho thứ nhất có 93 bao gạo, kho thứ hai có ít hơn kho thứ nhất 26 bao gạo. Kho thứ hai có bao nhiêu bao gạo?
\shortans{67}
\loigiai{
Kho thứ hai có: $93 - 26 = 67$ bao gạo.
}
\end{ex}

\begin{ex}
Một xe buýt chở 54 hành khách, tại trạm dừng có 28 người xuống xe và không có ai lên. Trên xe buýt còn lại bao nhiêu người?
\shortans{26}
\loigiai{
Còn lại: $54 - 28 = 26$ người.
}
\end{ex}

\begin{ex}
Tính giá trị biểu thức: $82 - 37 - 15$.
\shortans{30}
\loigiai{
Ta có $82 - 37 = 45$, rồi $45 - 15 = 30$.
}
\end{ex}

\newpage

\TaoPhieuBaiTap{Chủ đề 4 -- Phép cộng, phép trừ (có nhớ) trong phạm vi 100}{24}{Luyện tập chung}{https://loptoanthaylinh.vn/bai-tap/toan-2/bai-24}
\begin{ex}
Tính: $35 + 47$.
\shortans{82}
\loigiai{
Ta có $35 + 47 = 82$.
}
\end{ex}

\begin{ex}
Tính: $82 - 35$.
\shortans{47}
\loigiai{
Ta có $82 - 35 = 47$.
}
\end{ex}

\begin{ex}
Tính: $82 - 47$.
\shortans{35}
\loigiai{
Ta có $82 - 47 = 35$.
}
\end{ex}

\begin{ex}
Tính: $46 + 28$.
\shortans{74}
\loigiai{
Ta có $46 + 28 = 74$.
}
\end{ex}

\begin{ex}
Tính: $74 - 46$.
\shortans{28}
\loigiai{
Ta có $74 - 46 = 28$.
}
\end{ex}

\begin{ex}
Tính: $53 + 39$.
\shortans{92}
\loigiai{
Ta có $53 + 39 = 92$.
}
\end{ex}

\begin{ex}
Tính: $92 - 53$.
\shortans{39}
\loigiai{
Ta có $92 - 53 = 39$.
}
\end{ex}

\begin{ex}
Tính: $60 - 24$.
\shortans{36}
\loigiai{
Ta có $60 - 24 = 36$.
}
\end{ex}

\begin{ex}
Tính: $100 - 45$.
\shortans{55}
\loigiai{
Ta có $100 - 45 = 55$.
}
\end{ex}

\begin{ex}
Tính: $100 - 68$.
\shortans{32}
\loigiai{
Ta có $100 - 68 = 32$.
}
\end{ex}

\begin{ex}
Tính giá trị biểu thức: $38 + 26 - 19$.
\shortans{45}
\loigiai{
Ta có $38 + 26 = 64$, rồi $64 - 19 = 45$.
}
\end{ex}

\begin{ex}
Tính giá trị biểu thức: $72 - 35 + 18$.
\shortans{55}
\loigiai{
Ta có $72 - 35 = 37$, rồi $37 + 18 = 55$.
}
\end{ex}

\begin{ex}
Tính giá trị biểu thức: $81 - 46 - 15$.
\shortans{20}
\loigiai{
Ta có $81 - 46 = 35$, rồi $35 - 15 = 20$.
}
\end{ex}

\begin{ex}
Điền số thích hợp: $45 + \dots = 81$.
\shortans{36}
\loigiai{
Số cần điền: $81 - 45 = 36$.
}
\end{ex}

\begin{ex}
Điền số thích hợp: $90 - \dots = 37$.
\shortans{53}
\loigiai{
Số cần điền: $90 - 37 = 53$.
}
\end{ex}

\begin{ex}
Điền số thích hợp: $\dots - 28 = 46$.
\shortans{74}
\loigiai{
Số cần điền: $46 + 28 = 74$.
}
\end{ex}

\begin{ex}
Bác nông dân có 62 quả dưa hấu, bác đã bán đi 38 quả. Bác còn lại bao nhiêu quả dưa hấu?
\shortans{24}
\loigiai{
Còn lại: $62 - 38 = 24$ quả.
}
\end{ex}

\begin{ex}
Một bể nước có 85 lít nước, người ta xả bớt 39 lít nước. Trong bể còn lại bao nhiêu lít nước?
\shortans{46}
\loigiai{
Còn lại: $85 - 39 = 46$ lít.
}
\end{ex}

\begin{ex}
Đàn gia súc có 73 con gồm bò và dê, trong đó có 36 con bò. Hỏi có bao nhiêu con dê?
\shortans{37}
\loigiai{
Số dê là: $73 - 36 = 37$ con.
}
\end{ex}

\begin{ex}
Hai số có tổng bằng 91, số hạng thứ nhất là 45. Số hạng thứ hai là bao nhiêu?
\shortans{46}
\loigiai{
Số hạng thứ hai: $91 - 45 = 46$.
}
\end{ex}

\begin{ex}
Hai số có hiệu bằng 27, số bị trừ là 63. Số trừ là bao nhiêu?
\shortans{36}
\loigiai{
Số trừ: $63 - 27 = 36$.
}
\end{ex}

\begin{ex}
Một cuộn dây thép dài 100 mét, bác thợ cắt đi 35 mét. Cuộn dây còn lại bao nhiêu mét?
\shortans{65}
\loigiai{
Còn lại: $100 - 35 = 65$ mét.
}
\end{ex}

\begin{ex}
Quyển sách dày 94 trang, An đã đọc được 58 trang. An còn phải đọc bao nhiêu trang nữa để hết cuốn sách?
\shortans{36}
\loigiai{
Còn phải đọc: $94 - 58 = 36$ trang.
}
\end{ex}

\begin{ex}
Bao ngô nặng 52 kg, bao thóc nhẹ hơn bao ngô 16 kg. Hỏi bao thóc nặng bao nhiêu ki-lô-gam?
\shortans{36}
\loigiai{
Bao thóc nặng: $52 - 16 = 36$ kg.
}
\end{ex}

\begin{ex}
Điền số thích hợp vào chỗ chấm: $65 - 28 = 15 + \dots$
\shortans{22}
\loigiai{
Vế trái: $65 - 28 = 37$. Số cần điền là $37 - 15 = 22$.
}
\end{ex}

\newpage


% =========================================================================
% Làm quen với hình phẳng
% =========================================================================
\section{Làm quen với hình phẳng}

\TaoPhieuBaiTap{Chủ đề 5 -- Làm quen với hình phẳng}{25}{Điểm, đoạn thẳng}{https://loptoanthaylinh.vn/bai-tap/toan-2/bai-25}
\begin{ex}
Một đoạn thẳng được giới hạn bởi bao nhiêu điểm đầu mút?
\shortans{2}
\loigiai{
Đoạn thẳng bị giới hạn bởi 2 điểm đầu mút.
}
\end{ex}

\begin{ex}
Qua 2 điểm phân biệt, ta vẽ được bao nhiêu đường thẳng?
\shortans{1}
\loigiai{
Qua 2 điểm phân biệt vẽ được duy nhất 1 đường thẳng.
}
\end{ex}

\begin{ex}
Cho 3 điểm A, B, C cùng nằm trên một đường thẳng. Hỏi có bao nhiêu điểm thẳng hàng?
\shortans{3}
\loigiai{
Có 3 điểm thẳng hàng.
}
\end{ex}

\begin{ex}
Cho 3 điểm không thẳng hàng, nối từng cặp điểm ta được bao nhiêu đoạn thẳng?
\shortans{3}
\loigiai{
Nối được 3 đoạn thẳng tạo thành một hình tam giác.
}
\end{ex}

\begin{ex}
Cho 4 điểm trong đó không có 3 điểm nào thẳng hàng. Nối từng cặp điểm ta vẽ được tất cả bao nhiêu đoạn thẳng?
\shortans{6}
\loigiai{
Số đoạn thẳng: $4 \times 3 : 2 = 6$ đoạn thẳng.
}
\end{ex}

\begin{ex}
Trên một đường thẳng lấy 4 điểm phân biệt A, B, C, D theo thứ tự đó. Có bao nhiêu đoạn thẳng trên hình vẽ?
\shortans{6}
\loigiai{
Các đoạn thẳng: AB, AC, AD, BC, BD, CD. Có 6 đoạn thẳng.
}
\end{ex}

\begin{ex}
Đoạn thẳng AB dài 8 cm, điểm C là điểm nằm giữa A và B sao cho AC dài 3 cm. Độ dài đoạn thẳng CB là bao nhiêu xăng-ti-mét?
\shortans{5}
\loigiai{
Độ dài đoạn CB: $8 - 3 = 5$ cm.
}
\end{ex}

\begin{ex}
Cho ba điểm M, N, P thẳng hàng theo thứ tự đó. Biết MN dài 4 cm, NP dài 5 cm. Độ dài đoạn thẳng MP là bao nhiêu xăng-ti-mét?
\shortans{9}
\loigiai{
Độ dài MP: $4 + 5 = 9$ cm.
}
\end{ex}

\begin{ex}
Có bao nhiêu điểm đầu mút trong 3 đoạn thẳng rời nhau?
\shortans{6}
\loigiai{
Mỗi đoạn có 2 đầu mút, 3 đoạn có: $3 \times 2 = 6$ đầu mút.
}
\end{ex}

\begin{ex}
Một hình tam giác có bao nhiêu đỉnh (điểm)?
\shortans{3}
\loigiai{
Hình tam giác có 3 đỉnh.
}
\end{ex}

\begin{ex}
Một hình tứ giác có bao nhiêu đỉnh (điểm)?
\shortans{4}
\loigiai{
Hình tứ giác có 4 đỉnh.
}
\end{ex}

\begin{ex}
Trên đoạn thẳng AB dài 12 cm, người ta đánh dấu thêm điểm O sao cho OA dài 6 cm. Điểm O chia đoạn AB thành 2 đoạn bằng nhau. Độ dài đoạn OB là bao nhiêu xăng-ti-mét?
\shortans{6}
\loigiai{
Độ dài OB: $12 - 6 = 6$ cm.
}
\end{ex}

\begin{ex}
Cho 5 điểm thẳng hàng liên tiếp cách đều nhau 2 cm. Khoảng cách từ điểm thứ nhất đến điểm thứ năm là bao nhiêu xăng-ti-mét?
\shortans{8}
\loigiai{
Có 4 khoảng cách, mỗi khoảng 2 cm: $2 \times 4 = 8$ cm.
}
\end{ex}

\begin{ex}
Có bao nhiêu đường cong khép kín trong hình vẽ gồm 2 hình tròn và 3 hình vuông?
\shortans{2}
\loigiai{
Hình tròn là đường cong khép kín, có 2 hình tròn.
}
\end{ex}

\begin{ex}
Một đường gấp khúc gồm 2 đoạn thẳng nối tiếp nhau có bao nhiêu điểm đỉnh (tính cả 2 đầu mút)?
\shortans{3}
\loigiai{
Có 3 điểm.
}
\end{ex}

\begin{ex}
Một đường gấp khúc gồm 3 đoạn thẳng nối tiếp nhau có bao nhiêu điểm?
\shortans{4}
\loigiai{
Có 4 điểm.
}
\end{ex}

\begin{ex}
Đoạn thẳng CD dài 15 cm. Điểm K nằm giữa C và D sao cho KD dài 7 cm. Đoạn CK dài bao nhiêu xăng-ti-mét?
\shortans{8}
\loigiai{
Đoạn CK dài: $15 - 7 = 8$ cm.
}
\end{ex}

\begin{ex}
Có 3 điểm thẳng hàng X, Y, Z. Biết XY dài 6 cm, YZ dài 8 cm. Hỏi đoạn XZ dài bao nhiêu xăng-ti-mét nếu Y nằm giữa X và Z?
\shortans{14}
\loigiai{
Đoạn XZ dài: $6 + 8 = 14$ cm.
}
\end{ex}

\begin{ex}
Trên một sợi dây thẳng dài 20 cm, người ta thắt 3 nút chia sợi dây thành 4 phần bằng nhau. Mỗi phần dài bao nhiêu xăng-ti-mét?
\shortans{5}
\loigiai{
Mỗi phần dài: $20 : 4 = 5$ cm.
}
\end{ex}

\begin{ex}
Số góc vuông trong một hình chữ nhật là bao nhiêu?
\shortans{4}
\loigiai{
Hình chữ nhật có 4 góc vuông.
}
\end{ex}

\begin{ex}
Trong các chữ cái in hoa: A, B, C, D, O, chữ cái nào được tạo bởi một đường cong khép kín tròn? Chữ O gồm bao nhiêu nét cong khép kín?
\shortans{1}
\loigiai{
Chữ O gồm 1 đường cong khép kín.
}
\end{ex}

\begin{ex}
Một người cắm 4 cọc tiêu thẳng hàng, mỗi cọc cách nhau 3 mét. Khoảng cách từ cọc đầu tiên đến cọc cuối cùng là bao nhiêu mét?
\shortans{9}
\loigiai{
Có 3 khoảng cách: $3 \times 3 = 9$ mét.
}
\end{ex}

\begin{ex}
Để nối 5 điểm thẳng hàng liên tiếp cần bao nhiêu đoạn thẳng nối tiếp nhau?
\shortans{4}
\loigiai{
Cần 4 đoạn thẳng nối tiếp.
}
\end{ex}

\begin{ex}
Một tam giác có chu vi là 18 cm, hai cạnh dài lần lượt là 6 cm và 5 cm. Cạnh thứ ba dài bao nhiêu xăng-ti-mét?
\shortans{7}
\loigiai{
Cạnh thứ ba dài: $18 - 6 - 5 = 7$ cm.
}
\end{ex}

\begin{ex}
Cho đoạn thẳng PQ dài 10 cm. Điểm I là trung điểm của PQ. Đoạn PI dài bao nhiêu xăng-ti-mét?
\shortans{5}
\loigiai{
Đoạn PI dài: $10 : 2 = 5$ cm.
}
\end{ex}

\newpage

\TaoPhieuBaiTap{Chủ đề 5 -- Làm quen với hình phẳng}{26}{Đường thẳng, đường cong, đường gấp khúc}{https://loptoanthaylinh.vn/bai-tap/toan-2/bai-26}
\begin{ex}
Đường gấp khúc ABC gồm bao nhiêu đoạn thẳng?
\shortans{2}
\loigiai{
Gồm 2 đoạn thẳng là AB và BC.
}
\end{ex}

\begin{ex}
Đường gấp khúc ABCD gồm bao nhiêu đoạn thẳng?
\shortans{3}
\loigiai{
Gồm 3 đoạn thẳng: AB, BC, CD.
}
\end{ex}

\begin{ex}
Đường gấp khúc MNPD gồm 3 đoạn thẳng có độ dài lần lượt là 3 cm, 4 cm, 5 cm. Tính độ dài đường gấp khúc MNPD (đơn vị cm).
\shortans{12}
\loigiai{
Độ dài đường gấp khúc: $3 + 4 + 5 = 12$ cm.
}
\end{ex}

\begin{ex}
Đường gấp khúc ABC có AB dài 12 cm, BC dài 15 cm. Tính độ dài đường gấp khúc ABC (đơn vị cm).
\shortans{27}
\loigiai{
Độ dài: $12 + 15 = 27$ cm.
}
\end{ex}

\begin{ex}
Đường gấp khúc gồm 4 đoạn thẳng bằng nhau, mỗi đoạn dài 5 cm. Tính độ dài đường gấp khúc đó (đơn vị cm).
\shortans{20}
\loigiai{
Độ dài: $5 \times 4 = 20$ cm.
}
\end{ex}

\begin{ex}
Đường gấp khúc dài 30 cm gồm 3 đoạn thẳng, biết hai đoạn đầu dài lần lượt là 8 cm và 12 cm. Đoạn thứ ba dài bao nhiêu xăng-ti-mét?
\shortans{10}
\loigiai{
Đoạn thứ ba dài: $30 - 8 - 12 = 10$ cm.
}
\end{ex}

\begin{ex}
Một hình tứ giác có bao nhiêu cạnh?
\shortans{4}
\loigiai{
Hình tứ giác có đúng 4 cạnh.
}
\end{ex}

\begin{ex}
Một hình tứ giác có bao nhiêu đỉnh?
\shortans{4}
\loigiai{
Hình tứ giác có đúng 4 đỉnh.
}
\end{ex}

\begin{ex}
Một hình tứ giác có độ dài các cạnh lần lượt là 3 cm, 4 cm, 5 cm, 6 cm. Tính chu vi hình tứ giác đó (đơn vị cm).
\shortans{18}
\loigiai{
Chu vi: $3 + 4 + 5 + 6 = 18$ cm.
}
\end{ex}

\begin{ex}
Một hình vuông có cạnh dài 5 cm. Chu vi hình vuông đó là bao nhiêu xăng-ti-mét?
\shortans{20}
\loigiai{
Chu vi: $5 \times 4 = 20$ cm.
}
\end{ex}

\begin{ex}
Một hình chữ nhật có chiều dài 6 cm, chiều rộng 4 cm. Chu vi hình chữ nhật đó là bao nhiêu xăng-ti-mét?
\shortans{20}
\loigiai{
Chu vi: $(6 + 4) \times 2 = 20$ cm.
}
\end{ex}

\begin{ex}
Đường gấp khúc khép kín gồm 4 đoạn thẳng tạo thành hình gì có 4 cạnh? Hình đó có bao nhiêu đỉnh?
\shortans{4}
\loigiai{
Đó là hình tứ giác có 4 đỉnh.
}
\end{ex}

\begin{ex}
Đường gấp khúc khép kín gồm 3 đoạn thẳng tạo thành hình gì? Hình tam giác có bao nhiêu cạnh?
\shortans{3}
\loigiai{
Hình tam giác có 3 cạnh.
}
\end{ex}

\begin{ex}
Một sợi dây đồng uốn thành hình tứ giác có độ dài các cạnh là 10 cm, 12 cm, 8 cm, 15 cm. Chiều dài sợi dây đồng là bao nhiêu xăng-ti-mét?
\shortans{45}
\loigiai{
Chiều dài sợi dây: $10 + 12 + 8 + 15 = 45$ cm.
}
\end{ex}

\begin{ex}
Một đường gấp khúc gồm 2 đoạn thẳng dài 24 cm, đoạn thứ nhất dài 10 cm. Đoạn thứ hai dài bao nhiêu xăng-ti-mét?
\shortans{14}
\loigiai{
Đoạn thứ hai dài: $24 - 10 = 14$ cm.
}
\end{ex}

\begin{ex}
Một hình tứ giác có 4 cạnh bằng nhau, mỗi cạnh dài 6 cm. Chu vi hình tứ giác đó là bao nhiêu xăng-ti-mét?
\shortans{24}
\loigiai{
Chu vi: $6 \times 4 = 24$ cm.
}
\end{ex}

\begin{ex}
Một con ốc sên bò theo đường gấp khúc gồm 3 đoạn: đoạn 1 dài 15 cm, đoạn 2 dài 20 cm, đoạn 3 dài 25 cm. Con ốc sên đã bò được quãng đường dài bao nhiêu xăng-ti-mét?
\shortans{60}
\loigiai{
Quãng đường: $15 + 20 + 25 = 60$ cm.
}
\end{ex}

\begin{ex}
Chu vi một hình tam giác là 24 cm. Nếu 3 cạnh bằng nhau thì mỗi cạnh dài bao nhiêu xăng-ti-mét?
\shortans{8}
\loigiai{
Mỗi cạnh dài: $24 : 3 = 8$ cm.
}
\end{ex}

\begin{ex}
Chu vi một hình tứ giác là 36 cm, tổng độ dài ba cạnh là 27 cm. Cạnh thứ tư dài bao nhiêu xăng-ti-mét?
\shortans{9}
\loigiai{
Cạnh thứ tư dài: $36 - 27 = 9$ cm.
}
\end{ex}

\begin{ex}
Có bao nhiêu hình tứ giác trong một hình vuông được chia làm 4 ô vuông nhỏ bởi 2 đường thẳng chữ thập?
\shortans{9}
\loigiai{
Có 4 hình vuông nhỏ, 4 hình chữ nhật ghép đôi, và 1 hình vuông lớn. Tổng cộng có 9 hình tứ giác.
}
\end{ex}

\begin{ex}
Một đường gấp khúc có 5 đoạn thẳng bằng nhau, mỗi đoạn dài 4 cm. Độ dài đường gấp khúc là bao nhiêu xăng-ti-mét?
\shortans{20}
\loigiai{
Độ dài: $4 \times 5 = 20$ cm.
}
\end{ex}

\begin{ex}
Một khung ảnh hình chữ nhật có chu vi là 30 cm, nửa chu vi khung ảnh là bao nhiêu xăng-ti-mét?
\shortans{15}
\loigiai{
Nửa chu vi: $30 : 2 = 15$ cm.
}
\end{ex}

\begin{ex}
Đường gấp khúc ABCDE gồm bao nhiêu đoạn thẳng?
\shortans{4}
\loigiai{
Gồm 4 đoạn thẳng: AB, BC, CD, DE.
}
\end{ex}

\begin{ex}
Đường gấp khúc đó có bao nhiêu điểm đỉnh?
\shortans{5}
\loigiai{
Có 5 điểm đỉnh: A, B, C, D, E.
}
\end{ex}

\begin{ex}
Tính độ dài đường gấp khúc gồm 3 đoạn thẳng dài lần lượt 14 cm, 16 cm và 20 cm.
\shortans{50}
\loigiai{
Độ dài: $14 + 16 + 20 = 50$ cm.
}
\end{ex}

\newpage

\TaoPhieuBaiTap{Chủ đề 5 -- Làm quen với hình phẳng}{27}{Ba điểm thẳng hàng}{https://loptoanthaylinh.vn/bai-tap/toan-2/bai-27}
\begin{ex}
Để vẽ đoạn thẳng AB dài 7 cm, ta đặt thước kẻ sao cho điểm A trùng với vạch số mấy của thước?
\shortans{0}
\loigiai{
Điểm A trùng vạch 0 cm.
}
\end{ex}

\begin{ex}
Điểm B sẽ trùng với vạch số mấy của thước kẻ?
\shortans{7}
\loigiai{
Điểm B trùng vạch số 7 cm.
}
\end{ex}

\begin{ex}
Cắt một tờ giấy hình vuông thành hai hình tam giác bằng nhau bởi một đường chéo. Cần bao nhiêu nhát kéo cắt thẳng?
\shortans{1}
\loigiai{
Cần 1 nhát kéo cắt theo đường chéo.
}
\end{ex}

\begin{ex}
Gấp đôi một tờ giấy hình chữ nhật theo chiều ngang, ta nhận được 2 hình gì bằng nhau? Mỗi phần là hình chữ nhật nhỏ. Có bao nhiêu hình chữ nhật nhỏ?
\shortans{2}
\loigiai{
Có 2 hình chữ nhật nhỏ.
}
\end{ex}

\begin{ex}
Dùng 4 que tính bằng nhau xếp được 1 hình vuông. Muốn xếp 3 hình vuông rời nhau cần bao nhiêu que tính?
\shortans{12}
\loigiai{
Cần: $4 \times 3 = 12$ que tính.
}
\end{ex}

\begin{ex}
Dùng 3 que tính bằng nhau xếp được 1 hình tam giác. Muốn xếp 4 hình tam giác rời nhau cần bao nhiêu que tính?
\shortans{12}
\loigiai{
Cần: $3 \times 4 = 12$ que tính.
}
\end{ex}

\begin{ex}
Dùng 7 que tính bằng nhau xếp được 2 hình vuông chung 1 cạnh. Cần bao nhiêu que tính?
\shortans{7}
\loigiai{
Cần đúng 7 que tính.
}
\end{ex}

\begin{ex}
Dùng 5 que tính bằng nhau xếp được 2 hình tam giác chung 1 cạnh. Cần bao nhiêu que tính?
\shortans{5}
\loigiai{
Cần đúng 5 que tính.
}
\end{ex}

\begin{ex}
Một miếng bìa hình vuông có chu vi 16 cm. Cạnh của miếng bìa hình vuông dài bao nhiêu xăng-ti-mét?
\shortans{4}
\loigiai{
Cạnh dài: $16 : 4 = 4$ cm.
}
\end{ex}

\begin{ex}
Ghép 2 hình tam giác vuông cân bằng nhau có thể tạo thành hình gì có 4 cạnh bằng nhau? Hình vuông đó có bao nhiêu cạnh?
\shortans{4}
\loigiai{
Tạo thành hình vuông có 4 cạnh.
}
\end{ex}

\begin{ex}
Ghép 6 hình tam giác đều nhỏ bằng nhau có thể tạo thành một hình lục giác đều. Hình lục giác đều có bao nhiêu cạnh?
\shortans{6}
\loigiai{
Hình lục giác có 6 cạnh.
}
\end{ex}

\begin{ex}
Cắt bỏ 1 góc của hình vuông bằng 1 đường cắt thẳng. Hình mới nhận được có bao nhiêu cạnh?
\shortans{5}
\loigiai{
Hình mới có 5 cạnh.
}
\end{ex}

\begin{ex}
Cắt đôi một hình vuông bằng một đường thẳng đi qua trung điểm hai cạnh đối diện, ta nhận được bao nhiêu hình chữ nhật?
\shortans{2}
\loigiai{
Nhận được 2 hình chữ nhật.
}
\end{ex}

\begin{ex}
Một hình chữ nhật được ghép từ 3 hình vuông nhỏ cạnh 2 cm xếp thành hàng ngang. Chiều dài hình chữ nhật đó là bao nhiêu xăng-ti-mét?
\shortans{6}
\loigiai{
Chiều dài: $2 \times 3 = 6$ cm.
}
\end{ex}

\begin{ex}
Chiều rộng của hình chữ nhật đó là bao nhiêu xăng-ti-mét?
\shortans{2}
\loigiai{
Chiều rộng bằng cạnh hình vuông là 2 cm.
}
\end{ex}

\begin{ex}
Chu vi của hình chữ nhật ghép từ 3 hình vuông đó là bao nhiêu xăng-ti-mét?
\shortans{16}
\loigiai{
Chu vi: $(6 + 2) \times 2 = 16$ cm.
}
\end{ex}

\begin{ex}
Có 8 que tính dài bằng nhau, xếp được nhiều nhất bao nhiêu hình vuông rời nhau?
\shortans{2}
\loigiai{
Xếp được: $8 : 4 = 2$ hình vuông.
}
\end{ex}

\begin{ex}
Có 9 que tính dài bằng nhau, xếp được nhiều nhất bao nhiêu hình tam giác rời nhau?
\shortans{3}
\loigiai{
Xếp được: $9 : 3 = 3$ hình tam giác.
}
\end{ex}

\begin{ex}
Một tờ giấy hình chữ nhật có chiều dài 15 cm, chiều rộng 10 cm. Cắt một hình vuông có cạnh 10 cm từ tờ giấy đó. Mảnh giấy còn lại là hình chữ nhật có chiều dài 10 cm và chiều rộng bao nhiêu xăng-ti-mét?
\shortans{5}
\loigiai{
Chiều rộng mảnh còn lại: $15 - 10 = 5$ cm.
}
\end{ex}

\begin{ex}
Bác thợ mộc có 10 thanh gỗ dài bằng nhau để đóng các khung tranh hình tam giác rời nhau. Bác đóng được nhiều nhất bao nhiêu khung tranh?
\shortans{3}
\loigiai{
Vì $10 : 3 = 3$ (dư 1) nên đóng được nhiều nhất 3 khung tranh.
}
\end{ex}

\begin{ex}
Lan dùng các que tính xếp thành chữ số 0 trên bảng điện tử. Chữ số 0 cần bao nhiêu que tính?
\shortans{6}
\loigiai{
Cần 6 đoạn que tính (2 ngang, 4 dọc).
}
\end{ex}

\begin{ex}
Chữ số 8 trên bảng điện tử cần bao nhiêu que tính?
\shortans{7}
\loigiai{
Cần 7 đoạn que tính.
}
\end{ex}

\begin{ex}
Chữ số 7 trên bảng điện tử cần bao nhiêu que tính?
\shortans{3}
\loigiai{
Cần 3 đoạn que tính.
}
\end{ex}

\begin{ex}
Xếp một hình ngôi nhà gồm 1 hình tam giác làm mái và 1 hình vuông làm thân. Ngôi nhà có tất cả bao nhiêu cạnh bên ngoài bao quanh?
\shortans{5}
\loigiai{
Bao quanh gồm 2 cạnh mái và 3 cạnh thân nhà, tổng là 5 cạnh.
}
\end{ex}

\begin{ex}
Để vẽ đoạn thẳng CD dài 14 cm, điểm C ở vạch 2 cm thì điểm D ở vạch số mấy của thước kẻ?
\shortans{16}
\loigiai{
Điểm D ở vạch: $2 + 14 = 16$ cm.
}
\end{ex}

\newpage

\TaoPhieuBaiTap{Chủ đề 5 -- Làm quen với hình phẳng}{28}{Luyện tập chung}{https://loptoanthaylinh.vn/bai-tap/toan-2/bai-28}
\begin{ex}
Đường gấp khúc ABCD gồm bao nhiêu đoạn thẳng?
\shortans{3}
\loigiai{
Gồm 3 đoạn: AB, BC, CD.
}
\end{ex}

\begin{ex}
Hình tứ giác có bao nhiêu góc?
\shortans{4}
\loigiai{
Hình tứ giác có 4 góc.
}
\end{ex}

\begin{ex}
Tính độ dài đường gấp khúc gồm 3 đoạn thẳng có độ dài là 15 cm, 20 cm, 25 cm (đơn vị cm).
\shortans{60}
\loigiai{
Độ dài: $15 + 20 + 25 = 60$ cm.
}
\end{ex}

\begin{ex}
Tính chu vi hình tam giác có độ dài các cạnh lần lượt là 7 cm, 8 cm, 9 cm (đơn vị cm).
\shortans{24}
\loigiai{
Chu vi: $7 + 8 + 9 = 24$ cm.
}
\end{ex}

\begin{ex}
Tính chu vi hình tứ giác có độ dài các cạnh lần lượt là 10 cm, 12 cm, 14 cm, 16 cm (đơn vị cm).
\shortans{52}
\loigiai{
Chu vi: $10 + 12 + 14 + 16 = 52$ cm.
}
\end{ex}

\begin{ex}
Đoạn thẳng AB dài 18 cm. Điểm M là điểm nằm giữa A và B sao cho AM dài 8 cm. Đoạn MB dài bao nhiêu xăng-ti-mét?
\shortans{10}
\loigiai{
Đoạn MB dài: $18 - 8 = 10$ cm.
}
\end{ex}

\begin{ex}
Ba điểm thẳng hàng A, B, C có AB dài 6 cm, BC dài 9 cm (B nằm giữa A và C). Đoạn AC dài bao nhiêu xăng-ti-mét?
\shortans{15}
\loigiai{
Đoạn AC dài: $6 + 9 = 15$ cm.
}
\end{ex}

\begin{ex}
Đường gấp khúc gồm 4 đoạn thẳng bằng nhau, mỗi đoạn dài 6 cm. Độ dài đường gấp khúc là bao nhiêu xăng-ti-mét?
\shortans{24}
\loigiai{
Độ dài: $6 \times 4 = 24$ cm.
}
\end{ex}

\begin{ex}
Một hình vuông có cạnh dài 7 cm. Chu vi hình vuông đó là bao nhiêu xăng-ti-mét?
\shortans{28}
\loigiai{
Chu vi: $7 \times 4 = 28$ cm.
}
\end{ex}

\begin{ex}
Một hình chữ nhật có chiều dài 9 cm, chiều rộng 5 cm. Nửa chu vi hình chữ nhật là bao nhiêu xăng-ti-mét?
\shortans{14}
\loigiai{
Nửa chu vi: $9 + 5 = 14$ cm.
}
\end{ex}

\begin{ex}
Chu vi hình chữ nhật đó là bao nhiêu xăng-ti-mét?
\shortans{28}
\loigiai{
Chu vi: $14 \times 2 = 28$ cm.
}
\end{ex}

\begin{ex}
Cho 4 điểm không có 3 điểm nào thẳng hàng. Nối tất cả các cặp điểm được bao nhiêu đoạn thẳng?
\shortans{6}
\loigiai{
Có 6 đoạn thẳng.
}
\end{ex}

\begin{ex}
Một sợi dây kẽm dài 50 cm uốn thành một hình tứ giác có chu vi 42 cm. Đoạn dây thừa dài bao nhiêu xăng-ti-mét?
\shortans{8}
\loigiai{
Đoạn dây thừa: $50 - 42 = 8$ cm.
}
\end{ex}

\begin{ex}
Một đường gấp khúc dài 45 cm gồm 3 đoạn thẳng. Biết đoạn thứ nhất dài 12 cm, đoạn thứ hai dài 18 cm. Đoạn thứ ba dài bao nhiêu xăng-ti-mét?
\shortans{15}
\loigiai{
Đoạn thứ ba dài: $45 - 12 - 18 = 15$ cm.
}
\end{ex}

\begin{ex}
Có bao nhiêu đoạn thẳng trong hình gồm 1 tam giác và một đường cao kẻ từ đỉnh xuống đáy?
\shortans{6}
\loigiai{
Có 3 cạnh tam giác, 1 đường cao, và 2 đoạn nhỏ ở đáy, tổng là 6 đoạn thẳng.
}
\end{ex}

\begin{ex}
Một tam giác có chu vi 30 cm, biết hai cạnh dài 10 cm và 12 cm. Cạnh còn lại dài bao nhiêu xăng-ti-mét?
\shortans{8}
\loigiai{
Cạnh còn lại: $30 - 10 - 12 = 8$ cm.
}
\end{ex}

\begin{ex}
Hình tứ giác có chu vi 40 cm, có 4 cạnh bằng nhau. Mỗi cạnh dài bao nhiêu xăng-ti-mét?
\shortans{10}
\loigiai{
Mỗi cạnh dài: $40 : 4 = 10$ cm.
}
\end{ex}

\begin{ex}
Một con kiến bò quanh mép một chiếc điện thoại hình chữ nhật có chiều dài 15 cm, chiều rộng 8 cm. Quãng đường kiến bò một vòng là bao nhiêu xăng-ti-mét?
\shortans{46}
\loigiai{
Quãng đường là chu vi: $(15 + 8) \times 2 = 46$ cm.
}
\end{ex}

\begin{ex}
Cần bao nhiêu que tính để xếp được 5 hình tam giác rời nhau?
\shortans{15}
\loigiai{
Cần: $3 \times 5 = 15$ que tính.
}
\end{ex}

\begin{ex}
Cần bao nhiêu que tính để xếp được 5 hình vuông rời nhau?
\shortans{20}
\loigiai{
Cần: $4 \times 5 = 20$ que tính.
}
\end{ex}

\begin{ex}
Một đoạn thẳng dài 24 cm được chia làm 3 đoạn bằng nhau. Mỗi đoạn dài bao nhiêu xăng-ti-mét?
\shortans{8}
\loigiai{
Mỗi đoạn dài: $24 : 3 = 8$ cm.
}
\end{ex}

\begin{ex}
Một đoạn thẳng dài 30 cm được chia làm 5 đoạn bằng nhau. Mỗi đoạn dài bao nhiêu xăng-ti-mét?
\shortans{6}
\loigiai{
Mỗi đoạn dài: $30 : 5 = 6$ cm.
}
\end{ex}

\begin{ex}
Hai hình vuông bằng nhau cạnh 4 cm ghép sát cạnh tạo thành hình chữ nhật. Chu vi hình chữ nhật đó là bao nhiêu xăng-ti-mét?
\shortans{24}
\loigiai{
Chiều dài 8 cm, chiều rộng 4 cm. Chu vi: $(8 + 4) \times 2 = 24$ cm.
}
\end{ex}

\begin{ex}
Một hình tứ giác có các cạnh dài 8 cm, 10 cm, 12 cm, 15 cm. Cạnh dài nhất hơn cạnh ngắn nhất bao nhiêu xăng-ti-mét?
\shortans{7}
\loigiai{
Hơn: $15 - 8 = 7$ cm.
}
\end{ex}

\begin{ex}
Số đỉnh của một hình ngũ giác (hình 5 cạnh) là bao nhiêu?
\shortans{5}
\loigiai{
Hình ngũ giác có 5 đỉnh.
}
\end{ex}

\newpage


% =========================================================================
% Ngày -- giờ, giờ -- phút, ngày -- tháng
% =========================================================================
\section{Ngày -- giờ, giờ -- phút, ngày -- tháng}

\TaoPhieuBaiTap{Chủ đề 6 -- Ngày -- giờ, giờ -- phút, ngày -- tháng}{29}{Ngày - Giờ}{https://loptoanthaylinh.vn/bai-tap/toan-2/bai-29}
\begin{ex}
Một ngày có bao nhiêu giờ?
\shortans{24}
\loigiai{
Một ngày có 24 giờ.
}
\end{ex}

\begin{ex}
Một giờ có bao nhiêu phút?
\shortans{60}
\loigiai{
Một giờ có 60 phút.
}
\end{ex}

\begin{ex}
Đồng hồ chỉ kim ngắn ở giữa số 2 và số 3, kim dài chỉ số 6. Đồng hồ chỉ 2 giờ bao nhiêu phút?
\shortans{30}
\loigiai{
Kim dài chỉ số 6 là 30 phút (hoặc 2 giờ rưỡi).
}
\end{ex}

\begin{ex}
Lúc 8 giờ 30 phút, kim dài của đồng hồ chỉ vào số mấy?
\shortans{6}
\loigiai{
Kim dài chỉ số 6.
}
\end{ex}

\begin{ex}
Lúc 9 giờ đúng, kim dài của đồng hồ chỉ vào số mấy?
\shortans{12}
\loigiai{
Kim dài chỉ số 12.
}
\end{ex}

\begin{ex}
Lúc 4 giờ 15 phút, kim dài chỉ vào số mấy trên mặt đồng hồ?
\shortans{3}
\loigiai{
Mỗi khoảng là 5 phút, 15 phút kim dài chỉ số 3.
}
\end{ex}

\begin{ex}
Từ 7 giờ sáng đến 10 giờ sáng cùng ngày có bao nhiêu giờ?
\shortans{3}
\loigiai{
Thời gian: $10 - 7 = 3$ giờ.
}
\end{ex}

\begin{ex}
Từ 2 giờ chiều đến 5 giờ chiều cùng ngày có bao nhiêu giờ?
\shortans{3}
\loigiai{
Thời gian: $5 - 2 = 3$ giờ.
}
\end{ex}

\begin{ex}
2 giờ chiều còn gọi là bao nhiêu giờ trong hệ 24 giờ?
\shortans{14}
\loigiai{
2 giờ chiều là $12 + 2 = 14$ giờ.
}
\end{ex}

\begin{ex}
7 giờ tối còn gọi là bao nhiêu giờ trong hệ 24 giờ?
\shortans{19}
\loigiai{
7 giờ tối là $12 + 7 = 19$ giờ.
}
\end{ex}

\begin{ex}
9 giờ đêm còn gọi là bao nhiêu giờ trong hệ 24 giờ?
\shortans{21}
\loigiai{
9 giờ đêm là $12 + 9 = 21$ giờ.
}
\end{ex}

\begin{ex}
12 giờ đêm (kết thúc một ngày) tương ứng với bao nhiêu giờ?
\shortans{24}
\loigiai{
Tương ứng với 24 giờ (hoặc 0 giờ ngày mới).
}
\end{ex}

\begin{ex}
Bạn Lan học bài từ 8 giờ tối đến 9 giờ tối. Bạn Lan học trong bao nhiêu giờ?
\shortans{1}
\loigiai{
Thời gian học: $9 - 8 = 1$ giờ.
}
\end{ex}

\begin{ex}
Một tiết học kéo dài 40 phút, giờ ra chơi kéo dài 20 phút. Tổng thời gian một tiết học và ra chơi là bao nhiêu phút?
\shortans{60}
\loigiai{
Tổng thời gian: $40 + 20 = 60$ phút (đúng 1 giờ).
}
\end{ex}

\begin{ex}
Mẹ đi chợ từ 7 giờ sáng đến 8 giờ sáng. Mẹ đi chợ hết bao nhiêu phút?
\shortans{60}
\loigiai{
Đi hết 1 giờ, tức 60 phút.
}
\end{ex}

\begin{ex}
Bạn Nam ngủ từ 9 giờ tối (21 giờ) đến 6 giờ sáng hôm sau. Nam đã ngủ bao nhiêu giờ?
\shortans{9}
\loigiai{
Từ 21 giờ đến 24 giờ là 3 giờ, từ 0 giờ đến 6 giờ là 6 giờ. Tổng cộng: $3 + 6 = 9$ giờ.
}
\end{ex}

\begin{ex}
Một trận bóng đá mini kéo dài 40 phút gồm 2 hiệp bằng nhau. Mỗi hiệp kéo dài bao nhiêu phút?
\shortans{20}
\loigiai{
Mỗi hiệp kéo dài: $40 : 2 = 20$ phút.
}
\end{ex}

\begin{ex}
Đồng hồ đang chỉ 3 giờ. Sau 30 phút nữa đồng hồ chỉ 3 giờ bao nhiêu phút?
\shortans{30}
\loigiai{
Chỉ 3 giờ 30 phút.
}
\end{ex}

\begin{ex}
Đồng hồ đang chỉ 10 giờ 30 phút. Kim ngắn nằm ở giữa hai số nào?
\shortans{10}
\loigiai{
Nằm giữa số 10 và số 11. Số đầu tiên là 10.
}
\end{ex}

\begin{ex}
Bác thợ làm việc ca sáng từ 7 giờ đến 11 giờ. Bác thợ làm việc trong bao nhiêu giờ?
\shortans{4}
\loigiai{
Thời gian làm việc: $11 - 7 = 4$ giờ.
}
\end{ex}

\begin{ex}
Buổi chiều bắt đầu từ 13 giờ đến 18 giờ. Buổi chiều kéo dài bao nhiêu giờ?
\shortans{5}
\loigiai{
Thời gian: $18 - 13 = 5$ giờ.
}
\end{ex}

\begin{ex}
Hiện tại là 15 giờ. Sau 2 giờ nữa là mấy giờ?
\shortans{17}
\loigiai{
Sau 2 giờ là: $15 + 2 = 17$ giờ (5 giờ chiều).
}
\end{ex}

\begin{ex}
Hiện tại là 20 giờ. Trước đó 3 giờ là mấy giờ?
\shortans{17}
\loigiai{
Trước đó là: $20 - 3 = 17$ giờ (5 giờ chiều).
}
\end{ex}

\begin{ex}
Kim phút quay được nửa vòng đồng hồ (từ số 12 đến số 6) là bao nhiêu phút?
\shortans{30}
\loigiai{
Nửa vòng là 30 phút.
}
\end{ex}

\begin{ex}
Trong 2 giờ có tất cả bao nhiêu phút?
\shortans{120}
\loigiai{
Trong 2 giờ có: $60 + 60 = 120$ phút.
}
\end{ex}

\newpage

\TaoPhieuBaiTap{Chủ đề 6 -- Ngày -- giờ, giờ -- phút, ngày -- tháng}{30}{Giờ - Phút}{https://loptoanthaylinh.vn/bai-tap/toan-2/bai-30}
\begin{ex}
Một năm có bao nhiêu tháng?
\shortans{12}
\loigiai{
Một năm có 12 tháng.
}
\end{ex}

\begin{ex}
Tháng 1 có bao nhiêu ngày?
\shortans{31}
\loigiai{
Tháng 1 có 31 ngày.
}
\end{ex}

\begin{ex}
Tháng 4 có bao nhiêu ngày?
\shortans{30}
\loigiai{
Tháng 4 có 30 ngày.
}
\end{ex}

\begin{ex}
Tháng 6 có bao nhiêu ngày?
\shortans{30}
\loigiai{
Tháng 6 có 30 ngày.
}
\end{ex}

\begin{ex}
Tháng 7 có bao nhiêu ngày?
\shortans{31}
\loigiai{
Tháng 7 có 31 ngày.
}
\end{ex}

\begin{ex}
Tháng 8 có bao nhiêu ngày?
\shortans{31}
\loigiai{
Tháng 8 có 31 ngày.
}
\end{ex}

\begin{ex}
Tháng 9 có bao nhiêu ngày?
\shortans{30}
\loigiai{
Tháng 9 có 30 ngày.
}
\end{ex}

\begin{ex}
Tháng 11 có bao nhiêu ngày?
\shortans{30}
\loigiai{
Tháng 11 có 30 ngày.
}
\end{ex}

\begin{ex}
Tháng 12 có bao nhiêu ngày?
\shortans{31}
\loigiai{
Tháng 12 có 31 ngày.
}
\end{ex}

\begin{ex}
Trong một năm có bao nhiêu tháng có 31 ngày?
\shortans{7}
\loigiai{
Có 7 tháng: Tháng 1, 3, 5, 7, 8, 10, 12.
}
\end{ex}

\begin{ex}
Trong một năm có bao nhiêu tháng có 30 ngày? (Không tính tháng 2).
\shortans{4}
\loigiai{
Có 4 tháng: Tháng 4, 6, 9, 11.
}
\end{ex}

\begin{ex}
Tháng 2 năm nhuận có bao nhiêu ngày?
\shortans{29}
\loigiai{
Năm nhuận tháng 2 có 29 ngày.
}
\end{ex}

\begin{ex}
Tháng 2 năm không nhuận có bao nhiêu ngày?
\shortans{28}
\loigiai{
Năm thường tháng 2 có 28 ngày.
}
\end{ex}

\begin{ex}
Một tuần lễ có bao nhiêu ngày?
\shortans{7}
\loigiai{
Một tuần lễ có 7 ngày.
}
\end{ex}

\begin{ex}
Ngày Quốc khánh nước ta là ngày 2 tháng 9. Ngày này thuộc tháng thứ mấy trong năm?
\shortans{9}
\loigiai{
Thuộc tháng 9.
}
\end{ex}

\begin{ex}
Ngày Quốc tế Thiếu nhi là ngày 1 tháng 6. Ngày này thuộc tháng thứ mấy trong năm?
\shortans{6}
\loigiai{
Thuộc tháng 6.
}
\end{ex}

\begin{ex}
Ngày Nhà giáo Việt Nam là ngày 20 tháng 11. Tháng 11 có bao nhiêu ngày?
\shortans{30}
\loigiai{
Tháng 11 có 30 ngày.
}
\end{ex}

\begin{ex}
Nếu ngày 15 tháng 3 là Thứ Sáu, thì ngày 22 tháng 3 (sau đúng 1 tuần) là Thứ mấy?
\shortans{6}
\loigiai{
Sau đúng 7 ngày vẫn là Thứ Sáu.
}
\end{ex}

\begin{ex}
Nếu ngày 8 tháng 10 là Thứ Ba, thì ngày 1 tháng 10 (trước 1 tuần) là Thứ mấy?
\shortans{3}
\loigiai{
Trước 7 ngày vẫn là Thứ Ba.
}
\end{ex}

\begin{ex}
Nếu hôm nay là ngày 28 tháng 4 thì còn bao nhiêu ngày nữa là hết tháng 4?
\shortans{2}
\loigiai{
Tháng 4 có 30 ngày, còn: $30 - 28 = 2$ ngày.
}
\end{ex}

\begin{ex}
Nếu hôm nay là ngày 29 tháng 5 thì còn bao nhiêu ngày nữa là hết tháng 5?
\shortans{2}
\loigiai{
Tháng 5 có 31 ngày, còn: $31 - 29 = 2$ ngày.
}
\end{ex}

\begin{ex}
Bác nông dân đi nghỉ mát từ ngày 1 tháng 7 đến hết ngày 10 tháng 7. Bác đi nghỉ trong bao nhiêu ngày?
\shortans{10}
\loigiai{
Từ ngày 1 đến hết ngày 10 là 10 ngày.
}
\end{ex}

\begin{ex}
Trong 2 tuần lễ có bao nhiêu ngày?
\shortans{14}
\loigiai{
Hai tuần có: $7 \times 2 = 14$ ngày.
}
\end{ex}

\begin{ex}
Trong 4 tuần lễ có bao nhiêu ngày?
\shortans{28}
\loigiai{
Bốn tuần có: $7 \times 4 = 28$ ngày.
}
\end{ex}

\begin{ex}
Tháng cuối cùng của một năm dương lịch là tháng mấy?
\shortans{12}
\loigiai{
Là tháng 12.
}
\end{ex}

\newpage

\TaoPhieuBaiTap{Chủ đề 6 -- Ngày -- giờ, giờ -- phút, ngày -- tháng}{31}{Ngày - Tháng}{https://loptoanthaylinh.vn/bai-tap/toan-2/bai-31}
\begin{ex}
Đồng hồ chỉ 8 giờ 15 phút. Kim dài chỉ vào số mấy?
\shortans{3}
\loigiai{
Kim dài chỉ số 3.
}
\end{ex}

\begin{ex}
Đồng hồ chỉ 10 giờ 30 phút. Kim dài chỉ vào số mấy?
\shortans{6}
\loigiai{
Kim dài chỉ số 6.
}
\end{ex}

\begin{ex}
Đồng hồ chỉ 12 giờ đúng. Hai kim đồng hồ cùng chỉ vào số mấy?
\shortans{12}
\loigiai{
Cùng chỉ vào số 12.
}
\end{ex}

\begin{ex}
Lớp học bắt đầu lúc 7 giờ 30 phút và tan học lúc 10 giờ 30 phút. Buổi học kéo dài bao nhiêu giờ?
\shortans{3}
\loigiai{
Kéo dài: $10 - 7 = 3$ giờ.
}
\end{ex}

\begin{ex}
Bé Nam đi ngủ lúc 21 giờ và dậy lúc 6 giờ sáng. Nam đã ngủ bao nhiêu giờ?
\shortans{9}
\loigiai{
Ngủ 9 giờ.
}
\end{ex}

\begin{ex}
Tờ lịch tháng 10 có ngày 20 là Thứ Bảy. Hỏi ngày 21 tháng 10 là ngày gì trong tuần? (Chủ Nhật). Ngày 21 cách ngày 20 bao nhiêu ngày?
\shortans{1}
\loigiai{
Cách nhau 1 ngày.
}
\end{ex}

\begin{ex}
Nếu ngày 5 của tháng là Thứ Hai, thì các ngày Thứ Hai tiếp theo của tháng đó là ngày 12, ngày 19 và ngày bao nhiêu?
\shortans{26}
\loigiai{
Ngày tiếp theo: $19 + 7 = 26$.
}
\end{ex}

\begin{ex}
Tháng đó có tất cả bao nhiêu ngày Thứ Hai (ngày 5, 12, 19, 26)?
\shortans{4}
\loigiai{
Có 4 ngày Thứ Hai.
}
\end{ex}

\begin{ex}
Một tháng có 5 ngày Chủ Nhật, ngày Chủ Nhật đầu tiên là ngày 2. Ngày Chủ Nhật cuối cùng của tháng đó là ngày bao nhiêu?
\shortans{30}
\loigiai{
Các ngày Chủ Nhật: 2, 9, 16, 23, 30. Ngày cuối là ngày 30.
}
\end{ex}

\begin{ex}
Chuyến tàu xuất phát lúc 6 giờ sáng và đến ga lúc 2 giờ chiều (14 giờ). Chuyến tàu đi hết bao nhiêu giờ?
\shortans{8}
\loigiai{
Hết: $14 - 6 = 8$ giờ.
}
\end{ex}

\begin{ex}
Một buổi chiếu phim kéo dài 1 giờ 30 phút. Đổi 1 giờ 30 phút ra phút là bao nhiêu phút?
\shortans{90}
\loigiai{
Ta có: $60 + 30 = 90$ phút.
}
\end{ex}

\begin{ex}
Một tiết học kéo dài 35 phút. Sau 2 tiết học như thế, học sinh đã học tất cả bao nhiêu phút?
\shortans{70}
\loigiai{
Thời gian: $35 + 35 = 70$ phút.
}
\end{ex}

\begin{ex}
Bố đi công tác trong 1 tuần lễ và 3 ngày. Bố đi công tác bao nhiêu ngày?
\shortans{10}
\loigiai{
Đi: $7 + 3 = 10$ ngày.
}
\end{ex}

\begin{ex}
Mẹ đi du lịch trong 2 tuần lễ và 5 ngày. Mẹ đi du lịch tất cả bao nhiêu ngày?
\shortans{19}
\loigiai{
Đi: $14 + 5 = 19$ ngày.
}
\end{ex}

\begin{ex}
Nếu ngày 30 tháng 4 là Thứ Ba, thì ngày 1 tháng 5 tiếp theo là Thứ mấy?
\shortans{4}
\loigiai{
Sau Thứ Ba là Thứ Tư.
}
\end{ex}

\begin{ex}
Nếu ngày 31 tháng 12 là Thứ Sáu, thì ngày 1 tháng 1 năm mới là Thứ mấy?
\shortans{7}
\loigiai{
Sau Thứ Sáu là Thứ Bảy.
}
\end{ex}

\begin{ex}
Đồng hồ điện tử hiển thị 14:30. Đây là mấy giờ rưỡi chiều?
\shortans{2}
\loigiai{
Là 2 giờ rưỡi chiều.
}
\end{ex}

\begin{ex}
Đồng hồ điện tử hiển thị 19:00. Đây là mấy giờ tối?
\shortans{7}
\loigiai{
Là 7 giờ tối.
}
\end{ex}

\begin{ex}
Đồng hồ điện tử hiển thị 22:15. Đây là 10 giờ 15 phút đêm. Số giờ đêm là mấy?
\shortans{10}
\loigiai{
Là 10 giờ đêm.
}
\end{ex}

\begin{ex}
Khoảng thời gian từ 8 giờ sáng đến 12 giờ trưa là bao nhiêu giờ?
\shortans{4}
\loigiai{
Là: $12 - 8 = 4$ giờ.
}
\end{ex}

\begin{ex}
Khoảng thời gian từ 13 giờ đến 17 giờ là bao nhiêu giờ?
\shortans{4}
\loigiai{
Là: $17 - 13 = 4$ giờ.
}
\end{ex}

\begin{ex}
Tháng 8 có 31 ngày. Nếu hôm nay là ngày 25 tháng 8 thì còn bao nhiêu ngày nữa là hết tháng 8?
\shortans{6}
\loigiai{
Còn: $31 - 25 = 6$ ngày.
}
\end{ex}

\begin{ex}
Nếu hôm nay là ngày 10, hôm qua là ngày mấy?
\shortans{9}
\loigiai{
Hôm qua là ngày 9.
}
\end{ex}

\begin{ex}
Nếu hôm nay là ngày 10, ngày kia là ngày mấy?
\shortans{12}
\loigiai{
Ngày kia là ngày 12.
}
\end{ex}

\begin{ex}
Một ngày bác sĩ trực bệnh viện 24 giờ. Bác sĩ trực 2 ngày liên tiếp là bao nhiêu giờ?
\shortans{48}
\loigiai{
Hai ngày là: $24 + 24 = 48$ giờ.
}
\end{ex}

\newpage

\TaoPhieuBaiTap{Chủ đề 6 -- Ngày -- giờ, giờ -- phút, ngày -- tháng}{32}{Luyện tập chung thời gian}{https://loptoanthaylinh.vn/bai-tap/toan-2/bai-32}
\begin{ex}
Một tuần lễ có bao nhiêu ngày?
\shortans{7}
\loigiai{
Có 7 ngày.
}
\end{ex}

\begin{ex}
Một năm có bao nhiêu tháng?
\shortans{12}
\loigiai{
Có 12 tháng.
}
\end{ex}

\begin{ex}
Một ngày có bao nhiêu giờ?
\shortans{24}
\loigiai{
Có 24 giờ.
}
\end{ex}

\begin{ex}
Một giờ có bao nhiêu phút?
\shortans{60}
\loigiai{
Có 60 phút.
}
\end{ex}

\begin{ex}
Lúc 6 giờ 30 phút sáng, kim dài chỉ vào số mấy?
\shortans{6}
\loigiai{
Kim dài chỉ số 6.
}
\end{ex}

\begin{ex}
Lúc 3 giờ 15 phút chiều, kim dài chỉ vào số mấy?
\shortans{3}
\loigiai{
Kim dài chỉ số 3.
}
\end{ex}

\begin{ex}
Lúc 9 giờ đúng tối, kim ngắn chỉ vào số mấy?
\shortans{9}
\loigiai{
Kim ngắn chỉ số 9.
}
\end{ex}

\begin{ex}
Tháng 10 có bao nhiêu ngày?
\shortans{31}
\loigiai{
Tháng 10 có 31 ngày.
}
\end{ex}

\begin{ex}
Tháng 9 có bao nhiêu ngày?
\shortans{30}
\loigiai{
Tháng 9 có 30 ngày.
}
\end{ex}

\begin{ex}
Tháng 12 có bao nhiêu ngày?
\shortans{31}
\loigiai{
Tháng 12 có 31 ngày.
}
\end{ex}

\begin{ex}
Nếu ngày 12 là Thứ Tư, thì ngày 19 cùng tháng đó là Thứ mấy?
\shortans{4}
\loigiai{
Cách 7 ngày vẫn là Thứ Tư.
}
\end{ex}

\begin{ex}
Nếu ngày 25 là Chủ Nhật, thì ngày 18 cùng tháng đó (trước 1 tuần) là ngày gì? (Chủ Nhật). Khoảng cách giữa 2 ngày là bao nhiêu ngày?
\shortans{7}
\loigiai{
Khoảng cách là 7 ngày.
}
\end{ex}

\begin{ex}
Hôm nay là Thứ Sáu ngày 10. Thứ Sáu tuần sau là ngày bao nhiêu?
\shortans{17}
\loigiai{
Ngày đó là: $10 + 7 = 17$.
}
\end{ex}

\begin{ex}
Hôm nay là Thứ Hai ngày 20. Thứ Hai tuần trước là ngày bao nhiêu?
\shortans{13}
\loigiai{
Ngày đó là: $20 - 7 = 13$.
}
\end{ex}

\begin{ex}
15 giờ bằng mấy giờ chiều?
\shortans{3}
\loigiai{
15 giờ là 3 giờ chiều.
}
\end{ex}

\begin{ex}
20 giờ bằng mấy giờ tối?
\shortans{8}
\loigiai{
20 giờ là 8 giờ tối.
}
\end{ex}

\begin{ex}
17 giờ bằng mấy giờ chiều?
\shortans{5}
\loigiai{
17 giờ là 5 giờ chiều.
}
\end{ex}

\begin{ex}
Từ 9 giờ sáng đến 2 giờ chiều (14 giờ) có bao nhiêu giờ?
\shortans{5}
\loigiai{
Có: $14 - 9 = 5$ giờ.
}
\end{ex}

\begin{ex}
Bác bảo vệ trực từ 22 giờ đêm đến 6 giờ sáng hôm sau. Bác trực trong bao nhiêu giờ?
\shortans{8}
\loigiai{
Từ 22h đến 24h là 2h, từ 0h đến 6h là 6h. Tổng cộng: $2 + 6 = 8$ giờ.
}
\end{ex}

\begin{ex}
Một trận bóng đá bắt đầu lúc 15 giờ và kết thúc lúc 17 giờ. Trận đấu kéo dài bao nhiêu giờ?
\shortans{2}
\loigiai{
Kéo dài: $17 - 15 = 2$ giờ.
}
\end{ex}

\begin{ex}
Tháng 4 có 30 ngày. Nếu đã qua 20 ngày thì tháng 4 còn lại bao nhiêu ngày?
\shortans{10}
\loigiai{
Còn lại: $30 - 20 = 10$ ngày.
}
\end{ex}

\begin{ex}
Trong 3 tuần lễ có tất cả bao nhiêu ngày?
\shortans{21}
\loigiai{
Ba tuần có: $7 \times 3 = 21$ ngày.
}
\end{ex}

\begin{ex}
Kim phút quay được 1 vòng và thêm nửa vòng nữa là bao nhiêu phút?
\shortans{90}
\loigiai{
Ta có: $60 + 30 = 90$ phút.
}
\end{ex}

\begin{ex}
Thời gian nghỉ hè từ ngày 1 tháng 6 đến hết ngày 31 tháng 7. Tháng 6 có 30 ngày, tháng 7 có 31 ngày. Kỳ nghỉ hè kéo dài bao nhiêu ngày?
\shortans{61}
\loigiai{
Kéo dài: $30 + 31 = 61$ ngày.
}
\end{ex}

\begin{ex}
Một con tàu đi trên biển trong 1 tuần lễ và 5 ngày. Con tàu đã đi tất cả bao nhiêu ngày?
\shortans{12}
\loigiai{
Tất cả: $7 + 5 = 12$ ngày.
}
\end{ex}

\newpage


% =========================================================================
% Ôn tập học kì 1
% =========================================================================
\section{Ôn tập học kì 1}

\TaoPhieuBaiTap{Chủ đề 7 -- Ôn tập học kì 1}{33}{Ôn tập các số trong phạm vi 100}{https://loptoanthaylinh.vn/bai-tap/toan-2/bai-33}
\begin{ex}
Tính: $9 + 7$.
\shortans{16}
\loigiai{
Ta có $9 + 7 = 16$.
}
\end{ex}

\begin{ex}
Tính: $15 - 8$.
\shortans{7}
\loigiai{
Ta có $15 - 8 = 7$.
}
\end{ex}

\begin{ex}
Tính: $8 + 6$.
\shortans{14}
\loigiai{
Ta có $8 + 6 = 14$.
}
\end{ex}

\begin{ex}
Tính: $13 - 9$.
\shortans{4}
\loigiai{
Ta có $13 - 9 = 4$.
}
\end{ex}

\begin{ex}
Tính: $35 + 28$.
\shortans{63}
\loigiai{
Ta có $35 + 28 = 63$.
}
\end{ex}

\begin{ex}
Tính: $72 - 36$.
\shortans{36}
\loigiai{
Ta có $72 - 36 = 36$.
}
\end{ex}

\begin{ex}
Tính: $47 + 35$.
\shortans{82}
\loigiai{
Ta có $47 + 35 = 82$.
}
\end{ex}

\begin{ex}
Tính: $91 - 54$.
\shortans{37}
\loigiai{
Ta có $91 - 54 = 37$.
}
\end{ex}

\begin{ex}
Tính: $60 - 27$.
\shortans{33}
\loigiai{
Ta có $60 - 27 = 33$.
}
\end{ex}

\begin{ex}
Tính: $100 - 45$.
\shortans{55}
\loigiai{
Ta có $100 - 45 = 55$.
}
\end{ex}

\begin{ex}
Tính giá trị biểu thức: $38 + 25 - 19$.
\shortans{44}
\loigiai{
Ta có $38 + 25 = 63$, rồi $63 - 19 = 44$.
}
\end{ex}

\begin{ex}
Tính giá trị biểu thức: $82 - 47 + 29$.
\shortans{64}
\loigiai{
Ta có $82 - 47 = 35$, rồi $35 + 29 = 64$.
}
\end{ex}

\begin{ex}
Điền số thích hợp: $28 + \dots = 61$.
\shortans{33}
\loigiai{
Số cần điền: $61 - 28 = 33$.
}
\end{ex}

\begin{ex}
Điền số thích hợp: $75 - \dots = 38$.
\shortans{37}
\loigiai{
Số cần điền: $75 - 38 = 37$.
}
\end{ex}

\begin{ex}
Điền số thích hợp: $\dots + 36 = 82$.
\shortans{46}
\loigiai{
Số cần điền: $82 - 36 = 46$.
}
\end{ex}

\begin{ex}
Điền số thích hợp: $\dots - 29 = 45$.
\shortans{74}
\loigiai{
Số cần điền: $45 + 29 = 74$.
}
\end{ex}

\begin{ex}
Lớp 2A có 19 bạn nam và 18 bạn nữ. Hỏi lớp 2A có tất cả bao nhiêu bạn học sinh?
\shortans{37}
\loigiai{
Cả lớp có: $19 + 18 = 37$ bạn.
}
\end{ex}

\begin{ex}
Bác nông dân có 64 con gà, bác bán bớt 28 con gà. Bác còn lại bao nhiêu con gà?
\shortans{36}
\loigiai{
Còn lại: $64 - 28 = 36$ con gà.
}
\end{ex}

\begin{ex}
Đoạn dây thứ nhất dài 45 cm, đoạn thứ hai dài hơn đoạn thứ nhất 18 cm. Đoạn dây thứ hai dài bao nhiêu xăng-ti-mét?
\shortans{63}
\loigiai{
Đoạn thứ hai dài: $45 + 18 = 63$ cm.
}
\end{ex}

\begin{ex}
Thùng thứ nhất có 80 lít nước, thùng thứ hai có ít hơn thùng thứ nhất 25 lít nước. Thùng thứ hai có bao nhiêu lít nước?
\shortans{55}
\loigiai{
Thùng thứ hai có: $80 - 25 = 55$ lít.
}
\end{ex}

\begin{ex}
Hai số có tổng bằng 92, số hạng thứ nhất là 47. Số hạng thứ hai là bao nhiêu?
\shortans{45}
\loigiai{
Số hạng thứ hai: $92 - 47 = 45$.
}
\end{ex}

\begin{ex}
Hai số có hiệu bằng 38, số trừ là 25. Số bị trừ là bao nhiêu?
\shortans{63}
\loigiai{
Số bị trừ: $38 + 25 = 63$.
}
\end{ex}

\begin{ex}
Hai số có hiệu bằng 29, số bị trừ là 83. Số trừ là bao nhiêu?
\shortans{54}
\loigiai{
Số trừ: $83 - 29 = 54$.
}
\end{ex}

\begin{ex}
Hiệu của số lớn nhất có hai chữ số và số 35 là bao nhiêu?
\shortans{64}
\loigiai{
Số lớn nhất có 2 chữ số là 99. Hiệu: $99 - 35 = 64$.
}
\end{ex}

\begin{ex}
Tính: $26 + 34 + 15$.
\shortans{75}
\loigiai{
Ta có $26 + 34 = 60$, rồi $60 + 15 = 75$.
}
\end{ex}

\newpage

\TaoPhieuBaiTap{Chủ đề 7 -- Ôn tập học kì 1}{34}{Ôn tập phép cộng, phép trừ trong phạm vi 100}{https://loptoanthaylinh.vn/bai-tap/toan-2/bai-34}
\begin{ex}
Đường gấp khúc ABC gồm bao nhiêu đoạn thẳng?
\shortans{2}
\loigiai{
Gồm 2 đoạn: AB và BC.
}
\end{ex}

\begin{ex}
Đường gấp khúc MNPQ gồm bao nhiêu đoạn thẳng?
\shortans{3}
\loigiai{
Gồm 3 đoạn: MN, NP, PQ.
}
\end{ex}

\begin{ex}
Hình tứ giác có bao nhiêu cạnh?
\shortans{4}
\loigiai{
Có 4 cạnh.
}
\end{ex}

\begin{ex}
Hình tứ giác có bao nhiêu đỉnh?
\shortans{4}
\loigiai{
Có 4 đỉnh.
}
\end{ex}

\begin{ex}
Hình tam giác có bao nhiêu cạnh?
\shortans{3}
\loigiai{
Có 3 cạnh.
}
\end{ex}

\begin{ex}
Tính độ dài đường gấp khúc gồm 3 đoạn thẳng có độ dài lần lượt là 12 cm, 15 cm, 18 cm (đơn vị cm).
\shortans{45}
\loigiai{
Độ dài: $12 + 15 + 18 = 45$ cm.
}
\end{ex}

\begin{ex}
Tính chu vi hình tam giác có các cạnh dài 8 cm, 10 cm, 12 cm (đơn vị cm).
\shortans{30}
\loigiai{
Chu vi: $8 + 10 + 12 = 30$ cm.
}
\end{ex}

\begin{ex}
Tính chu vi hình tứ giác có các cạnh dài 5 cm, 7 cm, 9 cm, 11 cm (đơn vị cm).
\shortans{32}
\loigiai{
Chu vi: $5 + 7 + 9 + 11 = 32$ cm.
}
\end{ex}

\begin{ex}
Một hình vuông có cạnh dài 8 cm. Chu vi hình vuông là bao nhiêu xăng-ti-mét?
\shortans{32}
\loigiai{
Chu vi: $8 \times 4 = 32$ cm.
}
\end{ex}

\begin{ex}
Một hình chữ nhật có chiều dài 10 cm, chiều rộng 6 cm. Nửa chu vi là bao nhiêu xăng-ti-mét?
\shortans{16}
\loigiai{
Nửa chu vi: $10 + 6 = 16$ cm.
}
\end{ex}

\begin{ex}
Chu vi hình chữ nhật đó là bao nhiêu xăng-ti-mét?
\shortans{32}
\loigiai{
Chu vi: $16 \times 2 = 32$ cm.
}
\end{ex}

\begin{ex}
Cho 3 điểm thẳng hàng A, B, C (B nằm giữa A và C). Biết AB dài 7 cm, BC dài 8 cm. Độ dài đoạn AC là bao nhiêu xăng-ti-mét?
\shortans{15}
\loigiai{
Độ dài AC: $7 + 8 = 15$ cm.
}
\end{ex}

\begin{ex}
Đoạn thẳng MN dài 20 cm, điểm P nằm giữa M và N sao cho MP dài 9 cm. Đoạn thẳng PN dài bao nhiêu xăng-ti-mét?
\shortans{11}
\loigiai{
Đoạn PN dài: $20 - 9 = 11$ cm.
}
\end{ex}

\begin{ex}
Cần ít nhất bao nhiêu que tính để xếp được 1 hình tứ giác?
\shortans{4}
\loigiai{
Cần 4 que tính.
}
\end{ex}

\begin{ex}
Cần ít nhất bao nhiêu que tính để xếp được 1 hình tam giác?
\shortans{3}
\loigiai{
Cần 3 que tính.
}
\end{ex}

\begin{ex}
Dùng 9 que tính bằng nhau xếp được nhiều nhất bao nhiêu hình tam giác rời nhau?
\shortans{3}
\loigiai{
Xếp được: $9 : 3 = 3$ hình tam giác.
}
\end{ex}

\begin{ex}
Dùng 12 que tính bằng nhau xếp được nhiều nhất bao nhiêu hình vuông rời nhau?
\shortans{3}
\loigiai{
Xếp được: $12 : 4 = 3$ hình vuông.
}
\end{ex}

\begin{ex}
Một đường gấp khúc dài 36 cm gồm 3 đoạn thẳng bằng nhau. Mỗi đoạn dài bao nhiêu xăng-ti-mét?
\shortans{12}
\loigiai{
Mỗi đoạn dài: $36 : 3 = 12$ cm.
}
\end{ex}

\begin{ex}
Một hình tứ giác có chu vi 48 cm, bốn cạnh có độ dài bằng nhau. Mỗi cạnh dài bao nhiêu xăng-ti-mét?
\shortans{12}
\loigiai{
Mỗi cạnh dài: $48 : 4 = 12$ cm.
}
\end{ex}

\begin{ex}
Cắt một góc của hình tứ giác bằng một nhát kéo thẳng. Hình mới nhận được có bao nhiêu đỉnh?
\shortans{5}
\loigiai{
Hình mới có 5 đỉnh.
}
\end{ex}

\begin{ex}
Gấp đôi một hình vuông theo đường chéo ta được hình gì? Hình tam giác mới có bao nhiêu cạnh?
\shortans{3}
\loigiai{
Hình tam giác có 3 cạnh.
}
\end{ex}

\begin{ex}
Một con sâu bò trên đường gấp khúc gồm 4 đoạn, mỗi đoạn dài 6 cm. Con sâu bò được quãng đường bao nhiêu xăng-ti-mét?
\shortans{24}
\loigiai{
Quãng đường: $6 \times 4 = 24$ cm.
}
\end{ex}

\begin{ex}
Hình chữ nhật có bao nhiêu góc vuông?
\shortans{4}
\loigiai{
Có 4 góc vuông.
}
\end{ex}

\begin{ex}
Hình vuông có bao nhiêu góc vuông?
\shortans{4}
\loigiai{
Có 4 góc vuông.
}
\end{ex}

\begin{ex}
Tổng số cạnh của 2 hình tam giác rời nhau là bao nhiêu?
\shortans{6}
\loigiai{
Tổng số cạnh: $3 + 3 = 6$ cạnh.
}
\end{ex}

\newpage

\TaoPhieuBaiTap{Chủ đề 7 -- Ôn tập học kì 1}{35}{Ôn tập hình học và đo lường}{https://loptoanthaylinh.vn/bai-tap/toan-2/bai-35}
\begin{ex}
Tính: $35\text{ kg} + 28\text{ kg}$. Kết quả là bao nhiêu ki-lô-gam?
\shortans{63}
\loigiai{
Ta có: $35 + 28 = 63$ kg.
}
\end{ex}

\begin{ex}
Tính: $72\text{ kg} - 35\text{ kg}$. Kết quả là bao nhiêu ki-lô-gam?
\shortans{37}
\loigiai{
Ta có: $72 - 35 = 37$ kg.
}
\end{ex}

\begin{ex}
Tính: $46\text{ }l + 27\text{ }l$. Kết quả là bao nhiêu lít?
\shortans{73}
\loigiai{
Ta có: $46 + 27 = 73$ lít.
}
\end{ex}

\begin{ex}
Tính: $81\text{ }l - 45\text{ }l$. Kết quả là bao nhiêu lít?
\shortans{36}
\loigiai{
Ta có: $81 - 45 = 36$ lít.
}
\end{ex}

\begin{ex}
Một ngày có bao nhiêu giờ?
\shortans{24}
\loigiai{
Một ngày có 24 giờ.
}
\end{ex}

\begin{ex}
Một giờ có bao nhiêu phút?
\shortans{60}
\loigiai{
Một giờ có 60 phút.
}
\end{ex}

\begin{ex}
Tháng 5 có bao nhiêu ngày?
\shortans{31}
\loigiai{
Tháng 5 có 31 ngày.
}
\end{ex}

\begin{ex}
Tháng 6 có bao nhiêu ngày?
\shortans{30}
\loigiai{
Tháng 6 có 30 ngày.
}
\end{ex}

\begin{ex}
Một tuần lễ có bao nhiêu ngày?
\shortans{7}
\loigiai{
Có 7 ngày.
}
\end{ex}

\begin{ex}
Bao gạo nặng 45 kg, bao ngô nặng hơn bao gạo 8 kg. Bao ngô nặng bao nhiêu ki-lô-gam?
\shortans{53}
\loigiai{
Bao ngô nặng: $45 + 8 = 53$ kg.
}
\end{ex}

\begin{ex}
Thùng dầu có 65 lít, người ta lấy ra 28 lít. Thùng còn lại bao nhiêu lít dầu?
\shortans{37}
\loigiai{
Còn lại: $65 - 28 = 37$ lít.
}
\end{ex}

\begin{ex}
Đồng hồ chỉ 9 giờ 30 phút. Kim dài chỉ vào số mấy?
\shortans{6}
\loigiai{
Kim dài chỉ số 6.
}
\end{ex}

\begin{ex}
Đồng hồ chỉ 14 giờ. Đây là mấy giờ chiều?
\shortans{2}
\loigiai{
Là 2 giờ chiều.
}
\end{ex}

\begin{ex}
Đoạn dây thứ nhất dài 38 cm, đoạn thứ hai dài 25 cm. Cả hai đoạn dài bao nhiêu xăng-ti-mét?
\shortans{63}
\loigiai{
Dài: $38 + 25 = 63$ cm.
}
\end{ex}

\begin{ex}
Sợi dây đồng dài 90 cm, cắt bớt 35 cm. Sợi dây còn lại dài bao nhiêu xăng-ti-mét?
\shortans{55}
\loigiai{
Còn lại: $90 - 35 = 55$ cm.
}
\end{ex}

\begin{ex}
Bé Nam nặng 26 kg, khi bé bế thêm con mèo nặng 3 kg thì cả bé và mèo nặng bao nhiêu ki-lô-gam?
\shortans{29}
\loigiai{
Nặng: $26 + 3 = 29$ kg.
}
\end{ex}

\begin{ex}
Can thứ nhất có 15 lít nước, can thứ hai có 18 lít nước. Cả hai can có bao nhiêu lít nước?
\shortans{33}
\loigiai{
Có: $15 + 18 = 33$ lít.
}
\end{ex}

\begin{ex}
Nếu Thứ Sáu là ngày 14 thì Thứ Sáu tuần sau là ngày bao nhiêu?
\shortans{21}
\loigiai{
Ngày đó là: $14 + 7 = 21$.
}
\end{ex}

\begin{ex}
Nếu Thứ Ba là ngày 24 thì Thứ Ba tuần trước là ngày bao nhiêu?
\shortans{17}
\loigiai{
Ngày đó là: $24 - 7 = 17$.
}
\end{ex}

\begin{ex}
Tính giá trị biểu thức: $24\text{ kg} + 18\text{ kg} - 15\text{ kg}$.
\shortans{27}
\loigiai{
Ta có $24 + 18 = 42$, rồi $42 - 15 = 27$ kg.
}
\end{ex}

\begin{ex}
Tính giá trị biểu thức: $50\text{ }l - 25\text{ }l + 18\text{ }l$.
\shortans{43}
\loigiai{
Ta có $50 - 25 = 25$, rồi $25 + 18 = 43$ lít.
}
\end{ex}

\begin{ex}
Một lớp học có buổi dã ngoại kéo dài từ 8 giờ sáng đến 3 giờ chiều (15 giờ). Buổi dã ngoại kéo dài bao nhiêu giờ?
\shortans{7}
\loigiai{
Thời gian: $15 - 8 = 7$ giờ.
}
\end{ex}

\begin{ex}
Đĩa cân bên trái có gói đường và quả cân 1 kg, đĩa bên phải có quả cân 5 kg thì cân thăng bằng. Gói đường nặng bao nhiêu ki-lô-gam?
\shortans{4}
\loigiai{
Gói đường nặng: $5 - 1 = 4$ kg.
}
\end{ex}

\begin{ex}
Có 5 can nước, mỗi can chứa 3 lít. Cả 5 can chứa bao nhiêu lít nước?
\shortans{15}
\loigiai{
Chứa: $3 \times 5 = 15$ lít.
}
\end{ex}

\begin{ex}
Một đoạn đường dài 100 mét, công nhân đã sửa được 65 mét. Còn lại bao nhiêu mét đường chưa sửa?
\shortans{35}
\loigiai{
Còn lại: $100 - 65 = 35$ mét.
}
\end{ex}

\newpage

\TaoPhieuBaiTap{Chủ đề 7 -- Ôn tập học kì 1}{36}{Ôn tập chung học kì 1}{https://loptoanthaylinh.vn/bai-tap/toan-2/bai-36}
\begin{ex}
Số liền sau của số 89 là số nào?
\shortans{90}
\loigiai{
Số liền sau 89 là 90.
}
\end{ex}

\begin{ex}
Số liền trước của số 100 là số nào?
\shortans{99}
\loigiai{
Số liền trước 100 là 99.
}
\end{ex}

\begin{ex}
Tính: $48 + 36$.
\shortans{84}
\loigiai{
Ta có $48 + 36 = 84$.
}
\end{ex}

\begin{ex}
Tính: $82 - 45$.
\shortans{37}
\loigiai{
Ta có $82 - 45 = 37$.
}
\end{ex}

\begin{ex}
Tính: $65 + 27$.
\shortans{92}
\loigiai{
Ta có $65 + 27 = 92$.
}
\end{ex}

\begin{ex}
Tính: $94 - 58$.
\shortans{36}
\loigiai{
Ta có $94 - 58 = 36$.
}
\end{ex}

\begin{ex}
Tính giá trị biểu thức: $42 + 29 - 35$.
\shortans{36}
\loigiai{
Ta có $42 + 29 = 71$, rồi $71 - 35 = 36$.
}
\end{ex}

\begin{ex}
Tính giá trị biểu thức: $85 - 38 + 17$.
\shortans{64}
\loigiai{
Ta có $85 - 38 = 47$, rồi $47 + 17 = 64$.
}
\end{ex}

\begin{ex}
Điền số thích hợp: $37 + \dots = 81$.
\shortans{44}
\loigiai{
Số cần điền: $81 - 37 = 44$.
}
\end{ex}

\begin{ex}
Điền số thích hợp: $90 - \dots = 43$.
\shortans{47}
\loigiai{
Số cần điền: $90 - 43 = 47$.
}
\end{ex}

\begin{ex}
Đường gấp khúc ABCD gồm 3 đoạn thẳng có độ dài 10 cm, 15 cm, 20 cm. Độ dài đường gấp khúc là bao nhiêu xăng-ti-mét?
\shortans{45}
\loigiai{
Độ dài: $10 + 15 + 20 = 45$ cm.
}
\end{ex}

\begin{ex}
Chu vi hình tam giác có 3 cạnh bằng nhau, mỗi cạnh dài 8 cm là bao nhiêu xăng-ti-mét?
\shortans{24}
\loigiai{
Chu vi: $8 \times 3 = 24$ cm.
}
\end{ex}

\begin{ex}
Chu vi hình chữ nhật có chiều dài 12 cm, chiều rộng 8 cm là bao nhiêu xăng-ti-mét?
\shortans{40}
\loigiai{
Chu vi: $(12 + 8) \times 2 = 40$ cm.
}
\end{ex}

\begin{ex}
Một thùng cam có 54 quả, thùng quýt có 38 quả. Cả hai thùng có bao nhiêu quả?
\shortans{92}
\loigiai{
Cả hai thùng có: $54 + 38 = 92$ quả.
}
\end{ex}

\begin{ex}
Bác nông dân nuôi 72 con vịt, bác bán đi 35 con. Bác còn lại bao nhiêu con vịt?
\shortans{37}
\loigiai{
Còn lại: $72 - 35 = 37$ con vịt.
}
\end{ex}

\begin{ex}
Một cửa hàng ngày thứ nhất bán được 46 mét vải, ngày thứ hai bán được nhiều hơn ngày thứ nhất 18 mét vải. Ngày thứ hai bán được bao nhiêu mét vải?
\shortans{64}
\loigiai{
Ngày thứ hai bán: $46 + 18 = 64$ mét.
}
\end{ex}

\begin{ex}
Bao gạo nặng 50 kg, bao ngô nhẹ hơn bao gạo 15 kg. Bao ngô nặng bao nhiêu ki-lô-gam?
\shortans{35}
\loigiai{
Bao ngô nặng: $50 - 15 = 35$ kg.
}
\end{ex}

\begin{ex}
Một can 20 lít nước mắm, người ta đã rót ra 8 lít. Trong can còn lại bao nhiêu lít nước mắm?
\shortans{12}
\loigiai{
Còn lại: $20 - 8 = 12$ lít.
}
\end{ex}

\begin{ex}
Một ngày có bao nhiêu giờ?
\shortans{24}
\loigiai{
Có 24 giờ.
}
\end{ex}

\begin{ex}
Tháng 12 có bao nhiêu ngày?
\shortans{31}
\loigiai{
Có 31 ngày.
}
\end{ex}

\begin{ex}
Đồng hồ chỉ 16 giờ. Đây là mấy giờ chiều?
\shortans{4}
\loigiai{
Là 4 giờ chiều.
}
\end{ex}

\begin{ex}
Nếu Thứ Năm là ngày 15 thì Thứ Năm tuần trước là ngày bao nhiêu?
\shortans{8}
\loigiai{
Ngày đó là: $15 - 7 = 8$.
}
\end{ex}

\begin{ex}
Hiệu của số lớn nhất có hai chữ số và số bé nhất có hai chữ số là bao nhiêu?
\shortans{89}
\loigiai{
Số lớn nhất là 99, số bé nhất là 10. Hiệu: $99 - 10 = 89$.
}
\end{ex}

\begin{ex}
Tìm số trừ biết số bị trừ là 95 và hiệu là 48.
\shortans{47}
\loigiai{
Số trừ: $95 - 48 = 47$.
}
\end{ex}

\begin{ex}
Tìm số bị trừ biết số trừ là 36 và hiệu là 47.
\shortans{83}
\loigiai{
Số bị trừ: $47 + 36 = 83$.
}
\end{ex}

\newpage


% =========================================================================
% Phép nhân, phép chia
% =========================================================================
\section{Phép nhân, phép chia}

\TaoPhieuBaiTap{Chủ đề 8 -- Phép nhân, phép chia}{37}{Phép nhân. Thừa số - Tích}{https://loptoanthaylinh.vn/bai-tap/toan-2/bai-37}
\begin{ex}
Chuyển tổng $2 + 2 + 2 + 2$ thành tích $2 \times 4$. Tính kết quả của phép nhân đó.
\shortans{8}
\loigiai{
Ta có: $2 \times 4 = 8$.
}
\end{ex}

\begin{ex}
Chuyển tổng $5 + 5 + 5$ thành tích $5 \times 3$. Tính kết quả của phép nhân đó.
\shortans{15}
\loigiai{
Ta có: $5 \times 3 = 15$.
}
\end{ex}

\begin{ex}
Chuyển tổng $3 + 3 + 3 + 3 + 3$ thành tích. Phép tính có kết quả bằng bao nhiêu?
\shortans{15}
\loigiai{
Ta có: $3 \times 5 = 15$.
}
\end{ex}

\begin{ex}
Chuyển tổng $4 + 4 + 4$ thành tích. Phép tính có kết quả bằng bao nhiêu?
\shortans{12}
\loigiai{
Ta có: $4 \times 3 = 12$.
}
\end{ex}

\begin{ex}
Chuyển tổng $6 + 6$ thành tích. Phép tính có kết quả bằng bao nhiêu?
\shortans{12}
\loigiai{
Ta có: $6 \times 2 = 12$.
}
\end{ex}

\begin{ex}
Số 7 được lấy 3 lần viết thành phép nhân là $7 \times 3$. Tính kết quả.
\shortans{21}
\loigiai{
Ta có: $7 \times 3 = 21$.
}
\end{ex}

\begin{ex}
Số 8 được lấy 2 lần viết thành phép nhân là $8 \times 2$. Tính kết quả.
\shortans{16}
\loigiai{
Ta có: $8 \times 2 = 16$.
}
\end{ex}

\begin{ex}
Số 10 được lấy 4 lần viết thành phép nhân là $10 \times 4$. Tính kết quả.
\shortans{40}
\loigiai{
Ta có: $10 \times 4 = 40$.
}
\end{ex}

\begin{ex}
Mỗi đĩa có 3 quả cam, có 4 đĩa như thế. Hỏi có tất cả bao nhiêu quả cam?
\shortans{12}
\loigiai{
Có: $3 \times 4 = 12$ quả cam.
}
\end{ex}

\begin{ex}
Mỗi bạn có 2 chiếc bút chì, 5 bạn có bao nhiêu chiếc bút chì?
\shortans{10}
\loigiai{
Có: $2 \times 5 = 10$ chiếc bút.
}
\end{ex}

\begin{ex}
Mỗi con gà có 2 cái chân. Hỏi 6 con gà có tất cả bao nhiêu cái chân?
\shortans{12}
\loigiai{
Có: $2 \times 6 = 12$ cái chân.
}
\end{ex}

\begin{ex}
Mỗi con bò có 4 cái chân. Hỏi 5 con bò có tất cả bao nhiêu cái chân?
\shortans{20}
\loigiai{
Có: $4 \times 5 = 20$ cái chân.
}
\end{ex}

\begin{ex}
Mỗi tuần lễ bạn Nam đi học 5 ngày. Hỏi trong 3 tuần lễ bạn Nam đi học bao nhiêu ngày?
\shortans{15}
\loigiai{
Đi học: $5 \times 3 = 15$ ngày.
}
\end{ex}

\begin{ex}
Mỗi gói bánh có 5 chiếc bánh. Hỏi 4 gói bánh như thế có bao nhiêu chiếc bánh?
\shortans{20}
\loigiai{
Có: $5 \times 4 = 20$ chiếc bánh.
}
\end{ex}

\begin{ex}
Một bàn học có 2 học sinh ngồi. Hỏi 8 bàn học như thế có bao nhiêu học sinh ngồi?
\shortans{16}
\loigiai{
Có: $2 \times 8 = 16$ học sinh.
}
\end{ex}

\begin{ex}
Mỗi hàng cây có 6 cây. Hỏi 3 hàng cây như thế có tất cả bao nhiêu cây?
\shortans{18}
\loigiai{
Có: $6 \times 3 = 18$ cây.
}
\end{ex}

\begin{ex}
Điền số thích hợp: $2 + 2 + 2 + 2 + 2 = 2 \times \dots$
\shortans{5}
\loigiai{
Số 2 được lấy 5 lần nên số cần điền là 5.
}
\end{ex}

\begin{ex}
Điền số thích hợp: $5 + 5 + 5 + 5 = 5 \times \dots$
\shortans{4}
\loigiai{
Số 5 được lấy 4 lần nên số cần điền là 4.
}
\end{ex}

\begin{ex}
Điền số thích hợp: $3 \times 4 = 3 + 3 + 3 + \dots$
\shortans{3}
\loigiai{
Số cần điền là 3.
}
\end{ex}

\begin{ex}
Tính tổng $4 + 4 + 4 + 4 + 4$.
\shortans{20}
\loigiai{
Ta có $4 \times 5 = 20$.
}
\end{ex}

\begin{ex}
Tính tổng $10 + 10 + 10 + 10 + 10$.
\shortans{50}
\loigiai{
Ta có $10 \times 5 = 50$.
}
\end{ex}

\begin{ex}
Một chiếc xe ô tô con có 4 bánh xe. Hỏi 6 chiếc xe ô tô con như thế có bao nhiêu bánh xe?
\shortans{24}
\loigiai{
Có: $4 \times 6 = 24$ bánh xe.
}
\end{ex}

\begin{ex}
Mỗi lọ hoa cắm 5 bông hoa hồng. Hỏi 6 lọ hoa như thế cắm tất cả bao nhiêu bông hoa hồng?
\shortans{30}
\loigiai{
Có: $5 \times 6 = 30$ bông hoa.
}
\end{ex}

\begin{ex}
Số 2 được lấy 9 lần có giá trị bằng bao nhiêu?
\shortans{18}
\loigiai{
Giá trị: $2 \times 9 = 18$.
}
\end{ex}

\begin{ex}
Số 5 được lấy 7 lần có giá trị bằng bao nhiêu?
\shortans{35}
\loigiai{
Giá trị: $5 \times 7 = 35$.
}
\end{ex}

\newpage

\TaoPhieuBaiTap{Chủ đề 8 -- Phép nhân, phép chia}{38}{Bảng nhân 2}{https://loptoanthaylinh.vn/bai-tap/toan-2/bai-38}
\begin{ex}
Trong phép nhân $2 \times 5 = 10$, số 10 được gọi là gì? Tính tích của 2 và 5.
\shortans{10}
\loigiai{
Tích của 2 và 5 là 10.
}
\end{ex}

\begin{ex}
Trong phép nhân $5 \times 4 = 20$, thừa số thứ nhất là 5, thừa số thứ hai bằng bao nhiêu?
\shortans{4}
\loigiai{
Thừa số thứ hai là 4.
}
\end{ex}

\begin{ex}
Biết hai thừa số là 2 và 7. Tính tích của hai số đó.
\shortans{14}
\loigiai{
Tích: $2 \times 7 = 14$.
}
\end{ex}

\begin{ex}
Biết hai thừa số là 5 và 6. Tính tích của hai số đó.
\shortans{30}
\loigiai{
Tích: $5 \times 6 = 30$.
}
\end{ex}

\begin{ex}
Biết hai thừa số là 2 và 9. Tính tích của hai số đó.
\shortans{18}
\loigiai{
Tích: $2 \times 9 = 18$.
}
\end{ex}

\begin{ex}
Biết hai thừa số là 5 và 8. Tính tích của hai số đó.
\shortans{40}
\loigiai{
Tích: $5 \times 8 = 40$.
}
\end{ex}

\begin{ex}
Trong phép nhân, nếu một thừa số là 0 thì tích bằng bao nhiêu?
\shortans{0}
\loigiai{
Bất kỳ số nào nhân với 0 cũng bằng 0.
}
\end{ex}

\begin{ex}
Trong phép nhân, nếu một thừa số là 1 thì tích bằng số nào? Cho phép nhân $7 \times 1$, kết quả bằng bao nhiêu?
\shortans{7}
\loigiai{
Ta có $7 \times 1 = 7$.
}
\end{ex}

\begin{ex}
Tính: $2 \times 3 + 4$.
\shortans{10}
\loigiai{
Ta có $2 \times 3 = 6$, rồi $6 + 4 = 10$.
}
\end{ex}

\begin{ex}
Tính: $5 \times 2 + 8$.
\shortans{18}
\loigiai{
Ta có $5 \times 2 = 10$, rồi $10 + 8 = 18$.
}
\end{ex}

\begin{ex}
Tính: $2 \times 8 - 6$.
\shortans{10}
\loigiai{
Ta có $2 \times 8 = 16$, rồi $16 - 6 = 10$.
}
\end{ex}

\begin{ex}
Tính: $5 \times 5 - 15$.
\shortans{10}
\loigiai{
Ta có $5 \times 5 = 25$, rồi $25 - 15 = 10$.
}
\end{ex}

\begin{ex}
Điền số thích hợp: $2 \times \dots = 12$.
\shortans{6}
\loigiai{
Số cần điền là: $12 : 2 = 6$.
}
\end{ex}

\begin{ex}
Điền số thích hợp: $5 \times \dots = 25$.
\shortans{5}
\loigiai{
Số cần điền là: $25 : 5 = 5$.
}
\end{ex}

\begin{ex}
Điền số thích hợp: $\dots \times 7 = 14$.
\shortans{2}
\loigiai{
Số cần điền là: $14 : 7 = 2$.
}
\end{ex}

\begin{ex}
Điền số thích hợp: $\dots \times 9 = 45$.
\shortans{5}
\loigiai{
Số cần điền là: $45 : 9 = 5$.
}
\end{ex}

\begin{ex}
Một nhóm có 2 bạn, 6 nhóm như thế có tất cả bao nhiêu bạn?
\shortans{12}
\loigiai{
Có: $2 \times 6 = 12$ bạn.
}
\end{ex}

\begin{ex}
Mỗi hộp có 5 chiếc bút bi, 7 hộp như thế có bao nhiêu chiếc bút bi?
\shortans{35}
\loigiai{
Có: $5 \times 7 = 35$ chiếc bút.
}
\end{ex}

\begin{ex}
Mỗi chuồng có 2 con thỏ, 8 chuồng như thế có tất cả bao nhiêu con thỏ?
\shortans{16}
\loigiai{
Có: $2 \times 8 = 16$ con thỏ.
}
\end{ex}

\begin{ex}
Bác An mua 5 túi gạo, mỗi túi nặng 5 kg. Bác An đã mua tất cả bao nhiêu ki-lô-gam gạo?
\shortans{25}
\loigiai{
Bác mua: $5 \times 5 = 25$ kg.
}
\end{ex}

\begin{ex}
Có 4 can nước, mỗi can chứa 5 lít. Hỏi cả 4 can chứa bao nhiêu lít nước?
\shortans{20}
\loigiai{
Chứa: $5 \times 4 = 20$ lít.
}
\end{ex}

\begin{ex}
Tìm tích của số bé nhất có hai chữ số và số 2.
\shortans{20}
\loigiai{
Số bé nhất có hai chữ số là 10. Tích: $10 \times 2 = 20$.
}
\end{ex}

\begin{ex}
Tìm tích của số liền sau số 1 và số liền trước số 6.
\shortans{10}
\loigiai{
Số liền sau 1 là 2, liền trước 6 là 5. Tích: $2 \times 5 = 10$.
}
\end{ex}

\begin{ex}
Tích của hai số là 18, biết một thừa số là 2. Thừa số kia là bao nhiêu?
\shortans{9}
\loigiai{
Thừa số kia là: $18 : 2 = 9$.
}
\end{ex}

\begin{ex}
Tích của hai số là 35, biết một thừa số là 5. Thừa số kia là bao nhiêu?
\shortans{7}
\loigiai{
Thừa số kia là: $35 : 5 = 7$.
}
\end{ex}

\newpage

\TaoPhieuBaiTap{Chủ đề 8 -- Phép nhân, phép chia}{39}{Bảng nhân 5}{https://loptoanthaylinh.vn/bai-tap/toan-2/bai-39}
\begin{ex}
Tính: $2 \times 1$.
\shortans{2}
\loigiai{
Ta có $2 \times 1 = 2$.
}
\end{ex}

\begin{ex}
Tính: $2 \times 2$.
\shortans{4}
\loigiai{
Ta có $2 \times 2 = 4$.
}
\end{ex}

\begin{ex}
Tính: $2 \times 3$.
\shortans{6}
\loigiai{
Ta có $2 \times 3 = 6$.
}
\end{ex}

\begin{ex}
Tính: $2 \times 4$.
\shortans{8}
\loigiai{
Ta có $2 \times 4 = 8$.
}
\end{ex}

\begin{ex}
Tính: $2 \times 5$.
\shortans{10}
\loigiai{
Ta có $2 \times 5 = 10$.
}
\end{ex}

\begin{ex}
Tính: $2 \times 6$.
\shortans{12}
\loigiai{
Ta có $2 \times 6 = 12$.
}
\end{ex}

\begin{ex}
Tính: $2 \times 7$.
\shortans{14}
\loigiai{
Ta có $2 \times 7 = 14$.
}
\end{ex}

\begin{ex}
Tính: $2 \times 8$.
\shortans{16}
\loigiai{
Ta có $2 \times 8 = 16$.
}
\end{ex}

\begin{ex}
Tính: $2 \times 9$.
\shortans{18}
\loigiai{
Ta có $2 \times 9 = 18$.
}
\end{ex}

\begin{ex}
Tính: $2 \times 10$.
\shortans{20}
\loigiai{
Ta có $2 \times 10 = 20$.
}
\end{ex}

\begin{ex}
Tính giá trị biểu thức: $2 \times 4 + 15$.
\shortans{23}
\loigiai{
Ta có $2 \times 4 = 8$, rồi $8 + 15 = 23$.
}
\end{ex}

\begin{ex}
Tính giá trị biểu thức: $2 \times 7 + 26$.
\shortans{40}
\loigiai{
Ta có $2 \times 7 = 14$, rồi $14 + 26 = 40$.
}
\end{ex}

\begin{ex}
Tính giá trị biểu thức: $2 \times 9 - 9$.
\shortans{9}
\loigiai{
Ta có $2 \times 9 = 18$, rồi $18 - 9 = 9$.
}
\end{ex}

\begin{ex}
Tính giá trị biểu thức: $2 \times 6 + 18$.
\shortans{30}
\loigiai{
Ta có $2 \times 6 = 12$, rồi $12 + 18 = 30$.
}
\end{ex}

\begin{ex}
Tính giá trị biểu thức: $2 \times 8 - 10$.
\shortans{6}
\loigiai{
Ta có $2 \times 8 = 16$, rồi $16 - 10 = 6$.
}
\end{ex}

\begin{ex}
Điền số thích hợp: $2 \times \dots = 16$.
\shortans{8}
\loigiai{
Số cần điền là: $16 : 2 = 8$.
}
\end{ex}

\begin{ex}
Điền số thích hợp: $2 \times \dots = 14$.
\shortans{7}
\loigiai{
Số cần điền là: $14 : 2 = 7$.
}
\end{ex}

\begin{ex}
Một đôi đũa có 2 chiếc đũa. Hỏi 7 đôi đũa có tất cả bao nhiêu chiếc đũa?
\shortans{14}
\loigiai{
Có: $2 \times 7 = 14$ chiếc đũa.
}
\end{ex}

\begin{ex}
Một đôi chim bồ câu có 2 con. Hỏi 9 đôi chim bồ câu có tất cả bao nhiêu con?
\shortans{18}
\loigiai{
Có: $2 \times 9 = 18$ con chim.
}
\end{ex}

\begin{ex}
Mỗi bạn học sinh được phát 2 quyển vở. Hỏi 10 bạn học sinh được phát bao nhiêu quyển vở?
\shortans{20}
\loigiai{
Phát: $2 \times 10 = 20$ quyển vở.
}
\end{ex}

\begin{ex}
Một con chim có 2 cái cánh. Hỏi 8 con chim có tất cả bao nhiêu cái cánh?
\shortans{16}
\loigiai{
Có: $2 \times 8 = 16$ cái cánh.
}
\end{ex}

\begin{ex}
Mỗi chiếc xe đạp cần 2 bánh xe. Hỏi để lắp ráp 6 chiếc xe đạp cần bao nhiêu bánh xe?
\shortans{12}
\loigiai{
Cần: $2 \times 6 = 12$ bánh xe.
}
\end{ex}

\begin{ex}
Trong vườn có 5 luống hoa, mỗi luống có 2 cây hoa hồng. Có tất cả bao nhiêu cây hoa hồng?
\shortans{10}
\loigiai{
Có: $2 \times 5 = 10$ cây hoa.
}
\end{ex}

\begin{ex}
So sánh: Tích $2 \times 8$ lớn hơn tích $2 \times 6$ bao nhiêu đơn vị?
\shortans{4}
\loigiai{
Ta có $16 - 12 = 4$.
}
\end{ex}

\begin{ex}
Điền số thích hợp vào ô trống: $2 \times 5 + 2 = 2 \times \dots$
\shortans{6}
\loigiai{
Vế trái là $10 + 2 = 12$, vậy $2 \times 6 = 12$. Số cần điền là 6.
}
\end{ex}

\newpage

\TaoPhieuBaiTap{Chủ đề 8 -- Phép nhân, phép chia}{40}{Phép chia. Số bị chia - Số chia - Thương}{https://loptoanthaylinh.vn/bai-tap/toan-2/bai-40}
\begin{ex}
Tính: $5 \times 1$.
\shortans{5}
\loigiai{
Ta có $5 \times 1 = 5$.
}
\end{ex}

\begin{ex}
Tính: $5 \times 2$.
\shortans{10}
\loigiai{
Ta có $5 \times 2 = 10$.
}
\end{ex}

\begin{ex}
Tính: $5 \times 3$.
\shortans{15}
\loigiai{
Ta có $5 \times 3 = 15$.
}
\end{ex}

\begin{ex}
Tính: $5 \times 4$.
\shortans{20}
\loigiai{
Ta có $5 \times 4 = 20$.
}
\end{ex}

\begin{ex}
Tính: $5 \times 5$.
\shortans{25}
\loigiai{
Ta có $5 \times 5 = 25$.
}
\end{ex}

\begin{ex}
Tính: $5 \times 6$.
\shortans{30}
\loigiai{
Ta có $5 \times 6 = 30$.
}
\end{ex}

\begin{ex}
Tính: $5 \times 7$.
\shortans{35}
\loigiai{
Ta có $5 \times 7 = 35$.
}
\end{ex}

\begin{ex}
Tính: $5 \times 8$.
\shortans{40}
\loigiai{
Ta có $5 \times 8 = 40$.
}
\end{ex}

\begin{ex}
Tính: $5 \times 9$.
\shortans{45}
\loigiai{
Ta có $5 \times 9 = 45$.
}
\end{ex}

\begin{ex}
Tính: $5 \times 10$.
\shortans{50}
\loigiai{
Ta có $5 \times 10 = 50$.
}
\end{ex}

\begin{ex}
Tính giá trị biểu thức: $5 \times 3 + 15$.
\shortans{30}
\loigiai{
Ta có $5 \times 3 = 15$, rồi $15 + 15 = 30$.
}
\end{ex}

\begin{ex}
Tính giá trị biểu thức: $5 \times 6 - 12$.
\shortans{18}
\loigiai{
Ta có $5 \times 6 = 30$, rồi $30 - 12 = 18$.
}
\end{ex}

\begin{ex}
Tính giá trị biểu thức: $5 \times 8 + 20$.
\shortans{60}
\loigiai{
Ta có $5 \times 8 = 40$, rồi $40 + 20 = 60$.
}
\end{ex}

\begin{ex}
Tính giá trị biểu thức: $5 \times 9 - 25$.
\shortans{20}
\loigiai{
Ta có $5 \times 9 = 45$, rồi $45 - 25 = 20$.
}
\end{ex}

\begin{ex}
Tính giá trị biểu thức: $5 \times 7 + 15$.
\shortans{50}
\loigiai{
Ta có $5 \times 7 = 35$, rồi $35 + 15 = 50$.
}
\end{ex}

\begin{ex}
Điền số thích hợp: $5 \times \dots = 35$.
\shortans{7}
\loigiai{
Số cần điền là: $35 : 5 = 7$.
}
\end{ex}

\begin{ex}
Điền số thích hợp: $5 \times \dots = 40$.
\shortans{8}
\loigiai{
Số cần điền là: $40 : 5 = 8$.
}
\end{ex}

\begin{ex}
Một bàn tay có 5 ngón tay. Hỏi 6 bàn tay có tất cả bao nhiêu ngón tay?
\shortans{30}
\loigiai{
Có: $5 \times 6 = 30$ ngón tay.
}
\end{ex}

\begin{ex}
Một ngôi sao có 5 cánh. Hỏi 8 ngôi sao có tất cả bao nhiêu cánh sao?
\shortans{40}
\loigiai{
Có: $5 \times 8 = 40$ cánh sao.
}
\end{ex}

\begin{ex}
Mỗi can chứa 5 lít dầu. Hỏi 7 can như thế chứa bao nhiêu lít dầu?
\shortans{35}
\loigiai{
Chứa: $5 \times 7 = 35$ lít dầu.
}
\end{ex}

\begin{ex}
Mỗi bao gạo nặng 5 kg. Hỏi 9 bao gạo như thế nặng tất cả bao nhiêu ki-lô-gam?
\shortans{45}
\loigiai{
Nặng: $5 \times 9 = 45$ kg.
}
\end{ex}

\begin{ex}
Một tuần lễ đi học 5 ngày. Hỏi trong 4 tuần lễ học sinh đi học bao nhiêu ngày?
\shortans{20}
\loigiai{
Đi học: $5 \times 4 = 20$ ngày.
}
\end{ex}

\begin{ex}
Mỗi gói kẹo có 5 chiếc kẹo. Hỏi 10 gói kẹo như thế có bao nhiêu chiếc kẹo?
\shortans{50}
\loigiai{
Có: $5 \times 10 = 50$ chiếc kẹo.
}
\end{ex}

\begin{ex}
So sánh: Tích $5 \times 7$ lớn hơn tích $5 \times 5$ bao nhiêu đơn vị?
\shortans{10}
\loigiai{
Ta có $35 - 25 = 10$.
}
\end{ex}

\begin{ex}
Điền số thích hợp: $5 \times 4 + 5 = 5 \times \dots$
\shortans{5}
\loigiai{
Vế trái là $20 + 5 = 25$, vậy $5 \times 5 = 25$. Số cần điền là 5.
}
\end{ex}

\newpage

\TaoPhieuBaiTap{Chủ đề 8 -- Phép nhân, phép chia}{41}{Bảng chia 2}{https://loptoanthaylinh.vn/bai-tap/toan-2/bai-41}
\begin{ex}
Có 6 quả táo chia đều cho 2 bạn. Hỏi mỗi bạn được bao nhiêu quả táo?
\shortans{3}
\loigiai{
Mỗi bạn được: $6 : 2 = 3$ quả táo.
}
\end{ex}

\begin{ex}
Có 8 viên bi chia đều cho 2 bạn. Hỏi mỗi bạn được bao nhiêu viên bi?
\shortans{4}
\loigiai{
Mỗi bạn được: $8 : 2 = 4$ viên bi.
}
\end{ex}

\begin{ex}
Có 10 cái kẹo chia đều vào 2 đĩa. Hỏi mỗi đĩa có bao nhiêu cái kẹo?
\shortans{5}
\loigiai{
Mỗi đĩa có: $10 : 2 = 5$ cái kẹo.
}
\end{ex}

\begin{ex}
Có 15 bông hoa cắm đều vào 5 lọ hoa. Hỏi mỗi lọ có bao nhiêu bông hoa?
\shortans{3}
\loigiai{
Mỗi lọ có: $15 : 5 = 3$ bông hoa.
}
\end{ex}

\begin{ex}
Có 20 chiếc bút chì chia đều cho 5 bạn. Hỏi mỗi bạn được bao nhiêu chiếc bút chì?
\shortans{4}
\loigiai{
Mỗi bạn được: $20 : 5 = 4$ chiếc bút.
}
\end{ex}

\begin{ex}
Có 12 cái bánh chia đều cho 2 bạn. Mỗi bạn được bao nhiêu cái bánh?
\shortans{6}
\loigiai{
Mỗi bạn được: $12 : 2 = 6$ cái bánh.
}
\end{ex}

\begin{ex}
Có 14 quyển vở chia đều cho 2 tổ. Mỗi tổ được bao nhiêu quyển vở?
\shortans{7}
\loigiai{
Mỗi tổ được: $14 : 2 = 7$ quyển vở.
}
\end{ex}

\begin{ex}
Có 25 lít dầu chia đều vào 5 can. Mỗi can có bao nhiêu lít dầu?
\shortans{5}
\loigiai{
Mỗi can có: $25 : 5 = 5$ lít dầu.
}
\end{ex}

\begin{ex}
Có 30 kg gạo chia đều vào 5 túi. Mỗi túi nặng bao nhiêu ki-lô-gam?
\shortans{6}
\loigiai{
Mỗi túi nặng: $30 : 5 = 6$ kg.
}
\end{ex}

\begin{ex}
Từ phép nhân $2 \times 6 = 12$, viết phép chia $12 : 2$. Kết quả bằng bao nhiêu?
\shortans{6}
\loigiai{
Ta có $12 : 2 = 6$.
}
\end{ex}

\begin{ex}
Từ phép nhân $5 \times 4 = 20$, viết phép chia $20 : 5$. Kết quả bằng bao nhiêu?
\shortans{4}
\loigiai{
Ta có $20 : 5 = 4$.
}
\end{ex}

\begin{ex}
Từ phép nhân $2 \times 8 = 16$, viết phép chia $16 : 8$. Kết quả bằng bao nhiêu?
\shortans{2}
\loigiai{
Ta có $16 : 8 = 2$.
}
\end{ex}

\begin{ex}
Từ phép nhân $5 \times 7 = 35$, viết phép chia $35 : 7$. Kết quả bằng bao nhiêu?
\shortans{5}
\loigiai{
Ta có $35 : 7 = 5$.
}
\end{ex}

\begin{ex}
Có 18 học sinh xếp thành các hàng, mỗi hàng có 2 bạn. Hỏi xếp được bao nhiêu hàng?
\shortans{9}
\loigiai{
Số hàng: $18 : 2 = 9$ hàng.
}
\end{ex}

\begin{ex}
Có 40 bạn học sinh xếp thành các nhóm, mỗi nhóm có 5 bạn. Hỏi có bao nhiêu nhóm?
\shortans{8}
\loigiai{
Số nhóm: $40 : 5 = 8$ nhóm.
}
\end{ex}

\begin{ex}
Có 16 chiếc ghế xếp đều vào 2 dãy. Mỗi dãy có bao nhiêu chiếc ghế?
\shortans{8}
\loigiai{
Mỗi dãy có: $16 : 2 = 8$ chiếc ghế.
}
\end{ex}

\begin{ex}
Có 35 củ cà rốt chia đều cho 5 chú thỏ. Mỗi chú thỏ được bao nhiêu củ cà rốt?
\shortans{7}
\loigiai{
Mỗi chú thỏ được: $35 : 5 = 7$ củ cà rốt.
}
\end{ex}

\begin{ex}
Một sợi dây dài 10 dm cắt thành 2 đoạn bằng nhau. Mỗi đoạn dài bao nhiêu đề-xi-mét?
\shortans{5}
\loigiai{
Mỗi đoạn dài: $10 : 2 = 5$ dm.
}
\end{ex}

\begin{ex}
Một thanh gỗ dài 45 cm cắt thành 5 đoạn bằng nhau. Mỗi đoạn dài bao nhiêu xăng-ti-mét?
\shortans{9}
\loigiai{
Mỗi đoạn dài: $45 : 5 = 9$ cm.
}
\end{ex}

\begin{ex}
Có 18 con cá thả đều vào 2 bể kính. Mỗi bể kính có bao nhiêu con cá?
\shortans{9}
\loigiai{
Mỗi bể có: $18 : 2 = 9$ con cá.
}
\end{ex}

\begin{ex}
Có 50 lít nước mắm đóng đều vào các can 5 lít. Hỏi đóng được bao nhiêu can?
\shortans{10}
\loigiai{
Đóng được: $50 : 5 = 10$ can.
}
\end{ex}

\begin{ex}
Có 20 cái kẹo chia cho mỗi bạn 2 cái. Hỏi chia được cho bao nhiêu bạn?
\shortans{10}
\loigiai{
Chia được cho: $20 : 2 = 10$ bạn.
}
\end{ex}

\begin{ex}
Điền số thích hợp: $14 : 2 = \dots$
\shortans{7}
\loigiai{
Ta có $14 : 2 = 7$.
}
\end{ex}

\begin{ex}
Điền số thích hợp: $45 : 5 = \dots$
\shortans{9}
\loigiai{
Ta có $45 : 5 = 9$.
}
\end{ex}

\begin{ex}
Tính: $20 : 2 - 3$.
\shortans{7}
\loigiai{
Ta có $20 : 2 = 10$, rồi $10 - 3 = 7$.
}
\end{ex}

\newpage

\TaoPhieuBaiTap{Chủ đề 8 -- Phép nhân, phép chia}{42}{Bảng chia 5}{https://loptoanthaylinh.vn/bai-tap/toan-2/bai-42}
\begin{ex}
Trong phép chia $12 : 2 = 6$, số 6 được gọi là gì? Tính thương của 12 và 2.
\shortans{6}
\loigiai{
Thương là 6.
}
\end{ex}

\begin{ex}
Trong phép chia $20 : 5 = 4$, số bị chia bằng bao nhiêu?
\shortans{20}
\loigiai{
Số bị chia là 20.
}
\end{ex}

\begin{ex}
Trong phép chia $35 : 5 = 7$, số chia bằng bao nhiêu?
\shortans{5}
\loigiai{
Số chia là 5.
}
\end{ex}

\begin{ex}
Biết số bị chia là 16, số chia là 2. Tính thương.
\shortans{8}
\loigiai{
Thương: $16 : 2 = 8$.
}
\end{ex}

\begin{ex}
Biết số bị chia là 40, số chia là 5. Tính thương.
\shortans{8}
\loigiai{
Thương: $40 : 5 = 8$.
}
\end{ex}

\begin{ex}
Biết số bị chia là 18, số chia là 2. Tính thương.
\shortans{9}
\loigiai{
Thương: $18 : 2 = 9$.
}
\end{ex}

\begin{ex}
Biết số bị chia là 45, số chia là 5. Tính thương.
\shortans{9}
\loigiai{
Thương: $45 : 5 = 9$.
}
\end{ex}

\begin{ex}
Biết thương là 7, số chia là 2. Tìm số bị chia.
\shortans{14}
\loigiai{
Số bị chia: $7 \times 2 = 14$.
}
\end{ex}

\begin{ex}
Biết thương là 6, số chia là 5. Tìm số bị chia.
\shortans{30}
\loigiai{
Số bị chia: $6 \times 5 = 30$.
}
\end{ex}

\begin{ex}
Biết thương là 9, số chia là 2. Tìm số bị chia.
\shortans{18}
\loigiai{
Số bị chia: $9 \times 2 = 18$.
}
\end{ex}

\begin{ex}
Biết thương là 8, số chia là 5. Tìm số bị chia.
\shortans{40}
\loigiai{
Số bị chia: $8 \times 5 = 40$.
}
\end{ex}

\begin{ex}
Biết số bị chia là 10, thương là 5. Tìm số chia.
\shortans{2}
\loigiai{
Số chia: $10 : 5 = 2$.
}
\end{ex}

\begin{ex}
Biết số bị chia là 25, thương là 5. Tìm số chia.
\shortans{5}
\loigiai{
Số chia: $25 : 5 = 5$.
}
\end{ex}

\begin{ex}
Tính giá trị biểu thức: $16 : 2 + 15$.
\shortans{23}
\loigiai{
Ta có $16 : 2 = 8$, rồi $8 + 15 = 23$.
}
\end{ex}

\begin{ex}
Tính giá trị biểu thức: $30 : 5 + 24$.
\shortans{30}
\loigiai{
Ta có $30 : 5 = 6$, rồi $6 + 24 = 30$.
}
\end{ex}

\begin{ex}
Tính giá trị biểu thức: $14 : 2 \times 3$.
\shortans{21}
\loigiai{
Ta có $14 : 2 = 7$, rồi $7 \times 3 = 21$.
}
\end{ex}

\begin{ex}
Tính giá trị biểu thức: $40 : 5 \times 2$.
\shortans{16}
\loigiai{
Ta có $40 : 5 = 8$, rồi $8 \times 2 = 16$.
}
\end{ex}

\begin{ex}
Điền số thích hợp: $18 : \dots = 9$.
\shortans{2}
\loigiai{
Số cần điền là: $18 : 9 = 2$.
}
\end{ex}

\begin{ex}
Điền số thích hợp: $35 : \dots = 7$.
\shortans{5}
\loigiai{
Số cần điền là: $35 : 7 = 5$.
}
\end{ex}

\begin{ex}
Điền số thích hợp: $\dots : 2 = 8$.
\shortans{16}
\loigiai{
Số cần điền là: $8 \times 2 = 16$.
}
\end{ex}

\begin{ex}
Điền số thích hợp: $\dots : 5 = 9$.
\shortans{45}
\loigiai{
Số cần điền là: $9 \times 5 = 45$.
}
\end{ex}

\begin{ex}
Trong một phép chia có số chia là 2, thương là số lớn nhất có một chữ số. Số bị chia là bao nhiêu?
\shortans{18}
\loigiai{
Số lớn nhất có 1 chữ số là 9. Số bị chia: $9 \times 2 = 18$.
}
\end{ex}

\begin{ex}
Trong một phép chia có số chia là 5, thương là số chẵn lớn nhất có một chữ số. Số bị chia là bao nhiêu?
\shortans{40}
\loigiai{
Số chẵn lớn nhất có 1 chữ số là 8. Số bị chia: $8 \times 5 = 40$.
}
\end{ex}

\begin{ex}
Một số chia cho 2 thì được 10. Lấy số đó trừ đi 5 thì được bao nhiêu?
\shortans{15}
\loigiai{
Số đó là: $10 \times 2 = 20$. Hiệu: $20 - 5 = 15$.
}
\end{ex}

\begin{ex}
Tính: $50 : 5 - 2$.
\shortans{8}
\loigiai{
Ta có $50 : 5 = 10$, rồi $10 - 2 = 8$.
}
\end{ex}

\newpage

\TaoPhieuBaiTap{Chủ đề 8 -- Phép nhân, phép chia}{43}{Số 1 và số 0 trong phép nhân, phép chia}{https://loptoanthaylinh.vn/bai-tap/toan-2/bai-43}
\begin{ex}
Tính: $2 : 2$.
\shortans{1}
\loigiai{
Ta có $2 : 2 = 1$.
}
\end{ex}

\begin{ex}
Tính: $4 : 2$.
\shortans{2}
\loigiai{
Ta có $4 : 2 = 2$.
}
\end{ex}

\begin{ex}
Tính: $6 : 2$.
\shortans{3}
\loigiai{
Ta có $6 : 2 = 3$.
}
\end{ex}

\begin{ex}
Tính: $8 : 2$.
\shortans{4}
\loigiai{
Ta có $8 : 2 = 4$.
}
\end{ex}

\begin{ex}
Tính: $10 : 2$.
\shortans{5}
\loigiai{
Ta có $10 : 2 = 5$.
}
\end{ex}

\begin{ex}
Tính: $12 : 2$.
\shortans{6}
\loigiai{
Ta có $12 : 2 = 6$.
}
\end{ex}

\begin{ex}
Tính: $14 : 2$.
\shortans{7}
\loigiai{
Ta có $14 : 2 = 7$.
}
\end{ex}

\begin{ex}
Tính: $16 : 2$.
\shortans{8}
\loigiai{
Ta có $16 : 2 = 8$.
}
\end{ex}

\begin{ex}
Tính: $18 : 2$.
\shortans{9}
\loigiai{
Ta có $18 : 2 = 9$.
}
\end{ex}

\begin{ex}
Tính: $20 : 2$.
\shortans{10}
\loigiai{
Ta có $20 : 2 = 10$.
}
\end{ex}

\begin{ex}
Tính giá trị biểu thức: $18 : 2 + 11$.
\shortans{20}
\loigiai{
Ta có $18 : 2 = 9$, rồi $9 + 11 = 20$.
}
\end{ex}

\begin{ex}
Tính giá trị biểu thức: $14 : 2 + 25$.
\shortans{32}
\loigiai{
Ta có $14 : 2 = 7$, rồi $7 + 25 = 32$.
}
\end{ex}

\begin{ex}
Tính giá trị biểu thức: $16 : 2 - 3$.
\shortans{5}
\loigiai{
Ta có $16 : 2 = 8$, rồi $8 - 3 = 5$.
}
\end{ex}

\begin{ex}
Tính giá trị biểu thức: $20 : 2 + 45$.
\shortans{55}
\loigiai{
Ta có $20 : 2 = 10$, rồi $10 + 45 = 55$.
}
\end{ex}

\begin{ex}
Tính giá trị biểu thức: $12 : 2 \times 5$.
\shortans{30}
\loigiai{
Ta có $12 : 2 = 6$, rồi $6 \times 5 = 30$.
}
\end{ex}

\begin{ex}
Điền số thích hợp: $10 : 2 = \dots - 3$.
\shortans{8}
\loigiai{
Vế trái: $10 : 2 = 5$. Số cần điền: $5 + 3 = 8$.
}
\end{ex}

\begin{ex}
Điền số thích hợp: $16 : 2 = \dots + 4$.
\shortans{4}
\loigiai{
Vế trái: $16 : 2 = 8$. Số cần điền: $8 - 4 = 4$.
}
\end{ex}

\begin{ex}
Có 12 chiếc kẹo chia đều cho 2 bạn. Mỗi bạn được bao nhiêu chiếc kẹo?
\shortans{6}
\loigiai{
Mỗi bạn được: $12 : 2 = 6$ chiếc kẹo.
}
\end{ex}

\begin{ex}
Có 18 lít dầu chia đều vào 2 can. Mỗi can chứa bao nhiêu lít dầu?
\shortans{7}
\loigiai{
Mỗi can chứa: $18 : 2 = 9$ lít dầu (kiểm tra $18 : 2 = 9$).
}
\end{ex}

\begin{ex}
Mẹ mua 14 kg khoai tây chia đều vào 2 túi. Mỗi túi nặng bao nhiêu ki-lô-gam?
\shortans{7}
\loigiai{
Mỗi túi nặng: $14 : 2 = 7$ kg.
}
\end{ex}

\begin{ex}
Một sợi dây dài 16 mét cắt làm 2 đoạn bằng nhau. Mỗi đoạn dài bao nhiêu mét?
\shortans{8}
\loigiai{
Mỗi đoạn dài: $16 : 2 = 8$ mét.
}
\end{ex}

\begin{ex}
Một đoạn đường dài 20 km chia làm 2 chặng bằng nhau. Mỗi chặng dài bao nhiêu ki-lô-mét?
\shortans{10}
\loigiai{
Mỗi chặng dài: $20 : 2 = 10$ km.
}
\end{ex}

\begin{ex}
Có 8 con vịt chia vào 2 chuồng. Mỗi chuồng có bao nhiêu con vịt?
\shortans{4}
\loigiai{
Mỗi chuồng có: $8 : 2 = 4$ con vịt.
}
\end{ex}

\begin{ex}
Lớp có 20 học sinh xếp thành 2 hàng đều nhau. Mỗi hàng có bao nhiêu học sinh?
\shortans{10}
\loigiai{
Mỗi hàng có: $20 : 2 = 10$ học sinh.
}
\end{ex}

\begin{ex}
Tính: $2 \times 8 : 2$.
\shortans{8}
\loigiai{
Ta có $2 \times 8 = 16$, rồi $16 : 2 = 8$.
}
\end{ex}

\newpage

\TaoPhieuBaiTap{Chủ đề 8 -- Phép nhân, phép chia}{44}{Luyện tập phép nhân và phép chia}{https://loptoanthaylinh.vn/bai-tap/toan-2/bai-44}
\begin{ex}
Tính: $5 : 5$.
\shortans{1}
\loigiai{
Ta có $5 : 5 = 1$.
}
\end{ex}

\begin{ex}
Tính: $10 : 5$.
\shortans{2}
\loigiai{
Ta có $10 : 5 = 2$.
}
\end{ex}

\begin{ex}
Tính: $15 : 5$.
\shortans{3}
\loigiai{
Ta có $15 : 5 = 3$.
}
\end{ex}

\begin{ex}
Tính: $20 : 5$.
\shortans{4}
\loigiai{
Ta có $20 : 5 = 4$.
}
\end{ex}

\begin{ex}
Tính: $25 : 5$.
\shortans{5}
\loigiai{
Ta có $25 : 5 = 5$.
}
\end{ex}

\begin{ex}
Tính: $30 : 5$.
\shortans{6}
\loigiai{
Ta có $30 : 5 = 6$.
}
\end{ex}

\begin{ex}
Tính: $35 : 5$.
\shortans{7}
\loigiai{
Ta có $35 : 5 = 7$.
}
\end{ex}

\begin{ex}
Tính: $40 : 5$.
\shortans{8}
\loigiai{
Ta có $40 : 5 = 8$.
}
\end{ex}

\begin{ex}
Tính: $45 : 5$.
\shortans{9}
\loigiai{
Ta có $45 : 5 = 9$.
}
\end{ex}

\begin{ex}
Tính: $50 : 5$.
\shortans{10}
\loigiai{
Ta có $50 : 5 = 10$.
}
\end{ex}

\begin{ex}
Tính giá trị biểu thức: $35 : 5 + 18$.
\shortans{25}
\loigiai{
Ta có $35 : 5 = 7$, rồi $7 + 18 = 25$.
}
\end{ex}

\begin{ex}
Tính giá trị biểu thức: $40 : 5 + 32$.
\shortans{40}
\loigiai{
Ta có $40 : 5 = 8$, rồi $8 + 32 = 40$.
}
\end{ex}

\begin{ex}
Tính giá trị biểu thức: $45 : 5 - 4$.
\shortans{5}
\loigiai{
Ta có $45 : 5 = 9$, rồi $9 - 4 = 5$.
}
\end{ex}

\begin{ex}
Tính giá trị biểu thức: $50 : 5 - 7$.
\shortans{3}
\loigiai{
Ta có $50 : 5 = 10$, rồi $10 - 7 = 3$.
}
\end{ex}

\begin{ex}
Tính giá trị biểu thức: $25 : 5 \times 4$.
\shortans{20}
\loigiai{
Ta có $25 : 5 = 5$, rồi $5 \times 4 = 20$.
}
\end{ex}

\begin{ex}
Tính giá trị biểu thức: $30 : 5 \times 3$.
\shortans{18}
\loigiai{
Ta có $30 : 5 = 6$, rồi $6 \times 3 = 18$.
}
\end{ex}

\begin{ex}
Điền số thích hợp: $20 : 5 = \dots - 6$.
\shortans{10}
\loigiai{
Vế trái: $20 : 5 = 4$. Số cần điền: $4 + 6 = 10$.
}
\end{ex}

\begin{ex}
Điền số thích hợp: $40 : 5 = \dots + 5$.
\shortans{3}
\loigiai{
Vế trái: $40 : 5 = 8$. Số cần điền: $8 - 5 = 3$.
}
\end{ex}

\begin{ex}
Có 30 bông hoa cắm đều vào 5 lọ hoa. Mỗi lọ có bao nhiêu bông hoa?
\shortans{6}
\loigiai{
Mỗi lọ có: $30 : 5 = 6$ bông hoa.
}
\end{ex}

\begin{ex}
Có 35 cái bánh chia đều cho 5 bạn. Mỗi bạn được bao nhiêu cái bánh?
\shortans{7}
\loigiai{
Mỗi bạn được: $35 : 5 = 7$ cái bánh.
}
\end{ex}

\begin{ex}
Có 40 kg gạo chia đều vào 5 bao. Mỗi bao nặng bao nhiêu ki-lô-gam?
\shortans{8}
\loigiai{
Mỗi bao nặng: $40 : 5 = 8$ kg.
}
\end{ex}

\begin{ex}
Có 45 lít dầu chia đều vào 5 can. Mỗi can chứa bao nhiêu lít dầu?
\shortans{9}
\loigiai{
Mỗi can chứa: $45 : 5 = 9$ lít dầu.
}
\end{ex}

\begin{ex}
Có 50 bạn học sinh chia đều thành 5 đội kéo co. Mỗi đội có bao nhiêu bạn?
\shortans{10}
\loigiai{
Mỗi đội có: $50 : 5 = 10$ bạn.
}
\end{ex}

\begin{ex}
Một sợi dây đồng dài 25 cm cắt thành 5 đoạn bằng nhau. Mỗi đoạn dài bao nhiêu xăng-ti-mét?
\shortans{5}
\loigiai{
Mỗi đoạn dài: $25 : 5 = 5$ cm.
}
\end{ex}

\begin{ex}
Tính: $5 \times 6 : 5$.
\shortans{6}
\loigiai{
Ta có $5 \times 6 = 30$, rồi $30 : 5 = 6$.
}
\end{ex}

\newpage

\TaoPhieuBaiTap{Chủ đề 8 -- Phép nhân, phép chia}{45}{Luyện tập chung}{https://loptoanthaylinh.vn/bai-tap/toan-2/bai-45}
\begin{ex}
Tính: $2 \times 9$.
\shortans{18}
\loigiai{
Ta có $2 \times 9 = 18$.
}
\end{ex}

\begin{ex}
Tính: $18 : 2$.
\shortans{9}
\loigiai{
Ta có $18 : 2 = 9$.
}
\end{ex}

\begin{ex}
Tính: $5 \times 8$.
\shortans{40}
\loigiai{
Ta có $5 \times 8 = 40$.
}
\end{ex}

\begin{ex}
Tính: $40 : 5$.
\shortans{8}
\loigiai{
Ta có $40 : 5 = 8$.
}
\end{ex}

\begin{ex}
Tính: $2 \times 7$.
\shortans{14}
\loigiai{
Ta có $2 \times 7 = 14$.
}
\end{ex}

\begin{ex}
Tính: $14 : 2$.
\shortans{7}
\loigiai{
Ta có $14 : 2 = 7$.
}
\end{ex}

\begin{ex}
Tính: $5 \times 6$.
\shortans{30}
\loigiai{
Ta có $5 \times 6 = 30$.
}
\end{ex}

\begin{ex}
Tính: $30 : 5$.
\shortans{6}
\loigiai{
Ta có $30 : 5 = 6$.
}
\end{ex}

\begin{ex}
Tính giá trị biểu thức: $2 \times 8 + 15$.
\shortans{31}
\loigiai{
Ta có $2 \times 8 = 16$, rồi $16 + 15 = 31$.
}
\end{ex}

\begin{ex}
Tính giá trị biểu thức: $5 \times 7 - 18$.
\shortans{17}
\loigiai{
Ta có $5 \times 7 = 35$, rồi $35 - 18 = 17$.
}
\end{ex}

\begin{ex}
Tính giá trị biểu thức: $16 : 2 + 32$.
\shortans{40}
\loigiai{
Ta có $16 : 2 = 8$, rồi $8 + 32 = 40$.
}
\end{ex}

\begin{ex}
Tính giá trị biểu thức: $45 : 5 + 26$.
\shortans{35}
\loigiai{
Ta có $45 : 5 = 9$, rồi $9 + 26 = 35$.
}
\end{ex}

\begin{ex}
Tính giá trị biểu thức: $2 \times 5 \times 3$.
\shortans{30}
\loigiai{
Ta có $2 \times 5 = 10$, rồi $10 \times 3 = 30$.
}
\end{ex}

\begin{ex}
Tính giá trị biểu thức: $20 : 2 : 5$.
\shortans{2}
\loigiai{
Ta có $20 : 2 = 10$, rồi $10 : 5 = 2$.
}
\end{ex}

\begin{ex}
Điền số thích hợp: $2 \times \dots = 20$.
\shortans{10}
\loigiai{
Số cần điền: $20 : 2 = 10$.
}
\end{ex}

\begin{ex}
Điền số thích hợp: $5 \times \dots = 45$.
\shortans{9}
\loigiai{
Số cần điền: $45 : 5 = 9$.
}
\end{ex}

\begin{ex}
Điền số thích hợp: $\dots : 2 = 7$.
\shortans{14}
\loigiai{
Số cần điền: $7 \times 2 = 14$.
}
\end{ex}

\begin{ex}
Điền số thích hợp: $\dots : 5 = 8$.
\shortans{40}
\loigiai{
Số cần điền: $8 \times 5 = 40$.
}
\end{ex}

\begin{ex}
Mỗi bạn có 2 quả bóng bay, 9 bạn có tất cả bao nhiêu quả bóng bay?
\shortans{18}
\loigiai{
Có: $2 \times 9 = 18$ quả bóng.
}
\end{ex}

\begin{ex}
Có 35 bông hoa chia đều vào các bó, mỗi bó 5 bông hoa. Hỏi bó được bao nhiêu bó hoa?
\shortans{7}
\loigiai{
Bó được: $35 : 5 = 7$ bó hoa.
}
\end{ex}

\begin{ex}
Bác nông dân có 16 lít mật ong rót đều vào 2 can. Mỗi can có bao nhiêu lít mật ong?
\shortans{8}
\loigiai{
Mỗi can có: $16 : 2 = 8$ lít.
}
\end{ex}

\begin{ex}
Một chuồng gà có 25 con gà chia đều vào 5 ngăn. Mỗi ngăn có bao nhiêu con gà?
\shortans{5}
\loigiai{
Mỗi ngăn có: $25 : 5 = 5$ con gà.
}
\end{ex}

\begin{ex}
Tìm một số biết lấy số đó nhân với 2 thì được 16.
\shortans{8}
\loigiai{
Số đó là: $16 : 2 = 8$.
}
\end{ex}

\begin{ex}
Tìm một số biết lấy số đó chia cho 5 thì được 6.
\shortans{30}
\loigiai{
Số đó là: $6 \times 5 = 30$.
}
\end{ex}

\begin{ex}
Tính: $5 \times 4 + 2 \times 5$.
\shortans{30}
\loigiai{
Ta có $5 \times 4 = 20$, $2 \times 5 = 10$. Tổng: $20 + 10 = 30$.
}
\end{ex}

\newpage


% =========================================================================
% Làm quen với hình khối
% =========================================================================
\section{Làm quen với hình khối}

\TaoPhieuBaiTap{Chủ đề 9 -- Làm quen với hình khối}{46}{Khối trụ, khối cầu}{https://loptoanthaylinh.vn/bai-tap/toan-2/bai-46}
\begin{ex}
Khối trụ có bao nhiêu mặt đáy phẳng hình tròn?
\shortans{2}
\loigiai{
Khối trụ có 2 mặt đáy phẳng hình tròn (đáy trên và đáy dưới).
}
\end{ex}

\begin{ex}
Mặt đáy của khối trụ có dạng hình gì? (Hình tròn). Số cạnh thẳng của hình tròn là bao nhiêu?
\shortans{0}
\loigiai{
Hình tròn không có cạnh thẳng nào.
}
\end{ex}

\begin{ex}
Quả bóng đá có dạng hình khối gì? (Khối cầu). Khối cầu có bao nhiêu đỉnh nhọn?
\shortans{0}
\loigiai{
Khối cầu tròn đều, không có đỉnh nhọn nào.
}
\end{ex}

\begin{ex}
Khối cầu có bao nhiêu cạnh thẳng?
\shortans{0}
\loigiai{
Khối cầu không có cạnh thẳng nào.
}
\end{ex}

\begin{ex}
Một lon nước ngọt có dạng khối gì? (Khối trụ). Lon nước ngọt có mấy mặt đáy?
\shortans{2}
\loigiai{
Lon nước ngọt có 2 mặt đáy.
}
\end{ex}

\begin{ex}
Quả bóng bàn có dạng khối gì? Khối cầu có thể lăn được về bao nhiêu hướng? (Mọi hướng). Số mặt phẳng của khối cầu là bao nhiêu?
\shortans{0}
\loigiai{
Khối cầu không có mặt phẳng nào.
}
\end{ex}

\begin{ex}
Một hộp sữa đặc có dạng khối trụ. Nếu đặt đứng, hộp tiếp xúc với mặt bàn bằng mấy mặt?
\shortans{1}
\loigiai{
Tiếp xúc bằng đúng 1 mặt đáy.
}
\end{ex}

\begin{ex}
Một khối rubik hình lập phương có 6 mặt phẳng. Một khối trụ có bao nhiêu mặt phẳng?
\shortans{2}
\loigiai{
Khối trụ có 2 mặt phẳng tròn ở đáy.
}
\end{ex}

\begin{ex}
Một viên bi ve có dạng khối gì? (Khối cầu). Viên bi ve có bao nhiêu đỉnh?
\shortans{0}
\loigiai{
Khối cầu có 0 đỉnh.
}
\end{ex}

\begin{ex}
Trong các vật sau: quả bóng, hộp sữa ông Thọ, viên xúc xắc, quả cam. Có bao nhiêu vật có dạng khối cầu?
\shortans{2}
\loigiai{
Quả bóng và quả cam có dạng khối cầu (2 vật).
}
\end{ex}

\begin{ex}
Trong các vật trên, có bao nhiêu vật có dạng khối trụ?
\shortans{1}
\loigiai{
Hộp sữa ông Thọ có dạng khối trụ (1 vật).
}
\end{ex}

\begin{ex}
Trong các vật trên, có bao nhiêu vật có dạng khối lập phương?
\shortans{1}
\loigiai{
Viên xúc xắc có dạng khối lập phương (1 vật).
}
\end{ex}

\begin{ex}
Trên bàn có 3 khối trụ và 4 khối cầu. Hỏi có tất cả bao nhiêu khối hình?
\shortans{7}
\loigiai{
Có: $3 + 4 = 7$ khối hình.
}
\end{ex}

\begin{ex}
Lan có 5 khối cầu, Mai có nhiều hơn Lan 3 khối cầu. Mai có bao nhiêu khối cầu?
\shortans{8}
\loigiai{
Mai có: $5 + 3 = 8$ khối cầu.
}
\end{ex}

\begin{ex}
Trong rổ đồ chơi có 10 khối hình gồm khối trụ và khối cầu, trong đó có 6 khối trụ. Hỏi có bao nhiêu khối cầu?
\shortans{4}
\loigiai{
Có: $10 - 6 = 4$ khối cầu.
}
\end{ex}

\begin{ex}
Dùng 2 khối trụ giống nhau đặt chồng khít lên nhau tạo thành một khối trụ cao hơn. Khối trụ mới có bao nhiêu mặt đáy phẳng ở hai đầu?
\shortans{2}
\loigiai{
Vẫn có 2 mặt đáy phẳng ở hai đầu ngoài.
}
\end{ex}

\begin{ex}
Một khúc gỗ tròn cắt ngang hai đầu tạo thành dạng khối gì? Khối trụ đó có bao nhiêu mặt đáy phẳng?
\shortans{2}
\loigiai{
Khối trụ có 2 mặt đáy phẳng.
}
\end{ex}

\begin{ex}
Một chiếc trống trường có thân cong tròn và hai mặt bịt da phẳng hình tròn. Chiếc trống có dạng giống khối gì? Hai mặt da là bao nhiêu mặt phẳng?
\shortans{2}
\loigiai{
Có 2 mặt phẳng.
}
\end{ex}

\begin{ex}
Một chiếc cột nhà hình tròn thẳng đứng có dạng khối gì? (Khối trụ). Cột nhà có bao nhiêu đỉnh nhọn?
\shortans{0}
\loigiai{
Khối trụ không có đỉnh nhọn nào.
}
\end{ex}

\begin{ex}
Có 8 khối cầu chia đều cho 2 bạn. Mỗi bạn được bao nhiêu khối cầu?
\shortans{4}
\loigiai{
Mỗi bạn được: $8 : 2 = 4$ khối cầu.
}
\end{ex}

\begin{ex}
Có 15 khối trụ chia đều vào 5 hộp. Mỗi hộp có bao nhiêu khối trụ?
\shortans{3}
\loigiai{
Mỗi hộp có: $15 : 5 = 3$ khối trụ.
}
\end{ex}

\begin{ex}
Trong một trò chơi xếp hình: bạn Nam dùng 4 khối trụ làm cột nhà, 2 khối cầu trang trí trên nóc. Nam đã dùng tất cả bao nhiêu khối hình?
\shortans{6}
\loigiai{
Nam dùng: $4 + 2 = 6$ khối.
}
\end{ex}

\begin{ex}
Một chiếc cốc uống nước hình trụ có một miệng hở và một đáy phẳng. Chiếc cốc có bao nhiêu đáy phẳng?
\shortans{1}
\loigiai{
Có 1 đáy phẳng.
}
\end{ex}

\begin{ex}
Khối nào sau đây có thể lăn được về mọi hướng? (Khối cầu). Khối cầu có bao nhiêu cạnh?
\shortans{0}
\loigiai{
Khối cầu có 0 cạnh.
}
\end{ex}

\begin{ex}
Tổng số mặt phẳng của 3 khối trụ rời nhau là bao nhiêu?
\shortans{6}
\loigiai{
Mỗi khối có 2 mặt phẳng, 3 khối có: $2 \times 3 = 6$ mặt phẳng.
}
\end{ex}

\newpage

\TaoPhieuBaiTap{Chủ đề 9 -- Làm quen với hình khối}{47}{Nhận biết hình khối thực tế}{https://loptoanthaylinh.vn/bai-tap/toan-2/bai-47}
\begin{ex}
Khối trụ có bao nhiêu mặt phẳng đáy?
\shortans{2}
\loigiai{
Có 2 mặt đáy phẳng.
}
\end{ex}

\begin{ex}
Khối cầu có bao nhiêu đỉnh?
\shortans{0}
\loigiai{
Có 0 đỉnh.
}
\end{ex}

\begin{ex}
Khối lập phương có bao nhiêu mặt vuông?
\shortans{6}
\loigiai{
Có 6 mặt vuông.
}
\end{ex}

\begin{ex}
Khối hộp chữ nhật có bao nhiêu đỉnh?
\shortans{8}
\loigiai{
Có 8 đỉnh.
}
\end{ex}

\begin{ex}
Khối hộp chữ nhật có bao nhiêu cạnh?
\shortans{12}
\loigiai{
Có 12 cạnh.
}
\end{ex}

\begin{ex}
Trong các khối: khối lập phương, khối hộp chữ nhật, khối trụ, khối cầu. Có bao nhiêu khối có đỉnh nhọn?
\shortans{2}
\loigiai{
Khối lập phương và khối hộp chữ nhật có đỉnh nhọn (2 khối).
}
\end{ex}

\begin{ex}
Trong các khối trên, có bao nhiêu khối có thể lăn được trên mặt bàn phẳng?
\shortans{2}
\loigiai{
Khối trụ và khối cầu có thể lăn được (2 khối).
}
\end{ex}

\begin{ex}
Trên giá đồ chơi có: 4 khối cầu, 3 khối trụ, 2 khối lập phương. Có tất cả bao nhiêu khối hình?
\shortans{9}
\loigiai{
Có tất cả: $4 + 3 + 2 = 9$ khối hình.
}
\end{ex}

\begin{ex}
Số khối cầu nhiều hơn số khối lập phương bao nhiêu khối?
\shortans{2}
\loigiai{
Nhiều hơn: $4 - 2 = 2$ khối.
}
\end{ex}

\begin{ex}
Xếp 3 khối lập phương chồng lên nhau được một cột cao. Xếp 2 khối trụ chồng lên nhau được một cột khác. Có tất cả bao nhiêu khối hình được dùng?
\shortans{5}
\loigiai{
Có: $3 + 2 = 5$ khối.
}
\end{ex}

\begin{ex}
Một người thợ tiện gỗ làm được 10 quả bóng gỗ dạng khối cầu. Bác chia đều vào 2 hộp. Mỗi hộp có bao nhiêu quả bóng?
\shortans{5}
\loigiai{
Mỗi hộp có: $10 : 2 = 5$ quả bóng.
}
\end{ex}

\begin{ex}
Một nhà máy sản xuất lon nước ngọt đóng thành từng lốc, mỗi lốc 6 lon hình khối trụ. Hỏi 3 lốc như thế có bao nhiêu lon nước ngọt?
\shortans{18}
\loigiai{
Có: $6 \times 3 = 18$ lon nước ngọt.
}
\end{ex}

\begin{ex}
Một hộp đựng bóng tennis dạng khối trụ chứa vừa vặn 3 quả bóng tennis dạng khối cầu. Hỏi 4 hộp như thế chứa được bao nhiêu quả bóng tennis?
\shortans{12}
\loigiai{
Chứa được: $3 \times 4 = 12$ quả bóng.
}
\end{ex}

\begin{ex}
Khối trụ có mấy mặt cong xung quanh?
\shortans{1}
\loigiai{
Khối trụ có 1 mặt cong xung quanh.
}
\end{ex}

\begin{ex}
Khối cầu có mấy mặt cong bao quanh?
\shortans{1}
\loigiai{
Khối cầu có 1 mặt cong khép kín bao quanh.
}
\end{ex}

\begin{ex}
Tổng số mặt đáy phẳng của 4 khối trụ là bao nhiêu?
\shortans{8}
\loigiai{
Có: $2 \times 4 = 8$ mặt đáy.
}
\end{ex}

\begin{ex}
Tổng số đỉnh của 2 khối hộp chữ nhật là bao nhiêu?
\shortans{16}
\loigiai{
Có: $8 + 8 = 16$ đỉnh.
}
\end{ex}

\begin{ex}
Tổng số mặt của 2 khối lập phương là bao nhiêu?
\shortans{12}
\loigiai{
Có: $6 + 6 = 12$ mặt.
}
\end{ex}

\begin{ex}
Một cây nến hình khối trụ có chiều cao 15 cm. Sau khi thắp sáng một lúc, cây nến bị cháy ngắn bớt 6 cm. Cây nến còn lại dài bao nhiêu xăng-ti-mét?
\shortans{9}
\loigiai{
Còn lại: $15 - 6 = 9$ cm.
}
\end{ex}

\begin{ex}
Bác thợ mộc có một khúc gỗ dài 20 dm dạng khối trụ. Bác cưa thành 4 khúc bằng nhau. Mỗi khúc dài bao nhiêu đề-xi-mét?
\shortans{5}
\loigiai{
Mỗi khúc dài: $20 : 4 = 5$ dm.
}
\end{ex}

\begin{ex}
Bạn Hùng có 14 viên bi ve khối cầu, bạn Dũng có ít hơn bạn Hùng 5 viên. Bạn Dũng có bao nhiêu viên bi ve?
\shortans{9}
\loigiai{
Dũng có: $14 - 5 = 9$ viên bi.
}
\end{ex}

\begin{ex}
Có bao nhiêu khối cầu trong một hộp chứa 8 quả bóng bàn và 6 chiếc vợt cầu lông?
\shortans{8}
\loigiai{
Có 8 quả bóng bàn dạng khối cầu.
}
\end{ex}

\begin{ex}
Xếp 4 khối trụ thành một hàng ngang, cần thêm bao nhiêu khối trụ nữa để có 10 khối trụ?
\shortans{6}
\loigiai{
Cần thêm: $10 - 4 = 6$ khối.
}
\end{ex}

\begin{ex}
Một chiếc bánh sinh nhật dạng khối trụ tròn có 2 tầng bằng nhau. Cả hai tầng có bao nhiêu mặt đáy phẳng tròn ngoài cùng?
\shortans{2}
\loigiai{
Có 2 mặt đáy phẳng ngoài cùng.
}
\end{ex}

\begin{ex}
Số cạnh của một khối cầu là bao nhiêu?
\shortans{0}
\loigiai{
Khối cầu có 0 cạnh.
}
\end{ex}

\newpage


% =========================================================================
% Các số trong phạm vi 1 000
% =========================================================================
\section{Các số trong phạm vi 1 000}

\TaoPhieuBaiTap{Chủ đề 10 -- Các số trong phạm vi 1 000}{48}{Đơn vị đo: Đề-xi-mét}{https://loptoanthaylinh.vn/bai-tap/toan-2/bai-48}
\begin{ex}
1 chục bằng bao nhiêu đơn vị?
\shortans{10}
\loigiai{
$1\text{ chục} = 10\text{ đơn vị}$.
}
\end{ex}

\begin{ex}
10 chục bằng một trăm. Một trăm viết bằng số nào?
\shortans{100}
\loigiai{
$10\text{ chục} = 100$.
}
\end{ex}

\begin{ex}
10 trăm bằng một nghìn. Một nghìn viết bằng chữ số 1 và mấy chữ số 0?
\shortans{3}
\loigiai{
Số 1000 có 3 chữ số 0.
}
\end{ex}

\begin{ex}
Số một nghìn viết bằng chữ số nào?
\shortans{1000}
\loigiai{
Một nghìn viết là 1000.
}
\end{ex}

\begin{ex}
Số 1000 là số có mấy chữ số?
\shortans{4}
\loigiai{
Số 1000 có 4 chữ số.
}
\end{ex}

\begin{ex}
Số bé nhất có bốn chữ số là số nào?
\shortans{1000}
\loigiai{
Số bé nhất có bốn chữ số là 1000.
}
\end{ex}

\begin{ex}
Số lớn nhất có ba chữ số là số nào?
\shortans{999}
\loigiai{
Số lớn nhất có ba chữ số là 999.
}
\end{ex}

\begin{ex}
Số liền sau của số 999 là số nào?
\shortans{1000}
\loigiai{
Số liền sau của 999 là 1000.
}
\end{ex}

\begin{ex}
Số liền trước của số 1000 là số nào?
\shortans{999}
\loigiai{
Số liền trước của 1000 là 999.
}
\end{ex}

\begin{ex}
Có bao nhiêu chục trong số 100?
\shortans{10}
\loigiai{
Có 10 chục.
}
\end{ex}

\begin{ex}
Có bao nhiêu trăm trong số 1000?
\shortans{10}
\loigiai{
Có 10 trăm.
}
\end{ex}

\begin{ex}
Số gồm 3 trăm và 0 chục, 0 đơn vị là số nào?
\shortans{300}
\loigiai{
Số đó là 300.
}
\end{ex}

\begin{ex}
Số gồm 7 trăm và 0 chục, 0 đơn vị là số nào?
\shortans{700}
\loigiai{
Số đó là 700.
}
\end{ex}

\begin{ex}
Đếm thêm 100: $100, 200, 300, \dots, 500$. Số còn thiếu là số nào?
\shortans{400}
\loigiai{
Số còn thiếu là 400.
}
\end{ex}

\begin{ex}
Đếm thêm 100: $600, 700, \dots, 900, 1000$. Số còn thiếu là số nào?
\shortans{800}
\loigiai{
Số còn thiếu là 800.
}
\end{ex}

\begin{ex}
Số liền sau của số 100 là số nào?
\shortans{101}
\loigiai{
Số liền sau của 100 là 101.
}
\end{ex}

\begin{ex}
Số liền trước của số 100 là số nào?
\shortans{99}
\loigiai{
Số liền trước của 100 là 99.
}
\end{ex}

\begin{ex}
Một xưởng sản xuất đóng 10 hộp bánh, mỗi hộp có 10 chiếc bánh. Xưởng đã đóng được tất cả bao nhiêu chiếc bánh?
\shortans{100}
\loigiai{
Có: $10 \times 10 = 100$ chiếc bánh.
}
\end{ex}

\begin{ex}
Một cửa hàng có 10 thùng sách, mỗi thùng có 100 cuốn sách. Cửa hàng có tất cả bao nhiêu cuốn sách?
\shortans{1000}
\loigiai{
Có: $100 \times 10 = 1000$ cuốn sách.
}
\end{ex}

\begin{ex}
Một nghìn trừ đi một trăm bằng bao nhiêu?
\shortans{900}
\loigiai{
Ta có: $1000 - 100 = 900$.
}
\end{ex}

\begin{ex}
Số 1000 gồm mấy chữ số khác nhau? (chữ số 1 và chữ số 0). Số lượng chữ số khác nhau là bao nhiêu?
\shortans{2}
\loigiai{
Có 2 chữ số khác nhau là 1 và 0.
}
\end{ex}

\begin{ex}
Tính: $300 + 200$.
\shortans{500}
\loigiai{
Ta có $300 + 200 = 500$.
}
\end{ex}

\begin{ex}
Tính: $800 - 500$.
\shortans{300}
\loigiai{
Ta có $800 - 500 = 300$.
}
\end{ex}

\begin{ex}
Số nào cộng với 600 thì được 1000?
\shortans{400}
\loigiai{
Số đó là: $1000 - 600 = 400$.
}
\end{ex}

\begin{ex}
Số nào trừ đi 300 thì được 700?
\shortans{1000}
\loigiai{
Số đó là: $700 + 300 = 1000$.
}
\end{ex}

\newpage

\TaoPhieuBaiTap{Chủ đề 10 -- Các số trong phạm vi 1 000}{49}{Đơn vị đo: Mét}{https://loptoanthaylinh.vn/bai-tap/toan-2/bai-49}
\begin{ex}
Số tròn trăm liền sau số 200 là số nào?
\shortans{300}
\loigiai{
Số tròn trăm liền sau 200 là 300.
}
\end{ex}

\begin{ex}
Số tròn trăm liền trước số 600 là số nào?
\shortans{500}
\loigiai{
Số tròn trăm liền trước 600 là 500.
}
\end{ex}

\begin{ex}
Số tròn trăm lớn nhất có ba chữ số là số nào?
\shortans{900}
\loigiai{
Số đó là 900.
}
\end{ex}

\begin{ex}
Số tròn trăm bé nhất có ba chữ số là số nào?
\shortans{100}
\loigiai{
Số đó là 100.
}
\end{ex}

\begin{ex}
Có bao nhiêu số tròn trăm có ba chữ số? (100, 200, ..., 900).
\shortans{9}
\loigiai{
Có đúng 9 số.
}
\end{ex}

\begin{ex}
Tính: $200 + 300$.
\shortans{500}
\loigiai{
Ta có $200 + 300 = 500$.
}
\end{ex}

\begin{ex}
Tính: $400 + 500$.
\shortans{900}
\loigiai{
Ta có $400 + 500 = 900$.
}
\end{ex}

\begin{ex}
Tính: $700 - 200$.
\shortans{500}
\loigiai{
Ta có $700 - 200 = 500$.
}
\end{ex}

\begin{ex}
Tính: $900 - 600$.
\shortans{300}
\loigiai{
Ta có $900 - 600 = 300$.
}
\end{ex}

\begin{ex}
Tính: $1000 - 400$.
\shortans{600}
\loigiai{
Ta có $1000 - 400 = 600$.
}
\end{ex}

\begin{ex}
Tính giá trị biểu thức: $300 + 400 - 200$.
\shortans{500}
\loigiai{
Ta có $300 + 400 = 700$, rồi $700 - 200 = 500$.
}
\end{ex}

\begin{ex}
Tính giá trị biểu thức: $800 - 500 + 400$.
\shortans{700}
\loigiai{
Ta có $800 - 500 = 300$, rồi $300 + 400 = 700$.
}
\end{ex}

\begin{ex}
Tính giá trị biểu thức: $1000 - 300 - 200$.
\shortans{500}
\loigiai{
Ta có $1000 - 300 = 700$, rồi $700 - 200 = 500$.
}
\end{ex}

\begin{ex}
Số tròn chục liền sau số 140 là số nào?
\shortans{150}
\loigiai{
Số tròn chục liền sau 140 là 150.
}
\end{ex}

\begin{ex}
Số tròn chục liền trước số 280 là số nào?
\shortans{270}
\loigiai{
Số tròn chục liền trước 280 là 270.
}
\end{ex}

\begin{ex}
Điền số thích hợp: $400 + \dots = 700$.
\shortans{300}
\loigiai{
Số cần điền: $700 - 400 = 300$.
}
\end{ex}

\begin{ex}
Điền số thích hợp: $800 - \dots = 300$.
\shortans{500}
\loigiai{
Số cần điền: $800 - 300 = 500$.
}
\end{ex}

\begin{ex}
Điền số thích hợp: $\dots + 500 = 1000$.
\shortans{500}
\loigiai{
Số cần điền: $1000 - 500 = 500$.
}
\end{ex}

\begin{ex}
Điền số thích hợp: $\dots - 200 = 600$.
\shortans{800}
\loigiai{
Số cần điền: $600 + 200 = 800$.
}
\end{ex}

\begin{ex}
Khối lớp Hai có 200 học sinh, khối lớp Ba có 300 học sinh. Cả hai khối có bao nhiêu học sinh?
\shortans{500}
\loigiai{
Cả hai khối có: $200 + 300 = 500$ học sinh.
}
\end{ex}

\begin{ex}
Một cửa hàng buổi sáng bán được 400 kg gạo, buổi chiều bán được 500 kg gạo. Cả ngày cửa hàng bán được bao nhiêu ki-lô-gam gạo?
\shortans{900}
\loigiai{
Cả ngày bán được: $400 + 500 = 900$ kg.
}
\end{ex}

\begin{ex}
Kho thóc có 800 bao thóc, người ta chuyển đi 300 bao thóc. Trong kho còn lại bao nhiêu bao thóc?
\shortans{500}
\loigiai{
Trong kho còn lại: $800 - 300 = 500$ bao.
}
\end{ex}

\begin{ex}
Vườn cây có 600 cây ăn quả, trong đó có 200 cây cam, còn lại là cây xoài. Vườn có bao nhiêu cây xoài?
\shortans{400}
\loigiai{
Số cây xoài: $600 - 200 = 400$ cây.
}
\end{ex}

\begin{ex}
Tổng của số tròn trăm lớn nhất có ba chữ số và số tròn trăm bé nhất có ba chữ số là bao nhiêu?
\shortans{1000}
\loigiai{
Tổng: $900 + 100 = 1000$.
}
\end{ex}

\begin{ex}
Hiệu của số tròn trăm lớn nhất có ba chữ số và số tròn trăm bé nhất có ba chữ số là bao nhiêu?
\shortans{800}
\loigiai{
Hiệu: $900 - 100 = 800$.
}
\end{ex}

\newpage

\TaoPhieuBaiTap{Chủ đề 10 -- Các số trong phạm vi 1 000}{50}{Đơn vị đo: Ki-lô-mét}{https://loptoanthaylinh.vn/bai-tap/toan-2/bai-50}
\begin{ex}
Trong hai số 300 và 500, số nào lớn hơn?
\shortans{500}
\loigiai{
Vì chữ số hàng trăm $5 > 3$ nên $500 > 300$.
}
\end{ex}

\begin{ex}
Trong hai số 700 và 400, số nào bé hơn?
\shortans{400}
\loigiai{
Số bé hơn là 400.
}
\end{ex}

\begin{ex}
Trong hai số 800 và 1000, số nào lớn hơn?
\shortans{1000}
\loigiai{
Số 1000 có 4 chữ số, lớn hơn 800 có 3 chữ số.
}
\end{ex}

\begin{ex}
Trong các số $200, 700, 400, 900$, số nào lớn nhất?
\shortans{900}
\loigiai{
Số lớn nhất là 900.
}
\end{ex}

\begin{ex}
Trong các số $600, 100, 800, 300$, số nào bé nhất?
\shortans{100}
\loigiai{
Số bé nhất là 100.
}
\end{ex}

\begin{ex}
So sánh hai số tròn chục: 240 và 270. Số nào lớn hơn?
\shortans{270}
\loigiai{
Số lớn hơn là 270.
}
\end{ex}

\begin{ex}
So sánh hai số tròn chục: 350 và 310. Số nào bé hơn?
\shortans{310}
\loigiai{
Số bé hơn là 310.
}
\end{ex}

\begin{ex}
Trong các số $420, 480, 450, 410$, số nào lớn nhất?
\shortans{480}
\loigiai{
Số lớn nhất là 480.
}
\end{ex}

\begin{ex}
Trong các số $670, 620, 690, 650$, số nào bé nhất?
\shortans{620}
\loigiai{
Số bé nhất là 620.
}
\end{ex}

\begin{ex}
Có bao nhiêu số tròn trăm lớn hơn 300 và bé hơn 800?
\shortans{4}
\loigiai{
Các số: 400, 500, 600, 700. Có 4 số.
}
\end{ex}

\begin{ex}
Có bao nhiêu số tròn chục lớn hơn 120 và bé hơn 170?
\shortans{4}
\loigiai{
Các số: 130, 140, 150, 160. Có 4 số.
}
\end{ex}

\begin{ex}
Sắp xếp dãy số $400, 800, 200, 600$ từ bé đến lớn, số đứng ở vị trí thứ ba là số nào?
\shortans{600}
\loigiai{
Thứ tự: 200, 400, 600, 800. Số thứ ba là 600.
}
\end{ex}

\begin{ex}
Sắp xếp dãy số $900, 300, 700, 500$ từ lớn đến bé, số đứng ở vị trí thứ hai là số nào?
\shortans{700}
\loigiai{
Thứ tự: 900, 700, 500, 300. Số thứ hai là 700.
}
\end{ex}

\begin{ex}
Điền số tròn trăm thích hợp: $400 < \dots < 600$.
\shortans{500}
\loigiai{
Số tròn trăm ở giữa là 500.
}
\end{ex}

\begin{ex}
Điền số tròn chục thích hợp: $340 < \dots < 360$.
\shortans{350}
\loigiai{
Số tròn chục ở giữa là 350.
}
\end{ex}

\begin{ex}
Trường Tiểu học Kim Đồng có 500 học sinh, trường Tiểu học Lê Lợi có 600 học sinh. Trường nào có nhiều học sinh hơn? Trường Lê Lợi nhiều hơn trường Kim Đồng bao nhiêu học sinh?
\shortans{100}
\loigiai{
Nhiều hơn: $600 - 500 = 100$ học sinh.
}
\end{ex}

\begin{ex}
Kho A có 700 bao xi măng, kho B có 400 bao xi măng. Kho A nhiều hơn kho B bao nhiêu bao xi măng?
\shortans{300}
\loigiai{
Nhiều hơn: $700 - 400 = 300$ bao.
}
\end{ex}

\begin{ex}
Một đoạn đường dài 800 mét, đoạn đường thứ hai dài 500 mét. Đoạn thứ hai ngắn hơn đoạn thứ nhất bao nhiêu mét?
\shortans{300}
\loigiai{
Ngắn hơn: $800 - 500 = 300$ mét.
}
\end{ex}

\begin{ex}
Tìm số tròn trăm lớn nhất bé hơn 750.
\shortans{700}
\loigiai{
Số tròn trăm lớn nhất bé hơn 750 là 700.
}
\end{ex}

\begin{ex}
Tìm số tròn trăm bé nhất lớn hơn 420.
\shortans{500}
\loigiai{
Số tròn trăm bé nhất lớn hơn 420 là 500.
}
\end{ex}

\begin{ex}
Có bao nhiêu số tròn trăm từ 100 đến 1000?
\shortans{10}
\loigiai{
Có 10 số (100, 200, ..., 1000).
}
\end{ex}

\begin{ex}
Trong các số tròn chục: $150, 240, 190, 310$, có bao nhiêu số lớn hơn 200?
\shortans{2}
\loigiai{
Các số lớn hơn 200 là 240 và 310. Có 2 số.
}
\end{ex}

\begin{ex}
Trong các số tròn chục: $420, 380, 510, 490$, có bao nhiêu số bé hơn 500?
\shortans{3}
\loigiai{
Các số: 420, 380, 490. Có 3 số.
}
\end{ex}

\begin{ex}
Hiệu của số tròn trăm liền sau 600 và số tròn trăm liền trước 600 là bao nhiêu?
\shortans{200}
\loigiai{
Số liền sau là 700, liền trước là 500. Hiệu: $700 - 500 = 200$.
}
\end{ex}

\begin{ex}
Tìm số tròn chục ở giữa hai số 580 và 600.
\shortans{590}
\loigiai{
Số ở giữa là 590.
}
\end{ex}

\newpage

\TaoPhieuBaiTap{Chủ đề 10 -- Các số trong phạm vi 1 000}{51}{Đơn vị đo: Mi-li-mét}{https://loptoanthaylinh.vn/bai-tap/toan-2/bai-51}
\begin{ex}
Số gồm 2 trăm, 3 chục và 5 đơn vị là số nào?
\shortans{235}
\loigiai{
Số đó là 235.
}
\end{ex}

\begin{ex}
Số gồm 4 trăm, 8 chục và 1 đơn vị là số nào?
\shortans{481}
\loigiai{
Số đó là 481.
}
\end{ex}

\begin{ex}
Số gồm 5 trăm, 0 chục và 7 đơn vị là số nào?
\shortans{507}
\loigiai{
Số đó là 507.
}
\end{ex}

\begin{ex}
Số gồm 6 trăm, 4 chục và 0 đơn vị là số nào?
\shortans{640}
\loigiai{
Số đó là 640.
}
\end{ex}

\begin{ex}
Số 358 có chữ số hàng trăm là bao nhiêu?
\shortans{3}
\loigiai{
Chữ số hàng trăm là 3.
}
\end{ex}

\begin{ex}
Số 724 có chữ số hàng chục là bao nhiêu?
\shortans{2}
\loigiai{
Chữ số hàng chục là 2.
}
\end{ex}

\begin{ex}
Số 896 có chữ số hàng đơn vị là bao nhiêu?
\shortans{6}
\loigiai{
Chữ số hàng đơn vị là 6.
}
\end{ex}

\begin{ex}
Số liền sau của số 123 là số nào?
\shortans{124}
\loigiai{
Số liền sau 123 là 124.
}
\end{ex}

\begin{ex}
Số liền sau của số 259 là số nào?
\shortans{260}
\loigiai{
Số liền sau 259 là 260.
}
\end{ex}

\begin{ex}
Số liền sau của số 399 là số nào?
\shortans{400}
\loigiai{
Số liền sau 399 là 400.
}
\end{ex}

\begin{ex}
Số liền trước của số 450 là số nào?
\shortans{449}
\loigiai{
Số liền trước 450 là 449.
}
\end{ex}

\begin{ex}
Số liền trước của số 600 là số nào?
\shortans{599}
\loigiai{
Số liền trước 600 là 599.
}
\end{ex}

\begin{ex}
Số liền trước của số 811 là số nào?
\shortans{810}
\loigiai{
Số liền trước 811 là 810.
}
\end{ex}

\begin{ex}
Số bé nhất có ba chữ số là số nào?
\shortans{100}
\loigiai{
Số bé nhất có ba chữ số là 100.
}
\end{ex}

\begin{ex}
Số lớn nhất có ba chữ số là số nào?
\shortans{999}
\loigiai{
Số lớn nhất có ba chữ số là 999.
}
\end{ex}

\begin{ex}
Số bé nhất có ba chữ số khác nhau là số nào?
\shortans{102}
\loigiai{
Số bé nhất có ba chữ số khác nhau là 102.
}
\end{ex}

\begin{ex}
Số lớn nhất có ba chữ số khác nhau là số nào?
\shortans{987}
\loigiai{
Số lớn nhất có ba chữ số khác nhau là 987.
}
\end{ex}

\begin{ex}
Số gồm 9 trăm, 9 chục và 9 đơn vị là số nào?
\shortans{999}
\loigiai{
Số đó là 999.
}
\end{ex}

\begin{ex}
Điền số thích hợp: $345, 346, \dots, 348, 349$.
\shortans{347}
\loigiai{
Số còn thiếu là 347.
}
\end{ex}

\begin{ex}
Điền số thích hợp: $578, 579, \dots, 581$.
\shortans{580}
\loigiai{
Số còn thiếu là 580.
}
\end{ex}

\begin{ex}
Từ 100 đến 110 có tất cả bao nhiêu số tự nhiên?
\shortans{11}
\loigiai{
Có $(110 - 100) + 1 = 11$ số.
}
\end{ex}

\begin{ex}
Có bao nhiêu số có ba chữ số giống nhau? (111, 222, ..., 999).
\shortans{9}
\loigiai{
Có đúng 9 số.
}
\end{ex}

\begin{ex}
Tính tổng các chữ số của số 243.
\shortans{9}
\loigiai{
Tổng là: $2 + 4 + 3 = 9$.
}
\end{ex}

\begin{ex}
Tính tổng các chữ số của số 512.
\shortans{8}
\loigiai{
Tổng là: $5 + 1 + 2 = 8$.
}
\end{ex}

\begin{ex}
Tìm số ở giữa hai số 789 và 791.
\shortans{790}
\loigiai{
Số ở giữa là 790.
}
\end{ex}

\newpage

\TaoPhieuBaiTap{Chủ đề 10 -- Các số trong phạm vi 1 000}{52}{Thực hành đo và ước lượng độ dài}{https://loptoanthaylinh.vn/bai-tap/toan-2/bai-52}
\begin{ex}
Viết số 345 thành tổng: $300 + 40 + \dots$ Số còn thiếu là bao nhiêu?
\shortans{5}
\loigiai{
Số 345 gồm 300, 40 và 5.
}
\end{ex}

\begin{ex}
Viết số 567 thành tổng: $500 + \dots + 7$. Số còn thiếu là bao nhiêu?
\shortans{60}
\loigiai{
Số 567 gồm 500, 60 và 7.
}
\end{ex}

\begin{ex}
Viết số 728 thành tổng: $\dots + 20 + 8$. Số còn thiếu là bao nhiêu?
\shortans{700}
\loigiai{
Số 728 gồm 700, 20 và 8.
}
\end{ex}

\begin{ex}
Số 408 viết thành tổng các trăm và đơn vị là $400 + \dots$ Số còn thiếu là bao nhiêu?
\shortans{8}
\loigiai{
$408 = 400 + 8$.
}
\end{ex}

\begin{ex}
Số 690 viết thành tổng các trăm và chục là $600 + \dots$ Số còn thiếu là bao nhiêu?
\shortans{90}
\loigiai{
$690 = 600 + 90$.
}
\end{ex}

\begin{ex}
Tính tổng: $200 + 50 + 6$.
\shortans{256}
\loigiai{
Ta có: $200 + 50 + 6 = 256$.
}
\end{ex}

\begin{ex}
Tính tổng: $400 + 80 + 3$.
\shortans{483}
\loigiai{
Ta có: $400 + 80 + 3 = 483$.
}
\end{ex}

\begin{ex}
Tính tổng: $600 + 70 + 9$.
\shortans{679}
\loigiai{
Ta có: $600 + 70 + 9 = 679$.
}
\end{ex}

\begin{ex}
Tính tổng: $800 + 30$.
\shortans{830}
\loigiai{
Ta có: $800 + 30 = 830$.
}
\end{ex}

\begin{ex}
Tính tổng: $900 + 5$.
\shortans{905}
\loigiai{
Ta có: $900 + 5 = 905$.
}
\end{ex}

\begin{ex}
Tính tổng: $700 + 40 + 2$.
\shortans{742}
\loigiai{
Ta có: $700 + 40 + 2 = 742$.
}
\end{ex}

\begin{ex}
Tính tổng: $100 + 90 + 9$.
\shortans{199}
\loigiai{
Ta có: $100 + 90 + 9 = 199$.
}
\end{ex}

\begin{ex}
Tính tổng: $300 + 8$.
\shortans{308}
\loigiai{
Ta có: $300 + 8 = 308$.
}
\end{ex}

\begin{ex}
Tính tổng: $500 + 50$.
\shortans{550}
\loigiai{
Ta có: $500 + 50 = 550$.
}
\end{ex}

\begin{ex}
Một số có tổng các trăm, chục, đơn vị là $600 + 30 + 4$. Số đó là số nào?
\shortans{634}
\loigiai{
Số đó là 634.
}
\end{ex}

\begin{ex}
Một số có tổng các trăm, chục, đơn vị là $800 + 0 + 9$. Số đó là số nào?
\shortans{809}
\loigiai{
Số đó là 809.
}
\end{ex}

\begin{ex}
Số 888 viết thành tổng: $800 + 80 + \dots$ Số cần điền là bao nhiêu?
\shortans{8}
\loigiai{
Số cần điền là 8.
}
\end{ex}

\begin{ex}
Giá trị của chữ số 5 trong số 524 là bao nhiêu?
\shortans{500}
\loigiai{
Chữ số 5 ở hàng trăm có giá trị là 500.
}
\end{ex}

\begin{ex}
Giá trị của chữ số 7 trong số 376 là bao nhiêu?
\shortans{70}
\loigiai{
Chữ số 7 ở hàng chục có giá trị là 70.
}
\end{ex}

\begin{ex}
Giá trị của chữ số 9 trong số 819 là bao nhiêu?
\shortans{9}
\loigiai{
Chữ số 9 ở hàng đơn vị có giá trị là 9.
}
\end{ex}

\begin{ex}
Hiệu giữa giá trị chữ số 6 ở hàng trăm và chữ số 6 ở hàng chục trong số 660 là bao nhiêu?
\shortans{540}
\loigiai{
Giá trị: $600 - 60 = 540$.
}
\end{ex}

\begin{ex}
Một số có chữ số hàng trăm là 4, hàng chục gấp đôi hàng trăm (bằng 8), hàng đơn vị bằng 1. Số đó là số nào?
\shortans{481}
\loigiai{
Số đó là 481.
}
\end{ex}

\begin{ex}
Một số có chữ số hàng trăm là 5, hàng chục là 0, hàng đơn vị lớn nhất có một chữ số (bằng 9). Số đó là số nào?
\shortans{509}
\loigiai{
Số đó là 509.
}
\end{ex}

\begin{ex}
Điền số thích hợp: $754 - 54 = \dots$
\shortans{700}
\loigiai{
Số còn lại là 700.
}
\end{ex}

\begin{ex}
Điền số thích hợp: $836 - 800 = \dots$
\shortans{36}
\loigiai{
Số còn lại là 36.
}
\end{ex}

\newpage

\TaoPhieuBaiTap{Chủ đề 10 -- Các số trong phạm vi 1 000}{53}{Tiền Việt Nam}{https://loptoanthaylinh.vn/bai-tap/toan-2/bai-53}
\begin{ex}
Trong hai số 345 và 435, số nào lớn hơn?
\shortans{435}
\loigiai{
So sánh hàng trăm $4 > 3$ nên $435 > 345$.
}
\end{ex}

\begin{ex}
Trong hai số 562 và 581, số nào bé hơn?
\shortans{562}
\loigiai{
Cùng hàng trăm là 5, so sánh hàng chục $6 < 8$ nên số bé hơn là 562.
}
\end{ex}

\begin{ex}
Trong hai số 748 và 743, số nào lớn hơn?
\shortans{748}
\loigiai{
Cùng hàng trăm và chục, so sánh hàng đơn vị $8 > 3$ nên số lớn hơn là 748.
}
\end{ex}

\begin{ex}
Trong các số $452, 618, 389, 524$, số nào lớn nhất?
\shortans{618}
\loigiai{
Số có hàng trăm lớn nhất là 618.
}
\end{ex}

\begin{ex}
Trong các số $735, 719, 782, 706$, số nào bé nhất?
\shortans{706}
\loigiai{
Số có hàng chục bé nhất là 706.
}
\end{ex}

\begin{ex}
Trong các số $831, 839, 834, 830$, số nào lớn nhất?
\shortans{839}
\loigiai{
Số có hàng đơn vị lớn nhất là 839.
}
\end{ex}

\begin{ex}
Có bao nhiêu số tự nhiên lớn hơn 245 và bé hơn 250?
\shortans{4}
\loigiai{
Các số: 246, 247, 248, 249. Có 4 số.
}
\end{ex}

\begin{ex}
Sắp xếp dãy số $324, 185, 460, 279$ từ bé đến lớn, số đứng ở vị trí thứ hai là số nào?
\shortans{279}
\loigiai{
Thứ tự: 185, 279, 324, 460. Số thứ hai là 279.
}
\end{ex}

\begin{ex}
Sắp xếp dãy số $672, 815, 549, 730$ từ lớn đến bé, số đứng ở vị trí thứ ba là số nào?
\shortans{672}
\loigiai{
Thứ tự: 815, 730, 672, 549. Số thứ ba là 672.
}
\end{ex}

\begin{ex}
Điền số tự nhiên thích hợp: $518 < \dots < 520$.
\shortans{519}
\loigiai{
Số ở giữa là 519.
}
\end{ex}

\begin{ex}
Tìm số lớn nhất trong các số: $908, 899, 910, 890$.
\shortans{910}
\loigiai{
Số lớn nhất là 910.
}
\end{ex}

\begin{ex}
Tìm số bé nhất trong các số: $654, 645, 564, 546$.
\shortans{546}
\loigiai{
Số bé nhất là 546.
}
\end{ex}

\begin{ex}
Số 500 lớn hơn số 450 bao nhiêu đơn vị?
\shortans{50}
\loigiai{
Lớn hơn: $500 - 450 = 50$ đơn vị.
}
\end{ex}

\begin{ex}
Số 780 bé hơn số 800 bao nhiêu đơn vị?
\shortans{20}
\loigiai{
Bé hơn: $800 - 780 = 20$ đơn vị.
}
\end{ex}

\begin{ex}
Trong các số $320, 450, 290, 510$, có bao nhiêu số lớn hơn 400?
\shortans{2}
\loigiai{
Các số lớn hơn 400 là 450 và 510 (2 số).
}
\end{ex}

\begin{ex}
Trong các số $610, 580, 720, 690$, có bao nhiêu số bé hơn 700?
\shortans{3}
\loigiai{
Các số bé hơn 700 là 610, 580, 690 (3 số).
}
\end{ex}

\begin{ex}
Trường Tiểu học A có 485 học sinh, trường Tiểu học B có 512 học sinh. Trường B nhiều hơn trường A bao nhiêu học sinh?
\shortans{27}
\loigiai{
Nhiều hơn: $512 - 485 = 27$ học sinh.
}
\end{ex}

\begin{ex}
Đoạn đường thứ nhất dài 350 m, đoạn thứ hai dài 420 m. Đoạn thứ hai dài hơn đoạn thứ nhất bao nhiêu mét?
\shortans{70}
\loigiai{
Dài hơn: $420 - 350 = 70$ m.
}
\end{ex}

\begin{ex}
Kho thóc thứ nhất có 650 bao, kho thứ hai có 600 bao. Kho thứ nhất nhiều hơn kho thứ hai bao nhiêu bao thóc?
\shortans{50}
\loigiai{
Nhiều hơn: $650 - 600 = 50$ bao.
}
\end{ex}

\begin{ex}
Tìm số liền trước của số bé nhất có ba chữ số khác nhau (số đó là 102).
\shortans{101}
\loigiai{
Số bé nhất có ba chữ số khác nhau là 102. Số liền trước là 101.
}
\end{ex}

\begin{ex}
Tìm số liền sau của số lớn nhất có ba chữ số (số đó là 999).
\shortans{1000}
\loigiai{
Số liền sau là 1000.
}
\end{ex}

\begin{ex}
So sánh: $300 + 40$ và $300 + 50$. Hiệu của hai tổng là bao nhiêu?
\shortans{10}
\loigiai{
Hiệu: $350 - 340 = 10$.
}
\end{ex}

\begin{ex}
So sánh: $500 + 80 + 2$ và $500 + 70 + 9$. Hai số này hơn kém nhau bao nhiêu đơn vị?
\shortans{3}
\loigiai{
Số thứ nhất là 582, số thứ hai là 579. Hơn kém: $582 - 579 = 3$ đơn vị.
}
\end{ex}

\begin{ex}
Cho 3 chữ số 2, 5, 8. Viết được số lớn nhất có ba chữ số khác nhau từ 3 chữ số này là số nào?
\shortans{852}
\loigiai{
Số lớn nhất là 852.
}
\end{ex}

\begin{ex}
Viết được số bé nhất có ba chữ số khác nhau từ 3 chữ số 2, 5, 8 là số nào?
\shortans{258}
\loigiai{
Số bé nhất là 258.
}
\end{ex}

\newpage

\TaoPhieuBaiTap{Chủ đề 10 -- Các số trong phạm vi 1 000}{54}{Luyện tập chung đại lượng}{https://loptoanthaylinh.vn/bai-tap/toan-2/bai-54}
\begin{ex}
Số liền sau của số 599 là số nào?
\shortans{600}
\loigiai{
Số liền sau 599 là 600.
}
\end{ex}

\begin{ex}
Số liền trước của số 700 là số nào?
\shortans{699}
\loigiai{
Số liền trước 700 là 699.
}
\end{ex}

\begin{ex}
Số gồm 6 trăm, 5 chục và 8 đơn vị là số nào?
\shortans{658}
\loigiai{
Số đó là 658.
}
\end{ex}

\begin{ex}
Số 904 gồm 9 trăm và mấy đơn vị?
\shortans{4}
\loigiai{
Số 904 gồm 9 trăm và 4 đơn vị.
}
\end{ex}

\begin{ex}
Viết số 472 thành tổng: $400 + \dots + 2$. Số cần điền là bao nhiêu?
\shortans{70}
\loigiai{
Số cần điền là 70.
}
\end{ex}

\begin{ex}
Tính: $300 + 500$.
\shortans{800}
\loigiai{
Ta có $300 + 500 = 800$.
}
\end{ex}

\begin{ex}
Tính: $900 - 400$.
\shortans{500}
\loigiai{
Ta có $900 - 400 = 500$.
}
\end{ex}

\begin{ex}
Tính: $250 + 50$.
\shortans{300}
\loigiai{
Ta có $250 + 50 = 300$.
}
\end{ex}

\begin{ex}
Tính: $400 - 80$.
\shortans{320}
\loigiai{
Ta có $400 - 80 = 320$.
}
\end{ex}

\begin{ex}
Trong các số $582, 614, 599, 608$, số nào lớn nhất?
\shortans{614}
\loigiai{
Số lớn nhất là 614.
}
\end{ex}

\begin{ex}
Trong các số $435, 428, 450, 419$, số nào bé nhất?
\shortans{419}
\loigiai{
Số bé nhất là 419.
}
\end{ex}

\begin{ex}
Sắp xếp dãy số $712, 589, 830, 645$ từ bé đến lớn, số đứng ở vị trí thứ ba là số nào?
\shortans{712}
\loigiai{
Thứ tự: 589, 645, 712, 830. Số thứ ba là 712.
}
\end{ex}

\begin{ex}
Điền số thích hợp: $789 < \dots < 791$.
\shortans{790}
\loigiai{
Số ở giữa là 790.
}
\end{ex}

\begin{ex}
Điền số thích hợp: $500 + \dots = 900$.
\shortans{400}
\loigiai{
Số cần điền là: $900 - 500 = 400$.
}
\end{ex}

\begin{ex}
Điền số thích hợp: $1000 - \dots = 300$.
\shortans{700}
\loigiai{
Số cần điền là: $1000 - 300 = 700$.
}
\end{ex}

\begin{ex}
Tính giá trị biểu thức: $400 + 300 - 200$.
\shortans{500}
\loigiai{
Ta có $400 + 300 = 700$, rồi $700 - 200 = 500$.
}
\end{ex}

\begin{ex}
Tính giá trị biểu thức: $800 - 600 + 500$.
\shortans{700}
\loigiai{
Ta có $800 - 600 = 200$, rồi $200 + 500 = 700$.
}
\end{ex}

\begin{ex}
Bác nông dân thu hoạch được 500 bắp ngô nếp và 400 bắp ngô ngọt. Bác thu hoạch được tất cả bao nhiêu bắp ngô?
\shortans{900}
\loigiai{
Bác thu hoạch được: $500 + 400 = 900$ bắp ngô.
}
\end{ex}

\begin{ex}
Một đoàn tàu chở 600 hành khách, đến ga có 200 hành khách xuống tàu. Đoàn tàu còn lại bao nhiêu hành khách?
\shortans{400}
\loigiai{
Đoàn tàu còn lại: $600 - 200 = 400$ hành khách.
}
\end{ex}

\begin{ex}
Kho gạo thứ nhất có 350 bao gạo, kho thứ hai có nhiều hơn kho thứ nhất 50 bao gạo. Kho thứ hai có bao nhiêu bao gạo?
\shortans{400}
\loigiai{
Kho thứ hai có: $350 + 50 = 400$ bao gạo.
}
\end{ex}

\begin{ex}
Tìm số lớn nhất có ba chữ số mà chữ số hàng trăm là 7.
\shortans{799}
\loigiai{
Số lớn nhất là 799.
}
\end{ex}

\begin{ex}
Tìm số bé nhất có ba chữ số mà chữ số hàng trăm là 4.
\shortans{400}
\loigiai{
Số bé nhất là 400.
}
\end{ex}

\begin{ex}
Số 1000 trừ đi số lớn nhất có ba chữ số thì được kết quả bằng bao nhiêu?
\shortans{1}
\loigiai{
Số lớn nhất có 3 chữ số là 999. Hiệu: $1000 - 999 = 1$.
}
\end{ex}

\begin{ex}
Có bao nhiêu số tròn trăm từ 200 đến 800?
\shortans{7}
\loigiai{
Các số: 200, 300, 400, 500, 600, 700, 800. Có 7 số.
}
\end{ex}

\begin{ex}
Tính tổng các chữ số của số bé nhất có ba chữ số khác nhau (102).
\shortans{3}
\loigiai{
Số 102 có tổng các chữ số: $1 + 0 + 2 = 3$.
}
\end{ex}

\newpage


% =========================================================================
% Độ dài và đơn vị đo độ dài. tiền việt nam
% =========================================================================
\section{Độ dài và đơn vị đo độ dài. tiền việt nam}

\TaoPhieuBaiTap{Chủ đề 11 -- Độ dài và đơn vị đo độ dài. tiền việt nam}{55}{Đơn vị, chục, trăm, nghìn}{https://loptoanthaylinh.vn/bai-tap/toan-2/bai-55}
\begin{ex}
1 mét bằng bao nhiêu đề-xi-mét?
\shortans{10}
\loigiai{
$1\text{ m} = 10\text{ dm}$.
}
\end{ex}

\begin{ex}
1 mét bằng bao nhiêu xăng-ti-mét?
\shortans{100}
\loigiai{
$1\text{ m} = 100\text{ cm}$.
}
\end{ex}

\begin{ex}
1 đề-xi-mét bằng bao nhiêu xăng-ti-mét?
\shortans{10}
\loigiai{
$1\text{ dm} = 10\text{ cm}$.
}
\end{ex}

\begin{ex}
1 ki-lô-mét bằng bao nhiêu mét?
\shortans{1000}
\loigiai{
$1\text{ km} = 1000\text{ m}$.
}
\end{ex}

\begin{ex}
Tính: $35\text{ m} + 20\text{ m}$. Kết quả là bao nhiêu mét?
\shortans{55}
\loigiai{
Ta có: $35 + 20 = 55$ m.
}
\end{ex}

\begin{ex}
Tính: $80\text{ m} - 35\text{ m}$. Kết quả là bao nhiêu mét?
\shortans{45}
\loigiai{
Ta có: $80 - 35 = 45$ m.
}
\end{ex}

\begin{ex}
Tính: $400\text{ m} + 500\text{ m}$. Kết quả là bao nhiêu mét?
\shortans{900}
\loigiai{
Ta có: $400 + 500 = 900$ m.
}
\end{ex}

\begin{ex}
Tính: $1000\text{ m} - 600\text{ m}$. Kết quả là bao nhiêu mét?
\shortans{400}
\loigiai{
Ta có: $1000 - 600 = 400$ m.
}
\end{ex}

\begin{ex}
Tính: $42\text{ km} + 18\text{ km}$. Kết quả là bao nhiêu ki-lô-mét?
\shortans{60}
\loigiai{
Ta có: $42 + 18 = 60$ km.
}
\end{ex}

\begin{ex}
Tính: $95\text{ km} - 45\text{ km}$. Kết quả là bao nhiêu ki-lô-mét?
\shortans{50}
\loigiai{
Ta có: $95 - 45 = 50$ km.
}
\end{ex}

\begin{ex}
Quãng đường từ nhà An đến trường dài 2 km, quãng đường từ nhà Bình đến trường dài 3 km. Quãng đường từ nhà Bình đến trường dài hơn nhà An bao nhiêu ki-lô-mét?
\shortans{1}
\loigiai{
Dài hơn: $3 - 2 = 1$ km.
}
\end{ex}

\begin{ex}
Một sợi dây thừng dài 10 m, người ta cắt bớt 4 m. Sợi dây thừng còn lại dài bao nhiêu mét?
\shortans{6}
\loigiai{
Còn lại: $10 - 4 = 6$ m.
}
\end{ex}

\begin{ex}
Cây cột cờ cao 8 m, cái cây bên cạnh cao 5 m. Cột cờ cao hơn cái cây bao nhiêu mét?
\shortans{3}
\loigiai{
Cao hơn: $8 - 5 = 3$ m.
}
\end{ex}

\begin{ex}
Đoạn đường từ tỉnh A đến tỉnh B dài 35 km, từ tỉnh B đến tỉnh C dài 45 km. Quãng đường từ tỉnh A đến tỉnh C đi qua B dài bao nhiêu ki-lô-mét?
\shortans{80}
\loigiai{
Dài: $35 + 45 = 80$ km.
}
\end{ex}

\begin{ex}
Bác thợ may có tấm vải dài 50 m, bác may rèm cửa hết 32 m. Tấm vải còn lại dài bao nhiêu mét?
\shortans{18}
\loigiai{
Còn lại: $50 - 32 = 18$ m.
}
\end{ex}

\begin{ex}
Điền số thích hợp: $1\text{ km} = \dots\text{ m}$.
\shortans{1000}
\loigiai{
Số cần điền là 1000.
}
\end{ex}

\begin{ex}
Điền số thích hợp: $2\text{ m} = \dots\text{ dm}$.
\shortans{20}
\loigiai{
$2\text{ m} = 20\text{ dm}$.
}
\end{ex}

\begin{ex}
Điền số thích hợp: $5\text{ dm} = \dots\text{ cm}$.
\shortans{50}
\loigiai{
$5\text{ dm} = 50\text{ cm}$.
}
\end{ex}

\begin{ex}
Điền số thích hợp: $300\text{ m} + \dots\text{ m} = 1000\text{ m}$.
\shortans{700}
\loigiai{
Số cần điền là: $1000 - 300 = 700$.
}
\end{ex}

\begin{ex}
Một đường chạy dài 100 m. Vận động viên đã chạy được 75 m. Vận động viên còn phải chạy bao nhiêu mét nữa để về đích?
\shortans{25}
\loigiai{
Còn phải chạy: $100 - 75 = 25$ m.
}
\end{ex}

\begin{ex}
Cuộn dây điện dài 48 m, thợ điện đã dùng 23 m. Cuộn dây còn lại dài bao nhiêu mét?
\shortans{25}
\loigiai{
Còn lại: $48 - 23 = 25$ m.
}
\end{ex}

\begin{ex}
Tòa nhà thứ nhất cao 28 m, tòa nhà thứ hai cao hơn tòa nhà thứ nhất 7 m. Tòa nhà thứ hai cao bao nhiêu mét?
\shortans{35}
\loigiai{
Tòa nhà 2 cao: $28 + 7 = 35$ m.
}
\end{ex}

\begin{ex}
Khoảng cách từ Hà Nội đến Phủ Lý là 60 km, từ Phủ Lý đến Nam Định là 30 km. Khoảng cách từ Hà Nội đến Nam Định đi qua Phủ Lý là bao nhiêu ki-lô-mét?
\shortans{90}
\loigiai{
Khoảng cách: $60 + 30 = 90$ km.
}
\end{ex}

\begin{ex}
Một con kênh dài 15 km, người ta đã nạo vét được 8 km. Còn lại bao nhiêu ki-lô-mét kênh chưa nạo vét?
\shortans{7}
\loigiai{
Còn lại: $15 - 8 = 7$ km.
}
\end{ex}

\begin{ex}
Tính: $250\text{ m} + 150\text{ m} + 100\text{ m}$.
\shortans{500}
\loigiai{
Ta có: $250 + 150 + 100 = 500$ m.
}
\end{ex}

\newpage

\TaoPhieuBaiTap{Chủ đề 11 -- Độ dài và đơn vị đo độ dài. tiền việt nam}{56}{Các số tròn chục, tròn trăm}{https://loptoanthaylinh.vn/bai-tap/toan-2/bai-56}
\begin{ex}
Tờ tiền 100 đồng cộng với tờ tiền 200 đồng bằng bao nhiêu đồng?
\shortans{300}
\loigiai{
Tổng: $100 + 200 = 300$ đồng.
}
\end{ex}

\begin{ex}
Tờ tiền 200 đồng cộng với tờ tiền 500 đồng bằng bao nhiêu đồng?
\shortans{700}
\loigiai{
Tổng: $200 + 500 = 700$ đồng.
}
\end{ex}

\begin{ex}
Có hai tờ tiền 500 đồng. Tổng số tiền là bao nhiêu đồng?
\shortans{1000}
\loigiai{
Tổng: $500 + 500 = 1000$ đồng.
}
\end{ex}

\begin{ex}
Có năm tờ tiền 200 đồng. Tổng số tiền là bao nhiêu đồng?
\shortans{1000}
\loigiai{
Tổng: $200 \times 5 = 1000$ đồng.
}
\end{ex}

\begin{ex}
Có mười tờ tiền 100 đồng. Tổng số tiền là bao nhiêu đồng?
\shortans{1000}
\loigiai{
Tổng: $100 \times 10 = 1000$ đồng.
}
\end{ex}

\begin{ex}
Bạn Mai có 1 tờ 500 đồng và 2 tờ 200 đồng. Mai có tất cả bao nhiêu đồng?
\shortans{900}
\loigiai{
Mai có: $500 + 200 + 200 = 900$ đồng.
}
\end{ex}

\begin{ex}
Bạn Nam có 1 tờ 1000 đồng, Nam mua một cục tẩy hết 600 đồng. Nam còn lại bao nhiêu đồng?
\shortans{400}
\loigiai{
Nam còn: $1000 - 600 = 400$ đồng.
}
\end{ex}

\begin{ex}
Một cây kem giá 800 đồng, bé đưa cho cô bán hàng 1000 đồng. Cô bán hàng phải trả lại bé bao nhiêu đồng?
\shortans{200}
\loigiai{
Trả lại: $1000 - 800 = 200$ đồng.
}
\end{ex}

\begin{ex}
Một phong bì thư giá 300 đồng, một con tem giá 500 đồng. Mua cả phong bì và tem hết bao nhiêu đồng?
\shortans{800}
\loigiai{
Hết: $300 + 500 = 800$ đồng.
}
\end{ex}

\begin{ex}
Lan có 700 đồng, mẹ cho thêm Lan 300 đồng. Lan có tất cả bao nhiêu đồng?
\shortans{1000}
\loigiai{
Lan có: $700 + 300 = 1000$ đồng.
}
\end{ex}

\begin{ex}
Hùng có 900 đồng, Hùng mua một chiếc kẹo mút hết 500 đồng. Hùng còn lại bao nhiêu đồng?
\shortans{400}
\loigiai{
Hùng còn: $900 - 500 = 400$ đồng.
}
\end{ex}

\begin{ex}
Có bao nhiêu tờ tiền 100 đồng để đổi được một tờ tiền 500 đồng?
\shortans{5}
\loigiai{
Cần: $500 : 100 = 5$ tờ.
}
\end{ex}

\begin{ex}
Có bao nhiêu tờ tiền 200 đồng và một tờ 100 đồng để đổi được 500 đồng? Cần mấy tờ 200 đồng?
\shortans{2}
\loigiai{
Vì $200 \times 2 + 100 = 500$ đồng nên cần 2 tờ 200 đồng.
}
\end{ex}

\begin{ex}
Một chiếc bút chì giá 700 đồng, một chiếc thước kẻ giá 300 đồng. Chiếc bút chì đắt hơn chiếc thước kẻ bao nhiêu đồng?
\shortans{400}
\loigiai{
Đắt hơn: $700 - 300 = 400$ đồng.
}
\end{ex}

\begin{ex}
Mua 2 chiếc bánh rán, mỗi chiếc giá 400 đồng. Hết tất cả bao nhiêu đồng?
\shortans{800}
\loigiai{
Hết: $400 \times 2 = 800$ đồng.
}
\end{ex}

\begin{ex}
Một quyển vở nhỏ giá 900 đồng, bạn An chỉ có 600 đồng. An còn thiếu bao nhiêu đồng để mua quyển vở?
\shortans{300}
\loigiai{
Thiếu: $900 - 600 = 300$ đồng.
}
\end{ex}

\begin{ex}
Điền số thích hợp: $400\text{ đồng} + \dots\text{ đồng} = 1000\text{ đồng}$.
\shortans{600}
\loigiai{
Số cần điền: $1000 - 400 = 600$.
}
\end{ex}

\begin{ex}
Điền số thích hợp: $1000\text{ đồng} - \dots\text{ đồng} = 300\text{ đồng}$.
\shortans{700}
\loigiai{
Số cần điền: $1000 - 300 = 700$.
}
\end{ex}

\begin{ex}
Trong ống tiết kiệm của bé có 3 tờ 200 đồng và 4 tờ 100 đồng. Bé đã tiết kiệm được bao nhiêu đồng?
\shortans{1000}
\loigiai{
Tiết kiệm được: $200 \times 3 + 100 \times 4 = 1000$ đồng.
}
\end{ex}

\begin{ex}
Một gói bim bim giá 1000 đồng, bạn Bình có 800 đồng. Bình cần thêm bao nhiêu đồng nữa?
\shortans{200}
\loigiai{
Cần thêm: $1000 - 800 = 200$ đồng.
}
\end{ex}

\begin{ex}
Bác nông dân bán 1 mớ rau muống được 800 đồng, bác mua 1 củ hành hết 300 đồng. Bác còn lại bao nhiêu đồng?
\shortans{500}
\loigiai{
Còn lại: $800 - 300 = 500$ đồng.
}
\end{ex}

\begin{ex}
Lấy số tiền lớn nhất trong các tờ tiền: 100 đồng, 200 đồng, 500 đồng, 1000 đồng. Tờ tiền đó có mệnh giá bao nhiêu đồng?
\shortans{1000}
\loigiai{
Mệnh giá lớn nhất là 1000 đồng.
}
\end{ex}

\begin{ex}
Tờ tiền có mệnh giá nhỏ nhất trong các tờ trên là bao nhiêu đồng?
\shortans{100}
\loigiai{
Mệnh giá nhỏ nhất là 100 đồng.
}
\end{ex}

\begin{ex}
Hiệu giữa tờ tiền 1000 đồng và tờ tiền 500 đồng là bao nhiêu đồng?
\shortans{500}
\loigiai{
Hiệu: $1000 - 500 = 500$ đồng.
}
\end{ex}

\begin{ex}
Tính: $200\text{ đồng} + 300\text{ đồng} + 500\text{ đồng}$.
\shortans{1000}
\loigiai{
Tổng: $200 + 300 + 500 = 1000$ đồng.
}
\end{ex}

\newpage

\TaoPhieuBaiTap{Chủ đề 11 -- Độ dài và đơn vị đo độ dài. tiền việt nam}{57}{Các số có ba chữ số}{https://loptoanthaylinh.vn/bai-tap/toan-2/bai-57}
\begin{ex}
Dùng thước dây đo chiều dài bảng lớp được 3 mét. Đổi 3 mét ra đề-xi-mét là bao nhiêu đề-xi-mét?
\shortans{30}
\loigiai{
$3\text{ m} = 30\text{ dm}$.
}
\end{ex}

\begin{ex}
Đo chiều dài lớp học được 8 mét. Đổi 8 mét ra đề-xi-mét là bao nhiêu đề-xi-mét?
\shortans{80}
\loigiai{
$8\text{ m} = 80\text{ dm}$.
}
\end{ex}

\begin{ex}
Chiều dài cuốn sách Toán 2 đo được 24 cm. Chiều rộng đo được 17 cm. Chiều dài hơn chiều rộng bao nhiêu xăng-ti-mét?
\shortans{7}
\loigiai{
Hơn: $24 - 17 = 7$ cm.
}
\end{ex}

\begin{ex}
Bạn Lan bước 5 bước chân, mỗi bước chân dài khoảng 40 cm. Quãng đường Lan bước dài bao nhiêu xăng-ti-mét?
\shortans{200}
\loigiai{
Dài: $40 \times 5 = 200$ cm.
}
\end{ex}

\begin{ex}
Đổi 200 cm ra mét được bao nhiêu mét?
\shortans{2}
\loigiai{
$200\text{ cm} = 2\text{ m}$.
}
\end{ex}

\begin{ex}
Một cái bàn học dài 120 cm. Đổi ra đề-xi-mét là bao nhiêu đề-xi-mét?
\shortans{12}
\loigiai{
$120\text{ cm} = 12\text{ dm}$.
}
\end{ex}

\begin{ex}
Một thanh nhôm dài 2 m, bác thợ cắt đi 8 dm. Thanh nhôm còn lại dài bao nhiêu đề-xi-mét?
\shortans{12}
\loigiai{
$2\text{ m} = 20\text{ dm}$. Còn lại: $20 - 8 = 12$ dm.
}
\end{ex}

\begin{ex}
Cây cau cao 9 m, cây chuối cao 3 m. Cây cau cao gấp mấy lần cây chuối? Cây cau cao hơn cây chuối bao nhiêu mét?
\shortans{6}
\loigiai{
Cao hơn: $9 - 3 = 6$ mét.
}
\end{ex}

\begin{ex}
Đoạn thẳng AB dài 15 cm, vẽ thêm đoạn BC dài 25 cm thẳng hàng với AB. Cả đoạn thẳng AC dài bao nhiêu xăng-ti-mét?
\shortans{40}
\loigiai{
Đoạn AC dài: $15 + 25 = 40$ cm.
}
\end{ex}

\begin{ex}
Đoạn đường quanh sân trường hình chữ nhật có chu vi là 300 m. Bạn Hùng chạy 2 vòng quanh sân trường thì chạy được bao nhiêu mét?
\shortans{600}
\loigiai{
Chạy được: $300 \times 2 = 600$ m.
}
\end{ex}

\begin{ex}
Một cuộn chỉ dài 100 m, người ta dùng hết 65 m. Cuộn chỉ còn lại bao nhiêu mét?
\shortans{35}
\loigiai{
Còn lại: $100 - 65 = 35$ m.
}
\end{ex}

\begin{ex}
Bạn An cao 128 cm, bạn Bình cao 135 cm. Bạn Bình cao hơn bạn An bao nhiêu xăng-ti-mét?
\shortans{7}
\loigiai{
Cao hơn: $135 - 128 = 7$ cm.
}
\end{ex}

\begin{ex}
Bác nông dân rào vườn hết 85 m lưới thép, bác mua cuộn lưới dài 100 m. Bác còn thừa bao nhiêu mét lưới thép?
\shortans{15}
\loigiai{
Thừa: $100 - 85 = 15$ m.
}
\end{ex}

\begin{ex}
Một chiếc cột điện cao 12 m, chiếc cột đèn cao 8 m. Cột điện cao hơn cột đèn bao nhiêu mét?
\shortans{4}
\loigiai{
Cao hơn: $12 - 8 = 4$ m.
}
\end{ex}

\begin{ex}
Đoạn đường từ cổng trường đến thư viện dài 45 m, từ thư viện đến nhà ăn dài 35 m. Quãng đường từ cổng trường đến nhà ăn đi qua thư viện dài bao nhiêu mét?
\shortans{80}
\loigiai{
Dài: $45 + 35 = 80$ m.
}
\end{ex}

\begin{ex}
Một con sâu đo một chiếc lá dài 12 cm, mỗi lần sâu vươn mình tiến được 2 cm. Sâu cần đo bao nhiêu lần để đi hết chiếc lá?
\shortans{6}
\loigiai{
Cần đo: $12 : 2 = 6$ lần.
}
\end{ex}

\begin{ex}
Đoạn thẳng MN dài 5 dm, đoạn thẳng PQ dài 30 cm. Đoạn MN dài hơn đoạn PQ bao nhiêu xăng-ti-mét?
\shortans{20}
\loigiai{
$5\text{ dm} = 50\text{ cm}$. Hơn: $50 - 30 = 20$ cm.
}
\end{ex}

\begin{ex}
Một sợi dây thừng dài 18 m được cắt làm 2 đoạn bằng nhau. Mỗi đoạn dài bao nhiêu mét?
\shortans{9}
\loigiai{
Mỗi đoạn dài: $18 : 2 = 9$ m.
}
\end{ex}

\begin{ex}
Một sợi dây đồng dài 20 dm được cắt thành các đoạn 5 dm. Cắt được bao nhiêu đoạn?
\shortans{4}
\loigiai{
Cắt được: $20 : 5 = 4$ đoạn.
}
\end{ex}

\begin{ex}
Đo chiều cao bục giảng được 20 cm, bục có 2 bậc bằng nhau. Mỗi bậc cao bao nhiêu xăng-ti-mét?
\shortans{10}
\loigiai{
Mỗi bậc cao: $20 : 2 = 10$ cm.
}
\end{ex}

\begin{ex}
Một chiếc thang có 10 bậc, khoảng cách giữa các bậc là 2 dm. Chiều dài chiếc thang tính từ bậc 1 đến bậc 10 (gồm 9 khoảng cách) là bao nhiêu đề-xi-mét?
\shortans{18}
\loigiai{
Chiều dài: $2 \times 9 = 18$ dm.
}
\end{ex}

\begin{ex}
Tính: $45\text{ m} + 35\text{ m} - 20\text{ m}$.
\shortans{60}
\loigiai{
Ta có $45 + 35 = 80$, rồi $80 - 20 = 60$ m.
}
\end{ex}

\begin{ex}
Tính: $100\text{ m} - 40\text{ m} - 30\text{ m}$.
\shortans{30}
\loigiai{
Ta có $100 - 40 = 60$, rồi $60 - 30 = 30$ m.
}
\end{ex}

\begin{ex}
Tính: $500\text{ m} + 500\text{ m}$. Kết quả đổi ra ki-lô-mét là bao nhiêu ki-lô-mét?
\shortans{1}
\loigiai{
Ta có $500 + 500 = 1000\text{ m} = 1$ km.
}
\end{ex}

\begin{ex}
Một người đi bộ 1 giờ được 4 km. Trong 2 giờ người đó đi được bao nhiêu ki-lô-mét?
\shortans{8}
\loigiai{
Đi được: $4 \times 2 = 8$ km.
}
\end{ex}

\newpage

\TaoPhieuBaiTap{Chủ đề 11 -- Độ dài và đơn vị đo độ dài. tiền việt nam}{58}{Viết số thành tổng các trăm, chục, đơn vị}{https://loptoanthaylinh.vn/bai-tap/toan-2/bai-58}
\begin{ex}
Tính: $45\text{ m} + 28\text{ m}$.
\shortans{73}
\loigiai{
Ta có: $45 + 28 = 73$ m.
}
\end{ex}

\begin{ex}
Tính: $82\text{ m} - 37\text{ m}$.
\shortans{45}
\loigiai{
Ta có: $82 - 37 = 45$ m.
}
\end{ex}

\begin{ex}
Tính: $35\text{ km} + 48\text{ km}$.
\shortans{83}
\loigiai{
Ta có: $35 + 48 = 83$ km.
}
\end{ex}

\begin{ex}
Tính: $90\text{ km} - 54\text{ km}$.
\shortans{36}
\loigiai{
Ta có: $90 - 54 = 36$ km.
}
\end{ex}

\begin{ex}
Tính: $300\text{ đồng} + 400\text{ đồng} + 300\text{ đồng}$.
\shortans{1000}
\loigiai{
Tổng: $300 + 400 + 300 = 1000$ đồng.
}
\end{ex}

\begin{ex}
Tính: $1000\text{ đồng} - 700\text{ đồng}$.
\shortans{300}
\loigiai{
Hiệu: $1000 - 700 = 300$ đồng.
}
\end{ex}

\begin{ex}
Điền số thích hợp: $1\text{ m} = \dots\text{ cm}$.
\shortans{100}
\loigiai{
$1\text{ m} = 100\text{ cm}$.
}
\end{ex}

\begin{ex}
Điền số thích hợp: $1\text{ km} = \dots\text{ m}$.
\shortans{1000}
\loigiai{
$1\text{ km} = 1000\text{ m}$.
}
\end{ex}

\begin{ex}
Điền số thích hợp: $50\text{ dm} = \dots\text{ m}$.
\shortans{5}
\loigiai{
$50\text{ dm} = 5\text{ m}$.
}
\end{ex}

\begin{ex}
Bạn Lan có 1000 đồng, Lan mua 1 cái bánh giá 500 đồng và 1 cái kẹo giá 200 đồng. Lan còn lại bao nhiêu đồng?
\shortans{300}
\loigiai{
Lan đã tiêu: $500 + 200 = 700$ đồng. Còn lại: $1000 - 700 = 300$ đồng.
}
\end{ex}

\begin{ex}
Quãng đường từ nhà bác An đến thị trấn dài 15 km, từ thị trấn đến thành phố dài 25 km. Quãng đường từ nhà bác An đến thành phố dài bao nhiêu ki-lô-mét?
\shortans{40}
\loigiai{
Dài: $15 + 25 = 40$ km.
}
\end{ex}

\begin{ex}
Một cuộn dây cáp dài 92 m, thợ điện đã dùng 45 m. Cuộn cáp còn lại bao nhiêu mét?
\shortans{47}
\loigiai{
Còn lại: $92 - 45 = 47$ m.
}
\end{ex}

\begin{ex}
Mẹ mua 1 chiếc khăn tay giá 800 đồng, mẹ đưa cô bán hàng 1000 đồng. Cô bán hàng trả lại bao nhiêu đồng?
\shortans{200}
\loigiai{
Trả lại: $1000 - 800 = 200$ đồng.
}
\end{ex}

\begin{ex}
Đoạn thẳng thứ nhất dài 38 cm, đoạn thứ hai dài 27 cm. Cả hai đoạn dài bao nhiêu xăng-ti-mét?
\shortans{65}
\loigiai{
Dài: $38 + 27 = 65$ cm.
}
\end{ex}

\begin{ex}
Một con tàu đi được 45 km trong ngày thứ nhất, ngày thứ hai đi được nhiều hơn ngày thứ nhất 18 km. Ngày thứ hai tàu đi được bao nhiêu ki-lô-mét?
\shortans{63}
\loigiai{
Ngày thứ hai đi: $45 + 18 = 63$ km.
}
\end{ex}

\begin{ex}
Đổi 4 m ra xăng-ti-mét được bao nhiêu xăng-ti-mét?
\shortans{400}
\loigiai{
$4\text{ m} = 400\text{ cm}$.
}
\end{ex}

\begin{ex}
Đổi 600 cm ra mét được bao nhiêu mét?
\shortans{6}
\loigiai{
$600\text{ cm} = 6\text{ m}$.
}
\end{ex}

\begin{ex}
Đổi 7 m ra đề-xi-mét được bao nhiêu đề-xi-mét?
\shortans{70}
\loigiai{
$7\text{ m} = 70\text{ dm}$.
}
\end{ex}

\begin{ex}
Tính giá trị biểu thức: $35\text{ m} + 45\text{ m} - 50\text{ m}$.
\shortans{30}
\loigiai{
Ta có $35 + 45 = 80$, rồi $80 - 50 = 30$ m.
}
\end{ex}

\begin{ex}
Tính giá trị biểu thức: $1000\text{ m} - 400\text{ m} - 200\text{ m}$.
\shortans{400}
\loigiai{
Ta có $1000 - 400 = 600$, rồi $600 - 200 = 400$ m.
}
\end{ex}

\begin{ex}
Có 5 tờ tiền mệnh giá 200 đồng. Tổng số tiền là bao nhiêu đồng?
\shortans{1000}
\loigiai{
Tổng: $200 \times 5 = 1000$ đồng.
}
\end{ex}

\begin{ex}
Có 2 tờ tiền mệnh giá 500 đồng. Tổng số tiền là bao nhiêu đồng?
\shortans{1000}
\loigiai{
Tổng: $500 \times 2 = 1000$ đồng.
}
\end{ex}

\begin{ex}
Một cây cầu dài 450 m, cây cầu thứ hai dài 550 m. Tổng chiều dài hai cây cầu là bao nhiêu mét?
\shortans{1000}
\loigiai{
Tổng chiều dài: $450 + 550 = 1000$ m.
}
\end{ex}

\begin{ex}
Tổng chiều dài hai cây cầu đó bằng bao nhiêu ki-lô-mét?
\shortans{1}
\loigiai{
$1000\text{ m} = 1\text{ km}$.
}
\end{ex}

\begin{ex}
Tìm một số biết lấy số đó cộng với 350 thì được 1000.
\shortans{650}
\loigiai{
Số đó là: $1000 - 350 = 650$.
}
\end{ex}

\newpage


% =========================================================================
% Phép cộng, phép trừ trong phạm vi 1 000
% =========================================================================
\section{Phép cộng, phép trừ trong phạm vi 1 000}

\TaoPhieuBaiTap{Chủ đề 12 -- Phép cộng, phép trừ trong phạm vi 1 000}{59}{So sánh các số có ba chữ số}{https://loptoanthaylinh.vn/bai-tap/toan-2/bai-59}
\begin{ex}
Tính kết quả của phép tính: $234 + 125$.
\shortans{359}
\loigiai{
Ta có $234 + 125 = 359$.
}
\end{ex}

\begin{ex}
Tính kết quả của phép tính: $342 + 256$.
\shortans{598}
\loigiai{
Ta có $342 + 256 = 598$.
}
\end{ex}

\begin{ex}
Tính kết quả của phép tính: $415 + 372$.
\shortans{787}
\loigiai{
Ta có $415 + 372 = 787$.
}
\end{ex}

\begin{ex}
Tính kết quả của phép tính: $520 + 268$.
\shortans{788}
\loigiai{
Ta có $520 + 268 = 788$.
}
\end{ex}

\begin{ex}
Tính kết quả của phép tính: $631 + 247$.
\shortans{878}
\loigiai{
Ta có $631 + 247 = 878$.
}
\end{ex}

\begin{ex}
Tính kết quả của phép tính: $714 + 185$.
\shortans{899}
\loigiai{
Ta có $714 + 185 = 899$.
}
\end{ex}

\begin{ex}
Tính kết quả của phép tính: $153 + 42$.
\shortans{195}
\loigiai{
Ta có $153 + 42 = 195$.
}
\end{ex}

\begin{ex}
Tính kết quả của phép tính: $324 + 5$.
\shortans{329}
\loigiai{
Ta có $324 + 5 = 329$.
}
\end{ex}

\begin{ex}
Tính kết quả của phép tính: $506 + 72$.
\shortans{578}
\loigiai{
Ta có $506 + 72 = 578$.
}
\end{ex}

\begin{ex}
Tính kết quả của phép tính: $812 + 63$.
\shortans{875}
\loigiai{
Ta có $812 + 63 = 875$.
}
\end{ex}

\begin{ex}
Tính kết quả của phép tính: $400 + 350$.
\shortans{750}
\loigiai{
Ta có $400 + 350 = 750$.
}
\end{ex}

\begin{ex}
Tính kết quả của phép tính: $230 + 145$.
\shortans{375}
\loigiai{
Ta có $230 + 145 = 375$.
}
\end{ex}

\begin{ex}
Tính kết quả của phép tính: $605 + 132$.
\shortans{737}
\loigiai{
Ta có $605 + 132 = 737$.
}
\end{ex}

\begin{ex}
Điền số thích hợp: $241 + \dots = 568$.
\shortans{327}
\loigiai{
Số cần điền: $568 - 241 = 327$.
}
\end{ex}

\begin{ex}
Điền số thích hợp: $352 + \dots = 786$.
\shortans{434}
\loigiai{
Số cần điền: $786 - 352 = 434$.
}
\end{ex}

\begin{ex}
Điền số thích hợp: $\dots + 123 = 457$.
\shortans{334}
\loigiai{
Số cần điền: $457 - 123 = 334$.
}
\end{ex}

\begin{ex}
Điền số thích hợp: $\dots + 240 = 685$.
\shortans{445}
\loigiai{
Số cần điền: $685 - 240 = 445$.
}
\end{ex}

\begin{ex}
Trường Tiểu học Lê Quý Đôn có 324 học sinh nam và 342 học sinh nữ. Trường có tất cả bao nhiêu học sinh?
\shortans{666}
\loigiai{
Có tất cả: $324 + 342 = 666$ học sinh.
}
\end{ex}

\begin{ex}
Một trang trại nuôi 420 con gà và 250 con vịt. Trang trại nuôi tất cả bao nhiêu con gia cầm?
\shortans{670}
\loigiai{
Có tất cả: $420 + 250 = 670$ con.
}
\end{ex}

\begin{ex}
Kho thóc thứ nhất có 315 bao thóc, kho thóc thứ hai có 260 bao thóc. Cả hai kho có bao nhiêu bao thóc?
\shortans{575}
\loigiai{
Cả hai kho có: $315 + 260 = 575$ bao.
}
\end{ex}

\begin{ex}
Bác nông dân thu hoạch được 542 kg bưởi vào đợt một và 325 kg bưởi vào đợt hai. Bác thu hoạch được tất cả bao nhiêu ki-lô-gam bưởi?
\shortans{867}
\loigiai{
Cả hai đợt thu hoạch: $542 + 325 = 867$ kg.
}
\end{ex}

\begin{ex}
Quãng đường từ A đến B dài 235 km, từ B đến C dài 142 km. Quãng đường từ A đến C đi qua B dài bao nhiêu ki-lô-mét?
\shortans{377}
\loigiai{
Dài: $235 + 142 = 377$ km.
}
\end{ex}

\begin{ex}
Tính giá trị biểu thức: $123 + 234 + 341$.
\shortans{698}
\loigiai{
Ta có $123 + 234 = 357$, rồi $357 + 341 = 698$.
}
\end{ex}

\begin{ex}
Tính giá trị biểu thức: $400 + 250 + 130$.
\shortans{780}
\loigiai{
Ta có $400 + 250 = 650$, rồi $650 + 130 = 780$.
}
\end{ex}

\begin{ex}
Tìm tổng của số 432 và số 355.
\shortans{787}
\loigiai{
Tổng: $432 + 355 = 787$.
}
\end{ex}

\newpage

\TaoPhieuBaiTap{Chủ đề 12 -- Phép cộng, phép trừ trong phạm vi 1 000}{60}{Luyện tập các số trong phạm vi 1000}{https://loptoanthaylinh.vn/bai-tap/toan-2/bai-60}
\begin{ex}
Tính kết quả của phép tính: $248 + 135$.
\shortans{383}
\loigiai{
Ta có $8 + 5 = 13$, viết 3 nhớ 1; $4 + 3 + 1 = 8$; $2 + 1 = 3$. Kết quả là 383.
}
\end{ex}

\begin{ex}
Tính kết quả của phép tính: $356 + 218$.
\shortans{574}
\loigiai{
Ta có $356 + 218 = 574$.
}
\end{ex}

\begin{ex}
Tính kết quả của phép tính: $427 + 345$.
\shortans{772}
\loigiai{
Ta có $427 + 345 = 772$.
}
\end{ex}

\begin{ex}
Tính kết quả của phép tính: $539 + 246$.
\shortans{785}
\loigiai{
Ta có $539 + 246 = 785$.
}
\end{ex}

\begin{ex}
Tính kết quả của phép tính: $618 + 175$.
\shortans{793}
\loigiai{
Ta có $618 + 175 = 793$.
}
\end{ex}

\begin{ex}
Tính kết quả của phép tính: $382 + 254$.
\shortans{636}
\loigiai{
Ta có $8 + 5 = 13$, viết 3 nhớ 1 sang hàng trăm. Kết quả là 636.
}
\end{ex}

\begin{ex}
Tính kết quả của phép tính: $475 + 162$.
\shortans{637}
\loigiai{
Ta có $475 + 162 = 637$.
}
\end{ex}

\begin{ex}
Tính kết quả của phép tính: $563 + 281$.
\shortans{844}
\loigiai{
Ta có $563 + 281 = 844$.
}
\end{ex}

\begin{ex}
Tính kết quả của phép tính: $684 + 153$.
\shortans{837}
\loigiai{
Ta có $684 + 153 = 837$.
}
\end{ex}

\begin{ex}
Tính kết quả của phép tính: $267 + 358$.
\shortans{625}
\loigiai{
Ta có $267 + 358 = 625$.
}
\end{ex}

\begin{ex}
Tính kết quả của phép tính: $485 + 247$.
\shortans{732}
\loigiai{
Ta có $485 + 247 = 732$.
}
\end{ex}

\begin{ex}
Tính kết quả của phép tính: $379 + 456$.
\shortans{835}
\loigiai{
Ta có $379 + 456 = 835$.
}
\end{ex}

\begin{ex}
Tính kết quả của phép tính: $594 + 238$.
\shortans{832}
\loigiai{
Ta có $594 + 238 = 832$.
}
\end{ex}

\begin{ex}
Tính kết quả của phép tính: $648 + 176$.
\shortans{824}
\loigiai{
Ta có $648 + 176 = 824$.
}
\end{ex}

\begin{ex}
Tính kết quả của phép tính: $365 + 478$.
\shortans{843}
\loigiai{
Ta có $365 + 478 = 843$.
}
\end{ex}

\begin{ex}
Điền số thích hợp: $258 + \dots = 482$.
\shortans{224}
\loigiai{
Số cần điền: $482 - 258 = 224$.
}
\end{ex}

\begin{ex}
Điền số thích hợp: $375 + \dots = 620$.
\shortans{245}
\loigiai{
Số cần điền: $620 - 375 = 245$.
}
\end{ex}

\begin{ex}
Điền số thích hợp: $\dots + 184 = 536$.
\shortans{352}
\loigiai{
Số cần điền: $536 - 184 = 352$.
}
\end{ex}

\begin{ex}
Điền số thích hợp: $\dots + 267 = 715$.
\shortans{448}
\loigiai{
Số cần điền: $715 - 267 = 448$.
}
\end{ex}

\begin{ex}
Trường Tiểu học Quang Trung có 458 học sinh nam và 465 học sinh nữ. Trường có tất cả bao nhiêu học sinh?
\shortans{923}
\loigiai{
Có tất cả: $458 + 465 = 923$ học sinh.
}
\end{ex}

\begin{ex}
Một cửa hàng ngày thứ nhất bán được 348 mét vải, ngày thứ hai bán được 275 mét vải. Cả hai ngày bán được bao nhiêu mét vải?
\shortans{623}
\loigiai{
Cả hai ngày bán: $348 + 275 = 623$ mét.
}
\end{ex}

\begin{ex}
Kho thóc thứ nhất có 468 bao thóc, kho thóc thứ hai có 385 bao thóc. Cả hai kho có bao nhiêu bao thóc?
\shortans{853}
\loigiai{
Có tất cả: $468 + 385 = 853$ bao.
}
\end{ex}

\begin{ex}
Bác nông dân thu hoạch đợt một được 376 kg cam, đợt hai được 289 kg cam. Bác thu hoạch được tất cả bao nhiêu ki-lô-gam cam?
\shortans{665}
\loigiai{
Tổng sản lượng: $376 + 289 = 665$ kg.
}
\end{ex}

\begin{ex}
Tính giá trị biểu thức: $245 + 178 + 326$.
\shortans{749}
\loigiai{
Ta có $245 + 178 = 423$, rồi $423 + 326 = 749$.
}
\end{ex}

\begin{ex}
Tính tổng của số 548 và số 276.
\shortans{824}
\loigiai{
Tổng: $548 + 276 = 824$.
}
\end{ex}

\newpage

\TaoPhieuBaiTap{Chủ đề 12 -- Phép cộng, phép trừ trong phạm vi 1 000}{61}{Phép cộng (không nhớ) trong phạm vi 1000}{https://loptoanthaylinh.vn/bai-tap/toan-2/bai-61}
\begin{ex}
Tính kết quả của phép tính: $458 - 235$.
\shortans{223}
\loigiai{
Ta có $458 - 235 = 223$.
}
\end{ex}

\begin{ex}
Tính kết quả của phép tính: $579 - 346$.
\shortans{233}
\loigiai{
Ta có $579 - 346 = 233$.
}
\end{ex}

\begin{ex}
Tính kết quả của phép tính: $684 - 251$.
\shortans{433}
\loigiai{
Ta có $684 - 251 = 433$.
}
\end{ex}

\begin{ex}
Tính kết quả của phép tính: $796 - 463$.
\shortans{333}
\loigiai{
Ta có $796 - 463 = 333$.
}
\end{ex}

\begin{ex}
Tính kết quả của phép tính: $875 - 542$.
\shortans{333}
\loigiai{
Ta có $875 - 542 = 333$.
}
\end{ex}

\begin{ex}
Tính kết quả của phép tính: $968 - 634$.
\shortans{334}
\loigiai{
Ta có $968 - 634 = 334$.
}
\end{ex}

\begin{ex}
Tính kết quả của phép tính: $354 - 23$.
\shortans{331}
\loigiai{
Ta có $354 - 23 = 331$.
}
\end{ex}

\begin{ex}
Tính kết quả của phép tính: $487 - 65$.
\shortans{422}
\loigiai{
Ta có $487 - 65 = 422$.
}
\end{ex}

\begin{ex}
Tính kết quả của phép tính: $569 - 4$.
\shortans{565}
\loigiai{
Ta có $569 - 4 = 565$.
}
\end{ex}

\begin{ex}
Tính kết quả của phép tính: $748 - 300$.
\shortans{448}
\loigiai{
Ta có $748 - 300 = 448$.
}
\end{ex}

\begin{ex}
Tính kết quả của phép tính: $890 - 450$.
\shortans{440}
\loigiai{
Ta có $890 - 450 = 440$.
}
\end{ex}

\begin{ex}
Tính kết quả của phép tính: $675 - 205$.
\shortans{470}
\loigiai{
Ta có $675 - 205 = 470$.
}
\end{ex}

\begin{ex}
Tính kết quả của phép tính: $987 - 787$.
\shortans{200}
\loigiai{
Ta có $987 - 787 = 200$.
}
\end{ex}

\begin{ex}
Điền số thích hợp: $568 - \dots = 324$.
\shortans{244}
\loigiai{
Số cần điền: $568 - 324 = 244$.
}
\end{ex}

\begin{ex}
Điền số thích hợp: $789 - \dots = 456$.
\shortans{333}
\loigiai{
Số cần điền: $789 - 456 = 333$.
}
\end{ex}

\begin{ex}
Điền số thích hợp: $\dots - 213 = 452$.
\shortans{665}
\loigiai{
Số cần điền: $452 + 213 = 665$.
}
\end{ex}

\begin{ex}
Điền số thích hợp: $\dots - 340 = 528$.
\shortans{868}
\loigiai{
Số cần điền: $528 + 340 = 868$.
}
\end{ex}

\begin{ex}
Một trường tiểu học có 768 học sinh, trong đó có 345 học sinh nam. Hỏi trường có bao nhiêu học sinh nữ?
\shortans{423}
\loigiai{
Số học sinh nữ: $768 - 345 = 423$ bạn.
}
\end{ex}

\begin{ex}
Kho thóc có 895 bao thóc, người ta đã xuất đi 462 bao thóc. Trong kho còn lại bao nhiêu bao thóc?
\shortans{433}
\loigiai{
Còn lại: $895 - 462 = 433$ bao.
}
\end{ex}

\begin{ex}
Một xí nghiệp may được 685 bộ quần áo, đã xuất bán 350 bộ. Xí nghiệp còn lại bao nhiêu bộ quần áo?
\shortans{335}
\loigiai{
Còn lại: $685 - 350 = 335$ bộ.
}
\end{ex}

\begin{ex}
Cuộn dây cáp dài 589 mét, thợ điện đã dùng 264 mét. Cuộn dây còn lại bao nhiêu mét?
\shortans{325}
\loigiai{
Còn lại: $589 - 264 = 325$ mét.
}
\end{ex}

\begin{ex}
Bác Ba thu hoạch được 758 kg ngô, bác bán đi 425 kg ngô. Bác còn lại bao nhiêu ki-lô-gam ngô?
\shortans{333}
\loigiai{
Còn lại: $758 - 425 = 333$ kg.
}
\end{ex}

\begin{ex}
Tính giá trị biểu thức: $789 - 235 - 142$.
\shortans{412}
\loigiai{
Ta có $789 - 235 = 554$, rồi $554 - 142 = 412$.
}
\end{ex}

\begin{ex}
Tính giá trị biểu thức: $896 - 452 - 210$.
\shortans{234}
\loigiai{
Ta có $896 - 452 = 444$, rồi $444 - 210 = 234$.
}
\end{ex}

\begin{ex}
Tìm hiệu của số 876 và số 543.
\shortans{333}
\loigiai{
Hiệu: $876 - 543 = 333$.
}
\end{ex}

\newpage

\TaoPhieuBaiTap{Chủ đề 12 -- Phép cộng, phép trừ trong phạm vi 1 000}{62}{Phép trừ (không nhớ) trong phạm vi 1000}{https://loptoanthaylinh.vn/bai-tap/toan-2/bai-62}
\begin{ex}
Tính kết quả của phép tính: $452 - 138$.
\shortans{314}
\loigiai{
Ta có 2 không trừ được 8, lấy $12 - 8 = 4$, nhớ 1; $3 + 1 = 4$, rồi $5 - 4 = 1$; $4 - 1 = 3$. Kết quả là 314.
}
\end{ex}

\begin{ex}
Tính kết quả của phép tính: $563 - 247$.
\shortans{316}
\loigiai{
Ta có $563 - 247 = 316$.
}
\end{ex}

\begin{ex}
Tính kết quả của phép tính: $674 - 328$.
\shortans{346}
\loigiai{
Ta có $674 - 328 = 346$.
}
\end{ex}

\begin{ex}
Tính kết quả của phép tính: $781 - 435$.
\shortans{346}
\loigiai{
Ta có $781 - 435 = 346$.
}
\end{ex}

\begin{ex}
Tính kết quả của phép tính: $890 - 546$.
\shortans{344}
\loigiai{
Ta có $890 - 546 = 344$.
}
\end{ex}

\begin{ex}
Tính kết quả của phép tính: $635 - 281$.
\shortans{354}
\loigiai{
Ta có 3 không trừ được 8, lấy $13 - 8 = 5$, nhớ 1; $2 + 1 = 3$, rồi $6 - 3 = 3$. Kết quả là 354.
}
\end{ex}

\begin{ex}
Tính kết quả của phép tính: $724 - 351$.
\shortans{373}
\loigiai{
Ta có $724 - 351 = 373$.
}
\end{ex}

\begin{ex}
Tính kết quả của phép tính: $846 - 472$.
\shortans{374}
\loigiai{
Ta có $846 - 472 = 374$.
}
\end{ex}

\begin{ex}
Tính kết quả của phép tính: $915 - 562$.
\shortans{353}
\loigiai{
Ta có $915 - 562 = 353$.
}
\end{ex}

\begin{ex}
Tính kết quả của phép tính: $500 - 245$.
\shortans{255}
\loigiai{
Ta có $500 - 245 = 255$.
}
\end{ex}

\begin{ex}
Tính kết quả của phép tính: $600 - 328$.
\shortans{272}
\loigiai{
Ta có $600 - 328 = 272$.
}
\end{ex}

\begin{ex}
Tính kết quả của phép tính: $700 - 465$.
\shortans{235}
\loigiai{
Ta có $700 - 465 = 235$.
}
\end{ex}

\begin{ex}
Tính kết quả của phép tính: $800 - 537$.
\shortans{263}
\loigiai{
Ta có $800 - 537 = 263$.
}
\end{ex}

\begin{ex}
Tính kết quả của phép tính: $1000 - 458$.
\shortans{542}
\loigiai{
Ta có $1000 - 458 = 542$.
}
\end{ex}

\begin{ex}
Tính kết quả của phép tính: $1000 - 675$.
\shortans{325}
\loigiai{
Ta có $1000 - 675 = 325$.
}
\end{ex}

\begin{ex}
Điền số thích hợp: $532 - \dots = 278$.
\shortans{254}
\loigiai{
Số cần điền: $532 - 278 = 254$.
}
\end{ex}

\begin{ex}
Điền số thích hợp: $715 - \dots = 369$.
\shortans{346}
\loigiai{
Số cần điền: $715 - 369 = 346$.
}
\end{ex}

\begin{ex}
Điền số thích hợp: $\dots - 185 = 447$.
\shortans{632}
\loigiai{
Số cần điền: $447 + 185 = 632$.
}
\end{ex}

\begin{ex}
Điền số thích hợp: $\dots - 294 = 528$.
\shortans{822}
\loigiai{
Số cần điền: $528 + 294 = 822$.
}
\end{ex}

\begin{ex}
Một cửa hàng có 625 kg đường, đã bán được 278 kg. Cửa hàng còn lại bao nhiêu ki-lô-gam đường?
\shortans{347}
\loigiai{
Còn lại: $625 - 278 = 347$ kg.
}
\end{ex}

\begin{ex}
Đoạn đường dài 750 mét, đội công nhân đã trải nhựa được 385 mét. Còn lại bao nhiêu mét đường chưa trải nhựa?
\shortans{365}
\loigiai{
Còn lại: $750 - 385 = 365$ mét.
}
\end{ex}

\begin{ex}
Một thư viện có 920 cuốn sách, học sinh đã mượn 465 cuốn. Trong thư viện còn lại bao nhiêu cuốn sách?
\shortans{455}
\loigiai{
Còn lại: $920 - 465 = 455$ cuốn sách.
}
\end{ex}

\begin{ex}
Kho thóc có 800 bao thóc, người ta xuất kho 342 bao. Kho còn lại bao nhiêu bao thóc?
\shortans{458}
\loigiai{
Còn lại: $800 - 342 = 458$ bao.
}
\end{ex}

\begin{ex}
Tính giá trị biểu thức: $735 - 268 - 189$.
\shortans{278}
\loigiai{
Ta có $735 - 268 = 467$, rồi $467 - 189 = 278$.
}
\end{ex}

\begin{ex}
Tìm hiệu của số 612 và số 248.
\shortans{364}
\loigiai{
Hiệu: $612 - 248 = 364$.
}
\end{ex}

\newpage

\TaoPhieuBaiTap{Chủ đề 12 -- Phép cộng, phép trừ trong phạm vi 1 000}{63}{Phép cộng (có nhớ) trong phạm vi 1000}{https://loptoanthaylinh.vn/bai-tap/toan-2/bai-63}
\begin{ex}
Tính: $356 + 278$.
\shortans{634}
\loigiai{
Ta có $356 + 278 = 634$.
}
\end{ex}

\begin{ex}
Tính: $634 - 356$.
\shortans{278}
\loigiai{
Ta có $634 - 356 = 278$.
}
\end{ex}

\begin{ex}
Tính: $634 - 278$.
\shortans{356}
\loigiai{
Ta có $634 - 278 = 356$.
}
\end{ex}

\begin{ex}
Tính: $485 + 367$.
\shortans{852}
\loigiai{
Ta có $485 + 367 = 852$.
}
\end{ex}

\begin{ex}
Tính: $852 - 485$.
\shortans{367}
\loigiai{
Ta có $852 - 485 = 367$.
}
\end{ex}

\begin{ex}
Tính: $500 - 186$.
\shortans{314}
\loigiai{
Ta có $500 - 186 = 314$.
}
\end{ex}

\begin{ex}
Tính: $1000 - 548$.
\shortans{452}
\loigiai{
Ta có $1000 - 548 = 452$.
}
\end{ex}

\begin{ex}
Tính giá trị biểu thức: $345 + 289 - 176$.
\shortans{458}
\loigiai{
Ta có $345 + 289 = 634$, rồi $634 - 176 = 458$.
}
\end{ex}

\begin{ex}
Tính giá trị biểu thức: $725 - 368 + 247$.
\shortans{604}
\loigiai{
Ta có $725 - 368 = 357$, rồi $357 + 247 = 604$.
}
\end{ex}

\begin{ex}
Tính giá trị biểu thức: $800 - 245 - 155$.
\shortans{400}
\loigiai{
Ta có $800 - 245 = 555$, rồi $555 - 155 = 400$.
}
\end{ex}

\begin{ex}
Điền số thích hợp: $275 + \dots = 612$.
\shortans{337}
\loigiai{
Số cần điền: $612 - 275 = 337$.
}
\end{ex}

\begin{ex}
Điền số thích hợp: $804 - \dots = 358$.
\shortans{446}
\loigiai{
Số cần điền: $804 - 358 = 446$.
}
\end{ex}

\begin{ex}
Điền số thích hợp: $\dots + 386 = 720$.
\shortans{334}
\loigiai{
Số cần điền: $720 - 386 = 334$.
}
\end{ex}

\begin{ex}
Điền số thích hợp: $\dots - 452 = 389$.
\shortans{841}
\loigiai{
Số cần điền: $389 + 452 = 841$.
}
\end{ex}

\begin{ex}
Lớp 2A quyên góp được 275 quyển vở, lớp 2B quyên góp được 318 quyển vở. Cả hai lớp quyên góp được bao nhiêu quyển vở?
\shortans{593}
\loigiai{
Cả hai lớp quyên góp: $275 + 318 = 593$ quyển vở.
}
\end{ex}

\begin{ex}
Kho gạo có 950 bao gạo, đã xuất bán 475 bao. Kho còn lại bao nhiêu bao gạo?
\shortans{475}
\loigiai{
Còn lại: $950 - 475 = 475$ bao gạo.
}
\end{ex}

\begin{ex}
Trang trại thu hoạch được 620 quả dưa hấu, đã bán 385 quả. Còn lại bao nhiêu quả dưa hấu?
\shortans{235}
\loigiai{
Còn lại: $620 - 385 = 235$ quả dưa hấu.
}
\end{ex}

\begin{ex}
Hai số có tổng bằng 762, số hạng thứ nhất là 385. Số hạng thứ hai là bao nhiêu?
\shortans{377}
\loigiai{
Số hạng thứ hai: $762 - 385 = 377$.
}
\end{ex}

\begin{ex}
Hai số có hiệu bằng 268, số trừ là 354. Số bị trừ là bao nhiêu?
\shortans{622}
\loigiai{
Số bị trừ: $268 + 354 = 622$.
}
\end{ex}

\begin{ex}
Hai số có hiệu bằng 345, số bị trừ là 812. Số trừ là bao nhiêu?
\shortans{467}
\loigiai{
Số trừ: $812 - 345 = 467$.
}
\end{ex}

\begin{ex}
Một xưởng sản xuất tháng thứ nhất được 450 sản phẩm, tháng thứ hai nhiều hơn tháng thứ nhất 75 sản phẩm. Tháng thứ hai sản xuất được bao nhiêu sản phẩm?
\shortans{525}
\loigiai{
Tháng 2: $450 + 75 = 525$ sản phẩm.
}
\end{ex}

\begin{ex}
Cả hai tháng xưởng sản xuất được bao nhiêu sản phẩm?
\shortans{975}
\loigiai{
Cả hai tháng: $450 + 525 = 975$ sản phẩm.
}
\end{ex}

\begin{ex}
Một tàu chở hàng chở 650 tấn hàng, sau khi bốc dỡ ở cảng A thì trên tàu còn 280 tấn hàng. Đã bốc dỡ ở cảng A bao nhiêu tấn hàng?
\shortans{370}
\loigiai{
Đã bốc dỡ: $650 - 280 = 370$ tấn.
}
\end{ex}

\begin{ex}
Tính: $1000 - 300 - 400 + 150$.
\shortans{450}
\loigiai{
Ta có $1000 - 300 - 400 = 300$, rồi $300 + 150 = 450$.
}
\end{ex}

\begin{ex}
Số lớn nhất có ba chữ số trừ đi số bé nhất có ba chữ số bằng bao nhiêu?
\shortans{899}
\loigiai{
Số lớn nhất là 999, bé nhất là 100. Hiệu: $999 - 100 = 899$.
}
\end{ex}

\newpage


% =========================================================================
% Làm quen với yếu tố thống kê, xác suất
% =========================================================================
\section{Làm quen với yếu tố thống kê, xác suất}

\TaoPhieuBaiTap{Chủ đề 13 -- Làm quen với yếu tố thống kê, xác suất}{64}{Phép trừ (có nhớ) trong phạm vi 1000}{https://loptoanthaylinh.vn/bai-tap/toan-2/bai-64}
\begin{ex}
Trong sân có: 5 con gà trống, 8 con gà mái và 4 con vịt. Có tất cả bao nhiêu con gà?
\shortans{13}
\loigiai{
Số con gà là: $5 + 8 = 13$ con.
}
\end{ex}

\begin{ex}
Trong sân có tất cả bao nhiêu con gia cầm (gà và vịt)?
\shortans{17}
\loigiai{
Có tất cả: $5 + 8 + 4 = 17$ con.
}
\end{ex}

\begin{ex}
Số gà mái nhiều hơn số gà trống bao nhiêu con?
\shortans{3}
\loigiai{
Nhiều hơn: $8 - 5 = 3$ con.
}
\end{ex}

\begin{ex}
Trong hộp bút của An có: 6 bút bi xanh, 4 bút bi đỏ và 2 bút chì. Có tất cả bao nhiêu chiếc bút bi?
\shortans{10}
\loigiai{
Số bút bi: $6 + 4 = 10$ chiếc.
}
\end{ex}

\begin{ex}
Hộp bút của An có tất cả bao nhiêu chiếc bút?
\shortans{12}
\loigiai{
Có tất cả: $6 + 4 + 2 = 12$ chiếc.
}
\end{ex}

\begin{ex}
Một rổ quả có: 10 quả cam, 8 quả táo và 6 quả xoài. Hỏi có tất cả bao nhiêu quả trong rổ?
\shortans{24}
\loigiai{
Có tất cả: $10 + 8 + 6 = 24$ quả.
}
\end{ex}

\begin{ex}
Trong rổ quả trên, số cam nhiều hơn số xoài bao nhiêu quả?
\shortans{4}
\loigiai{
Nhiều hơn: $10 - 6 = 4$ quả.
}
\end{ex}

\begin{ex}
Kiểm đếm số học sinh thích các môn thể thao của tổ Một: Bóng đá có 5 bạn, Cầu lông có 3 bạn, Bơi lội có 4 bạn. Có bao nhiêu bạn thích bóng đá hoặc cầu lông?
\shortans{8}
\loigiai{
Có: $5 + 3 = 8$ bạn.
}
\end{ex}

\begin{ex}
Tổ Một có tất cả bao nhiêu bạn tham gia khảo sát?
\shortans{12}
\loigiai{
Có: $5 + 3 + 4 = 12$ bạn.
}
\end{ex}

\begin{ex}
Môn thể thao nào được nhiều bạn thích nhất? Số bạn thích môn đó là bao nhiêu?
\shortans{5}
\loigiai{
Bóng đá nhiều bạn thích nhất với 5 bạn.
}
\end{ex}

\begin{ex}
Môn thể thao nào ít bạn thích nhất? Số bạn thích môn đó là bao nhiêu?
\shortans{3}
\loigiai{
Cầu lông ít bạn thích nhất với 3 bạn.
}
\end{ex}

\begin{ex}
Trong bể cá có: 7 con cá vàng, 5 con cá bảy màu và 3 con cá lau kính. Có tất cả bao nhiêu con cá trong bể?
\shortans{15}
\loigiai{
Có: $7 + 5 + 3 = 15$ con cá.
}
\end{ex}

\begin{ex}
Số cá vàng nhiều hơn số cá lau kính bao nhiêu con?
\shortans{4}
\loigiai{
Nhiều hơn: $7 - 3 = 4$ con.
}
\end{ex}

\begin{ex}
Kiểm đếm phương tiện giao thông đi qua cổng trường trong 5 phút: 12 xe máy, 4 xe ô tô con và 6 xe đạp. Có tất cả bao nhiêu phương tiện đã đi qua?
\shortans{22}
\loigiai{
Có: $12 + 4 + 6 = 22$ phương tiện.
}
\end{ex}

\begin{ex}
Số xe máy nhiều hơn số xe ô tô con bao nhiêu chiếc?
\shortans{8}
\loigiai{
Nhiều hơn: $12 - 4 = 8$ chiếc.
}
\end{ex}

\begin{ex}
Trong vườn có: 15 cây hoa hồng, 12 cây hoa cúc và 8 cây hoa đồng tiền. Có tất cả bao nhiêu cây hoa trong vườn?
\shortans{35}
\loigiai{
Có: $15 + 12 + 8 = 35$ cây hoa.
}
\end{ex}

\begin{ex}
Số hoa hồng nhiều hơn số hoa đồng tiền bao nhiêu cây?
\shortans{7}
\loigiai{
Nhiều hơn: $15 - 8 = 7$ cây.
}
\end{ex}

\begin{ex}
Quan sát giá sách lớp học: có 20 cuốn truyện tranh, 15 cuốn sách Toán, 10 cuốn sách Tiếng Việt. Có tất cả bao nhiêu cuốn sách giáo khoa (Toán và Tiếng Việt)?
\shortans{25}
\loigiai{
Có: $15 + 10 = 25$ cuốn sách giáo khoa.
}
\end{ex}

\begin{ex}
Tổng số sách và truyện trên giá là bao nhiêu cuốn?
\shortans{45}
\loigiai{
Có: $20 + 15 + 10 = 45$ cuốn.
}
\end{ex}

\begin{ex}
Khảo sát thời tiết trong tháng 4: có 18 ngày nắng, 8 ngày mưa và 4 ngày râm mát. Có bao nhiêu ngày không mưa trong tháng 4?
\shortans{22}
\loigiai{
Có: $18 + 4 = 22$ ngày không mưa.
}
\end{ex}

\begin{ex}
Số ngày nắng nhiều hơn số ngày mưa bao nhiêu ngày?
\shortans{10}
\loigiai{
Nhiều hơn: $18 - 8 = 10$ ngày.
}
\end{ex}

\begin{ex}
Trong đĩa hoa quả có: 4 quả chuối, 6 quả mận, 8 quả dâu tây. Tổng số quả là bao nhiêu?
\shortans{18}
\loigiai{
Có: $4 + 6 + 8 = 18$ quả.
}
\end{ex}

\begin{ex}
Có bao nhiêu loại quả trong đĩa hoa quả trên?
\shortans{3}
\loigiai{
Có 3 loại quả: chuối, mận, dâu tây.
}
\end{ex}

\begin{ex}
Số dâu tây gấp mấy lần số quả chuối?
\shortans{2}
\loigiai{
Gấp: $8 : 4 = 2$ lần.
}
\end{ex}

\begin{ex}
Nếu ăn hết 3 quả mận thì đĩa còn lại bao nhiêu quả?
\shortans{15}
\loigiai{
Còn lại: $18 - 3 = 15$ quả.
}
\end{ex}

\newpage

\TaoPhieuBaiTap{Chủ đề 13 -- Làm quen với yếu tố thống kê, xác suất}{65}{Luyện tập chung phép tính 1000}{https://loptoanthaylinh.vn/bai-tap/toan-2/bai-65}
\begin{ex}
Trên biểu đồ tranh, mỗi biểu tượng ngôi sao tương ứng với 1 điểm tốt. Bạn Mai có 8 ngôi sao. Bạn Mai có bao nhiêu điểm tốt?
\shortans{8}
\loigiai{
Mai có 8 điểm tốt.
}
\end{ex}

\begin{ex}
Bạn Lan có 6 ngôi sao biểu thị điểm tốt. Bạn Nam có 9 ngôi sao. Cả hai bạn có bao nhiêu điểm tốt?
\shortans{15}
\loigiai{
Cả hai bạn có: $6 + 9 = 15$ điểm tốt.
}
\end{ex}

\begin{ex}
Nam nhiều hơn Lan bao nhiêu điểm tốt?
\shortans{3}
\loigiai{
Nhiều hơn: $9 - 6 = 3$ điểm tốt.
}
\end{ex}

\begin{ex}
Trên biểu đồ tranh, mỗi hình tròn biểu thị 2 cây xanh trồng được. Lớp 2A có 5 hình tròn. Lớp 2A trồng được bao nhiêu cây xanh?
\shortans{10}
\loigiai{
Lớp 2A trồng: $2 \times 5 = 10$ cây xanh.
}
\end{ex}

\begin{ex}
Lớp 2B có 7 hình tròn biểu thị cây trồng được (mỗi hình tròn là 2 cây). Lớp 2B trồng được bao nhiêu cây xanh?
\shortans{14}
\loigiai{
Lớp 2B trồng: $2 \times 7 = 14$ cây xanh.
}
\end{ex}

\begin{ex}
Lớp 2B trồng nhiều hơn lớp 2A bao nhiêu cây xanh?
\shortans{4}
\loigiai{
Nhiều hơn: $14 - 10 = 4$ cây.
}
\end{ex}

\begin{ex}
Cả hai lớp trồng được tất cả bao nhiêu cây xanh?
\shortans{24}
\loigiai{
Cả hai lớp trồng: $10 + 14 = 24$ cây.
}
\end{ex}

\begin{ex}
Biểu đồ tranh số con vật nuôi: mỗi hình con vật biểu thị 5 con. Đàn gà có 6 hình con gà. Đàn gà có bao nhiêu con gà?
\shortans{30}
\loigiai{
Có: $5 \times 6 = 30$ con gà.
}
\end{ex}

\begin{ex}
Đàn vịt có 4 hình con vịt (mỗi hình là 5 con). Đàn vịt có bao nhiêu con vịt?
\shortans{20}
\loigiai{
Có: $5 \times 4 = 20$ con vịt.
}
\end{ex}

\begin{ex}
Đàn gà nhiều hơn đàn vịt bao nhiêu con?
\shortans{10}
\loigiai{
Nhiều hơn: $30 - 20 = 10$ con.
}
\end{ex}

\begin{ex}
Cả gà và vịt có tất cả bao nhiêu con?
\shortans{50}
\loigiai{
Cả hai đàn có: $30 + 20 = 50$ con.
}
\end{ex}

\begin{ex}
Biểu đồ tranh hoa tặng cô giáo: Tổ Một có 4 bông hoa biểu tượng, Tổ Hai có 6 bông hoa biểu tượng, Tổ Ba có 5 bông hoa biểu tượng. Mỗi biểu tượng là 1 bông hoa thật. Cả 3 tổ tặng bao nhiêu bông hoa?
\shortans{15}
\loigiai{
Tặng: $4 + 6 + 5 = 15$ bông hoa.
}
\end{ex}

\begin{ex}
Tổ Hai tặng nhiều hơn Tổ Một bao nhiêu bông hoa?
\shortans{2}
\loigiai{
Nhiều hơn: $6 - 4 = 2$ bông hoa.
}
\end{ex}

\begin{ex}
Tổ nào tặng nhiều hoa nhất? Số hoa của tổ đó là bao nhiêu?
\shortans{6}
\loigiai{
Tổ Hai nhiều nhất với 6 bông hoa.
}
\end{ex}

\begin{ex}
Tổ nào tặng ít hoa nhất? Số hoa của tổ đó là bao nhiêu?
\shortans{4}
\loigiai{
Tổ Một ít nhất với 4 bông hoa.
}
\end{ex}

\begin{ex}
Biểu đồ số xe bán được trong tuần: Thứ Hai bán được 3 xe, Thứ Ba bán được 5 xe, Thứ Tư bán được 4 xe. Cả 3 ngày bán được bao nhiêu xe?
\shortans{12}
\loigiai{
Bán được: $3 + 5 + 4 = 12$ chiếc xe.
}
\end{ex}

\begin{ex}
Ngày Thứ Ba bán nhiều hơn ngày Thứ Hai bao nhiêu chiếc xe?
\shortans{2}
\loigiai{
Nhiều hơn: $5 - 3 = 2$ chiếc xe.
}
\end{ex}

\begin{ex}
Biểu đồ số bạn học sinh thích đọc sách: mỗi biểu tượng mặt cười là 2 bạn. Thể loại truyện cổ tích có 4 mặt cười. Có bao nhiêu bạn thích truyện cổ tích?
\shortans{8}
\loigiai{
Có: $2 \times 4 = 8$ bạn.
}
\end{ex}

\begin{ex}
Thể loại truyện tranh có 6 mặt cười (mỗi mặt cười là 2 bạn). Có bao nhiêu bạn thích truyện tranh?
\shortans{12}
\loigiai{
Có: $2 \times 6 = 12$ bạn.
}
\end{ex}

\begin{ex}
Số bạn thích truyện tranh nhiều hơn truyện cổ tích bao nhiêu bạn?
\shortans{4}
\loigiai{
Nhiều hơn: $12 - 8 = 4$ bạn.
}
\end{ex}

\begin{ex}
Tổng số bạn thích cả hai thể loại truyện trên là bao nhiêu bạn?
\shortans{20}
\loigiai{
Tổng số: $8 + 12 = 20$ bạn.
}
\end{ex}

\begin{ex}
Biểu đồ hoa quả bán trong ngày: Táo có 5 giỏ, Cam có 7 giỏ, Lê có 3 giỏ. Mỗi giỏ có 10 quả. Có bao nhiêu quả lê đã bán?
\shortans{30}
\loigiai{
Có: $10 \times 3 = 30$ quả lê.
}
\end{ex}

\begin{ex}
Có bao nhiêu quả cam đã bán?
\shortans{70}
\loigiai{
Có: $10 \times 7 = 70$ quả cam.
}
\end{ex}

\begin{ex}
Số quả cam bán nhiều hơn số quả táo bao nhiêu quả?
\shortans{20}
\loigiai{
Nhiều hơn: $70 - 50 = 20$ quả.
}
\end{ex}

\begin{ex}
Tổng số quả của cả 3 loại đã bán là bao nhiêu quả?
\shortans{150}
\loigiai{
Tổng số quả: $50 + 70 + 30 = 150$ quả.
}
\end{ex}

\newpage

\TaoPhieuBaiTap{Chủ đề 13 -- Làm quen với yếu tố thống kê, xác suất}{66}{Thu thập, phân loại, kiểm đếm số liệu}{https://loptoanthaylinh.vn/bai-tap/toan-2/bai-66}
\begin{ex}
Trong hộp chỉ có 5 viên bi màu đỏ. Nếu nhắm mắt lấy ra 1 viên bi, sự kiện 'Lấy được viên bi màu đỏ' là chắc chắn, có thể hay không thể? Quy ước: Chắc chắn = 1, Có thể = 2, Không thể = 3.
\shortans{1}
\loigiai{
Vì toàn bi đỏ nên chắc chắn lấy được bi đỏ. Đáp án là 1.
}
\end{ex}

\begin{ex}
Trong hộp đó chỉ có bi đỏ, sự kiện 'Lấy được viên bi màu xanh' xảy ra như thế nào? (Quy ước: Chắc chắn = 1, Có thể = 2, Không thể = 3).
\shortans{3}
\loigiai{
Không có bi xanh nên không thể xảy ra. Đáp án là 3.
}
\end{ex}

\begin{ex}
Trong túi có 3 viên bi đỏ và 3 viên bi xanh. Nhắm mắt lấy 1 viên bi, sự kiện 'Lấy được viên bi màu đỏ' là: Chắc chắn = 1, Có thể = 2, Không thể = 3.
\shortans{2}
\loigiai{
Có thể lấy được bi đỏ hoặc bi xanh. Đáp án là 2.
}
\end{ex}

\begin{ex}
Gieo một con xúc xắc 6 mặt (từ 1 chấm đến 6 chấm). Sự kiện 'Xuất hiện mặt có 7 chấm' xảy ra như thế nào? (Chắc chắn = 1, Có thể = 2, Không thể = 3).
\shortans{3}
\loigiai{
Xúc xắc chỉ có từ 1 đến 6 chấm nên không thể có mặt 7 chấm. Đáp án là 3.
}
\end{ex}

\begin{ex}
Gieo một con xúc xắc 6 mặt. Sự kiện 'Xuất hiện mặt có số chấm nhỏ hơn 7' là: (Chắc chắn = 1, Có thể = 2, Không thể = 3).
\shortans{1}
\loigiai{
Mọi mặt đều từ 1 đến 6 (nhỏ hơn 7) nên chắc chắn xảy ra. Đáp án là 1.
}
\end{ex}

\begin{ex}
Gieo một con xúc xắc. Sự kiện 'Xuất hiện mặt 6 chấm' là: (Chắc chắn = 1, Có thể = 2, Không thể = 3).
\shortans{2}
\loigiai{
Có thể xuất hiện mặt 6 chấm. Đáp án là 2.
}
\end{ex}

\begin{ex}
Ngày mai Mặt Trời mọc ở hướng Đông. Sự kiện này là: (Chắc chắn = 1, Có thể = 2, Không thể = 3).
\shortans{1}
\loigiai{
Mặt Trời luôn mọc ở hướng Đông nên chắc chắn xảy ra. Đáp án là 1.
}
\end{ex}

\begin{ex}
Ngày mai Mặt Trời lặn ở hướng Đông. Sự kiện này là: (Chắc chắn = 1, Có thể = 2, Không thể = 3).
\shortans{3}
\loigiai{
Mặt Trời lặn ở hướng Tây nên không thể xảy ra. Đáp án là 3.
}
\end{ex}

\begin{ex}
Một tuần lễ có 7 ngày. Sự kiện này là: (Chắc chắn = 1, Có thể = 2, Không thể = 3).
\shortans{1}
\loigiai{
Đúng quy luật tự nhiên nên chắc chắn xảy ra. Đáp án là 1.
}
\end{ex}

\begin{ex}
Trong đĩa có 5 quả cam và 4 quả táo. Nam lấy ra 6 quả. Sự kiện 'Trong số quả lấy ra có ít nhất 1 quả cam' là: (Chắc chắn = 1, Có thể = 2, Không thể = 3).
\shortans{1}
\loigiai{
Vì chỉ có 4 quả táo nên lấy 6 quả thì chắc chắn phải có quả cam. Đáp án là 1.
}
\end{ex}

\begin{ex}
Trong hộp có 10 chiếc bút chì màu đen. Lấy ra 2 chiếc bút, cả 2 chiếc đều màu đen là: (Chắc chắn = 1, Có thể = 2, Không thể = 3).
\shortans{1}
\loigiai{
Chắc chắn vì toàn bộ bút trong hộp đều màu đen. Đáp án là 1.
}
\end{ex}

\begin{ex}
Bé Nam 8 tuổi cao 3 mét. Sự kiện này là: (Chắc chắn = 1, Có thể = 2, Không thể = 3).
\shortans{3}
\loigiai{
Không thể xảy ra đối với con người. Đáp án là 3.
}
\end{ex}

\begin{ex}
Hôm nay là Thứ Bảy, ngày mai là Chủ Nhật. Sự kiện này là: (Chắc chắn = 1, Có thể = 2, Không thể = 3).
\shortans{1}
\loigiai{
Chắc chắn xảy ra theo thứ tự các ngày trong tuần. Đáp án là 1.
}
\end{ex}

\begin{ex}
Tung một đồng xu có hai mặt Sấp và Ngửa. Sự kiện 'Đồng xu rơi xuống xuất hiện mặt Sấp' là: (Chắc chắn = 1, Có thể = 2, Không thể = 3).
\shortans{2}
\loigiai{
Có thể xuất hiện mặt Sấp hoặc mặt Ngửa. Đáp án là 2.
}
\end{ex}

\begin{ex}
Trong túi có 5 quả bóng xanh và 5 quả bóng vàng. Lấy ra 6 quả bóng, sự kiện 'Có ít nhất 1 quả bóng xanh' là: (Chắc chắn = 1, Có thể = 2, Không thể = 3).
\shortans{1}
\loigiai{
Chắc chắn vì bóng vàng chỉ có 5 quả, lấy 6 quả chắc chắn có bóng xanh. Đáp án là 1.
}
\end{ex}

\begin{ex}
Lấy ra từ túi đó 1 quả bóng màu tím. Sự kiện này là: (Chắc chắn = 1, Có thể = 2, Không thể = 3).
\shortans{3}
\loigiai{
Trong túi không có bóng tím nên không thể. Đáp án là 3.
}
\end{ex}

\begin{ex}
Tháng Hai có 35 ngày. Sự kiện này là: (Chắc chắn = 1, Có thể = 2, Không thể = 3).
\shortans{3}
\loigiai{
Tháng Hai nhiều nhất chỉ có 29 ngày nên không thể có 35 ngày. Đáp án là 3.
}
\end{ex}

\begin{ex}
Một năm có 12 tháng. Sự kiện này là: (Chắc chắn = 1, Có thể = 2, Không thể = 3).
\shortans{1}
\loigiai{
Chắc chắn đúng. Đáp án là 1.
}
\end{ex}

\begin{ex}
Số 24 là số chẵn. Sự kiện này là: (Chắc chắn = 1, Có thể = 2, Không thể = 3).
\shortans{1}
\loigiai{
Chắc chắn đúng. Đáp án là 1.
}
\end{ex}

\begin{ex}
Tổng của hai số lẻ là một số lẻ. Sự kiện này là: (Chắc chắn = 1, Có thể = 2, Không thể = 3).
\shortans{3}
\loigiai{
Tổng hai số lẻ luôn là số chẵn nên không thể là số lẻ. Đáp án là 3.
}
\end{ex}

\begin{ex}
Chọn ngẫu nhiên một bạn học sinh trong lớp 2A toàn nữ. Bạn được chọn là học sinh nữ: (Chắc chắn = 1, Có thể = 2, Không thể = 3).
\shortans{1}
\loigiai{
Toàn nữ nên chắc chắn là học sinh nữ. Đáp án là 1.
}
\end{ex}

\begin{ex}
Bạn được chọn là học sinh nam: (Chắc chắn = 1, Có thể = 2, Không thể = 3).
\shortans{3}
\loigiai{
Không có bạn nam nào nên không thể. Đáp án là 3.
}
\end{ex}

\begin{ex}
Thời tiết ngày mai có mưa: (Chắc chắn = 1, Có thể = 2, Không thể = 3).
\shortans{2}
\loigiai{
Có thể mưa hoặc không mưa. Đáp án là 2.
}
\end{ex}

\begin{ex}
Có bao nhiêu sự kiện 'Chắc chắn' trong 3 sự kiện sau: (A) Cá sống dưới nước; (B) Con người bay bằng cánh tay; (C) 1 ngày có 24 giờ?
\shortans{2}
\loigiai{
Sự kiện A và C chắc chắn (2 sự kiện).
}
\end{ex}

\begin{ex}
Có bao nhiêu sự kiện 'Không thể' trong 3 sự kiện trên?
\shortans{1}
\loigiai{
Sự kiện B không thể (1 sự kiện).
}
\end{ex}

\newpage

\TaoPhieuBaiTap{Chủ đề 13 -- Làm quen với yếu tố thống kê, xác suất}{67}{Biểu đồ tranh}{https://loptoanthaylinh.vn/bai-tap/toan-2/bai-67}
\begin{ex}
Trong giờ ra chơi, kiểm đếm số bạn chơi các trò: Nhảy dây có 8 bạn, Đá cầu có 6 bạn, Đọc sách có 5 bạn. Có tất cả bao nhiêu bạn tham gia các trò chơi?
\shortans{19}
\loigiai{
Có tất cả: $8 + 6 + 5 = 19$ bạn.
}
\end{ex}

\begin{ex}
Số bạn chơi nhảy dây nhiều hơn số bạn chơi đá cầu bao nhiêu bạn?
\shortans{2}
\loigiai{
Nhiều hơn: $8 - 6 = 2$ bạn.
}
\end{ex}

\begin{ex}
Số bạn chơi nhảy dây nhiều hơn số bạn đọc sách bao nhiêu bạn?
\shortans{3}
\loigiai{
Nhiều hơn: $8 - 5 = 3$ bạn.
}
\end{ex}

\begin{ex}
Kiểm đếm số cây trồng trong vườn trường: Cây phượng có 6 cây, Cây bàng có 4 cây, Cây bằng lăng có 5 cây. Có bao nhiêu cây bóng mát (phượng và bàng)?
\shortans{10}
\loigiai{
Có: $6 + 4 = 10$ cây.
}
\end{ex}

\begin{ex}
Vườn trường có tất cả bao nhiêu cây?
\shortans{15}
\loigiai{
Có tất cả: $6 + 4 + 5 = 15$ cây.
}
\end{ex}

\begin{ex}
Khảo sát bữa sáng của 25 bạn học sinh: 10 bạn ăn bánh mì, 8 bạn ăn xôi, còn lại ăn phở. Có bao nhiêu bạn ăn phở?
\shortans{7}
\loigiai{
Số bạn ăn phở: $25 - 10 - 8 = 7$ bạn.
}
\end{ex}

\begin{ex}
Số bạn ăn bánh mì nhiều hơn số bạn ăn xôi bao nhiêu bạn?
\shortans{2}
\loigiai{
Nhiều hơn: $10 - 8 = 2$ bạn.
}
\end{ex}

\begin{ex}
Kiểm đếm dụng cụ học tập trong cặp của Lan: 5 quyển vở, 3 cuốn sách giáo khoa, 2 cuốn truyện tranh. Có tất cả bao nhiêu cuốn sách và vở?
\shortans{10}
\loigiai{
Có tất cả: $5 + 3 + 2 = 10$ cuốn.
}
\end{ex}

\begin{ex}
Lan có 4 chiếc bút chì, 2 chiếc bút mực và 1 chiếc thước kẻ. Lan có bao nhiêu chiếc bút các loại?
\shortans{6}
\loigiai{
Số bút: $4 + 2 = 6$ chiếc.
}
\end{ex}

\begin{ex}
Thống kê số điểm 10 môn Toán của tổ Hai trong tuần: Thứ Hai được 4 điểm 10, Thứ Ba được 5 điểm 10, Thứ Tư được 3 điểm 10, Thứ Năm được 6 điểm 10, Thứ Sáu được 7 điểm 10. Cả tuần tổ Hai đạt được bao nhiêu điểm 10?
\shortans{25}
\loigiai{
Tổng số điểm 10: $4 + 5 + 3 + 6 + 7 = 25$ điểm 10.
}
\end{ex}

\begin{ex}
Ngày nào đạt được nhiều điểm 10 nhất? Số điểm 10 ngày đó là bao nhiêu?
\shortans{7}
\loigiai{
Thứ Sáu nhiều nhất với 7 điểm 10.
}
\end{ex}

\begin{ex}
Ngày nào đạt ít điểm 10 nhất? Số điểm 10 ngày đó là bao nhiêu?
\shortans{3}
\loigiai{
Thứ Tư ít nhất với 3 điểm 10.
}
\end{ex}

\begin{ex}
Chênh lệch giữa ngày đạt nhiều điểm 10 nhất và ít điểm 10 nhất là bao nhiêu điểm?
\shortans{4}
\loigiai{
Chênh lệch: $7 - 3 = 4$ điểm.
}
\end{ex}

\begin{ex}
Thống kê số thú cưng nuôi tại nhà của 20 bạn: 8 bạn nuôi chó, 7 bạn nuôi mèo, 3 bạn nuôi chim, 2 bạn nuôi cá cảnh. Có bao nhiêu bạn nuôi chó hoặc mèo?
\shortans{15}
\loigiai{
Có: $8 + 7 = 15$ bạn.
}
\end{ex}

\begin{ex}
Số bạn nuôi mèo nhiều hơn số bạn nuôi chim bao nhiêu bạn?
\shortans{4}
\loigiai{
Nhiều hơn: $7 - 3 = 4$ bạn.
}
\end{ex}

\begin{ex}
Có bao nhiêu bạn nuôi chim hoặc cá cảnh?
\shortans{5}
\loigiai{
Có: $3 + 2 = 5$ bạn.
}
\end{ex}

\begin{ex}
Kiểm đếm các loại rác tái chế gom được: 15 vỏ lon nhôm, 20 chai nhựa và 10 tờ bìa carton. Có tất cả bao nhiêu đồ vật tái chế gom được?
\shortans{45}
\loigiai{
Có tất cả: $15 + 20 + 10 = 45$ đồ vật.
}
\end{ex}

\begin{ex}
Số chai nhựa nhiều hơn số vỏ lon bao nhiêu chiếc?
\shortans{5}
\loigiai{
Nhiều hơn: $20 - 15 = 5$ chiếc.
}
\end{ex}

\begin{ex}
Trong tủ thuốc gia đình có: 4 vỉ thuốc cảm, 3 cuộn băng gạc, 2 lọ thuốc ho. Có tất cả bao nhiêu đồ dùng y tế?
\shortans{9}
\loigiai{
Có tất cả: $4 + 3 + 2 = 9$ đồ dùng.
}
\end{ex}

\begin{ex}
Khảo sát mùa yêu thích của 30 bạn: 12 bạn thích mùa hè, 8 bạn thích mùa thu, 6 bạn thích mùa xuân, còn lại thích mùa đông. Có bao nhiêu bạn thích mùa đông?
\shortans{4}
\loigiai{
Số bạn thích mùa đông: $30 - 12 - 8 - 6 = 4$ bạn.
}
\end{ex}

\begin{ex}
Số bạn thích mùa hè nhiều hơn mùa thu bao nhiêu bạn?
\shortans{4}
\loigiai{
Nhiều hơn: $12 - 8 = 4$ bạn.
}
\end{ex}

\begin{ex}
Số bạn thích mùa xuân nhiều hơn mùa đông bao nhiêu bạn?
\shortans{2}
\loigiai{
Nhiều hơn: $6 - 4 = 2$ bạn.
}
\end{ex}

\begin{ex}
Trong vườn có 3 tổ chim: tổ 1 có 4 con non, tổ 2 có 3 con non, tổ 3 có 5 con non. Cả 3 tổ có bao nhiêu con chim non?
\shortans{12}
\loigiai{
Có: $4 + 3 + 5 = 12$ con chim non.
}
\end{ex}

\begin{ex}
Tổ chim đông con nhất có bao nhiêu con non?
\shortans{5}
\loigiai{
Tổ 3 đông nhất có 5 con non.
}
\end{ex}

\begin{ex}
Tổ chim ít con nhất có bao nhiêu con non?
\shortans{3}
\loigiai{
Tổ 2 ít nhất có 3 con non.
}
\end{ex}

\newpage


% =========================================================================
% Ôn tập cuối năm
% =========================================================================
\section{Ôn tập cuối năm}

\TaoPhieuBaiTap{Chủ đề 14 -- Ôn tập cuối năm}{68}{Chắc chắn, có thể, không thể}{https://loptoanthaylinh.vn/bai-tap/toan-2/bai-68}
\begin{ex}
Số liền sau của số 999 là số nào?
\shortans{1000}
\loigiai{
Số liền sau 999 là 1000.
}
\end{ex}

\begin{ex}
Số liền trước của số 1000 là số nào?
\shortans{999}
\loigiai{
Số liền trước 1000 là 999.
}
\end{ex}

\begin{ex}
Số gồm 7 trăm, 8 chục và 4 đơn vị là số nào?
\shortans{784}
\loigiai{
Số đó là 784.
}
\end{ex}

\begin{ex}
Số 605 gồm 6 trăm và mấy đơn vị?
\shortans{5}
\loigiai{
Số 605 gồm 6 trăm và 5 đơn vị.
}
\end{ex}

\begin{ex}
Số 830 có chữ số hàng chục là bao nhiêu?
\shortans{3}
\loigiai{
Chữ số hàng chục là 3.
}
\end{ex}

\begin{ex}
Số tròn trăm lớn nhất có ba chữ số là số nào?
\shortans{900}
\loigiai{
Số đó là 900.
}
\end{ex}

\begin{ex}
Số tròn trăm bé nhất có ba chữ số là số nào?
\shortans{100}
\loigiai{
Số đó là 100.
}
\end{ex}

\begin{ex}
Số lớn nhất có ba chữ số là số nào?
\shortans{999}
\loigiai{
Số đó là 999.
}
\end{ex}

\begin{ex}
Số bé nhất có ba chữ số là số nào?
\shortans{100}
\loigiai{
Số đó là 100.
}
\end{ex}

\begin{ex}
Số bé nhất có ba chữ số khác nhau là số nào?
\shortans{102}
\loigiai{
Số đó là 102.
}
\end{ex}

\begin{ex}
Số lớn nhất có ba chữ số khác nhau là số nào?
\shortans{987}
\loigiai{
Số đó là 987.
}
\end{ex}

\begin{ex}
Trong các số $485, 521, 498, 512$, số nào lớn nhất?
\shortans{521}
\loigiai{
Số lớn nhất là 521.
}
\end{ex}

\begin{ex}
Trong các số $735, 762, 719, 740$, số nào bé nhất?
\shortans{719}
\loigiai{
Số bé nhất là 719.
}
\end{ex}

\begin{ex}
Sắp xếp dãy số $352, 480, 299, 415$ từ bé đến lớn, số đứng ở vị trí thứ ba là số nào?
\shortans{415}
\loigiai{
Thứ tự: 299, 352, 415, 480. Số thứ ba là 415.
}
\end{ex}

\begin{ex}
Điền số thích hợp: $698, 699, \dots, 701, 702$.
\shortans{700}
\loigiai{
Số còn thiếu là 700.
}
\end{ex}

\begin{ex}
Điền số thích hợp: $850, 860, \dots, 880, 890$.
\shortans{870}
\loigiai{
Số còn thiếu là 870.
}
\end{ex}

\begin{ex}
Số 458 viết thành tổng các trăm, chục, đơn vị: $400 + \dots + 8$. Số cần điền là bao nhiêu?
\shortans{50}
\loigiai{
Số cần điền là 50.
}
\end{ex}

\begin{ex}
Tính tổng các chữ số của số 345.
\shortans{12}
\loigiai{
Tổng: $3 + 4 + 5 = 12$.
}
\end{ex}

\begin{ex}
Tính tổng các chữ số của số 809.
\shortans{17}
\loigiai{
Tổng: $8 + 0 + 9 = 17$.
}
\end{ex}

\begin{ex}
Tìm số tự nhiên ở giữa số 599 và số 601.
\shortans{600}
\loigiai{
Số ở giữa là 600.
}
\end{ex}

\begin{ex}
Có bao nhiêu số tròn chục lớn hơn 250 và nhỏ hơn 300?
\shortans{4}
\loigiai{
Các số: 260, 270, 280, 290. Có 4 số.
}
\end{ex}

\begin{ex}
Có bao nhiêu số tròn trăm từ 300 đến 700?
\shortans{5}
\loigiai{
Các số: 300, 400, 500, 600, 700. Có 5 số.
}
\end{ex}

\begin{ex}
Giá trị của chữ số 7 trong số 742 là bao nhiêu?
\shortans{700}
\loigiai{
Chữ số 7 ở hàng trăm có giá trị là 700.
}
\end{ex}

\begin{ex}
Giá trị của chữ số 4 trong số 742 là bao nhiêu?
\shortans{40}
\loigiai{
Chữ số 4 ở hàng chục có giá trị là 40.
}
\end{ex}

\begin{ex}
Số 1000 có bao nhiêu chữ số 0?
\shortans{3}
\loigiai{
Số 1000 có 3 chữ số 0.
}
\end{ex}

\newpage

\TaoPhieuBaiTap{Chủ đề 14 -- Ôn tập cuối năm}{69}{Luyện tập chung thống kê và xác suất}{https://loptoanthaylinh.vn/bai-tap/toan-2/bai-69}
\begin{ex}
Tính: $35 + 47$.
\shortans{82}
\loigiai{
Ta có $35 + 47 = 82$.
}
\end{ex}

\begin{ex}
Tính: $82 - 35$.
\shortans{47}
\loigiai{
Ta có $82 - 35 = 47$.
}
\end{ex}

\begin{ex}
Tính: $48 + 36$.
\shortans{84}
\loigiai{
Ta có $48 + 36 = 84$.
}
\end{ex}

\begin{ex}
Tính: $84 - 48$.
\shortans{36}
\loigiai{
Ta có $84 - 48 = 36$.
}
\end{ex}

\begin{ex}
Tính: $59 + 25$.
\shortans{84}
\loigiai{
Ta có $59 + 25 = 84$.
}
\end{ex}

\begin{ex}
Tính: $73 - 38$.
\shortans{35}
\loigiai{
Ta có $73 - 38 = 35$.
}
\end{ex}

\begin{ex}
Tính: $60 - 24$.
\shortans{36}
\loigiai{
Ta có $60 - 24 = 36$.
}
\end{ex}

\begin{ex}
Tính: $100 - 37$.
\shortans{63}
\loigiai{
Ta có $100 - 37 = 63$.
}
\end{ex}

\begin{ex}
Tính: $100 - 65$.
\shortans{35}
\loigiai{
Ta có $100 - 65 = 35$.
}
\end{ex}

\begin{ex}
Tính giá trị biểu thức: $38 + 27 - 45$.
\shortans{20}
\loigiai{
Ta có $38 + 27 = 65$, rồi $65 - 45 = 20$.
}
\end{ex}

\begin{ex}
Tính giá trị biểu thức: $91 - 46 + 25$.
\shortans{70}
\loigiai{
Ta có $91 - 46 = 45$, rồi $45 + 25 = 70$.
}
\end{ex}

\begin{ex}
Tính giá trị biểu thức: $54 + 38 - 62$.
\shortans{30}
\loigiai{
Ta có $54 + 38 = 92$, rồi $92 - 62 = 30$.
}
\end{ex}

\begin{ex}
Điền số thích hợp: $46 + \dots = 81$.
\shortans{35}
\loigiai{
Số cần điền: $81 - 46 = 35$.
}
\end{ex}

\begin{ex}
Điền số thích hợp: $72 - \dots = 38$.
\shortans{34}
\loigiai{
Số cần điền: $72 - 38 = 34$.
}
\end{ex}

\begin{ex}
Điền số thích hợp: $\dots + 29 = 65$.
\shortans{36}
\loigiai{
Số cần điền: $65 - 29 = 36$.
}
\end{ex}

\begin{ex}
Điền số thích hợp: $\dots - 37 = 48$.
\shortans{85}
\loigiai{
Số cần điền: $48 + 37 = 85$.
}
\end{ex}

\begin{ex}
Lớp 2A có 18 bạn nam và 19 bạn nữ. Hỏi lớp 2A có bao nhiêu bạn học sinh?
\shortans{37}
\loigiai{
Cả lớp có: $18 + 19 = 37$ bạn.
}
\end{ex}

\begin{ex}
Bác nông dân nuôi 52 con gà, bác bán đi 28 con gà. Bác còn lại bao nhiêu con gà?
\shortans{24}
\loigiai{
Còn lại: $52 - 28 = 24$ con gà.
}
\end{ex}

\begin{ex}
Đoạn dây dài 65 cm, người ta cắt đi 38 cm. Đoạn dây còn lại dài bao nhiêu xăng-ti-mét?
\shortans{27}
\loigiai{
Còn lại: $65 - 38 = 27$ cm.
}
\end{ex}

\begin{ex}
Một can dầu có 40 lít, sau khi dùng còn lại 16 lít. Đã dùng bao nhiêu lít dầu?
\shortans{24}
\loigiai{
Đã dùng: $40 - 16 = 24$ lít.
}
\end{ex}

\begin{ex}
Hai số có tổng bằng 83, số thứ nhất là 45. Số thứ hai là bao nhiêu?
\shortans{38}
\loigiai{
Số thứ hai là: $83 - 45 = 38$.
}
\end{ex}

\begin{ex}
Hai số có hiệu bằng 26, số trừ là 37. Số bị trừ là bao nhiêu?
\shortans{63}
\loigiai{
Số bị trừ là: $26 + 37 = 63$.
}
\end{ex}

\begin{ex}
Hiệu của số lớn nhất có hai chữ số và số 48 là bao nhiêu?
\shortans{51}
\loigiai{
Hiệu: $99 - 48 = 51$.
}
\end{ex}

\begin{ex}
Tính tổng của số tròn chục lớn nhất có hai chữ số và số 9.
\shortans{99}
\loigiai{
Số tròn chục lớn nhất có 2 chữ số là 90. Tổng: $90 + 9 = 99$.
}
\end{ex}

\begin{ex}
Tính: $28 + 17 + 35$.
\shortans{80}
\loigiai{
Ta có $28 + 17 = 45$, rồi $45 + 35 = 80$.
}
\end{ex}

\newpage

\TaoPhieuBaiTap{Chủ đề 14 -- Ôn tập cuối năm}{70}{Ôn tập các số trong phạm vi 1000}{https://loptoanthaylinh.vn/bai-tap/toan-2/bai-70}
\begin{ex}
Tính: $345 + 234$.
\shortans{579}
\loigiai{
Ta có $345 + 234 = 579$.
}
\end{ex}

\begin{ex}
Tính: $678 - 345$.
\shortans{333}
\loigiai{
Ta có $678 - 345 = 333$.
}
\end{ex}

\begin{ex}
Tính: $456 + 278$.
\shortans{734}
\loigiai{
Ta có $456 + 278 = 734$.
}
\end{ex}

\begin{ex}
Tính: $734 - 456$.
\shortans{278}
\loigiai{
Ta có $734 - 456 = 278$.
}
\end{ex}

\begin{ex}
Tính: $528 + 365$.
\shortans{893}
\loigiai{
Ta có $528 + 365 = 893$.
}
\end{ex}

\begin{ex}
Tính: $820 - 456$.
\shortans{364}
\loigiai{
Ta có $820 - 456 = 364$.
}
\end{ex}

\begin{ex}
Tính: $600 - 245$.
\shortans{355}
\loigiai{
Ta có $600 - 245 = 355$.
}
\end{ex}

\begin{ex}
Tính: $1000 - 368$.
\shortans{632}
\loigiai{
Ta có $1000 - 368 = 632$.
}
\end{ex}

\begin{ex}
Tính: $1000 - 750$.
\shortans{250}
\loigiai{
Ta có $1000 - 750 = 250$.
}
\end{ex}

\begin{ex}
Tính giá trị biểu thức: $425 + 285 - 310$.
\shortans{400}
\loigiai{
Ta có $425 + 285 = 710$, rồi $710 - 310 = 400$.
}
\end{ex}

\begin{ex}
Tính giá trị biểu thức: $850 - 360 + 150$.
\shortans{640}
\loigiai{
Ta có $850 - 360 = 490$, rồi $490 + 150 = 640$.
}
\end{ex}

\begin{ex}
Tính giá trị biểu thức: $1000 - 450 - 250$.
\shortans{300}
\loigiai{
Ta có $1000 - 450 = 550$, rồi $550 - 250 = 300$.
}
\end{ex}

\begin{ex}
Điền số thích hợp: $350 + \dots = 800$.
\shortans{450}
\loigiai{
Số cần điền: $800 - 350 = 450$.
}
\end{ex}

\begin{ex}
Điền số thích hợp: $920 - \dots = 460$.
\shortans{460}
\loigiai{
Số cần điền: $920 - 460 = 460$.
}
\end{ex}

\begin{ex}
Điền số thích hợp: $\dots + 248 = 615$.
\shortans{367}
\loigiai{
Số cần điền: $615 - 248 = 367$.
}
\end{ex}

\begin{ex}
Điền số thích hợp: $\dots - 375 = 425$.
\shortans{800}
\loigiai{
Số cần điền: $425 + 375 = 800$.
}
\end{ex}

\begin{ex}
Trường Tiểu học Chu Văn An có 540 học sinh, trường Tiểu học Trần Phú có 430 học sinh. Cả hai trường có bao nhiêu học sinh?
\shortans{970}
\loigiai{
Cả hai trường có: $540 + 430 = 970$ học sinh.
}
\end{ex}

\begin{ex}
Kho A có 725 bao gạo, kho B có ít hơn kho A 150 bao gạo. Kho B có bao nhiêu bao gạo?
\shortans{575}
\loigiai{
Kho B có: $725 - 150 = 575$ bao gạo.
}
\end{ex}

\begin{ex}
Bác nông dân thu hoạch được 850 kg lúa, bác đã bán đi 560 kg. Bác còn lại bao nhiêu ki-lô-gam lúa?
\shortans{290}
\loigiai{
Bác còn lại: $850 - 560 = 290$ kg.
}
\end{ex}

\begin{ex}
Quãng đường từ Hà Nội đến Thanh Hóa dài khoảng 160 km, từ Thanh Hóa đến Vinh dài khoảng 140 km. Quãng đường từ Hà Nội đến Vinh đi qua Thanh Hóa dài bao nhiêu ki-lô-mét?
\shortans{300}
\loigiai{
Dài: $160 + 140 = 300$ km.
}
\end{ex}

\begin{ex}
Hai số có tổng bằng 950, số hạng thứ nhất là 420. Số hạng thứ hai là bao nhiêu?
\shortans{530}
\loigiai{
Số hạng thứ hai: $950 - 420 = 530$.
}
\end{ex}

\begin{ex}
Hai số có hiệu bằng 320, số trừ là 280. Số bị trừ là bao nhiêu?
\shortans{600}
\loigiai{
Số bị trừ: $320 + 280 = 600$.
}
\end{ex}

\begin{ex}
Tính tổng của số tròn trăm lớn nhất có ba chữ số và số 85.
\shortans{985}
\loigiai{
Số tròn trăm lớn nhất có 3 chữ số là 900. Tổng: $900 + 85 = 985$.
}
\end{ex}

\begin{ex}
Lấy 1000 trừ đi số tròn trăm liền sau 400 thì được bao nhiêu?
\shortans{500}
\loigiai{
Số tròn trăm liền sau 400 là 500. Hiệu: $1000 - 500 = 500$.
}
\end{ex}

\begin{ex}
Tính: $250 + 350 + 400$.
\shortans{1000}
\loigiai{
Ta có $250 + 350 = 600$, rồi $600 + 400 = 1000$.
}
\end{ex}

\newpage

\TaoPhieuBaiTap{Chủ đề 14 -- Ôn tập cuối năm}{71}{Ôn tập phép cộng, phép trừ trong phạm vi 1000}{https://loptoanthaylinh.vn/bai-tap/toan-2/bai-71}
\begin{ex}
Tính: $2 \times 6$.
\shortans{12}
\loigiai{
Ta có $2 \times 6 = 12$.
}
\end{ex}

\begin{ex}
Tính: $2 \times 8$.
\shortans{16}
\loigiai{
Ta có $2 \times 8 = 16$.
}
\end{ex}

\begin{ex}
Tính: $2 \times 9$.
\shortans{18}
\loigiai{
Ta có $2 \times 9 = 18$.
}
\end{ex}

\begin{ex}
Tính: $5 \times 4$.
\shortans{20}
\loigiai{
Ta có $5 \times 4 = 20$.
}
\end{ex}

\begin{ex}
Tính: $5 \times 7$.
\shortans{35}
\loigiai{
Ta có $5 \times 7 = 35$.
}
\end{ex}

\begin{ex}
Tính: $5 \times 9$.
\shortans{45}
\loigiai{
Ta có $5 \times 9 = 45$.
}
\end{ex}

\begin{ex}
Tính: $14 : 2$.
\shortans{7}
\loigiai{
Ta có $14 : 2 = 7$.
}
\end{ex}

\begin{ex}
Tính: $18 : 2$.
\shortans{9}
\loigiai{
Ta có $18 : 2 = 9$.
}
\end{ex}

\begin{ex}
Tính: $35 : 5$.
\shortans{7}
\loigiai{
Ta có $35 : 5 = 7$.
}
\end{ex}

\begin{ex}
Tính: $45 : 5$.
\shortans{9}
\loigiai{
Ta có $45 : 5 = 9$.
}
\end{ex}

\begin{ex}
Tính giá trị biểu thức: $2 \times 7 + 16$.
\shortans{30}
\loigiai{
Ta có $2 \times 7 = 14$, rồi $14 + 16 = 30$.
}
\end{ex}

\begin{ex}
Tính giá trị biểu thức: $5 \times 8 - 15$.
\shortans{25}
\loigiai{
Ta có $5 \times 8 = 40$, rồi $40 - 15 = 25$.
}
\end{ex}

\begin{ex}
Tính giá trị biểu thức: $16 : 2 + 22$.
\shortans{30}
\loigiai{
Ta có $16 : 2 = 8$, rồi $8 + 22 = 30$.
}
\end{ex}

\begin{ex}
Tính giá trị biểu thức: $40 : 5 + 32$.
\shortans{40}
\loigiai{
Ta có $40 : 5 = 8$, rồi $8 + 32 = 40$.
}
\end{ex}

\begin{ex}
Tính giá trị biểu thức: $2 \times 5 \times 4$.
\shortans{40}
\loigiai{
Ta có $2 \times 5 = 10$, rồi $10 \times 4 = 40$.
}
\end{ex}

\begin{ex}
Tính giá trị biểu thức: $20 : 2 : 2$.
\shortans{5}
\loigiai{
Ta có $20 : 2 = 10$, rồi $10 : 2 = 5$.
}
\end{ex}

\begin{ex}
Điền số thích hợp: $2 \times \dots = 16$.
\shortans{8}
\loigiai{
Số cần điền là: $16 : 2 = 8$.
}
\end{ex}

\begin{ex}
Điền số thích hợp: $5 \times \dots = 30$.
\shortans{6}
\loigiai{
Số cần điền là: $30 : 5 = 6$.
}
\end{ex}

\begin{ex}
Điền số thích hợp: $\dots : 2 = 9$.
\shortans{18}
\loigiai{
Số cần điền là: $9 \times 2 = 18$.
}
\end{ex}

\begin{ex}
Điền số thích hợp: $\dots : 5 = 7$.
\shortans{35}
\loigiai{
Số cần điền là: $7 \times 5 = 35$.
}
\end{ex}

\begin{ex}
Có 5 can dầu, mỗi can chứa 5 lít. Hỏi cả 5 can chứa bao nhiêu lít dầu?
\shortans{25}
\loigiai{
Chứa: $5 \times 5 = 25$ lít.
}
\end{ex}

\begin{ex}
Có 18 viên bi chia đều cho 2 bạn. Mỗi bạn được bao nhiêu viên bi?
\shortans{9}
\loigiai{
Mỗi bạn được: $18 : 2 = 9$ viên bi.
}
\end{ex}

\begin{ex}
Có 30 kg gạo chia đều vào 5 túi. Mỗi túi nặng bao nhiêu ki-lô-gam?
\shortans{6}
\loigiai{
Mỗi túi nặng: $30 : 5 = 6$ kg.
}
\end{ex}

\begin{ex}
Mỗi bạn học sinh có 2 bàn tay. Hỏi 8 bạn học sinh có tất cả bao nhiêu bàn tay?
\shortans{16}
\loigiai{
Có: $2 \times 8 = 16$ bàn tay.
}
\end{ex}

\begin{ex}
Tính: $5 \times 6 : 5$.
\shortans{6}
\loigiai{
Ta có $5 \times 6 = 30$, rồi $30 : 5 = 6$.
}
\end{ex}

\newpage

\TaoPhieuBaiTap{Chủ đề 14 -- Ôn tập cuối năm}{72}{Ôn tập phép nhân, phép chia}{https://loptoanthaylinh.vn/bai-tap/toan-2/bai-72}
\begin{ex}
Đường gấp khúc ABCD gồm bao nhiêu đoạn thẳng?
\shortans{3}
\loigiai{
Gồm 3 đoạn thẳng: AB, BC, CD.
}
\end{ex}

\begin{ex}
Hình tứ giác có bao nhiêu đỉnh?
\shortans{4}
\loigiai{
Có 4 đỉnh.
}
\end{ex}

\begin{ex}
Hình tam giác có bao nhiêu cạnh?
\shortans{3}
\loigiai{
Có 3 cạnh.
}
\end{ex}

\begin{ex}
Khối trụ có bao nhiêu mặt phẳng đáy?
\shortans{2}
\loigiai{
Có 2 mặt đáy phẳng.
}
\end{ex}

\begin{ex}
Khối cầu có bao nhiêu đỉnh?
\shortans{0}
\loigiai{
Khối cầu có 0 đỉnh.
}
\end{ex}

\begin{ex}
Khối lập phương có bao nhiêu mặt vuông?
\shortans{6}
\loigiai{
Có 6 mặt vuông.
}
\end{ex}

\begin{ex}
Khối hộp chữ nhật có bao nhiêu đỉnh?
\shortans{8}
\loigiai{
Có 8 đỉnh.
}
\end{ex}

\begin{ex}
Khối hộp chữ nhật có bao nhiêu cạnh?
\shortans{12}
\loigiai{
Có 12 cạnh.
}
\end{ex}

\begin{ex}
Đường gấp khúc MNPQ gồm 3 đoạn thẳng dài lần lượt 10 cm, 20 cm, 30 cm. Độ dài đường gấp khúc là bao nhiêu xăng-ti-mét?
\shortans{60}
\loigiai{
Độ dài: $10 + 20 + 30 = 60$ cm.
}
\end{ex}

\begin{ex}
Chu vi hình tam giác có 3 cạnh dài lần lượt 12 cm, 15 cm, 18 cm là bao nhiêu xăng-ti-mét?
\shortans{45}
\loigiai{
Chu vi: $12 + 15 + 18 = 45$ cm.
}
\end{ex}

\begin{ex}
Chu vi hình tứ giác có độ dài các cạnh lần lượt là 10 cm, 15 cm, 20 cm, 25 cm là bao nhiêu xăng-ti-mét?
\shortans{70}
\loigiai{
Chu vi: $10 + 15 + 20 + 25 = 70$ cm.
}
\end{ex}

\begin{ex}
Một hình vuông có cạnh dài 9 cm. Chu vi hình vuông đó là bao nhiêu xăng-ti-mét?
\shortans{36}
\loigiai{
Chu vi: $9 \times 4 = 36$ cm.
}
\end{ex}

\begin{ex}
Một hình chữ nhật có chiều dài 15 cm, chiều rộng 10 cm. Chu vi hình chữ nhật đó là bao nhiêu xăng-ti-mét?
\shortans{50}
\loigiai{
Chu vi: $(15 + 10) \times 2 = 50$ cm.
}
\end{ex}

\begin{ex}
Cho 3 điểm A, B, C thẳng hàng theo thứ tự đó. Biết AB dài 14 cm, BC dài 16 cm. Độ dài đoạn AC là bao nhiêu xăng-ti-mét?
\shortans{30}
\loigiai{
Độ dài AC: $14 + 16 = 30$ cm.
}
\end{ex}

\begin{ex}
Đoạn thẳng MN dài 25 cm, điểm P nằm giữa M và N sao cho MP dài 10 cm. Đoạn PN dài bao nhiêu xăng-ti-mét?
\shortans{15}
\loigiai{
Đoạn PN dài: $25 - 10 = 15$ cm.
}
\end{ex}

\begin{ex}
Cần ít nhất bao nhiêu que tính để xếp được 1 hình vuông?
\shortans{4}
\loigiai{
Cần 4 que tính.
}
\end{ex}

\begin{ex}
Cần ít nhất bao nhiêu que tính để xếp được 1 hình tam giác?
\shortans{3}
\loigiai{
Cần 3 que tính.
}
\end{ex}

\begin{ex}
Dùng 15 que tính bằng nhau xếp được nhiều nhất bao nhiêu hình tam giác rời nhau?
\shortans{5}
\loigiai{
Xếp được: $15 : 3 = 5$ hình tam giác.
}
\end{ex}

\begin{ex}
Dùng 16 que tính bằng nhau xếp được nhiều nhất bao nhiêu hình vuông rời nhau?
\shortans{4}
\loigiai{
Xếp được: $16 : 4 = 4$ hình vuông.
}
\end{ex}

\begin{ex}
Hai hình tam giác rời nhau có tổng số cạnh là bao nhiêu?
\shortans{6}
\loigiai{
Có: $3 + 3 = 6$ cạnh.
}
\end{ex}

\begin{ex}
Hai hình vuông rời nhau có tổng số góc vuông là bao nhiêu?
\shortans{8}
\loigiai{
Có: $4 + 4 = 8$ góc vuông.
}
\end{ex}

\begin{ex}
Một hình tam giác có chu vi 36 cm, 3 cạnh có độ dài bằng nhau. Mỗi cạnh dài bao nhiêu xăng-ti-mét?
\shortans{12}
\loigiai{
Mỗi cạnh dài: $36 : 3 = 12$ cm.
}
\end{ex}

\begin{ex}
Một hình tứ giác có chu vi 60 cm, 4 cạnh có độ dài bằng nhau. Mỗi cạnh dài bao nhiêu xăng-ti-mét?
\shortans{15}
\loigiai{
Mỗi cạnh dài: $60 : 4 = 15$ cm.
}
\end{ex}

\begin{ex}
Tổng số mặt phẳng đáy của 2 khối trụ rời nhau là bao nhiêu?
\shortans{4}
\loigiai{
Mỗi khối trụ có 2 mặt đáy phẳng, hai khối có: $2 + 2 = 4$ mặt đáy phẳng.
}
\end{ex}

\begin{ex}
Hình chữ nhật có mấy cạnh?
\shortans{4}
\loigiai{
Hình chữ nhật có 4 cạnh.
}
\end{ex}

\newpage

\TaoPhieuBaiTap{Chủ đề 14 -- Ôn tập cuối năm}{73}{Ôn tập hình học và đo lường}{https://loptoanthaylinh.vn/bai-tap/toan-2/bai-73}
\begin{ex}
1 mét bằng bao nhiêu xăng-ti-mét?
\shortans{100}
\loigiai{
$1\text{ m} = 100\text{ cm}$.
}
\end{ex}

\begin{ex}
1 mét bằng bao nhiêu đề-xi-mét?
\shortans{10}
\loigiai{
$1\text{ m} = 10\text{ dm}$.
}
\end{ex}

\begin{ex}
1 ki-lô-mét bằng bao nhiêu mét?
\shortans{1000}
\loigiai{
$1\text{ km} = 1000\text{ m}$.
}
\end{ex}

\begin{ex}
Một ngày có bao nhiêu giờ?
\shortans{24}
\loigiai{
Một ngày có 24 giờ.
}
\end{ex}

\begin{ex}
Một giờ có bao nhiêu phút?
\shortans{60}
\loigiai{
Một giờ có 60 phút.
}
\end{ex}

\begin{ex}
Một năm có bao nhiêu tháng?
\shortans{12}
\loigiai{
Một năm có 12 tháng.
}
\end{ex}

\begin{ex}
Tháng 7 có bao nhiêu ngày?
\shortans{31}
\loigiai{
Tháng 7 có 31 ngày.
}
\end{ex}

\begin{ex}
Tháng 9 có bao nhiêu ngày?
\shortans{30}
\loigiai{
Tháng 9 có 30 ngày.
}
\end{ex}

\begin{ex}
Tính: $45\text{ kg} + 38\text{ kg}$.
\shortans{83}
\loigiai{
Ta có: $45 + 38 = 83$ kg.
}
\end{ex}

\begin{ex}
Tính: $92\text{ kg} - 45\text{ kg}$.
\shortans{47}
\loigiai{
Ta có: $92 - 45 = 47$ kg.
}
\end{ex}

\begin{ex}
Tính: $35\text{ }l + 28\text{ }l$.
\shortans{63}
\loigiai{
Ta có: $35 + 28 = 63$ lít.
}
\end{ex}

\begin{ex}
Tính: $80\text{ }l - 35\text{ }l$.
\shortans{45}
\loigiai{
Ta có: $80 - 35 = 45$ lít.
}
\end{ex}

\begin{ex}
Tính: $450\text{ m} + 350\text{ m}$.
\shortans{800}
\loigiai{
Ta có: $450 + 350 = 800$ m.
}
\end{ex}

\begin{ex}
Tính: $1000\text{ m} - 650\text{ m}$.
\shortans{350}
\loigiai{
Ta có: $1000 - 650 = 350$ m.
}
\end{ex}

\begin{ex}
Bác An có 1000 đồng, bác mua 1 mớ rau hết 700 đồng. Bác An còn lại bao nhiêu đồng?
\shortans{300}
\loigiai{
Còn lại: $1000 - 700 = 300$ đồng.
}
\end{ex}

\begin{ex}
Đồng hồ chỉ 15 giờ 30 phút. Đây là mấy giờ rưỡi chiều?
\shortans{3}
\loigiai{
Là 3 giờ rưỡi chiều.
}
\end{ex}

\begin{ex}
Đồng hồ chỉ 20 giờ. Đây là mấy giờ tối?
\shortans{8}
\loigiai{
Là 8 giờ tối.
}
\end{ex}

\begin{ex}
Nếu Thứ Ba tuần này là ngày 12 thì Thứ Ba tuần sau là ngày bao nhiêu?
\shortans{19}
\loigiai{
Ngày đó là: $12 + 7 = 19$.
}
\end{ex}

\begin{ex}
Bao gạo nặng 45 kg, bao ngô nặng 38 kg. Cả hai bao nặng bao nhiêu ki-lô-gam?
\shortans{83}
\loigiai{
Cả hai bao nặng: $45 + 38 = 83$ kg.
}
\end{ex}

\begin{ex}
Thùng to có 55 lít nước ngọt, thùng nhỏ có ít hơn thùng to 18 lít. Thùng nhỏ có bao nhiêu lít nước ngọt?
\shortans{37}
\loigiai{
Thùng nhỏ có: $55 - 18 = 37$ lít.
}
\end{ex}

\begin{ex}
Một tấm vải dài 80 m, người ta cắt đi 35 m. Tấm vải còn lại dài bao nhiêu mét?
\shortans{45}
\loigiai{
Còn lại: $80 - 35 = 45$ m.
}
\end{ex}

\begin{ex}
Quãng đường từ A đến B dài 48 km, từ B đến C dài 35 km. Quãng đường từ A đến C đi qua B dài bao nhiêu ki-lô-mét?
\shortans{83}
\loigiai{
Dài: $48 + 35 = 83$ km.
}
\end{ex}

\begin{ex}
Một ngày bé ngủ 9 giờ. Thời gian bé thức trong một ngày là bao nhiêu giờ?
\shortans{15}
\loigiai{
Thời gian thức: $24 - 9 = 15$ giờ.
}
\end{ex}

\begin{ex}
Tính giá trị biểu thức: $25\text{ kg} + 35\text{ kg} - 20\text{ kg}$.
\shortans{40}
\loigiai{
Ta có $25 + 35 = 60$, rồi $60 - 20 = 40$ kg.
}
\end{ex}

\begin{ex}
Có 5 can nước loại 5 lít. Cả 5 can chứa tất cả bao nhiêu lít nước?
\shortans{25}
\loigiai{
Chứa: $5 \times 5 = 25$ lít.
}
\end{ex}

\newpage

\TaoPhieuBaiTap{Chủ đề 14 -- Ôn tập cuối năm}{74}{Ôn tập yếu tố thống kê, xác suất}{https://loptoanthaylinh.vn/bai-tap/toan-2/bai-74}
\begin{ex}
Trong sân có: 6 con gà trống, 9 con gà mái, 5 con vịt. Có tất cả bao nhiêu con gia cầm?
\shortans{20}
\loigiai{
Có: $6 + 9 + 5 = 20$ con.
}
\end{ex}

\begin{ex}
Số gà mái nhiều hơn số gà trống bao nhiêu con?
\shortans{3}
\loigiai{
Nhiều hơn: $9 - 6 = 3$ con.
}
\end{ex}

\begin{ex}
Có tất cả bao nhiêu con gà?
\shortans{15}
\loigiai{
Số gà: $6 + 9 = 15$ con.
}
\end{ex}

\begin{ex}
Trên biểu đồ tranh, mỗi ngôi sao tương ứng với 2 điểm tốt. Bạn An có 6 ngôi sao. Bạn An có bao nhiêu điểm tốt?
\shortans{12}
\loigiai{
An có: $2 \times 6 = 12$ điểm tốt.
}
\end{ex}

\begin{ex}
Bạn Bình có 8 ngôi sao (mỗi ngôi sao là 2 điểm tốt). Bình có bao nhiêu điểm tốt?
\shortans{16}
\loigiai{
Bình có: $2 \times 8 = 16$ điểm tốt.
}
\end{ex}

\begin{ex}
Bạn Bình nhiều hơn bạn An bao nhiêu điểm tốt?
\shortans{4}
\loigiai{
Nhiều hơn: $16 - 12 = 4$ điểm tốt.
}
\end{ex}

\begin{ex}
Cả hai bạn có tất cả bao nhiêu điểm tốt?
\shortans{28}
\loigiai{
Có: $12 + 16 = 28$ điểm tốt.
}
\end{ex}

\begin{ex}
Trong hộp chỉ có các viên bi màu vàng. Lấy ngẫu nhiên 1 viên bi, sự kiện 'Lấy được viên bi màu vàng' là: Chắc chắn = 1, Có thể = 2, Không thể = 3.
\shortans{1}
\loigiai{
Toàn bi vàng nên chắc chắn. Đáp án là 1.
}
\end{ex}

\begin{ex}
Trong hộp đó chỉ có bi vàng, sự kiện 'Lấy được viên bi màu đỏ' là: Chắc chắn = 1, Có thể = 2, Không thể = 3.
\shortans{3}
\loigiai{
Không có bi đỏ nên không thể. Đáp án là 3.
}
\end{ex}

\begin{ex}
Trong hộp có cả bi xanh và bi đỏ. Lấy 1 viên bi, sự kiện 'Lấy được bi đỏ' là: Chắc chắn = 1, Có thể = 2, Không thể = 3.
\shortans{2}
\loigiai{
Có thể lấy được bi đỏ hoặc xanh. Đáp án là 2.
}
\end{ex}

\begin{ex}
Gieo một con xúc xắc 6 mặt. Mặt xuất hiện có số chấm lớn hơn 6 là: Chắc chắn = 1, Có thể = 2, Không thể = 3.
\shortans{3}
\loigiai{
Không thể có quá 6 chấm. Đáp án là 3.
}
\end{ex}

\begin{ex}
Mặt xuất hiện có số chấm từ 1 đến 6 là: Chắc chắn = 1, Có thể = 2, Không thể = 3.
\shortans{1}
\loigiai{
Chắc chắn vì xúc xắc có 6 mặt từ 1 đến 6. Đáp án là 1.
}
\end{ex}

\begin{ex}
Khảo sát loài hoa yêu thích của 30 bạn: 14 bạn thích hoa hồng, 9 bạn thích hoa cúc, còn lại thích hoa sen. Có bao nhiêu bạn thích hoa sen?
\shortans{7}
\loigiai{
Số bạn thích hoa sen: $30 - 14 - 9 = 7$ bạn.
}
\end{ex}

\begin{ex}
Số bạn thích hoa hồng nhiều hơn hoa cúc bao nhiêu bạn?
\shortans{5}
\loigiai{
Nhiều hơn: $14 - 9 = 5$ bạn.
}
\end{ex}

\begin{ex}
Kiểm đếm phương tiện trong bãi đỗ: 15 xe ô tô, 25 xe máy, 10 xe đạp. Có tất cả bao nhiêu phương tiện?
\shortans{50}
\loigiai{
Tổng số phương tiện: $15 + 25 + 10 = 50$ chiếc.
}
\end{ex}

\begin{ex}
Số xe máy nhiều hơn số xe ô tô bao nhiêu chiếc?
\shortans{10}
\loigiai{
Nhiều hơn: $25 - 15 = 10$ chiếc.
}
\end{ex}

\begin{ex}
Số xe đạp ít hơn số xe máy bao nhiêu chiếc?
\shortans{15}
\loigiai{
Ít hơn: $25 - 10 = 15$ chiếc.
}
\end{ex}

\begin{ex}
Trong túi có 4 quả bóng xanh và 1 quả bóng đỏ. Nam lấy ra 2 quả bóng. Sự kiện 'Trong đó có ít nhất 1 quả bóng xanh' là: Chắc chắn = 1, Có thể = 2, Không thể = 3.
\shortans{1}
\loigiai{
Vì chỉ có 1 bóng đỏ nên lấy 2 quả chắc chắn có bóng xanh. Đáp án là 1.
}
\end{ex}

\begin{ex}
Một năm có 365 ngày (năm thường) hoặc 366 ngày (năm nhuận). Sự kiện 'Một năm có 400 ngày' là: Chắc chắn = 1, Có thể = 2, Không thể = 3.
\shortans{3}
\loigiai{
Không thể xảy ra. Đáp án là 3.
}
\end{ex}

\begin{ex}
Hôm nay là Thứ Hai, ngày mai là Thứ Ba. Sự kiện này là: Chắc chắn = 1, Có thể = 2, Không thể = 3.
\shortans{1}
\loigiai{
Chắc chắn đúng. Đáp án là 1.
}
\end{ex}

\begin{ex}
Một bạn học sinh ném bóng vào rổ, bóng lọt vào rổ: Chắc chắn = 1, Có thể = 2, Không thể = 3.
\shortans{2}
\loigiai{
Có thể trúng hoặc trượt. Đáp án là 2.
}
\end{ex}

\begin{ex}
Trong rổ có 10 quả cam và 12 quả xoài. Tổng số quả trong rổ là bao nhiêu?
\shortans{22}
\loigiai{
Có: $10 + 12 = 22$ quả.
}
\end{ex}

\begin{ex}
Số xoài nhiều hơn số cam bao nhiêu quả?
\shortans{2}
\loigiai{
Nhiều hơn: $12 - 10 = 2$ quả.
}
\end{ex}

\begin{ex}
Có bao nhiêu sự kiện 'Chắc chắn' trong 2 sự kiện: (A) Nước sôi ở 100 độ C; (B) Con voi biết bay?
\shortans{1}
\loigiai{
Chỉ có A chắc chắn (1 sự kiện).
}
\end{ex}

\begin{ex}
Sự kiện 'Con voi biết bay' là: Chắc chắn = 1, Có thể = 2, Không thể = 3.
\shortans{3}
\loigiai{
Không thể xảy ra. Đáp án là 3.
}
\end{ex}

\newpage

\TaoPhieuBaiTap{Chủ đề 14 -- Ôn tập cuối năm}{75}{Ôn tập chung cuối năm}{https://loptoanthaylinh.vn/bai-tap/toan-2/bai-75}
\begin{ex}
Số lớn nhất có ba chữ số là số nào?
\shortans{999}
\loigiai{
Số đó là 999.
}
\end{ex}

\begin{ex}
Số bé nhất có ba chữ số là số nào?
\shortans{100}
\loigiai{
Số đó là 100.
}
\end{ex}

\begin{ex}
Số liền sau của số lớn nhất có ba chữ số là số nào?
\shortans{1000}
\loigiai{
Số liền sau của 999 là 1000.
}
\end{ex}

\begin{ex}
Tính: $348 + 275$.
\shortans{623}
\loigiai{
Ta có $348 + 275 = 623$.
}
\end{ex}

\begin{ex}
Tính: $752 - 384$.
\shortans{368}
\loigiai{
Ta có $752 - 384 = 368$.
}
\end{ex}

\begin{ex}
Tính: $5 \times 8$.
\shortans{40}
\loigiai{
Ta có $5 \times 8 = 40$.
}
\end{ex}

\begin{ex}
Tính: $18 : 2$.
\shortans{9}
\loigiai{
Ta có $18 : 2 = 9$.
}
\end{ex}

\begin{ex}
Tính: $45 : 5$.
\shortans{9}
\loigiai{
Ta có $45 : 5 = 9$.
}
\end{ex}

\begin{ex}
Tính giá trị biểu thức: $2 \times 9 + 35$.
\shortans{53}
\loigiai{
Ta có $2 \times 9 = 18$, rồi $18 + 35 = 53$.
}
\end{ex}

\begin{ex}
Tính giá trị biểu thức: $40 : 5 + 48$.
\shortans{56}
\loigiai{
Ta có $40 : 5 = 8$, rồi $8 + 48 = 56$.
}
\end{ex}

\begin{ex}
Tính giá trị biểu thức: $500 + 350 - 250$.
\shortans{600}
\loigiai{
Ta có $500 + 350 = 850$, rồi $850 - 250 = 600$.
}
\end{ex}

\begin{ex}
Điền số thích hợp: $456 + \dots = 789$.
\shortans{333}
\loigiai{
Số cần điền: $789 - 456 = 333$.
}
\end{ex}

\begin{ex}
Điền số thích hợp: $800 - \dots = 345$.
\shortans{455}
\loigiai{
Số cần điền: $800 - 345 = 455$.
}
\end{ex}

\begin{ex}
Hình chữ nhật có bao nhiêu góc vuông?
\shortans{4}
\loigiai{
Có 4 góc vuông.
}
\end{ex}

\begin{ex}
Hình tam giác có bao nhiêu cạnh?
\shortans{3}
\loigiai{
Có 3 cạnh.
}
\end{ex}

\begin{ex}
Khối lập phương có bao nhiêu đỉnh?
\shortans{8}
\loigiai{
Có 8 đỉnh.
}
\end{ex}

\begin{ex}
Khối trụ có bao nhiêu mặt phẳng đáy?
\shortans{2}
\loigiai{
Có 2 mặt phẳng đáy.
}
\end{ex}

\begin{ex}
Chu vi hình chữ nhật có chiều dài 20 cm, chiều rộng 10 cm là bao nhiêu xăng-ti-mét?
\shortans{60}
\loigiai{
Chu vi: $(20 + 10) \times 2 = 60$ cm.
}
\end{ex}

\begin{ex}
Chu vi hình tam giác có 3 cạnh bằng nhau, mỗi cạnh dài 9 cm là bao nhiêu xăng-ti-mét?
\shortans{27}
\loigiai{
Chu vi: $9 \times 3 = 27$ cm.
}
\end{ex}

\begin{ex}
Một tuần lễ có bao nhiêu ngày?
\shortans{7}
\loigiai{
Có 7 ngày.
}
\end{ex}

\begin{ex}
Một ngày có bao nhiêu giờ?
\shortans{24}
\loigiai{
Có 24 giờ.
}
\end{ex}

\begin{ex}
Tháng 1 có bao nhiêu ngày?
\shortans{31}
\loigiai{
Có 31 ngày.
}
\end{ex}

\begin{ex}
Lớp 2A có 35 học sinh, lớp 2B có 38 học sinh. Cả hai lớp có bao nhiêu học sinh?
\shortans{73}
\loigiai{
Cả hai lớp có: $35 + 38 = 73$ học sinh.
}
\end{ex}

\begin{ex}
Quãng đường từ nhà An đến trường dài 250 m, từ trường đến nhà bà dài 350 m. Quãng đường từ nhà An đến nhà bà đi qua trường dài bao nhiêu mét?
\shortans{600}
\loigiai{
Dài: $250 + 350 = 600$ m.
}
\end{ex}

\begin{ex}
Tính: $1000 - 250 - 350 - 100$.
\shortans{300}
\loigiai{
Ta có $1000 - 250 = 750$, $750 - 350 = 400$, $400 - 100 = 300$.
}
\end{ex}

\newpage

\end{document}
